\documentclass{amsart}
\usepackage{amssymb}
\newtheorem{theorem}{Theorem}[section]
\newtheorem{proposition}[theorem]{Proposition}
\newtheorem{lemma}[theorem]{Lemma}
\newtheorem{corollary}[theorem]{Corollary}
\theoremstyle{definition}
\newtheorem{definition}[theorem]{Definition}
\theoremstyle{remark}
\newtheorem*{remark}{Remark}
\numberwithin{equation}{section}
\makeatletter
\def\subsection{\@startsection{subsection}{2}\z@{.5\linespacing\@plus.7\linespacing}{-.5em}{\normalfont\bfseries\boldmath}}
\makeatother
\newcommand{\RunningHead}{Parking functions by maximum displacement}
\uppercasenonmath\RunningHead
\begin{document}
\title{Unimodality of parking functions counted by maximum displacement}\subjclass[2020]{Primary 05A20; Secondary 05A15, 05A16}\keywords{Parking functions, maximum displacement, unimodality, linear probing}
\begin{abstract}
We prove that for every $n\ge1$ the numbers of parking functions of length $n$
with maximum displacement $0,1,\dots,n-1$ form a unimodal sequence, as
conjectured by Celano, Elder, Hadaway, Harris, Martin, Priestley and Udell.
More precisely, for $n\ge2$ the sequence rises strictly up to a term, its
peak, and falls strictly from there on, except that the term after the peak
may be equal to it. The position of the peak in the sequence for $n+1$ is
the same as in the sequence for $n$ or larger by one, and in the second
case no two neighbouring terms of the sequence for $n+1$ are equal. The
proof starts from a recursion for the numbers of parking functions whose
maximum displacement is bounded, and it combines analytic estimates with
exact computations of these numbers for $n\le300$.
\end{abstract}
\maketitle
\sbox0{\normalfont\scriptsize\RunningHead}\ifdim\wd0>.85\textwidth\markboth{}{}\else\markboth{\RunningHead}{\RunningHead}\fi

\section{Introduction}\label{sec:intro}
\subsection{Parking functions and their maximum displacement}\label{subsec:intro-problem}

Konheim and Weiss \cite[Sections~3 and~6]{konheimweiss1966} studied an
occupancy discipline in which balls are placed one after another into cells
that are arranged in a circle, each ball into the first free cell at or after
a cell that is assigned to it, and they applied it to cars that park along a
street. The sequences of assignments under which
every car finds a spot in the street are the parking functions. Hashing with
linear probing places the items of a table by the same discipline, and the
distance between the place that a car or an item takes and the place assigned
to it, its displacement, has been studied in both settings
\cite{celano2025,pappe2024,petersson2009,seligzhu2024}. This paper is about
the largest displacement of a car in a parking function, and about the way in
which the parking functions of a given length are distributed over its
values.

Let $n\ge1$ be an integer. A \emph{preference sequence} of length $n$ is a
sequence $\alpha=(\alpha_1,\dots,\alpha_n)\in\{1,\dots,n\}^n$. It is read as
follows. Spots $1,\dots,n$ lie in a row along a street and are all empty at
the start. Cars
$1,\dots,n$ arrive in this order, and car $i$ prefers spot $\alpha_i$: it
looks from spot $\alpha_i$ onward, in the direction of increasing spot
numbers, and parks in the first empty spot it meets. If the spots
$\alpha_i,\alpha_i+1,\dots,n$ are all occupied, car $i$ cannot park. This is
the \emph{parking process}.
If every car parks, $\alpha$ is called a
\emph{parking function} of length $n$, and the set of all parking functions of
length $n$ is written $\mathrm{PF}_n$. For example, with the preference
sequence $(1,2,1)$ the three cars park in spots $1$, $2$ and $3$, whereas
with $(2,2,3)$ the third car finds spot $3$ occupied and cannot park. The
number of parking functions of length $n$ is $(n+1)^{n-1}$; this count
follows from Lemma~2 and formula~(6.1) of the article of Konheim and Weiss
\cite{konheimweiss1966}, and
Foata and Riordan \cite[Section~2]{foatariordan1974} give a proof by an
argument of Pollak.

For $\alpha\in\mathrm{PF}_n$ let $\operatorname{pos}_i(\alpha)$ be the spot in
which car $i$ parks. The \emph{displacement} of car $i$ is the nonnegative
integer $\operatorname{pos}_i(\alpha)-\alpha_i$, the distance between the spot
that the car takes and the spot that it prefers, and the \emph{maximum
displacement} of $\alpha$ is
\[
M(\alpha)=\max_{1\le i\le n}\bigl(\operatorname{pos}_i(\alpha)-\alpha_i\bigr)
\in\{0,1,\dots,n-1\}.
\]
For $0\le\ell\le n-1$ we put
\[
b_{n,\ell}=\#\{\alpha\in\mathrm{PF}_n:\ M(\alpha)=\ell\}.
\]
The numbers $b_{n,0},\dots,b_{n,n-1}$ add up to $(n+1)^{n-1}$. The two ends of
this sequence are known \cite[Example~2.8 and Section~6]{celano2025}: the
parking functions with maximum displacement $0$ are the permutations, so
that $b_{n,0}=n!$, and $b_{n,n-1}=n^{n-2}$ for $n\ge2$. The sequences for
$n\le7$ are as follows \cite[Table~1]{celano2025}.
\begin{equation}\label{eq:intro-rows}
\begin{array}{c|l}
n & b_{n,0},\ b_{n,1},\ \dots,\ b_{n,n-1}\\
\hline
1 & 1\\
2 & 2,\ 1\\
3 & 6,\ 7,\ 3\\
4 & 24,\ 51,\ 34,\ 16\\
5 & 120,\ 421,\ 377,\ 253,\ 125\\
6 & 720,\ 3963,\ 4594,\ 3688,\ 2546,\ 1296\\
7 & 5040,\ 42253,\ 62145,\ 57398,\ 46142,\ 32359,\ 16807
\end{array}
\end{equation}

Each of these sequences rises to a largest term and falls after it. A finite
sequence $(v_0,\dots,v_m)$ of numbers is \emph{unimodal} if there is
$j^*\in\{0,\dots,m\}$ such that $v_{j-1}\le v_j$ for $1\le j\le j^*$ and
$v_j\ge v_{j+1}$ for $j^*\le j\le m-1$; equality is allowed. For the total
displacement, the sum of the displacements of the cars, a result on the
shape of a sequence of counts is known. The total displacement is
essentially equivalent to the sum $\alpha_1+\dots+\alpha_n$ of the
preferences, as is noted in \cite{pappe2024}, and the generating polynomial
of this sum over the parking functions of length $n$ has log-concave
coefficients \cite[Theorem~3.1]{pappe2024}. For the maximum displacement,
Celano, Elder, Hadaway, Harris, Martin, Priestley and Udell conjectured that
the sequence $(b_{n,0},\dots,b_{n,n-1})$ is unimodal for every $n\ge1$
\cite[Conjecture~6.5]{celano2025}. In this paper we prove their conjecture.

\begin{theorem}\label{thm:main}
For every integer $n\ge1$ the sequence
\[
(b_{n,0},b_{n,1},\dots,b_{n,n-1})
\]
is unimodal.
\end{theorem}

Our second result describes these sequences more precisely and relates each
of them to the next: for $n\ge2$ the sequence rises strictly up to a term
$b_{n,p}$ and falls strictly from there on, except that $b_{n,p+1}$ may be
equal to $b_{n,p}$, and the index $p$ stays the same or increases by one
when $n$ is replaced by $n+1$. To state it we arrange the numbers $b_{n,k}$
in an array $(b_{n,k})$ whose rows are indexed by the length $n$ and whose
columns are indexed by the displacement, for which we write $k$ when it is
the index of a column:
\emph{row} $n$ of the array is $(b_{n,0},\dots,b_{n,n-1})$, and
\emph{column} $k$ consists of the entries $b_{n,k}$ with $n\ge k+1$.

It will be convenient to have the notions involved for arbitrary arrays. By
an \emph{array of natural numbers} we mean a family $f$ of nonnegative integers $f(n,k)$,
indexed by the integers $n\ge0$ and $k\ge0$; its row $n$ is
$(f(n,0),\dots,f(n,n-1))$. Let $n\ge1$ and $p\ge0$ be integers. Row $n$ of
$f$ \emph{has its peak at position $p$} if $p\le n-1$ and
\[
f(n,0)<\dots<f(n,p)\ \ge\ f(n,p+1)>\dots>f(n,n-1),
\]
and it \emph{has a strict peak at position $p$} if $p\le n-1$ and
\[
f(n,0)<\dots<f(n,p)>f(n,p+1)>\dots>f(n,n-1),
\]
where in both cases the terms after $f(n,p)$ are absent if $p=n-1$.

\begin{theorem}\label{thm:peaks}
For every integer $n\ge2$ there is an integer $p\ge0$ with the following two
properties.
\begin{enumerate}
\item Row $n$ of the array $(b_{n,k})$ has its peak at position $p$.
\item Row $n+1$ has its peak at position $p$ as well, or it has a
strict peak at position $p+1$.
\end{enumerate}
\end{theorem}

For $n\ge2$ we write $p_n$ for the least integer $p$ such that row $n$
has its peak at position $p$; by Theorem~\ref{thm:peaks} there is one. In
\eqref{eq:intro-rows} one reads off $p_2=0$, $p_3=p_4=p_5=1$ and
$p_6=p_7=2$, so that both alternatives of the second property occur.
Theorem~\ref{thm:main} follows from Theorem~\ref{thm:peaks}, and both
theorems are proved in Section~\ref{sec:staircase}.

Theorem~\ref{thm:peaks} leaves two questions about the rows open. It does
not give $p_n$ as a function of $n$. And it does not decide whether the
equality
$b_{n,p_n+1}=b_{n,p_n}$ ever occurs.

\subsection{Earlier work}\label{subsec:intro-earlier}

The paper \cite{celano2025} in which the conjecture is posed calls a parking
function an $\ell$-interval parking function when its maximum displacement
is at most $\ell$ (Definition~2.7 there). The number $b_{n,\ell}$ is
therefore the number of $\ell$-interval parking functions of length $n$ that
are not $(\ell-1)$-interval. The authors tabulate these numbers for $n\le9$
(Table~1) and give the two end values quoted above. Before the conjecture
they state (Section~6), as consequences of their enumerative results,
inequalities between the first three entries and between the last two
entries of a row, and they add that experimentally the pattern seems to
continue. Throughout, we cite this paper
in its second version, and we refer to its statements, sections and tables
by the numbers that they have in that version.

Exact descriptions are known for the \emph{cumulative counts}, which are also
the numbers that our proof works with. For integers $n\ge0$ and $\ell\ge0$
let $A_{n,\ell}$ be the number of parking functions of length $n$ whose
maximum displacement is at most $\ell$. Here the length $n=0$ is
admitted: the empty sequence is the one parking function of length $0$, and
its maximum displacement is $0$. We extend the definition of $b_{n,\ell}$
accordingly to all integers $n\ge0$ and $\ell\ge0$, as the number of parking
functions of length $n$ whose maximum displacement is $\ell$. Then
$b_{n,\ell}=0$ for $\ell\ge n\ge1$, while $b_{0,0}=1$ and $b_{0,\ell}=0$ for
$\ell\ge1$, and
\[
A_{n,\ell}=b_{n,0}+b_{n,1}+\dots+b_{n,\ell},\qquad A_{0,\ell}=1.
\]

Celano et al.\ write $A_{n,\ell}$ as a sum over outcomes. The outcome of a
parking function of length $n$ is the permutation $\sigma$ of
$\{1,\dots,n\}$ whose $i$-th entry $\sigma_i$ is the number of the car that
stands in spot $i$. For a permutation $\sigma$ and a spot $i$ let
$L_i(\sigma)$ be the length of the longest run of consecutive spots that ends
at spot $i$, spot $i$ included, and whose cars all have numbers at most
$\sigma_i$. Then the number of parking functions with outcome $\sigma$ and
maximum displacement at most $\ell$ is
\[
\prod_{i=1}^{n}\min\bigl(\ell+1,\,L_i(\sigma)\bigr),
\]
and $A_{n,\ell}$ is the sum of these products over all permutations $\sigma$
\cite[Section~4.1, (6)--(8)]{celano2025}. Aguilar-Fraga et al.\ obtain a
recursion in $n$ by fixing the spot of the car that arrives last
\cite[Theorem~3.7]{aguilarfraga2024}. With $n$ replaced by $n+1$ it reads
\begin{equation}\label{eq:intro-lastcar}
A_{n+1,\ell}=\sum_{j=0}^{n}\min(j+1,\ell+1)\binom{n}{j}A_{j,\ell}\,A_{n-j,\ell}
\end{equation}
for all $n\ge0$ and $\ell\ge0$. In the term of $j$ the last car parks in spot
$j+1$: the factor $\binom nj$ chooses the $j$ cars that stand to its left,
the factors $A_{j,\ell}$ and $A_{n-j,\ell}$ count the ways in which the cars
to its left and to its right park, and $\min(j+1,\ell+1)$ is the number of
preferences from which the last car reaches spot $j+1$ with a displacement of
at most $\ell$. We refer to \eqref{eq:intro-lastcar} as the \emph{last-car
recursion}. It is the starting point of our proof, and it is quoted in
Section~\ref{sec:counts} as Proposition~\ref{prop:lastcar}.

A limit theorem for the maximum displacement is due to Petersson
\cite[Theorem~2.1]{petersson2009}. It is stated for hashing with linear
probing, for a table with $m$ cells and $m-1$ items, conditioned on
its last cell remaining empty: if $M_m$ denotes the maximum displacement of
the items, then for every real number $u>0$ the probability that
$(m-M_m)/m^{2/3}\le u$ converges to $1-\psi(u^{3/2})$ as $m$ tends to
infinity. Here $\psi$ is an explicit power series, and it is the function
$\psi(s)=\mathbb{E}\,e^{-sJ}$, where $J$ is the integral over the interval
$[0,1]$ of the running maximum of a standard Brownian motion
\cite[Corollary~2.4]{jansonpetersson2009}. For $m=n+1$ such a table is a
parking function of length $n$ chosen uniformly at random. In terms of the
numbers $b_{n,\ell}$ the theorem therefore says that all but an arbitrarily
small fraction of the sum $(n+1)^{n-1}$ of row $n$ lies at a distance of
order $n^{2/3}$ from the end $\ell=n-1$ of the row.

There are also results on subfamilies of the parking functions. Selig and Zhu showed that for
non-decreasing parking functions the maximum displacement is the height of
the corresponding {\L}ukasiewicz path \cite[Theorem~2.1]{seligzhu2024}. For
arbitrary parking functions the two quantities can differ
\cite[Remark~2.3]{seligzhu2024}.
Ager-Hart, Beerbower, Harris, Jakovleski and McClinton count the
$\ell$-interval Fubini rankings, the Fubini rankings in which at most
$\ell+1$ competitors tie at a rank \cite[Corollary~4.1]{agerhart2026}; as
they recall, Fubini rankings are parking functions.

\subsection{Outline of the proof}\label{subsec:intro-outline}

Of the results described in the preceding subsection, our proof uses one,
the last-car recursion \eqref{eq:intro-lastcar}. We do not prove
Theorem~\ref{thm:main} directly. The shape of the rows is obtained by an
induction on $n$, and what this induction passes from row $n$ to row $n+1$
is the signs of the differences of neighbouring entries. The statement that
it yields is therefore Theorem~\ref{thm:peaks}, which describes these signs
in a row and relates each row to the next.

For integers $n\ge1$ and $k\ge1$ let
$d_{n,k}=b_{n,k}-b_{n,k-1}$ be the \emph{step} of column $k$ in row $n$; we
say that column $k$ is \emph{positive} in row $n$ if $d_{n,k}>0$, and
\emph{negative} if $d_{n,k}<0$. A \emph{position} is a pair $(n,k)$ of
integers with $1\le k\le n-1$, and the \emph{diagonal neighbour} of position
$(n,k)$ is position $(n+1,k+1)$. Theorem~\ref{thm:peaks} follows from two
rules about the signs of the steps when the length grows by one. The
\emph{column rule} says that $d_{n,k}>0$ implies $d_{n+1,k}>0$, and the
\emph{diagonal rule} says that $d_{n,k}\le0$ implies $d_{n+1,k+1}<0$, in
both cases for all integers $n\ge2$ and $1\le k\le n-1$.

We use the words column, step and position, and the names of the two rules,
for an arbitrary array $f$ of natural numbers as well, with
$d_{n,k}$ replaced by $f(n,k)-f(n,k-1)$. For every such array the two
rules imply the conclusion of Theorem~\ref{thm:peaks}, by an induction on
the row (Proposition~\ref{prop:rowinduction}): the column rule carries the
rising part of row $n$ to row $n+1$, the diagonal rule carries the rest one
column to the right and makes it strictly falling, and the one step of row
$n+1$ that neither rule reaches decides whether the peak stays or moves.

The parking process enters the proof of the two rules only through the
last-car recursion \eqref{eq:intro-lastcar}. The recursion holds for the
cumulative counts because the condition $M(\alpha)\le\ell$ is a condition on
each car separately, which is inherited by the two groups of cars into which
the last car divides a parking function; the condition $M(\alpha)=\ell$ is
not of this kind. The steps are recovered from the cumulative counts as
second differences in the bound on the displacement,
\[
d_{n,k}=A_{n,k}-2A_{n,k-1}+A_{n,k-2}\qquad(k\ge2),
\]
so that the two rules are statements on the sign of a second difference of
the solutions of a family of quadratic recursions. We write the recursion as
a differential equation for the generating series of the sequence
$(A_{n,\ell}/n!)_{n\ge0}$ and derive from it the linear relations on which the rest of
the proof is built.

One of these relations divides the array into two regions. Rows
$k+1\le n\le2k-1$ are called the \emph{first layer} of column $k$; it is
empty for $k=1$. Rows $n\ge2k$ lie \emph{beyond the first layer}. In the
first layer the step satisfies a linear recursion along the column
(Proposition~\ref{prop:columnequation}), in which $d_{n,k}$ is a
combination, with positive coefficients, of the entry $b_{n-1,k-1}$ of the
neighbouring column and of the steps of column $k$ in the rows from $k$ to
$n-1$. In every row $n\ge k+1$, and so also beyond the first layer, an
inequality of this form holds (Proposition~\ref{prop:columninequality}).

The two regions behave differently. Beyond the first layer the step is
positive at every position except eleven, the positions $(2k,k)$ with
$1\le k\le7$ and $(2k+1,k)$ with $2\le k\le5$
(Proposition~\ref{prop:beyondlayer}). At each of these eleven positions the
step is negative, at $(2,1)$ by row $2$ of \eqref{eq:intro-rows} and at the
other ten by Lemma~\ref{lem:fortyfive}, and both rules follow at once in
this region when one knows in addition that the step is negative at certain
positions next to the eleven. In the first layer a column $k\ge9$ turns positive in some row
of the layer and stays positive from the first such row to the end of the
layer, whereas columns $2\le k\le8$ are negative in their whole first layer.
We call the least row $n\ge k+1$ with $d_{n,k}>0$ the \emph{first positive
row} of column $k$. In the first layer the column rule says that a positive
step in a row of the layer is followed by a positive step in the next row of
the layer (Proposition~\ref{prop:layerpersistence}). The diagonal rule needs
more there, namely that a step of column $k$ which is not positive is
followed by a negative step of column $k+1$ one row below.

Each rule is proved region by region. Within a region the small columns are
settled by exact computation, entirely or up to a certain row from which an
analytic argument propagates the positive sign down the column, and the
large columns by analytic estimates; in the columns between them the two are
combined. The dividing columns are $82$, for the column rule in the first
layer and for the signs beyond the first layer, and $265$, for the diagonal
rule in the first layer. The results of the computations are stated as five
lemmas in Section~\ref{sec:tables}; they are obtained from the integers
$A_{n,\ell}$ in rows $n\le300$.

The difficulty of the proof lies in three places. The first is the diagonal
rule in the first layer of columns $k\ge265$. We call $r=n-k$ the
\emph{deficit} of position $(n,k)$, and we regard the position as the point
$(k,r)$ of the plane. By \emph{the curve} we mean the curve $k^2=r^3$ in
this plane, for real numbers $k$ and $r$. A
position with deficit $r$ lies \emph{behind the curve} if $r^3\le k^2$, and
\emph{beyond the curve} if $r^3>k^2$; it lies \emph{at most one row beyond
the curve} if $(r-1)^3\le k^2$. The diagonal rule in this region is obtained
from two analytic statements that have to meet.
Proposition~\ref{prop:diagonalstep} is the diagonal rule at the positions
behind the curve in columns $k\ge265$. Proposition~\ref{prop:curveincrease}
says that every column $k\ge265$ is positive in a row $j>k$ with
$(j-k-1)^3\le k^2$, that is, at a position $(j,k)$ that lies at most one
row beyond the curve. In Section~\ref{sec:staircase} the two propositions
are combined with the column rule in the first layer, which places every
step of the first layer of such a column that is not positive in a row
before row $j$, and so behind the curve, where
Proposition~\ref{prop:diagonalstep} applies.

Both propositions start from Proposition~\ref{prop:frozenstep}, which
writes the number $d_{n,k}/n!$ at every position of the first layer as
the difference of two terms: a second difference of a polynomial in a real
variable $y$, the curvature term, and a multiple of a sum over compositions,
the boundary term. The proposition also gives, in the same row and with the
same polynomial, a formula of this form for the step of column $k+1$. Read
for column $k-1$, that formula is a second decomposition of $d_{n,k}/n!$
into a curvature term and a boundary term, with the polynomial of column
$k-1$, and it holds for $k\ge3$ in rows $k\le n\le2k-3$. The sums in
question are the chain sums of
Section~\ref{sec:compositions}. For an integer $m\ge1$ and real numbers $y$
and $N$ with $N\ge m$ they are
\[
\Lambda_m(y;N)=\sum_{q\ge1}\ \sum_{0=s_0<s_1<\dots<s_q=m}\
\prod_{i=1}^{q}a_{s_i-s_{i-1}}\ \prod_{i=1}^{q-1}\frac{y}{N-s_i},
\]
where $a_i=i^{i-1}/i!$: a composition of $m$ is weighted by the factor $a_i$
for each part of size $i$ and by the factor $y/(N-s)$ for each cut at
position $s$.

In the proof of Proposition~\ref{prop:diagonalstep} the first of the two
decompositions is used at a
position behind the curve and at its diagonal neighbour, and the growth of
its two terms along the diagonal is compared. The curvature term is a sum,
and each of its terms is the product of a coefficient formed from the
numbers $a_i$ and of a sum over compositions. The sum over compositions
grows by at most an explicit factor, by H\"older's and Jensen's
inequalities (Proposition~\ref{prop:curvaturegrowth}), and so does the
coefficient, by Stirling's formula and the fact that the numbers
$e^{-i}a_i$, $i\ge1$, which are the probabilities of the Borel distribution
with parameter $1$, add up to one (Proposition~\ref{prop:deficitgrowth}).
Each term of the curvature term therefore grows by at most an explicit
factor, the product of these two, and the boundary term grows by at least
an explicit
factor, by an instance of the correlation inequality of Fortuin, Kasteleyn
and Ginibre \cite[Proposition~1]{fortuin1971} on the sets of cuts of the
compositions. That the first factor is smaller than the second is an
inequality between closed forms in $k$, $r$ and the index of the term
(Lemma~\ref{lem:closedforms}). The two conditions $r^3\le k^2$ and
$k\ge265$ of Proposition~\ref{prop:diagonalstep} are the conditions under
which we prove this inequality, and they enter it in different ways. After
logarithms, the inequality compares the bounds for the growth of a term of
the curvature term with the sum of $\log(1+1/k)$, which comes from a factor
$k/(k+1)$ in the first factor, and of the logarithm of the gain of the chain
sum in the boundary term. Its proof is a comparison of explicit constants.
The condition $r^3\le k^2$, under which $r^{3/2}$ is at most
$k$, is the range in which we carry out this comparison, and this is why
the proposition is restricted to the positions behind the curve. The bound
$k\ge265$ enters through the numerical constants in the proof of the
inequality, which are chosen for the columns from $265$ on; for the columns
up to $264$ the diagonal rule in the first layer is proved with the help of
the computations of Section~\ref{sec:tables} instead
(Proposition~\ref{prop:rulesmall}).

Proposition~\ref{prop:curveincrease} is a separate estimate, of the step at
a single position: the curvature term of the second decomposition is
bounded from
below and its boundary
term from above, and at the positions near the curve the two bounds are
compared with explicit numerical constants, for which we use Stirling's
formula.

The second place of difficulty is the column rule in the first layer of
columns $k\ge82$, and with it the positive sign of these columns beyond the
first layer. Both start from a positive row early in the first layer, which
is computed for $k\le264$ and is the row $j$ of
Proposition~\ref{prop:curveincrease} for $k\ge265$; the proof of that
proposition does not use the column rule. The obstacle is what the rows
before the first positive row contribute to the recursion along the column
at a later row: in these rows the step is not positive, and nothing is known
about its size. In the first layer we meet it by two arguments, each for the
rows that the other cannot reach. On a little less than the first half of
the layer, the recursion at row $n$ is compared with the recursion at the
first positive row. On the remaining rows
the step is bounded from below by the entries of column $k-1$, and the terms
of this bound whose coefficients may be negative are controlled by a
comparison with the renewal sequence of a positive log-convex kernel. This
comparison rests on a theorem of de Bruijn and Erd\H{o}s
\cite[Theorem~1]{debruijn1953}, by which such a renewal sequence is
log-convex. Their theorem is stated for a kernel whose sum is one, which is
not the case for the kernel of the comparison, and we include a proof of
the form that we need. Beyond the
first layer the positive sign is carried from row to row by an induction,
in which the contributions of the rows before the first positive row are
not bounded one by one; we bound a weighted sum of them, by a summation by
parts that rests on a convexity along the
diagonal of the cumulative counts, taken with suitable explicit weights.

The third place of difficulty is strictness. Both rules conclude with a
strict inequality, and the diagonal rule starts from a step that may be zero.
This is more than Theorem~\ref{thm:main} asks, since unimodality allows equal
neighbours everywhere in a row, while Theorem~\ref{thm:peaks} allows them in
one place only. Wherever a rule is read off from the signs of the steps at
certain positions, a step equal to zero has to be excluded there, and where
no argument excludes it, the signs are computed. This is the reason why
every sign of a step that Lemmas~\ref{lem:fortyfive}, \ref{lem:junction},
\ref{lem:signpair} and~\ref{lem:nextcolumn} state is a strict inequality,
$d_{n,k}<0$ or $d_{n,k}>0$.

The inequalities of H\"older and Jensen, Stirling's formula, the total mass
of the Borel distribution, the correlation
inequality of Fortuin, Kasteleyn and Ginibre and the theorem of de Bruijn
and Erd\H{o}s, which the arguments for the first two places of difficulty
use, are classical. What is particular to the proof is where they are
applied: to the two terms of the formulas of
Proposition~\ref{prop:frozenstep}, which are followed along a diagonal and
compared near the curve, and to the renewal sequence of the second place.
Proposition~\ref{prop:curveincrease} is the
statement on which both rules depend throughout columns $k\ge265$: the
diagonal rule in their first layer needs its bound for the first positive
row in the form in which it is stated, and the column rule in the first
layer and the positive signs beyond the first layer are derived from the
positive row that it supplies.

The paper is organised as follows. Section~\ref{sec:prelim} fixes
conventions and collects the facts about the numbers $a_i$ and related
sequences that several later sections use. Section~\ref{sec:counts}
quotes the last-car recursion, and Section~\ref{sec:series} derives from it the
recursion along a column in its first layer. Section~\ref{sec:propagation}
proves two inequalities for a column that hold in every row and the
propagation of a positive step down a column, and
Section~\ref{sec:tables} states the results of the exact computations.
Sections~\ref{sec:diagonal} to~\ref{sec:far} prove that a column $k\ge82$
which turns positive early in its first layer stays positive to the end of
the layer (Proposition~\ref{prop:staypositive}).
Sections~\ref{sec:compositions} to~\ref{sec:curve} prove
Proposition~\ref{prop:curveincrease} and draw its consequences for the first
layer. Sections~\ref{sec:kernel} to~\ref{sec:beyond} determine the signs
beyond the first layer. Sections~\ref{sec:gain} to~\ref{sec:step} prove
Proposition~\ref{prop:diagonalstep}. In Section~\ref{sec:staircase} the two
rules are assembled from these statements, and Theorems~\ref{thm:peaks}
and~\ref{thm:main} are proved.

\section{Preliminaries}\label{sec:prelim}
Throughout the paper an empty sum is $0$, an empty product is $1$, and
$0^0=1$. Later proofs use the following elementary inequalities. For every
real number $u\ge0$ the Taylor polynomials of $\log(1+u)$ bound it
alternately, those of even degree from below and those of odd degree from
above; we use this in the forms
\begin{gather*}
u-\frac{u^2}2\le\log(1+u)\le u-\frac{u^2}2+\frac{u^3}3,\\
u-\frac{u^2}2+\frac{u^3}3-\frac{u^4}4\le\log(1+u).
\end{gather*}
The exponential function satisfies
\[
e^{y}\le1+y+y^2\qquad(0\le y\le1.7),
\]
as is easily verified from the exponential series. By comparison with
integrals, for integers $n$ and real numbers $y$,
\begin{gather*}
\sum_{i=1}^{n}\frac1{\sqrt i}\le2\sqrt n-1\quad(n\ge1),\qquad
\sum_{i=0}^{n-1}\sqrt i\le\frac23\,n\sqrt n\quad(n\ge0),\\
\frac1{(y+1)\sqrt{y+1}}\le2\Bigl(\frac1{\sqrt y}-\frac1{\sqrt{y+1}}\Bigr)
\le\frac1{y\sqrt y}\quad(y>0).
\end{gather*}
Finally, we use Jensen's inequality for the convex function
$t\mapsto\theta^{t}$, where $\theta>0$ is a real number, with weights that
need not be normalised: if $I$ is a finite set, if $\rho_i\ge0$
and $t_i$ are real numbers for every $i\in I$, and if $\bar t$ is a real
number such that $\bar t\,\sum_{i\in I}\rho_i=\sum_{i\in I}\rho_i\,t_i$, then
\[
\theta^{\bar t}\,\sum_{i\in I}\rho_i\le\sum_{i\in I}\rho_i\,\theta^{t_i}.
\]

\subsection{The Cayley coefficients}\label{subsec:prelim-cayley}

The \emph{Cayley coefficients} are the real numbers
\[
a_0=0,\qquad a_j=\frac{j^{\,j-1}}{j!}\quad(j\ge1).
\]
By Cayley's formula, $j^{j-1}$ is the number of rooted labelled trees on $j$
vertices. The first value is
$a_1=1$, and $a_j>0$ for every $j\ge1$. The
value $a_0=0$ is chosen so that a sum over the pairs of indices with a given
sum, such as $\sum_{i=0}^{t}a_i\,a_{t-i}$, can run over all $i$ with
$0\le i\le t$: the two extreme terms contain $a_0$ and vanish.

Beside the Cayley coefficients we use the same numbers with the index shifted
by one,
\[
c_n=a_{n+1}=\frac{(n+1)^{n}}{(n+1)!}=\frac{(n+1)^{n-1}}{n!}\qquad(n\ge0).
\]
Thus $c_0=1$. By the count of the parking functions recalled in
Section~\ref{sec:intro}, there are $(n+1)^{n-1}$ parking functions of
length $n$ for every $n\ge1$. Since the empty sequence is
the one parking function of length $0$, the number $c_n$ is, for every
$n\ge0$, the number of parking functions of length $n$ divided by $n!$.

Consecutive Cayley coefficients satisfy, for every integer $j\ge1$,
\begin{equation}\label{eq:prelim-quotient}
a_{j+1}=\Bigl(1+\frac1j\Bigr)^{j-1}a_j,
\end{equation}
since both sides are equal to $(j+1)^{j-1}/j!$. The next lemma shows that
the ratio $(1+1/j)^{j-1}$ does not decrease in $j$. It is derived from the
classical fact
that the sequence $(1+1/i)^{i}$, whose limit is $e$, does not decrease: for
all integers $i$ and $j$ with $1\le i\le j$,
\begin{equation}\label{eq:prelim-euler}
\Bigl(1+\frac1i\Bigr)^{i}\le\Bigl(1+\frac1j\Bigr)^{j}.
\end{equation}

\begin{lemma}\label{lem:powermonotone}
For all integers $i$ and $j$ with $1\le i\le j$,
\[
\Bigl(1+\frac1i\Bigr)^{i-1}\le\Bigl(1+\frac1j\Bigr)^{j-1}.
\]
\end{lemma}

\begin{proof}
By induction on $j$ it suffices to prove the inequality for $j=i+1$. In this
case it follows from \eqref{eq:prelim-euler} on dividing the left-hand side
by $1+1/i$ and the right-hand side by the smaller number $1+1/(i+1)$.
\end{proof}

Together with Lemma~\ref{lem:powermonotone}, the identity
\eqref{eq:prelim-quotient} gives the
log-convexity of the Cayley coefficients on indices $j\ge1$, in the
following form; for $m=1$ and $j=i+1$ it is the usual inequality
$a_{i+1}^2\le a_i\,a_{i+2}$.

\begin{lemma}\label{lem:cayleyconvex}
For all integers $i$, $j$ and $m$ with $1\le i\le j$ and $m\ge0$, the Cayley
coefficients satisfy
\[
a_{i+m}\,a_j\le a_i\,a_{j+m}.
\]
\end{lemma}

\begin{proof}
We argue by induction on $m$, for fixed $i\le j$; the case $m=0$ is trivial.
Suppose that $a_{i+m}\,a_j\le a_i\,a_{j+m}$. Since $1\le i+m\le j+m$,
\eqref{eq:prelim-quotient} and Lemma~\ref{lem:powermonotone} give
\[
\frac{a_{i+m+1}}{a_{i+m}}=\Bigl(1+\frac1{i+m}\Bigr)^{i+m-1}
\le\Bigl(1+\frac1{j+m}\Bigr)^{j+m-1}=\frac{a_{j+m+1}}{a_{j+m}}.
\]
Multiplying this inequality by the induction hypothesis gives
$a_{i+m+1}\,a_j\le a_i\,a_{j+m+1}$.
\end{proof}

For the Cayley coefficients the inequality with a general shift came directly
from the explicit ratio. For other sequences, in Sections~\ref{sec:near}
and~\ref{sec:renewal}, only the inequality between three consecutive terms
will be known, and the following lemma passes from it to the inequality with
a general shift.

\begin{lemma}\label{lem:crossform}
Let $(v_n)_{n\ge0}$ be a sequence of real numbers such that $v_n>0$ and
$v_{n+1}^2\le v_n\,v_{n+2}$ for every integer $n\ge0$. Then for all integers
$i$, $j$ and $m$ with $0\le i\le j$ and $m\ge0$,
\[
v_{i+m}\,v_j\le v_i\,v_{j+m}.
\]
\end{lemma}

\begin{proof}
Since the numbers $v_n$ are positive, the hypothesis says that the ratio
$v_{n+1}/v_n$ does not decrease in $n$. Hence, for $i\le j$,
\[
\frac{v_{i+m}}{v_i}=\prod_{t=0}^{m-1}\frac{v_{i+t+1}}{v_{i+t}}
\le\prod_{t=0}^{m-1}\frac{v_{j+t+1}}{v_{j+t}}=\frac{v_{j+m}}{v_j},
\]
which is the claim.
\end{proof}

\subsection{Abel's identity and a recursion for the sequence $(c_n)$}
\label{subsec:prelim-abel}

The recursion for the numbers $c_n$ that we need is a case of a classical
identity, Abel's generalisation of the binomial theorem. For a
real number $x$ and an integer $k\ge0$ let
\[
w_0(x)=1,\qquad w_k(x)=x\,(x+k)^{k-1}\quad(k\ge1).
\]
In Abel's identity these polynomials in $x$ play the role of the powers
$x^k$ of the binomial theorem, and the powers $(y-k)^{n-k}$ that of
$y^{n-k}$.

\begin{lemma}[{Abel \cite[p.~159]{abel1826}}]\label{lem:abel}
For every integer $n\ge0$ and all real numbers $x$ and $y$,
\[
\sum_{k=0}^{n}\binom nk\,w_k(x)\,(y-k)^{n-k}=(x+y)^n.
\]
\end{lemma}

Abel states his identity for three arbitrary quantities $x$, $\alpha$ and
$\beta$ and a positive integer exponent, and his proof covers the exponent
$0$ as well; Lemma~\ref{lem:abel} is its case $\beta=-1$, with Abel's
$\alpha$ and $x$ replaced by our $x$ and $y$.

For $x=1$ the weights of Abel's identity are the numbers $c_k$ up to a
factorial: for every integer $k\ge0$,
\begin{equation}\label{eq:prelim-cabel}
c_k=\frac{w_k(1)}{k!}.
\end{equation}
Indeed, $w_k(1)=(k+1)^{k-1}$ for every $k\ge0$.
With \eqref{eq:prelim-cabel}, Abel's identity becomes a recursion for the
numbers $c_n$, with which Section~\ref{sec:counts} compares the last-car
recursion of Proposition~\ref{prop:lastcar}.

\begin{lemma}\label{lem:crecursion}
For every integer $n\ge0$, the numbers $c_n$ satisfy
\[
(n+1)\,c_{n+1}=\sum_{j=0}^{n}(j+1)\,c_j\,c_{n-j}.
\]
\end{lemma}

\begin{proof}
Since $(j+1)\,c_j=(j+1)^{j}/j!$ for every $j\ge0$, the term of index $j=n-k$
of the sum is, by \eqref{eq:prelim-cabel},
\[
(n-k+1)\,c_{n-k}\,c_k=\frac{(n+1-k)^{n-k}}{(n-k)!}\cdot\frac{w_k(1)}{k!}
=\frac1{n!}\binom nk\,w_k(1)\,(n+1-k)^{n-k}.
\]
By Lemma~\ref{lem:abel} with $x=1$ and $y=n+1$, the sum of these terms over
$0\le k\le n$ is $(n+2)^n/n!$, which is $(n+1)\,c_{n+1}$.
\end{proof}

Written for the Cayley coefficients themselves, the recursion of
Lemma~\ref{lem:crecursion} becomes the following convolution identity, which
is used in Section~\ref{sec:frozen}.

\begin{lemma}\label{lem:cayleyconvolution}
For every integer $t\ge0$,
\[
\sum_{i=0}^{t-1}i\,a_i\,a_{t-i}=(t-1)\,a_t.
\]
\end{lemma}

\begin{proof}
For $t\le1$ both sides are zero, since $a_0=0$. Let $t=n+2$ with $n\ge0$.
The term $i=0$ vanishes, and with $i=j+1$ the other terms are the numbers
$(j+1)\,c_j\,c_{n-j}$ with $0\le j\le n$, whose sum is
$(n+1)\,c_{n+1}=(t-1)\,a_t$ by Lemma~\ref{lem:crecursion}.
\end{proof}

\subsection{The Poisson mass and Stirling's formula}
\label{subsec:prelim-stirling}

For an integer $n\ge0$ let
\[
\nu_n=\frac{n^n}{n!},\qquad P(n)=e^{-n}\,\nu_n=e^{-n}\,\frac{n^n}{n!},
\]
so that $\nu_0=P(0)=1$. The number $P(n)$, the \emph{Poisson mass}, is the
probability that a Poisson variable with mean $n$ takes the value $n$. The
sequences of this section are estimated through it, and by Stirling's formula
$P(n)$ is close to $1/\sqrt{2\pi n}$. Note that $P(n)=(n/e)^n/n!$ for every
$n\ge0$, and that in a product of two Poisson masses whose arguments add up
to $m$ the exponential factors combine to $e^{-m}$, whatever the splitting:
\[
P(g)\,P(m-g)=e^{-m}\,\nu_g\,\nu_{m-g}\qquad(0\le g\le m).
\]

We need the approximation of $P(n)$ by $1/\sqrt{2\pi n}$ with
explicit bounds on both sides that hold for every $n\ge1$. Both bounds come
from Stirling's formula in the form that Robbins gave it
\cite[p.~26, (1) and~(2)]{robbins1955}: for every integer $n\ge1$,
\begin{equation}\label{eq:prelim-robbins}
\sqrt{2\pi n}\,\Bigl(\frac ne\Bigr)^{n}e^{1/(12n+1)}<n!
<\sqrt{2\pi n}\,\Bigl(\frac ne\Bigr)^{n}e^{1/(12n)}.
\end{equation}

Here and for the rest of the paper we write
\[
c_{\mathrm S}=\frac1{\sqrt{2\pi}}
\]
for the constant of Stirling's formula. Since $P(n)=(n/e)^n/n!$, the bounds
\eqref{eq:prelim-robbins} give, for every integer $n\ge1$,
\begin{equation}\label{eq:prelim-stirling}
c_{\mathrm S}\Bigl(1-\frac1{12n}\Bigr)\frac1{\sqrt n}\le P(n)
\le\frac{c_{\mathrm S}}{\sqrt n}.
\end{equation}
The upper bound follows from the left inequality of
\eqref{eq:prelim-robbins}, in which $e^{1/(12n+1)}\ge1$. For the lower
bound, the right inequality is used together with
$e^{-1/(12n)}\ge1-1/(12n)$, which holds because $e^{t}\ge1+t$ for every
real number $t$.

Where a numerical comparison needs the value of the constant, we use the
enclosure
\[
0.39894\le c_{\mathrm S}\le0.39894229.
\]

\section{The cumulative counts and the last-car recursion}\label{sec:counts}
The proof starts from the last-car recursion for the cumulative counts
$A_{n,\ell}$, which we recall from the literature
(Proposition~\ref{prop:lastcar}); it is through this recursion that
the parking process enters the later arguments. We then pass to the numbers
$A_{n,\ell}/n!$ and to the normalised counts and steps formed from them,
with which the rest of the paper works.

\subsection{The last-car recursion}\label{subsec:counts-lastcar}

Recall from Section~\ref{sec:intro} that the cumulative counts $A_{n,\ell}$
are defined for all integers $n\ge0$ and $\ell\ge0$, and that
$A_{0,\ell}=1$.
The last-car recursion expresses $A_{n+1,\ell}$ by the numbers $A_{j,\ell}$
with $j\le n$; it results from sorting the parking functions of length
$n+1$ by the spot in which the last car parks. It is stated by
Aguilar-Fraga et al.\ \cite[Theorem~3.7]{aguilarfraga2024}, with $n$ instead
of $n+1$, for the number of parking functions of length $n$ in which every
car parks at most $\ell$ spots after the spot that it prefers, which is
$A_{n,\ell}$.

\begin{proposition}[{\cite[Theorem~3.7]{aguilarfraga2024}}]\label{prop:lastcar}
For all integers $n\ge0$ and $\ell\ge0$, the cumulative counts satisfy
\[
A_{n+1,\ell}=\sum_{j=0}^{n}\min(j+1,\ell+1)\binom{n}{j}
A_{j,\ell}\,A_{n-j,\ell}.
\]
\end{proposition}

\subsection{The steps as second differences}\label{subsec:counts-steps}

The numbers $b_{n,\ell}$ enter the proof only through the cumulative counts.
For all integers $n\ge0$ and $\ell\ge0$,
\begin{equation}\label{eq:counts-differences}
b_{n,0}=A_{n,0},\qquad b_{n,\ell+1}=A_{n,\ell+1}-A_{n,\ell},
\end{equation}
and $A_{n,\ell}\le A_{n,\ell'}$ whenever $\ell\le\ell'$. These relations are
immediate from the definitions.

The theorems of this paper are statements on the signs of the steps
$d_{n,k}=b_{n,k}-b_{n,k-1}$, which Section~\ref{sec:intro} defines for
integers $n\ge1$ and $k\ge1$. The numbers $b_{0,\ell}$ are defined as well,
and we extend the definition of the step to length $n=0$ by the same
formula, so that $d_{n,k}$ is defined for all integers $n\ge0$ and
$k\ge1$. Since $b_{n,k}$ is a first difference of the cumulative counts
in the bound on the displacement, a step is a second difference: for all
integers $k\ge2$ and $n\ge0$, the relations
\eqref{eq:counts-differences}, applied to $b_{n,k}$ and, since $k-1\ge1$, to
$b_{n,k-1}$, give
\begin{equation}\label{eq:counts-second}
b_{n,k}-b_{n,k-1}=A_{n,k}-2A_{n,k-1}+A_{n,k-2}.
\end{equation}
It is from this second difference that the exact computations of
Section~\ref{sec:tables} read the sign of a step.

\subsection{The normalised counts and their recursion}
\label{subsec:counts-normalised}

From now on we work with the counts of length $n$ divided by $n!$. The
reason is that the last-car recursion contains a binomial coefficient, which
this division removes, since
\[
\frac{1}{n!}\binom{n}{j}=\frac{1}{j!\,(n-j)!}\qquad(0\le j\le n).
\]
The recursion that results can be written as a differential equation for a
generating series, and it can be compared with the recursion of
Lemma~\ref{lem:crecursion} for the sequence $(c_n)$. For integers $y\ge1$
and $n\ge0$ we put
\[
z^{(y)}_n=\frac{A_{n,y-1}}{n!},
\]
the \emph{normalised count of threshold $y$}: it is
the number of parking functions of length $n$ all of whose displacements are
smaller than $y$, divided by $n!$. The index is the threshold $y$ and not the
bound $y-1$ on the maximum displacement. With this choice the factor
$\min(j+1,\ell+1)$ of the last-car recursion is $\min(j+1,y)$, and
$z^{(y)}_n=c_n$ for the lengths $n\le y$ (Proposition~\ref{prop:freerange}). For
integers $k\ge0$ and $n\ge0$ we put
\[
\beta^{(k)}_n=\frac{b_{n,k}}{n!},
\]
the \emph{normalised count of column $k$}, and for integers $k\ge1$ and
$n\ge0$
\[
\delta^{(k)}_n=\beta^{(k)}_n-\beta^{(k-1)}_n=\frac{d_{n,k}}{n!},
\]
the \emph{normalised step} of column $k$ in row $n$. When column
$k$ is fixed, we write $\delta_n$ for $\delta^{(k)}_n$.

Clearly the numbers
$z^{(y)}_n$ and $\beta^{(k)}_n$ are nonnegative, $z^{(y)}_n\le z^{(y')}_n$
whenever $y\le y'$, and $z^{(y)}_0=1$. Dividing
\eqref{eq:counts-differences} by $n!$, we obtain
$\beta^{(0)}_n=z^{(1)}_n$ and
\begin{equation}\label{eq:counts-column}
\beta^{(k)}_n=z^{(k+1)}_n-z^{(k)}_n\qquad(k\ge1),
\end{equation}
and \eqref{eq:counts-second} becomes
\begin{equation}\label{eq:counts-step}
\delta^{(k)}_n=z^{(k+1)}_n-2\,z^{(k)}_n+z^{(k-1)}_n\qquad(k\ge2,\ n\ge0).
\end{equation}
The quantities of one column $k$ are therefore formed from the normalised
counts of the three thresholds $k+1$, $k$ and $k-1$. Finally, the real
number $\delta^{(k)}_n$ is positive, nonnegative, negative or nonpositive
exactly when the integer $d_{n,k}$ is. The analytic sections prove
inequalities for the real numbers $\delta^{(k)}_n$, whereas the theorems and
the exact computations speak of the integers $d_{n,k}$, and these equivalences
are the passage between the two.

The recursion that the numbers $z^{(y)}_n$ satisfy belongs to the following
family of recursions, which is used with a general weight in
Section~\ref{sec:series}.

\begin{definition}\label{def:weightedrecursion}
Let $\rho=(\rho(t))_{t\ge1}$ and $W=(W_n)_{n\ge0}$ be sequences of real
numbers. The sequence $W$ is a \emph{solution of the weighted last-car
recursion with weight $\rho$} if $W_0=1$ and
\[
(n+1)\,W_{n+1}=\sum_{j=0}^{n}\rho(j+1)\,W_j\,W_{n-j}
\qquad\text{for every integer $n\ge0$.}
\]
\end{definition}

The recursion gives $W_{n+1}$ in terms of $W_0,\dots,W_n$ and of
$\rho(1),\dots,\rho(n+1)$. The sequence $c=(c_n)_{n\ge0}$ is a solution
for the weight $\rho(t)=t$: we have $c_0=1$,
and the recursion is Lemma~\ref{lem:crecursion}.

\begin{corollary}\label{cor:normalised}
Let $y\ge1$ be an integer. Then the sequence $z^{(y)}=(z^{(y)}_n)_{n\ge0}$
of the normalised counts of threshold $y$ is a solution of the weighted
last-car recursion with the weight $\rho(t)=\min(t,y)$, in the sense of
Definition~\ref{def:weightedrecursion}.
\end{corollary}

\begin{proof}
We have $z^{(y)}_0=1$. For every integer $n\ge0$, dividing the recursion of
Proposition~\ref{prop:lastcar} for $\ell=y-1$ by $n!$ gives
\begin{equation}\label{eq:counts-normalised}
(n+1)\,z^{(y)}_{n+1}=\sum_{j=0}^{n}\min(j+1,y)\,z^{(y)}_j\,z^{(y)}_{n-j},
\end{equation}
which is the recursion of Definition~\ref{def:weightedrecursion} for the
weight $\rho(t)=\min(t,y)$.
\end{proof}

\subsection{The free range of a threshold}\label{subsec:counts-freerange}

For $t\le y$ the weight $\min(t,y)$ is the weight $t$ of the sequence $c$,
and so the two solutions agree up to index $y$. The lengths $n\le y$ form
the \emph{free range} of threshold $y$.

\begin{proposition}\label{prop:freerange}
Let $y\ge1$ and $n\ge0$ be integers with $n\le y$. Then $z^{(y)}_n=c_n$.
\end{proposition}

\begin{proof}
We prove the claim by induction on $n$. For $n=0$ we have
$z^{(y)}_0=1=c_0$. Let $n+1\le y$, and suppose that $z^{(y)}_j=c_j$ for
$0\le j\le n$. Since $\min(j+1,y)=j+1$ for these $j$,
Corollary~\ref{cor:normalised} and Lemma~\ref{lem:crecursion} give
\[
(n+1)\,z^{(y)}_{n+1}=\sum_{j=0}^{n}(j+1)\,c_j\,c_{n-j}=(n+1)\,c_{n+1},
\]
and therefore $z^{(y)}_{n+1}=c_{n+1}$.
\end{proof}

In terms of the counts,
Proposition~\ref{prop:freerange} says that
\[
A_{n,\ell}=(n+1)^{n-1}\qquad\text{whenever } n\le\ell+1,
\]
which is the classical count of parking functions of length $n$ given in
Section~\ref{sec:intro}.

The free range is the reason why a column $k\ge1$ begins in row $k+1$.
Indeed, a length $n\le k$ lies in the free range of both thresholds
$k$ and $k+1$, and \eqref{eq:counts-column} and
Proposition~\ref{prop:freerange} give
\[
\beta^{(k)}_n=c_n-c_n=0\qquad(k\ge1,\ n\le k).
\]
Consequently
\[
\delta^{(k)}_n=0\qquad(k\ge2,\ n<k).
\]
In row $n=k$ the step of a column $k\ge2$ is $\delta^{(k)}_k=-a_k/k$, by
Lemma~\ref{lem:initialstep} of Section~\ref{sec:series}, with the Cayley
coefficient $a_k$ of Section~\ref{sec:prelim}.

\subsection{The tilted recursion}\label{subsec:counts-tilts}

Several later arguments multiply the term of index $n$ of a sequence by the
$n$-th power of a real number. Such a number $x$ is called a \emph{tilt},
and to \emph{tilt} a sequence $(s_n)_{n\ge0}$ of real numbers by $x$ is to
pass to the sequence $(s_n\,x^n)_{n\ge0}$. The numbers $c_N\,x^N$ are called
the \emph{tilted Cayley masses}, and $z^{(y)}_N\,x^N$ is the \emph{tilted
normalised count} of threshold $y$. By Proposition~\ref{prop:freerange} the
tilted normalised count of threshold $y$ is the tilted Cayley mass
$c_N\,x^N$ for $N\le y$, and for a tilt $x\ge0$ the tilted normalised counts
are nonnegative and do not decrease in the threshold.

Multiplying the recursion \eqref{eq:counts-normalised}, taken for $n=N$, by
$x^{N+1}$ and writing $x^{N+1}=x\cdot x^{i}\cdot x^{N-i}$ in the term of
index $i$, we obtain the tilted recursion
\begin{equation}\label{eq:counts-tilted}
(N+1)\,z^{(y)}_{N+1}\,x^{N+1}
=x\sum_{i=0}^{N}\min(i+1,y)\,\bigl(z^{(y)}_i\,x^i\bigr)
\bigl(z^{(y)}_{N-i}\,x^{N-i}\bigr)
\end{equation}
for every integer $y\ge1$, every real number $x$ and every integer $N\ge0$.
The later sections apply it at the critical tilt of a column, which is
defined in Section~\ref{sec:propagation}.

\section{Generating series, the prime kernel and the column equation}\label{sec:series}
We now write the normalised last-car recursion of
Corollary~\ref{cor:normalised} as a differential equation for a formal power
series and derive from it, by the algebra of power series, the linear
relations on which the later sections rest: one for the solutions of a
weighted last-car recursion (Propositions~\ref{prop:chainforward}
and~\ref{prop:chainquotient}), and the \emph{column equation}, which is the
recursion that the normalised step satisfies along a column in its first
layer (Proposition~\ref{prop:columnequation}).

\subsection{The series of a threshold}\label{subsec:series-threshold}

All series in this section are formal power series in one variable $X$ with
real coefficients. For a series $\phi$
we write $[X^n]\phi$ for the coefficient of $X^n$ in $\phi$ and $\phi'$ for
the formal derivative of $\phi$, so that
$[X^n]\phi'=(n+1)\,[X^{n+1}]\phi$; the derivative obeys the product rule. For
two series $\phi$ and $\psi$ with the coefficients $\phi_j=[X^j]\phi$ and
$\psi_j=[X^j]\psi$ the coefficients of the product are
\begin{equation}\label{eq:series-product}
[X^n](\phi\psi)=\sum_{j=0}^{n}\phi_j\,\psi_{n-j}
=\phi_0\,\psi_n+\sum_{m=1}^{n}\phi_m\,\psi_{n-m}.
\end{equation}
A series $\phi$ with $\phi_0\ne0$ has an inverse
$\phi^{-1}$, which we also write $1/\phi$, and the constant coefficient of
$\phi^{-1}$ is $1/\phi_0$.

Let $y\ge1$ be an integer, and let $z^{(y)}_n=A_{n,y-1}/n!$ be the normalised
count of threshold $y$ (Section~\ref{sec:counts}). By
Corollary~\ref{cor:normalised} the sequence $z^{(y)}$ is the solution of the
weighted last-car recursion with the weight $\min(t,y)$: we have
$z^{(y)}_0=1$ and
\begin{equation}\label{eq:series-normalised}
(n+1)\,z^{(y)}_{n+1}=\sum_{j=0}^{n}\min(j+1,y)\,z^{(y)}_j\,z^{(y)}_{n-j}
\qquad(n\ge0).
\end{equation}
The \emph{series of the threshold} $y$ is the formal power series
\[
Z_y=\sum_{n\ge0}z^{(y)}_n\,X^n ,
\]
whose constant coefficient is $1$. Beside it we consider the series
\[
\Gamma_y=\sum_{t\ge1}\frac{\min(t,y)}{t}\,z^{(y)}_{t-1}\,X^t ,
\]
which is made so that the coefficients of its derivative are the first factors
on the right of \eqref{eq:series-normalised}:
\begin{equation}\label{eq:series-log-coefficients}
[X^j]\,\Gamma_y'=\min(j+1,y)\,z^{(y)}_j\qquad(j\ge0).
\end{equation}
The same two series can be formed for an arbitrary weight, and we need them in
this generality in Section~\ref{subsec:series-chain}. For real sequences
$\rho=(\rho(t))_{t\ge1}$ and $W=(W_n)_{n\ge0}$ we put
\[
\mathbf W=\sum_{n\ge0}W_n\,X^n,\qquad
\Gamma=\sum_{t\ge1}\frac{\rho(t)}{t}\,W_{t-1}\,X^t ,
\]
so that $[X^j]\,\Gamma'=\rho(j+1)\,W_j$ for $j\ge0$. The notation does not
show the two sequences from which $\mathbf W$ and $\Gamma$ are formed; they
will be named wherever these series are used. For the weight
$\rho(t)=\min(t,y)$ and the sequence $W=z^{(y)}$ the series $\mathbf W$ and
$\Gamma$ are $Z_y$ and $\Gamma_y$.

The next lemma is the last-car recursion in the language of series.

\begin{lemma}\label{lem:seriesode}
For every integer $y\ge1$, the series $Z_y$ of threshold $y$ satisfies
\[
Z_y'=\Gamma_y'\,Z_y .
\]
\end{lemma}

\begin{proof}
Let $\rho=(\rho(t))_{t\ge1}$ and $W=(W_n)_{n\ge0}$ be real sequences such
that $W$ is a solution of the weighted last-car recursion with weight $\rho$
in the sense of Definition~\ref{def:weightedrecursion}, so that
\[
(n+1)\,W_{n+1}=\sum_{j=0}^{n}\rho(j+1)\,W_j\,W_{n-j}
\]
for every $n\ge0$, and let $\mathbf W$ and $\Gamma$ be the series formed from
them. Since $[X^j]\,\Gamma'=\rho(j+1)\,W_j$, the right-hand side of the
recursion is the coefficient of $X^n$ in $\Gamma'\,\mathbf W$, by
\eqref{eq:series-product}, and the left-hand side is the coefficient of $X^n$
in $\mathbf W'$. Hence
\begin{equation}\label{eq:series-weighted-ode}
\mathbf W'=\Gamma'\,\mathbf W .
\end{equation}
By Corollary~\ref{cor:normalised} this applies to the sequence $W=z^{(y)}$ and
the weight $\rho(t)=\min(t,y)$, for which $\mathbf W=Z_y$ and
$\Gamma=\Gamma_y$.
\end{proof}

The proof gives more than the lemma states: the equation
\eqref{eq:series-weighted-ode} holds whenever $W$ is a solution of the
weighted last-car recursion with weight $\rho$ in the sense of
Definition~\ref{def:weightedrecursion}.

Since the constant coefficient of $Z_y$ is $1$, the series $Z_y$ has an
inverse. The \emph{prime kernel series} of threshold $y$ is
\[
K_y=1-Z_y^{-1},
\]
and the \emph{prime kernel} of threshold $y$ is the sequence of its
coefficients, $\kappa^{(y)}_m=[X^m]K_y$ for $m\ge0$. By definition
$Z_y^{-1}=1-K_y$, so that
\begin{equation}\label{eq:series-inverse}
Z_y\,(1-K_y)=1 .
\end{equation}
Since the
constant coefficient of $Z_y^{-1}$ is $1$, we have $\kappa^{(y)}_0=0$, and
by \eqref{eq:series-product}, for every series $\phi$ and every $n\ge0$,
\begin{equation}\label{eq:series-kernel-product}
[X^n]\bigl(\phi\,(1-K_y)\bigr)
=[X^n]\phi-\sum_{m=1}^{n}\kappa^{(y)}_m\,[X^{n-m}]\phi .
\end{equation}
For $\phi=Z_y$ and $n\ge1$ the left-hand side is zero by
\eqref{eq:series-inverse}, and we obtain
\begin{equation}\label{eq:series-renewal}
z^{(y)}_n=\sum_{m=1}^{n}\kappa^{(y)}_m\,z^{(y)}_{n-m}\qquad(n\ge1).
\end{equation}
An equation of this form, in which a sequence is expressed by its own earlier
terms through a convolution, is called a renewal equation, and the sequence
$\kappa^{(y)}$ is its kernel. As $z^{(y)}_0=1$, the equation
\eqref{eq:series-renewal} determines $\kappa^{(y)}_n$ from the normalised
counts up to index $n$.

The name prime kernel refers to the prime parking functions, a notion due
to Gessel, for which we refer to the survey of Yan \cite{yan2015}. The
properties of the kernel that we need follow from the differential equation
of Lemma~\ref{lem:seriesode}, which we first carry over to the series
$K_y$.

\begin{corollary}\label{cor:kernelode}
For every integer $y\ge1$, the prime kernel series $K_y$ of threshold $y$
satisfies
\[
K_y'=(1-K_y)\,\Gamma_y' .
\]
\end{corollary}

\begin{proof}
Differentiating \eqref{eq:series-inverse} by the product rule gives
$Z_y'\,(1-K_y)=Z_y\,K_y'$. By Lemma~\ref{lem:seriesode} and
\eqref{eq:series-inverse} the left-hand side is
$\Gamma_y'\,Z_y\,(1-K_y)=\Gamma_y'$. Hence $\Gamma_y'=Z_y\,K_y'$, and
multiplying by $1-K_y$ gives $(1-K_y)\,\Gamma_y'=K_y'$, again by
\eqref{eq:series-inverse}.
\end{proof}

\subsection{The series of all parking functions and the free prime kernel}
\label{subsec:series-free}

When the displacement is not bounded, the inverse of the series can be
computed. Recall from Section~\ref{sec:prelim} the Cayley coefficients $a_j$
and the numbers $c_n=a_{n+1}$; in particular $a_0=0$ and $c_0=1$. We put
\[
F=\sum_{n\ge0}c_n\,X^n=1+X+\tfrac32X^2+\tfrac83X^3+\cdots .
\]
We write $E=1/F$ for the inverse of $F$; thus $EF=1$, and the constant
coefficient of $E$ is $1$. The \emph{tree series} is $T=XF$. Since
$c_n=a_{n+1}$ and $a_0=0$, we have $[X^n]\,T=a_n$ for every $n\ge0$.

The sequence $c$ is the solution of the weighted last-car recursion with the
weight $\rho(t)=t$: we have $c_0=1$, and the recursion is
Lemma~\ref{lem:crecursion}. For this weight and the sequence $W=c$ the series
$\mathbf W$ is $F$, and the series $\Gamma$ is $T$, because its coefficient of
$X^t$ is $c_{t-1}=a_t$ for $t\ge1$. Thus \eqref{eq:series-weighted-ode} reads
\begin{equation}\label{eq:series-tree-ode}
F'=T'F .
\end{equation}
The inverse $E$ satisfies a similar equation, which is used in
Section~\ref{subsec:series-chain}.
Differentiating $EF=1$ gives $E'F=-EF'$, and $EF'=E\,T'F$ by
\eqref{eq:series-tree-ode}. Multiplying by $E$ and using $EF=1$, we obtain
\begin{equation}\label{eq:series-inverse-ode}
E'=-E\,T' .
\end{equation}

Since the constant terms of $F$ and $T$ are $1$ and $0$, the equation
\eqref{eq:series-tree-ode} says that $F=e^{T}$, and since $T=XF$ it follows
that $T=Xe^{T}$. The series $T$ is thus the Cayley tree function, the
exponential generating series of the rooted labelled trees
\cite[Section~II.5.1]{flajoletsedgewick2009}.

The \emph{free prime kernel} is the sequence $(\varepsilon_m)_{m\ge0}$ with
\[
\varepsilon_0=0,\qquad
\varepsilon_m=\frac{(m-1)^{m-1}}{m!}\quad(m\ge1),
\]
so that $\varepsilon_1=1$; the next values are
$\varepsilon_2=1/2$, $\varepsilon_3=2/3$ and $\varepsilon_4=9/8$. Clearly
$\varepsilon_m>0$ for every $m\ge1$, and
\[
m\,\varepsilon_m=\nu_{m-1}\qquad(m\ge1),
\]
where $\nu_n=n^n/n!$ (Section~\ref{sec:prelim}). The next lemma shows that
the coefficient of $X^m$ in $E$ is $-\varepsilon_m$ for every $m\ge1$, and
it expresses the coefficients of a quotient $\phi/F=E\phi$ through the
kernel $\varepsilon$.

\begin{lemma}\label{lem:inverseseries}
For every formal power series $\phi$ with real coefficients
$\phi_j=[X^j]\phi$ and every integer $n\ge0$, the inverse $E$ of the series
$F$ satisfies
\[
[X^n](E\phi)=\phi_n-\sum_{m=1}^{n}\varepsilon_m\,\phi_{n-m} .
\]
\end{lemma}

\begin{proof}
Since $F=e^{T}$, we have $E=e^{-T}$, and $T$ satisfies the equation
$T=Xe^{T}$. For this equation and a power series $H(u)$ the Lagrange
inversion formula \cite[Theorem~A.2]{flajoletsedgewick2009} reads
\[
[X^m]\,H(T)=\frac1m\,[u^{m-1}]\bigl(H'(u)\,e^{mu}\bigr)\qquad(m\ge1).
\]
We apply it with $H(u)=e^{-u}$. Then $H(T)=E$ and
$H'(u)\,e^{mu}=-e^{(m-1)u}$, and the coefficient of $u^{m-1}$ in
$e^{(m-1)u}$ is $(m-1)^{m-1}/(m-1)!$. Hence
\[
[X^m]E=-\frac1m\cdot\frac{(m-1)^{m-1}}{(m-1)!}=-\varepsilon_m
\]
for every $m\ge1$. Since $[X^0]E=1$, the assertion is the second form of
\eqref{eq:series-product} for the product $E\phi$.
\end{proof}

For $\phi=F$ the lemma gives a renewal equation for the sequence $(c_n)$,
with the kernel $\varepsilon$.

\begin{corollary}\label{cor:crenewal}
For every integer $n\ge1$,
\[
\sum_{m=1}^{n}\varepsilon_m\,c_{n-m}=c_n .
\]
\end{corollary}

\begin{proof}
By Lemma~\ref{lem:inverseseries} with $\phi=F$, the difference of the two
sides of the asserted equation is the coefficient of $X^n$ in $EF$. Since
$EF=1$, this coefficient is zero for $n\ge1$.
\end{proof}

We shall also need a renewal equation for the numbers $\nu_n$, with
the Cayley coefficients as the kernel; the term $j=0$ of its sum is zero,
because $a_0=0$.

\begin{lemma}\label{lem:nurenewal}
For every integer $n\ge1$,
\[
\sum_{j=0}^{n}a_j\,\nu_{n-j}=\nu_n .
\]
\end{lemma}

\begin{proof}
For the tree function $T$, the series $1/(1-T)$ is the exponential
generating series of the mappings of a finite set into itself, of which
there are $i^i$ on $i$ points
\cite[Section~II.5.2, (48) and~(49)]{flajoletsedgewick2009}. Thus
\[
\frac{1}{1-T}=\sum_{i\ge0}\nu_i\,X^i .
\]
The coefficients of $T$ are the numbers $a_j$. By \eqref{eq:series-product}
the difference of the two sides of the asserted equation is therefore the
coefficient of $X^n$ in the product of $1-T$ and $1/(1-T)$, which is $1$,
and this coefficient is zero because $n\ge1$.
\end{proof}

\subsection{The prime kernel in and beyond the free range}
\label{subsec:series-kernel}

We return to a threshold $y\ge1$. In its free range $n\le y$ the normalised
counts are those of all parking functions, $z^{(y)}_n=c_n$
(Proposition~\ref{prop:freerange}), so that the series $Z_y$ and $F$ agree up
to degree $y$. The same then holds for their inverses.

\begin{lemma}\label{lem:kernelfree}
Let $y\ge1$ and $m$ be integers with $1\le m\le y$. Then the prime kernel
of threshold $y$ satisfies
\[
\kappa^{(y)}_m=\varepsilon_m .
\]
\end{lemma}

\begin{proof}
We argue by induction on $m$ and suppose that $\kappa^{(y)}_i=\varepsilon_i$
for $1\le i<m$. In the renewal equation \eqref{eq:series-renewal} for $n=m$
all the indices $m$ and $m-i$ are at most $y$, so that by
Proposition~\ref{prop:freerange}
\[
c_m=\sum_{i=1}^{m}\kappa^{(y)}_i\,c_{m-i} .
\]
Subtracting the equation of Corollary~\ref{cor:crenewal}, we obtain
\[
\sum_{i=1}^{m}\bigl(\kappa^{(y)}_i-\varepsilon_i\bigr)\,c_{m-i}=0 .
\]
The terms with $i<m$ vanish by the induction hypothesis, and the term $i=m$ is
$(\kappa^{(y)}_m-\varepsilon_m)\,c_0=\kappa^{(y)}_m-\varepsilon_m$.
\end{proof}

Beyond the free range, for $y\ge2$, the kernel satisfies a linear recursion
with explicit coefficients. The
recursion comes from the fact that the weight $\min(j+1,y)$ of the last-car
recursion differs from the constant $y$ only for $j\le y-2$, where the
normalised counts are the free ones.

\begin{lemma}\label{lem:kernelrecursion}
Let $y\ge2$ and $N\ge y$ be integers. Then the prime kernel of threshold $y$
satisfies
\[
N\,\kappa^{(y)}_N=\sum_{\gamma=1}^{y-1}(y-\gamma)\,a_\gamma\,
\kappa^{(y)}_{N-\gamma} .
\]
\end{lemma}

\begin{proof}
Let $C_y$ be the polynomial
\[
C_y=\sum_{j=0}^{y-2}(y-1-j)\,c_j\,X^j .
\]
Then $\Gamma_y'=y\,Z_y-C_y$. Indeed, by \eqref{eq:series-log-coefficients}
the coefficient of $X^j$ in $\Gamma_y'$ is $y\,z^{(y)}_j$ for $j\ge y-1$,
and for $j\le y-2$ it is $(j+1)\,c_j=y\,c_j-(y-1-j)\,c_j$ by
Proposition~\ref{prop:freerange}. By Corollary~\ref{cor:kernelode} and
\eqref{eq:series-inverse} we therefore have
\[
K_y'=(1-K_y)\,(y\,Z_y-C_y)=y-(1-K_y)\,C_y .
\]
We compare the coefficients of $X^{N-1}$. On the left it is
$N\,\kappa^{(y)}_N$. On the right the constant $y$ contributes nothing,
because $N-1\ge1$, and the coefficient of $X^{N-1}$ in $C_y$ is zero,
because $N-1>y-2$. By
\eqref{eq:series-kernel-product} with $\phi=C_y$ we obtain
\[
N\,\kappa^{(y)}_N=\sum_{m=1}^{N-1}\kappa^{(y)}_m\,[X^{N-1-m}]\,C_y .
\]
We substitute $\gamma=N-m$. The coefficient of $X^{\gamma-1}$ in $C_y$ is
$(y-\gamma)\,c_{\gamma-1}=(y-\gamma)\,a_\gamma$ for $1\le\gamma\le y-1$ and
zero for $\gamma\ge y$, and since $N-1\ge y-1$ this is the assertion.
\end{proof}

For $y=2$ the recursion is $N\,\kappa^{(2)}_N=\kappa^{(2)}_{N-1}$ for
$N\ge2$.

Let $x$ be a tilt (Section~\ref{sec:counts}).
Multiplying the recursion by $x^N=x^\gamma x^{N-\gamma}$, we see that the
tilted prime kernel $\kappa^{(y)}_m\,x^m$ satisfies the same recursion with
the tilted weights: for $N\ge y\ge2$,
\begin{equation}\label{eq:series-tilted-recursion}
N\,\kappa^{(y)}_N\,x^N=\sum_{\gamma=1}^{y-1}(y-\gamma)\,a_\gamma x^\gamma
\cdot\kappa^{(y)}_{N-\gamma}\,x^{N-\gamma} .
\end{equation}
For $x\ge0$ the tilted weights $(y-\gamma)\,a_\gamma x^\gamma$ are
nonnegative. In the free range,
Lemma~\ref{lem:kernelfree} gives $\kappa^{(y)}_m\,x^m=\varepsilon_m\,x^m$ for
$1\le m\le y$, where $\varepsilon_m\,x^m$ is the tilted free prime kernel.

The coefficients of the inverse of a series have no evident sign. For the
prime kernel, tilted or not, the sign follows from the recursion.

\begin{lemma}\label{lem:kernelsign}
For every integer $y\ge2$, every real number $x\ge0$ and every integer
$m\ge0$, the prime kernel of threshold $y$ satisfies
\[
\kappa^{(y)}_m\,x^m\ge0 .
\]
\end{lemma}

\begin{proof}
We argue by induction on $m$. For $m=0$ we have $\kappa^{(y)}_0=0$. For
$1\le m\le y$ we have $\kappa^{(y)}_m\,x^m=\varepsilon_m\,x^m\ge0$ by
Lemma~\ref{lem:kernelfree}. For $m>y$ every term on the right-hand side of
\eqref{eq:series-tilted-recursion} with $N=m$ is nonnegative, by the
induction hypothesis and because the tilted weights are nonnegative. Hence
$m\,\kappa^{(y)}_m\,x^m\ge0$.
\end{proof}

For $x=1$ the lemma gives $\kappa^{(y)}_m\ge0$ for every $y\ge2$ and every
$m\ge0$.

\subsection{A linear relation for the weighted last-car recursion}
\label{subsec:series-chain}

The last-car recursion is quadratic. We now show that the solution of a
weighted last-car recursion satisfies a linear relation over a range of
indices, as soon as its weight agrees on an initial range with the weight
$\rho(t)=t$, whose solution is the sequence $c$. The column equations of
Section~\ref{subsec:series-column} are obtained from this relation.

Let $W$ be a solution of the weighted last-car recursion with weight $\rho$,
let $\mathbf W$ and $\Gamma$ be the series formed from $\rho$ and $W$, and
suppose that $\rho(t)=t$ for $1\le t\le k$. The proofs below compare the
quotient $\mathbf W/F$ with the series $T-\Gamma$, whose coefficients are
linear in the numbers $W_{t-1}$, and they determine the coefficients of the
quotient up to degree $2k+1$ from the differential equation of the quotient.

The resulting relation has two forms. The first is the relation multiplied
by $F$.

\begin{proposition}\label{prop:chainforward}
Let $\rho=(\rho(t))_{t\ge1}$ and $W=(W_n)_{n\ge0}$ be real sequences such
that $W$ is a solution of the weighted last-car recursion with weight $\rho$
in the sense of Definition~\ref{def:weightedrecursion}, and let $k\ge0$ be an
integer such that $\rho(t)=t$ for every $t$ with $1\le t\le k$. Then for
every integer $n$ with $0\le n\le2k+1$,
\[
c_n-W_n=\sum_{t=k+1}^{n}c_{n-t}\,
\Bigl(c_{t-1}-\frac{\rho(t)}{t}\,W_{t-1}\Bigr) .
\]
For $n\le k$ the sum is empty and the statement is $W_n=c_n$.
\end{proposition}

\begin{proof}
Let $\mathbf W$ and $\Gamma$ be the series formed from $\rho$ and $W$ as in
Section~\ref{subsec:series-threshold}, so that
$\mathbf W'=\Gamma'\,\mathbf W$ by \eqref{eq:series-weighted-ode}, and put
\[
R=T-\Gamma,\qquad Q=E\,\mathbf W .
\]
We write $R_i$ and $Q_i$ for the coefficients of $X^i$ in $R$ and in
$Q$. The coefficient of $X^t$ in $T$ is $a_t=c_{t-1}$ for $t\ge1$, and that in
$\Gamma$ is $(\rho(t)/t)\,W_{t-1}$. Hence $R_0=0$ and
\begin{equation}\label{eq:series-difference}
R_t=c_{t-1}-\frac{\rho(t)}{t}\,W_{t-1}\qquad(t\ge1).
\end{equation}

We first note that $W_n=c_n$ for $0\le n\le k$. This follows easily by
induction on $n$: we have $W_0=1=c_0$, and as long as $n+1\le k$ the
recursion that gives $W_{n+1}$ and the recursion that gives $c_{n+1}$
(Lemma~\ref{lem:crecursion}) have the same coefficients, because
$\rho(j+1)=j+1$ for $j+1\le n+1\le k$.
Consequently, for $1\le t\le k$ we have $\rho(t)/t=1$ and
$W_{t-1}=c_{t-1}$, and \eqref{eq:series-difference} shows that
\begin{equation}\label{eq:series-difference-zero}
R_t=0\qquad(0\le t\le k).
\end{equation}

Next, by the product rule, \eqref{eq:series-inverse-ode} and
\eqref{eq:series-weighted-ode},
\[
Q'=E'\,\mathbf W+E\,\mathbf W'=-E\,T'\,\mathbf W+E\,\Gamma'\,\mathbf W
=-R'\,Q .
\]
Clearly $Q_0=1$, and the coefficient of $X^{i-1}$ in the last equation is
\begin{equation}\label{eq:series-quotient-recursion}
i\,Q_i=-\sum_{j=0}^{i-1}(j+1)\,R_{j+1}\,Q_{i-1-j}\qquad(i\ge1).
\end{equation}
For $1\le i\le k$ every $R_{j+1}$ on the right has $j+1\le i\le k$ and
vanishes by \eqref{eq:series-difference-zero}; hence $Q_i=0$ for
$1\le i\le k$. Now let $1\le i\le2k+1$, and consider a term of
\eqref{eq:series-quotient-recursion} with $j<i-1$. If $j+1\le k$, then
$R_{j+1}=0$. If $j\ge k$, then $1\le i-1-j\le k$, and $Q_{i-1-j}=0$ by
what was just shown. So only the term $j=i-1$ remains, and it is
$i\,R_i\,Q_0=i\,R_i$. Altogether,
\begin{equation}\label{eq:series-quotient}
Q_0=1\qquad\text{and}\qquad Q_i=-R_i\quad(1\le i\le2k+1).
\end{equation}

Finally, $Q\,F=E\,F\,\mathbf W=\mathbf W$. Let $n\le2k+1$. By
\eqref{eq:series-product}, \eqref{eq:series-quotient} and
\eqref{eq:series-difference-zero}, the coefficient of $X^n$ in this identity
is
\[
W_n=Q_0\,c_n+\sum_{t=1}^{n}Q_t\,c_{n-t}
=c_n-\sum_{t=k+1}^{n}R_t\,c_{n-t},
\]
and with \eqref{eq:series-difference} this is the assertion.
\end{proof}

Proposition~\ref{prop:chainforward} gives $W_n$ through one earlier value at
a time, whereas the recursion of $W$ gives a value through products of two
earlier values. In Section~\ref{sec:frozen} it is applied repeatedly to the
normalised counts of the thresholds, which leads to the chain sums of
Section~\ref{sec:compositions}.

For a threshold $y\ge1$ the proposition applies to the sequence $W=z^{(y)}$
and the weight $\rho(t)=\min(t,y)$, by Corollary~\ref{cor:normalised}, with
$k=y$. Multiplied by $n!$, it then reads, for $n\le2y+1$,
\begin{multline*}
(n+1)^{n-1}-A_{n,y-1}\\
=\sum_{t=y+1}^{n}\binom nt\,(n-t+1)^{n-t-1}\,
\bigl(t^{t-1}-y\,A_{t-1,y-1}\bigr) .
\end{multline*}
The left-hand side is the number of parking functions of length $n$ in which
some car has a displacement of at least $y$, and on the right the cumulative
counts occur for smaller lengths only.

The second form is the same relation before the multiplication by $F$. Here
the kernel $\varepsilon$ appears, through Lemma~\ref{lem:inverseseries}.

\begin{proposition}\label{prop:chainquotient}
Let $\rho=(\rho(t))_{t\ge1}$ and $W=(W_n)_{n\ge0}$ be real sequences such
that $W$ is a solution of the weighted last-car recursion with weight $\rho$
in the sense of Definition~\ref{def:weightedrecursion}, and let $k\ge0$ be an
integer such that $\rho(t)=t$ for every $t$ with $1\le t\le k$. Then for
every integer $n$ with $1\le n\le2k+1$,
\[
W_n-\sum_{m=1}^{n}\varepsilon_m\,W_{n-m}=
\begin{cases}
0, & n\le k,\\[2pt]
\dfrac{\rho(n)}{n}\,W_{n-1}-c_{n-1}, & n>k.
\end{cases}
\]
\end{proposition}

\begin{proof}
The hypotheses are those of Proposition~\ref{prop:chainforward}, and we use
the notation of its proof. By Lemma~\ref{lem:inverseseries}, applied to the
series $\mathbf W$, the left-hand side is the coefficient $Q_n$ of $X^n$ in
$Q=E\,\mathbf W$, and $Q_n=-R_n$ by \eqref{eq:series-quotient}. For
$n\le k$ this is zero by \eqref{eq:series-difference-zero}, and for $n>k$ it
is $(\rho(n)/n)\,W_{n-1}-c_{n-1}$ by \eqref{eq:series-difference}.
\end{proof}

For the thresholds Proposition~\ref{prop:chainquotient} reads as follows. Let
$y\ge1$ and $n$ be integers with $y<n\le2y+1$. Then
\begin{equation}\label{eq:series-chain-threshold}
z^{(y)}_n-\sum_{m=1}^{n}\varepsilon_m\,z^{(y)}_{n-m}
=\frac yn\,z^{(y)}_{n-1}-c_{n-1} .
\end{equation}
Indeed, by Corollary~\ref{cor:normalised} the sequence $z^{(y)}$ is the
solution of the weighted last-car recursion with the weight $\min(t,y)$, which
is equal to $t$ for $t\le y$, so that Proposition~\ref{prop:chainquotient}
applies with $k=y$. It gives \eqref{eq:series-chain-threshold}, because
$\min(n,y)=y$.

\subsection{The column equations}\label{subsec:series-column}

We now fix a column and combine the relations
\eqref{eq:series-chain-threshold} of neighbouring thresholds. Recall from
Section~\ref{sec:counts} the normalised count $\beta^{(k)}_n=b_{n,k}/n!$ of
column $k$ and the normalised step
$\delta^{(k)}_n=\beta^{(k)}_n-\beta^{(k-1)}_n=d_{n,k}/n!$. They are a first
and a second difference of the normalised counts in the threshold, the
latter being the identity \eqref{eq:counts-step}, which holds for $k\ge2$,
and they vanish in the free
range of the thresholds involved: for all $n\ge0$,
\begin{align}
\beta^{(k)}_n&=z^{(k+1)}_n-z^{(k)}_n&&(k\ge1),
\label{eq:series-level-difference}\\
\beta^{(k)}_n&=0&&(k\ge1,\ n\le k),
\label{eq:series-level-zero}\\
\delta^{(k)}_n&=0&&(k\ge2,\ n<k).
\label{eq:series-step-zero}
\end{align}
The step of a column $k\ge2$ in row $n=k$ is given by the next lemma.

\begin{lemma}\label{lem:initialstep}
For every integer $k\ge2$, the normalised step of column $k$ in row $k$ is
\[
\delta^{(k)}_k=-\frac{a_k}{k} .
\]
\end{lemma}

\begin{proof}
By Corollary~\ref{cor:normalised} the sequence $z^{(k-1)}$ is the solution of
the weighted last-car recursion with the weight $\min(t,k-1)$, which is equal
to $t$ for $t\le k-1$. We apply Proposition~\ref{prop:chainforward} to it,
with $k$ replaced by $k-1$ and with $n=k$, which is allowed because
$k\le2(k-1)+1$. The sum has the single term $t=k$, and since
$z^{(k-1)}_{k-1}=c_{k-1}$ by Proposition~\ref{prop:freerange}, we obtain
\[
c_k-z^{(k-1)}_k=c_{k-1}-\frac{k-1}{k}\,c_{k-1}=\frac{c_{k-1}}{k}
=\frac{a_k}{k} .
\]
On the other hand, $\beta^{(k)}_k=0$ by
\eqref{eq:series-level-zero}, and
$\beta^{(k-1)}_k=z^{(k)}_k-z^{(k-1)}_k=c_k-z^{(k-1)}_k$ by
\eqref{eq:series-level-difference} and Proposition~\ref{prop:freerange}.
Therefore
\[
\delta^{(k)}_k=\beta^{(k)}_k-\beta^{(k-1)}_k
=-\bigl(c_k-z^{(k-1)}_k\bigr)=-\frac{a_k}{k}. \qedhere
\]
\end{proof}

Multiplied by $k!$, the lemma says that $d_{k,k}=-k^{k-2}$.

In the first layer of the column, from row $k+1$ to row $2k-1$, the
normalised step satisfies a linear recursion along the column, the column
equation, which starts from the value of Lemma~\ref{lem:initialstep}. Its
proof takes the second difference in the threshold of the relations
\eqref{eq:series-chain-threshold}, whose left-hand sides are linear in the
normalised counts with coefficients that do not depend on the threshold.

\begin{proposition}\label{prop:columnequation}
Let $k\ge2$ and $n$ be integers with $k+1\le n\le2k-1$, and write
$\delta_i=\delta^{(k)}_i$ for the normalised step of column $k$ in row $i$.
Then
\[
\delta_n=\frac{k+1}{n}\,\delta_{n-1}+\frac{2}{n}\,\beta^{(k-1)}_{n-1}
+\sum_{m=1}^{n-k}\varepsilon_m\,\delta_{n-m} .
\]
\end{proposition}

\begin{proof}
The hypothesis $k+1\le n\le2k-1$ gives $y<n\le2y+1$ for the two thresholds
$y=k$ and $y=k-1$, and $k-1\ge1$; the conditions that bind are $k<n$ for
$y=k$ and $n\le2(k-1)+1$ for $y=k-1$. For $y=k+1$ it gives $y<n\le2y+1$
unless $n=k+1$, and for $n=k+1$ the relation
\eqref{eq:series-chain-threshold} holds with $y=k+1$ as well: by
Corollary~\ref{cor:normalised} and Proposition~\ref{prop:chainquotient},
applied to the sequence $z^{(k+1)}$ and the weight $\min(t,k+1)$ with $k$
replaced by $k+1$, its left-hand side is zero, and its right-hand side is
$z^{(k+1)}_k-c_k$, which is zero by Proposition~\ref{prop:freerange}. We take
\eqref{eq:series-chain-threshold} for $y=k+1$, subtract twice
\eqref{eq:series-chain-threshold} for $y=k$, and add
\eqref{eq:series-chain-threshold} for $y=k-1$. The result is
\begin{equation}\label{eq:series-column-full}
\delta_n-\sum_{m=1}^{n}\varepsilon_m\,\delta_{n-m}
=\frac{k+1}{n}\,\delta_{n-1}+\frac{2}{n}\,\beta^{(k-1)}_{n-1} .
\end{equation}
Indeed, the left-hand sides combine to the left-hand side of
\eqref{eq:series-column-full} by \eqref{eq:counts-step}. On the
right the terms $c_{n-1}$ cancel, and what remains is $1/n$ times
\[
(k+1)\,z^{(k+1)}_{n-1}-2k\,z^{(k)}_{n-1}+(k-1)\,z^{(k-1)}_{n-1}
=(k+1)\,\delta_{n-1}+2\,\beta^{(k-1)}_{n-1},
\]
where the equality holds by \eqref{eq:counts-step} and by
\eqref{eq:series-level-difference} for column $k-1$. Finally, for
$m>n-k$ we have $n-m<k$ and $\delta_{n-m}=0$ by \eqref{eq:series-step-zero},
so that the sum in \eqref{eq:series-column-full} is the sum over
$1\le m\le n-k$. Moving it to the right gives the assertion.
\end{proof}

Multiplied by $n!$, the column equation is an identity between integers,
\begin{multline*}
d_{n,k}=(k+1)\,d_{n-1,k}+2\,b_{n-1,k-1}\\
+\sum_{m=1}^{n-k}\binom nm\,(m-1)^{m-1}\,d_{n-m,k} ,
\end{multline*}
which can be checked against the rows displayed in Section~\ref{sec:intro}:
for $k=3$ and $n=5$ it reads
$-124=4\cdot(-18)+2\cdot34+5\cdot(-18)+10\cdot(-3)$. Beyond the first
layer the equation is replaced by the inequality of
Proposition~\ref{prop:columninequality}, which has the prime kernel
$\kappa^{(k)}$ of threshold $k$ instead of $\varepsilon$. In the
first layer the two kernels are evaluated only at indices
$m\le n-k\le k-1$, where they agree by Lemma~\ref{lem:kernelfree}.

The normalised counts of a column satisfy a column equation of their own, of
the same kind as that of the normalised step and on a slightly longer range
of rows. It is used in Section~\ref{sec:far}, where it is applied to
column $k-1$, whose normalised counts
$\beta^{(k-1)}_{n-1}$ occur in Proposition~\ref{prop:columnequation}.

\begin{lemma}\label{lem:levelequation}
Let $k\ge1$ and $n$ be integers with $k+1\le n\le2k+1$. Then
\[
\beta^{(k)}_n=\frac{z^{(k+1)}_{n-1}}{n}
+\frac{k}{n}\,\beta^{(k)}_{n-1}
+\sum_{m=1}^{n-k-1}\varepsilon_m\,\beta^{(k)}_{n-m} .
\]
\end{lemma}

\begin{proof}
We apply \eqref{eq:series-chain-threshold} with $y=k+1$, which is allowed
because $k+1<n\le2k+3$ when $n>k+1$, and with $y=k$, which is allowed because
$k<n\le2k+1$. For $n=k+1$ the relation holds with $y=k+1$ as well: by
Corollary~\ref{cor:normalised} and Proposition~\ref{prop:chainquotient},
applied to the sequence $z^{(k+1)}$ and the weight $\min(t,k+1)$ with $k$
replaced by $k+1$, its left-hand side is zero, and its right-hand side is
$z^{(k+1)}_k-c_k$, which is zero by Proposition~\ref{prop:freerange}. We
subtract the second relation from the first. The
terms $c_{n-1}$ cancel, and by \eqref{eq:series-level-difference}, used on
the left for every index and on the right for index $n-1$, we obtain
\[
\beta^{(k)}_n-\sum_{m=1}^{n}\varepsilon_m\,\beta^{(k)}_{n-m}
=\frac{k+1}{n}\,z^{(k+1)}_{n-1}-\frac{k}{n}\,z^{(k)}_{n-1}
=\frac{z^{(k+1)}_{n-1}}{n}+\frac{k}{n}\,\beta^{(k)}_{n-1} .
\]
For $m>n-k-1$ we have $n-m\le k$ and $\beta^{(k)}_{n-m}=0$ by
\eqref{eq:series-level-zero}, so that the sum on the left is the sum over
$1\le m\le n-k-1$. Moving it to the right gives the assertion.
\end{proof}

The equation gives the normalised counts of column $k$ row by row. It starts
in row $n=k+1$, where the sum is empty.

\section{Two inequalities for a column and the propagation of a positive step}\label{sec:propagation}
In the first layer of a column the normalised step satisfies the column
equation of Proposition~\ref{prop:columnequation}, which is an identity.
Here we prove two inequalities for a column $k\ge2$ that are not restricted
to the first layer and that replace this identity beyond it. The first,
Proposition~\ref{prop:onetree}, holds in every row. The second,
Proposition~\ref{prop:columninequality}, holds in the rows of the first
layer and in all rows beyond it. From
Proposition~\ref{prop:onetree} we derive the propagation of a positive step
down a column (Proposition~\ref{prop:propagation}), which gives the positive
signs beyond the first layer for columns $2\le k\le81$ in
Section~\ref{sec:beyond}. The propagation requires a tilt $x$ with $0<x<1$
at which the tilted normalised counts of threshold $k$ do not increase. In
the last subsection we define the critical tilt of a column and show that
at the critical tilt of a column $k\ge2$ the tilted normalised counts of
threshold $k$ do not increase (Proposition~\ref{prop:tilteddecrease}).

\subsection{The tree curvature of a column}\label{subsec:propagation-curvature}

We work with the normalised quantities of Section~\ref{sec:counts}: the
normalised counts $z^{(y)}_n$ of a threshold $y\ge1$, the normalised counts
$\beta^{(\ell)}_n$ of a column $\ell\ge0$, and the normalised step
$\delta^{(k)}_n$ of a column $k\ge1$, which is positive, nonnegative,
negative or nonpositive exactly when the integer $d_{n,k}$ is.
Recall that the numbers $z^{(y)}_n$ are nonnegative and do not decrease in
the threshold $y$, that $z^{(y)}_0=1$, and that $z^{(y)}_n=c_n$ in the free
range $n\le y$ of the threshold (Proposition~\ref{prop:freerange}). Recall
also that $\beta^{(\ell)}_n=z^{(\ell+1)}_n-z^{(\ell)}_n$ for a column
$\ell\ge1$ and that, by \eqref{eq:counts-step}, the normalised
step $\delta^{(k)}_n$ of a column $k\ge2$ is the second difference, in the
threshold, of the normalised counts of the three thresholds $k+1$, $k$ and
$k-1$, and that this step vanishes for $n<k$. The following sequence is
formed from the same three thresholds.

\begin{definition}\label{def:treecurvature}
Let $k\ge2$ be an integer. The \emph{tree curvature} of column $k$ is the
sequence of real numbers $\xi^{(k)}_t$, $t\ge0$, defined by $\xi^{(k)}_0=0$
and, for $t\ge1$,
\[
\xi^{(k)}_t=\frac1t\Bigl(\min(t,k+1)\,z^{(k+1)}_{t-1}
-2\,\min(t,k)\,z^{(k)}_{t-1}+\min(t,k-1)\,z^{(k-1)}_{t-1}\Bigr).
\]
When column $k$ is fixed, we write $\xi_t$ for $\xi^{(k)}_t$.
\end{definition}

In terms of the series of Section~\ref{sec:series} the definition says that
$\xi^{(k)}_t$ is the coefficient of $X^t$ in
$\Gamma_{k+1}-2\,\Gamma_k+\Gamma_{k-1}$, where
\[
\Gamma_y=\sum_{t\ge1}\frac{\min(t,y)}{t}\,z^{(y)}_{t-1}\,X^t .
\]
The series $Z_y$ of the threshold, whose coefficients are the numbers
$z^{(y)}_n$, satisfies $Z_y'=\Gamma_y'Z_y$ (Lemma~\ref{lem:seriesode}). By
\eqref{eq:counts-step} the step $\delta^{(k)}_n$ is the
coefficient of $X^n$ in $Z_{k+1}-2Z_k+Z_{k-1}$. Thus the step is the second
difference, in the threshold, of the series $Z_y$, and the tree curvature is
the same second difference of the series $\Gamma_y$.

Let $k\ge2$. In the rest of this section we also write $\delta_n$ for
$\delta^{(k)}_n$ when column $k$ is fixed. Below index $k$
the tree curvature vanishes:
\begin{equation}\label{eq:propagation-vanishing}
\xi_t=0\qquad(0\le t<k).
\end{equation}
Indeed, for $1\le t\le k-1$ the three terms of the definition cancel: the
three minima are equal to $t$, and the three normalised counts are equal to
$c_{t-1}$, because index $t-1$ lies in the free range of each of the
thresholds $k+1$, $k$ and $k-1$.
Beyond index $k$ the minima are the thresholds themselves: for every
$n\ge k$,
\begin{equation}\label{eq:propagation-curvature-step}
\begin{split}
(n+1)\,\xi_{n+1}&=(k+1)\,z^{(k+1)}_n-2k\,z^{(k)}_n+(k-1)\,z^{(k-1)}_n\\
&=(k+1)\,\delta_n+2\,\beta^{(k-1)}_n ,
\end{split}
\end{equation}
where the second equality follows from \eqref{eq:counts-step}
and from $\beta^{(k-1)}_n=z^{(k)}_n-z^{(k-1)}_n$.

The right-hand side of \eqref{eq:propagation-curvature-step}, with $n$
replaced by $n-1$, is $n$ times the sum of the first two terms on the right
of the column equation. It is also the left-hand side of the inequality of
Proposition~\ref{prop:columninequality} below, which holds in every row
$n\ge k+1$ and in which the prime kernel of threshold $k$ stands where the
column equation has the free prime kernel.

\subsection{Two inequalities between the step and the tree
curvature}\label{subsec:propagation-inequalities}

The first inequality bounds the step of column $k$ from below by a sum over
the tree curvature of the column, in which the weights are normalised counts
of threshold $k$.

\begin{proposition}\label{prop:onetree}
Let $k\ge2$ and $n\ge0$ be integers. Then the tree curvature of column $k$
satisfies
\[
\sum_{t=0}^{n}\xi^{(k)}_t\,z^{(k)}_{n-t}\le\delta^{(k)}_n .
\]
\end{proposition}

The second inequality is the one announced at the end of
Section~\ref{subsec:propagation-curvature}. From it the
positive sign of columns $k\ge82$ beyond the first layer is derived in
Section~\ref{sec:beyond}.

\begin{proposition}\label{prop:columninequality}
Let $k\ge2$ and $n\ge k+1$ be integers. Then
\[
(k+1)\,\delta^{(k)}_{n-1}+2\,\beta^{(k-1)}_{n-1}
\le n\,\Bigl(\delta^{(k)}_n
-\sum_{m=1}^{n-k}\kappa^{(k)}_m\,\delta^{(k)}_{n-m}\Bigr),
\]
where $\kappa^{(k)}$ is the prime kernel of threshold $k$.
\end{proposition}

The two propositions are proved together, since both are obtained from one
relation between power series.
Let $k\ge2$. For a threshold $y\ge1$ we use the series $Z_y$, $K_y$ and
$\Gamma_y$ of Section~\ref{sec:series}. The prime kernel series
$K_y=1-Z_y^{-1}$ satisfies the relation $Z_y\,(1-K_y)=1$ of
\eqref{eq:series-inverse};
the coefficient of $X^0$ in $1-K_y$ is $1$, and the coefficient of $X^m$ is
$-\kappa^{(y)}_m$ for $m\ge1$. The derivative of $\Gamma_y$ has
coefficients $[X^j]\,\Gamma_y'=\min(j+1,y)\,z^{(y)}_j$, and
\begin{equation}\label{eq:propagation-ode}
Z_y'=\Gamma_y'\,Z_y,\qquad K_y'=(1-K_y)\,\Gamma_y'
\end{equation}
by Lemma~\ref{lem:seriesode} and Corollary~\ref{cor:kernelode}. We put
\[
\Delta=Z_{k+1}-2\,Z_k+Z_{k-1},\qquad
\Xi=\Gamma_{k+1}-2\,\Gamma_k+\Gamma_{k-1},
\]
so that the coefficient of $X^n$ in $\Delta$ is the step $\delta_n$ and the
coefficient of $X^t$ in $\Xi$ is the tree curvature $\xi_t$, as was said in
Section~\ref{subsec:propagation-curvature}. For formal power series $S$ and
$S^*$ with real coefficients we write $S\succeq0$ if every coefficient of $S$
is nonnegative, and $S\succeq S^*$ if $S-S^*\succeq0$. Sums and products of
series that are $\succeq0$ are $\succeq0$. With this notation the relation
is $\Delta\,(1-K_k)\succeq\Xi$. The proof establishes it and then reads it
at the coefficient of $X^n$, for Proposition~\ref{prop:onetree} after
multiplication by $Z_k$.

To see where this relation comes from, we compare thresholds
$k+1$ and $k-1$ with threshold $k$ twice, through the series $Z_y$ and
through the series $\Gamma_y$. For $y\ge2$ let $L_y=\Gamma_y-\Gamma_{y-1}$, and put
\begin{gather*}
L_+=L_{k+1},\qquad L_-=L_k,\\
Y_+=(Z_{k+1}-Z_k)\,(1-K_k),\qquad Y_-=(Z_k-Z_{k-1})\,(1-K_k).
\end{gather*}
Then $\Xi=L_+-L_-$ and $\Delta\,(1-K_k)=Y_+-Y_-$, so that
the relation follows from the two comparisons $Y_+\succeq L_+$ and
$L_-\succeq Y_-$.
By \eqref{eq:series-inverse}, $Y_+$ and $Y_-$ are the quotients
$Z_{k+1}/Z_k-1$ and $1-Z_{k-1}/Z_k$. The proof obtains the two comparisons
from the differential equations \eqref{eq:propagation-ode} and from a
principle on series with nonnegative coefficients, with which it begins.

\begin{proof}[Proof of Propositions~\ref{prop:onetree} and~\ref{prop:columninequality}]
Let $S$, $S_1$ and $S_2$ be formal power series with real coefficients such
that
\begin{equation}\label{eq:propagation-principle}
S'=S_1\,S+S_2,\qquad S_1\succeq0,\qquad S_2\succeq0,\qquad S(0)\ge0,
\end{equation}
where $S(0)$ denotes the constant term of $S$. Then $S\succeq0$. Indeed, the
coefficient of $X^n$ in the equation is
\[
(n+1)\,[X^{n+1}]\,S=\sum_{j=0}^{n}[X^j]\,S_1\cdot[X^{n-j}]\,S+[X^n]\,S_2 ,
\]
and the claim follows by induction on the index of the coefficient. For
$S_1=0$ this principle says that a series with $S'\succeq0$ and $S(0)\ge0$ is
$\succeq0$.

Let $y\ge2$, and let
\[
J_y=\sum_{j\ge y-1}z^{(y-1)}_j\,X^j .
\]
We show that
\begin{equation}\label{eq:propagation-adjacent}
L_y'=y\,(Z_y-Z_{y-1})+J_y,\qquad L_y'\succeq0,\qquad L_y\succeq0 .
\end{equation}
The coefficient of $X^j$ in $L_y'=\Gamma_y'-\Gamma_{y-1}'$ is
\[
\min(j+1,y)\,z^{(y)}_j-\min(j+1,y-1)\,z^{(y-1)}_j .
\]
For $j\ge y-1$ this is the sum of $y\,(z^{(y)}_j-z^{(y-1)}_j)$ and
$z^{(y-1)}_j$. For $j\le y-2$ it is $0$, as is the coefficient of $X^j$ on
the right-hand side, because both minima are $j+1$ and
$z^{(y)}_j=z^{(y-1)}_j=c_j$ in the free range of both thresholds. This
proves the identity. The series $Z_y-Z_{y-1}$ and $J_y$ are $\succeq0$,
because the normalised counts are nonnegative and do not decrease in the
threshold. Hence $L_y'\succeq0$, and since $L_y$ has constant term $0$,
the principle \eqref{eq:propagation-principle} with $S_1=0$ gives
$L_y\succeq0$. We use \eqref{eq:propagation-adjacent} for $L_+=L_{k+1}$ and
for $L_-=L_k$.

Next we compute the derivatives of $Y_+$ and $Y_-$. By
\eqref{eq:series-inverse}, $Y_+=Z_{k+1}\,(1-K_k)-1$ and
$Y_-=1-Z_{k-1}\,(1-K_k)$, and both series have constant term $0$. By
\eqref{eq:propagation-ode},
\begin{align*}
Y_+'&=Z_{k+1}'\,(1-K_k)-Z_{k+1}\,K_k'\\
&=Z_{k+1}\,(1-K_k)\,(\Gamma_{k+1}'-\Gamma_k')=(1+Y_+)\,L_+',\\
Y_-'&=-Z_{k-1}'\,(1-K_k)+Z_{k-1}\,K_k'\\
&=Z_{k-1}\,(1-K_k)\,(\Gamma_k'-\Gamma_{k-1}')=(1-Y_-)\,L_-'.
\end{align*}

The first equation gives
\[
(Y_+-L_+)'=Y_+\,L_+'=L_+'\,(Y_+-L_+)+L_+\,L_+' .
\]
Here $L_+'\succeq0$ and $L_+\,L_+'\succeq0$ by
\eqref{eq:propagation-adjacent}, and $Y_+-L_+$ has constant term $0$. By
\eqref{eq:propagation-principle}, $Y_+\succeq L_+$.

The second equation gives $(L_--Y_-)'=Y_-\,L_-'$, so that the comparison
$L_-\succeq Y_-$ will follow once we know that $Y_-\succeq0$. The equation
$Y_-'=(1-Y_-)\,L_-'$ does not show this as it stands: written as
$Y_-'=-L_-'\,Y_-+L_-'$, it multiplies $Y_-$ by the series $-L_-'$, whose
coefficients are not positive, and \eqref{eq:propagation-principle} does not
apply. We rewrite it. Multiplying $Y_-=(Z_k-Z_{k-1})\,(1-K_k)$ by $Z_k$ we
get $Z_k-Z_{k-1}=Y_-\,Z_k$ by \eqref{eq:series-inverse}, and therefore
\[
(1-Y_-)\,(Z_k-Z_{k-1})=Z_{k-1}\,(1-K_k)\cdot Y_-\,Z_k=Z_{k-1}\,Y_- .
\]
Inserting the identity of \eqref{eq:propagation-adjacent} for $y=k$ into
$Y_-'=(1-Y_-)\,L_-'$, we obtain
\[
Y_-'=k\,(1-Y_-)\,(Z_k-Z_{k-1})+(1-Y_-)\,J_k=(k\,Z_{k-1}-J_k)\,Y_-+J_k .
\]
In this equation the series that multiplies $Y_-$ is $\succeq0$: its
coefficient of $X^j$ is $k\,z^{(k-1)}_j$ for $j<k-1$ and
$(k-1)\,z^{(k-1)}_j$ for $j\ge k-1$. Since $J_k\succeq0$ and $Y_-$ has
constant term $0$, the principle \eqref{eq:propagation-principle} gives
$Y_-\succeq0$. Hence $(L_--Y_-)'=Y_-\,L_-'\succeq0$, and as $L_--Y_-$ has
constant term $0$, we conclude that $L_-\succeq Y_-$.

Since $\Xi=L_+-L_-$ and $\Delta\,(1-K_k)=Y_+-Y_-$, the two comparisons
$Y_+\succeq L_+$ and $L_-\succeq Y_-$ give
\begin{equation}\label{eq:propagation-key}
\Delta\,(1-K_k)\succeq\Xi .
\end{equation}

For Proposition~\ref{prop:onetree}, let $n\ge0$. Multiplying the
series $\Delta\,(1-K_k)-\Xi\succeq0$ by $Z_k\succeq0$ and using
\eqref{eq:series-inverse}, we find that $\Delta\succeq\Xi\,Z_k$. The
coefficient of $X^n$ in $\Delta$ is $\delta_n$, and the coefficient of $X^n$
in $\Xi\,Z_k$ is the left-hand side of the inequality of
Proposition~\ref{prop:onetree}.

For Proposition~\ref{prop:columninequality}, let $n\ge k+1$. By
\eqref{eq:propagation-key} the
coefficient of $X^n$ in $\Delta\,(1-K_k)$ is at least the coefficient of
$X^n$ in $\Xi$, which is $\xi_n$. The former is
\[
\delta_n-\sum_{m=1}^{n}\kappa^{(k)}_m\,\delta_{n-m}
=\delta_n-\sum_{m=1}^{n-k}\kappa^{(k)}_m\,\delta_{n-m} ,
\]
where the terms with $m>n-k$ vanish because $\delta_j=0$ for $j<k$. By
\eqref{eq:propagation-curvature-step}, applied with $n-1\ge k$ instead of
$n$, we have $n\,\xi_n=(k+1)\,\delta_{n-1}+2\,\beta^{(k-1)}_{n-1}$, and the
inequality of Proposition~\ref{prop:columninequality} follows on multiplying
by $n$.
\end{proof}

Divided by $n$, the inequality of Proposition~\ref{prop:columninequality} is
a lower bound for the step of row $n$ by quantities of the earlier rows,
\[
\delta_n\ge\frac{k+1}{n}\,\delta_{n-1}+\frac2n\,\beta^{(k-1)}_{n-1}
+\sum_{m=1}^{n-k}\kappa^{(k)}_m\,\delta_{n-m} .
\]
The term with $\beta^{(k-1)}_{n-1}$ is a nonnegative source, and the other
terms feed the earlier steps of the column back with nonnegative
coefficients, since $\kappa^{(k)}_m\ge0$ (Lemma~\ref{lem:kernelsign}). On rows
$k+1\le n\le2k-1$ of the first layer the kernel of the sum is the free prime
kernel: there $m\le n-k\le k-1$, so that $\kappa^{(k)}_m=\varepsilon_m$ by
Lemma~\ref{lem:kernelfree}.

\subsection{The propagation of a positive step}\label{subsec:propagation-step}

We now turn Proposition~\ref{prop:onetree} into a statement on the sign of
the step. The two rules from which
Theorem~\ref{thm:peaks} is derived conclude with strict inequalities, as was
explained in Section~\ref{sec:intro}, so that we need a lower bound for the
step that can be positive. The following lemma gives such a bound at a tilt
$x>0$ at which the tilted normalised counts $z^{(k)}_Nx^N$ do not increase
in $N$, under a condition on sums of the numbers $\xi_tx^t$. Its proof is a
summation by parts in the inequality of Proposition~\ref{prop:onetree},
after the weight that the term $t=n$ of the sum carries has been lowered.

\begin{lemma}\label{lem:partialsums}
Let $k\ge2$ and $n\ge k$ be integers and $x>0$ a real number. Assume that
\begin{enumerate}
\item $z^{(k)}_{N+1}\,x^{N+1}\le z^{(k)}_N\,x^N$ for every integer $N\ge0$;
\item for every integer $s$ with $k\le s\le n$,
\[
\sum_{t=s}^{n}\xi^{(k)}_t\,x^t\ge0 .
\]
\end{enumerate}
Then
\[
(1-x)\,\xi^{(k)}_n\le\delta^{(k)}_n .
\]
\end{lemma}

\begin{proof}
We multiply the inequality of Proposition~\ref{prop:onetree} by $x^n$. Since
the terms with $t<k$ vanish by \eqref{eq:propagation-vanishing}, this gives
\begin{equation}\label{eq:propagation-tilted}
\sum_{t=k}^{n}\xi_tx^t\cdot z^{(k)}_{n-t}\,x^{n-t}\le\delta_n\,x^n .
\end{equation}
For $k\le t\le n$ put $\sigma_t=\xi_tx^t$, and let
$\rho_t=z^{(k)}_{n-t}\,x^{n-t}$ for $k\le t<n$ and $\rho_n=z^{(k)}_1\,x$.
Thus $\rho_t$ is the weight of $\sigma_t$ in \eqref{eq:propagation-tilted},
except for $t=n$, where the weight $z^{(k)}_0=1$ has been replaced by
$z^{(k)}_1\,x$. Since $z^{(k)}_1=c_1=1$, we have
\begin{equation}\label{eq:propagation-lowered}
\sum_{t=k}^{n}\sigma_t\,z^{(k)}_{n-t}\,x^{n-t}
=\sum_{t=k}^{n}\sigma_t\,\rho_t+(1-x)\,\xi_n\,x^n .
\end{equation}
The numbers $\rho_t$ are nonnegative, and they do not decrease in $t$: for
$k\le s<n-1$, hypothesis~(1) with $N=n-s-1$ is the inequality
$\rho_s\le\rho_{s+1}$, and if $n>k$, then $\rho_{n-1}=z^{(k)}_1\,x=\rho_n$.

Summation by parts gives
\[
\sum_{t=k}^{n}\sigma_t\,\rho_t=\rho_k\sum_{t=k}^{n}\sigma_t
+\sum_{s=k+1}^{n}(\rho_s-\rho_{s-1})\sum_{t=s}^{n}\sigma_t .
\]
The right-hand side is nonnegative by hypothesis~(2), because $\rho_k\ge0$
and $\rho_{s-1}\le\rho_s$. With \eqref{eq:propagation-tilted} and
\eqref{eq:propagation-lowered} we get
$(1-x)\,\xi_n\,x^n\le\delta_n\,x^n$, and the claim follows.
\end{proof}

Lemma~\ref{lem:partialsums} makes the step $\delta_n$ positive when the tree
curvature $\xi_n$ is positive and $x<1$. In the other direction,
\eqref{eq:propagation-curvature-step} expresses the next tree curvature
$\xi_{n+1}$ by the step $\delta_n$ of a row $n\ge k$ and by the number
$\beta^{(k-1)}_n$, which is nonnegative. Applied alternately, the two
relations turn finitely many inequalities at one row into a sign in
infinitely many rows.

\begin{proposition}\label{prop:propagation}
Let $k\ge2$ and $m\ge k$ be integers and $x$ a real number with $0<x<1$.
Assume that
\begin{enumerate}
\item $z^{(k)}_{N+1}\,x^{N+1}\le z^{(k)}_N\,x^N$ for every integer $N\ge0$;
\item for every integer $s$ with $k\le s\le m+1$,
\[
\sum_{t=s}^{m+1}\xi^{(k)}_t\,x^t\ge0 ;
\]
\item $\delta^{(k)}_m>0$.
\end{enumerate}
Then $\delta^{(k)}_n>0$ for every integer $n\ge m$.
\end{proposition}

\begin{proof}
We say that the sums are nonnegative at a row $n\ge k$ if hypothesis~(2) of
Lemma~\ref{lem:partialsums} holds for this $n$, and we show by induction on
$n\ge m$ that $\delta_n>0$ and that the sums are nonnegative at row
$n+1$. For $n=m$ these are the hypotheses (3) and~(2).

Let $n\ge m$, and suppose that $\delta_n>0$ and that the sums are nonnegative
at row $n+1$. Since $n\ge k$ and $\beta^{(k-1)}_n\ge0$, the right-hand
side of \eqref{eq:propagation-curvature-step} is positive, and hence
$\xi_{n+1}>0$. We apply Lemma~\ref{lem:partialsums} at row $n+1$; its
hypotheses are hypothesis~(1) and the nonnegativity of the sums at row
$n+1$. As $1-x>0$, it gives
\[
\delta_{n+1}\ge(1-x)\,\xi_{n+1}>0 .
\]
By \eqref{eq:propagation-curvature-step} again, now for row $n+1$, we
have $\xi_{n+2}>0$. Hence the sums are nonnegative at row $n+2$ as well:
each of them is one of the sums at row $n+1$, or the empty sum, plus the
positive term $\xi_{n+2}\,x^{n+2}$. This completes the induction.
\end{proof}

In Section~\ref{sec:beyond} the proposition is applied to columns
$2\le k\le81$; hypotheses (2) and~(3) are then finitely many inequalities,
which Lemmas~\ref{lem:smallcolumns} and~\ref{lem:junction} supply.

\subsection{The critical tilt}\label{subsec:propagation-tilt}

It remains to find a tilt $x<1$ at which the tilted normalised counts
$z^{(k)}_Nx^N$ of threshold $k$ do not increase in $N$. In the free range
$N\le k$ these are the tilted Cayley masses $c_Nx^N$, and the ratio of two
consecutive masses is explicit: for every real number $x$ and every integer
$j\ge0$,
\begin{equation}\label{eq:propagation-ratio}
c_{j+1}\,x^{j+1}=x\,\Bigl(1+\frac{1}{j+1}\Bigr)^{j}\,c_j\,x^j .
\end{equation}
This follows from $c_j=a_{j+1}$ and from \eqref{eq:prelim-quotient}. By
Lemma~\ref{lem:powermonotone} the factor $(1+1/(j+1))^{j}$ does not decrease
in $j$. Hence, among the factors $(1+1/(j+1))^{j}$ with $j+1\le k$, the
largest is the factor $(1+1/k)^{k-1}$ of $j=k-1$.
We take the tilt $x$ for which $x\,(1+1/k)^{k-1}=1$.

For an integer $k\ge1$ the \emph{critical tilt} of column $k$ is
\[
x_k=\Bigl(\frac{k}{k+1}\Bigr)^{k-1} .
\]
It is the real number determined by
\begin{equation}\label{eq:propagation-critical}
x_k\,\Bigl(1+\frac1k\Bigr)^{k-1}=1 .
\end{equation}
Clearly $0<x_k\le1$, with $x_1=1$. The first values
are $x_2=2/3$, $x_3=9/16$ and $x_4=64/125$. By \eqref{eq:propagation-ratio}
for $j=k-1$, the equation \eqref{eq:propagation-critical} says that
at the critical tilt the tilted Cayley
masses of indices $k-1$ and $k$ are equal. When a column $k$ is fixed, we
write
\[
\omega_n=c_n\,x_k^{\,n}\qquad\text{and}\qquad u^{(y)}_N=z^{(y)}_N\,x_k^{\,N}
\]
for the tilted Cayley masses and for the tilted normalised counts of a
threshold $y\ge1$ at the critical tilt of this column, so that
$u^{(y)}_N=\omega_N$ for $N\le y$ by Proposition~\ref{prop:freerange}.

The first lemma is the monotonicity of the masses for which the critical tilt
was chosen.

\begin{lemma}\label{lem:massesdecrease}
Let $k$ and $n$ be integers with $n\ge0$ and $n+1\le k$. Then the critical
tilt $x_k$ of column $k$ satisfies
\[
c_{n+1}\,x_k^{\,n+1}\le c_n\,x_k^{\,n} .
\]
\end{lemma}

\begin{proof}
By \eqref{eq:propagation-ratio} the left-hand side is the right-hand side
multiplied by $x_k\,(1+1/(n+1))^{n}$, and this factor is at most
$x_k\,(1+1/k)^{k-1}=1$, by Lemma~\ref{lem:powermonotone} applied to the
integers $n+1\le k$ and by \eqref{eq:propagation-critical}.
\end{proof}

Since $u^{(k)}_N=\omega_N$ for $N\le k$, the lemma says for a column
$k\ge2$ that at the critical tilt the tilted normalised counts
of threshold $k$ do not increase in their free range.
Proposition~\ref{prop:tilteddecrease} below states that at this tilt they do
not increase anywhere. Besides Lemma~\ref{lem:massesdecrease}, its proof
needs one further inequality in the free range, namely
\begin{equation}\label{eq:propagation-mass}
x_k\sum_{j=1}^{k}c_j\,x_k^{\,j}\le1 .
\end{equation}
Since $x\cdot c_j\,x^j=a_{j+1}\,x^{j+1}$ for every real number $x$, the
left-hand side can also be written with the Cayley coefficients:
\begin{equation}\label{eq:propagation-cayley}
x_k\sum_{j=1}^{k}c_j\,x_k^{\,j}=\sum_{n=2}^{k+1}a_n\,x_k^{\,n} .
\end{equation}
For $k\ge12$ we prove
\eqref{eq:propagation-mass} by an estimate (Lemma~\ref{lem:cayleysum}); for
$2\le k\le11$ the sum is evaluated, in the proof of
Proposition~\ref{prop:tilteddecrease}. The estimate rests on Stirling's
formula, which controls the numbers $a_ne^{-n}$, the terms $a_n\,x^n$ at
the tilt $x=1/e$. The next lemma bounds the quotient of the critical tilt by
this tilt. It will be used in the form $(e\,x_k)^n\le e^{3n/(2k)}$, a bound
that stays below a fixed number as long as $n$ is at most a fixed multiple
of $k$.

\begin{lemma}\label{lem:tiltupper}
For every integer $k\ge1$ the critical tilt $x_k$ satisfies
\[
e\,x_k\le\exp\Bigl(\frac{3}{2k}\Bigr).
\]
\end{lemma}

\begin{proof}
For every real number $t\ge0$ one has $\log(1+t)\ge2t/(2+t)$, since the two
sides agree at $t=0$ and the derivative of their difference is
$t^2/((1+t)(2+t)^2)\ge0$. For $t=1/k$ this gives
\[
(k-1)\,\log\Bigl(1+\frac1k\Bigr)\ge\frac{2(k-1)}{2k+1}\ge1-\frac{3}{2k} ,
\]
where the second inequality is equivalent to $(2k-3)(2k+1)\le4k(k-1)$, and
the left-hand side of this is $4k(k-1)-3$. Hence
$(1+1/k)^{k-1}\ge\exp(1-3/(2k))$, and the claim follows from
\eqref{eq:propagation-critical}.
\end{proof}

We now prove \eqref{eq:propagation-mass} for $k\ge12$.

\begin{lemma}\label{lem:cayleysum}
For every integer $k\ge12$ the critical tilt $x_k$ satisfies
\[
x_k\sum_{j=1}^{k}c_j\,x_k^{\,j}\le1 .
\]
\end{lemma}

\begin{proof}
Write $x=x_k$. By \eqref{eq:propagation-cayley} the left-hand side is the sum
of $a_n\,x^n$ over $2\le n\le k+1$. We first bound a single term.
For $n\ge1$,
\[
a_n\,x^n=\frac{n^{n-1}}{n!}\,x^n=\frac{(e\,x)^n}{n}\,P(n)
\le\frac{c_{\mathrm S}\,(e\,x)^n}{n\sqrt n} ,
\]
by the definition of the Poisson mass and by the upper bound of
\eqref{eq:prelim-stirling}.
Lemma~\ref{lem:tiltupper} gives $(e\,x)^n\le e^{3n/(2k)}$. For
$2\le n\le k+1$ the exponent satisfies
$0\le3n/(2k)\le3(k+1)/(2k)\le1.625$, because $k\ge12$, and the inequality
$e^{\theta}\le1+\theta+\theta^2$, which holds for $0\le\theta\le1.7$
(Section~\ref{sec:prelim}), applies with $\theta=3n/(2k)$. Hence, for
$2\le n\le k+1$,
\begin{equation}\label{eq:propagation-term}
a_n\,x^n\le c_{\mathrm S}\,\Bigl(\frac{1}{n\sqrt n}
+\frac{3}{2k}\cdot\frac{1}{\sqrt n}+\frac{9}{4k^2}\,\sqrt n\Bigr).
\end{equation}

Next we sum the three terms over $2\le n\le k+1$. For the first term we use
that for every real number $t\ge1$,
\[
\frac{1}{t\sqrt t}\le\int_{t-1/2}^{t+1/2}\frac{d\sigma}{\sigma\sqrt\sigma}
=\frac{2}{\sqrt{t-1/2}}-\frac{2}{\sqrt{t+1/2}} ,
\]
because the function $\sigma\mapsto\sigma^{-3/2}$ is convex. The sum of the
right-hand side over $t=2,\dots,k+1$ telescopes, so that
\[
\sum_{n=2}^{k+1}\frac{1}{n\sqrt n}
\le\frac{2}{\sqrt{3/2}}-\frac{2}{\sqrt{k+3/2}} .
\]
For the other two terms we use the bounds
\[
\sum_{i=1}^{m}\frac{1}{\sqrt i}\le2\sqrt m-1\quad(m\ge1)
\qquad\text{and}\qquad
\sum_{i=0}^{m-1}\sqrt i\le\frac23\,m\sqrt m\quad(m\ge0)
\]
of Section~\ref{sec:prelim}, the first with $m=k+1$ and the second with
$m=k+2$. Leaving out the term $i=1$ of the first sum and the terms $i=0$ and
$i=1$ of the second, we get
\[
\sum_{n=2}^{k+1}\frac{1}{\sqrt n}\le2\sqrt{k+1}-2,\qquad
\sum_{n=2}^{k+1}\sqrt n\le\frac23\,(k+2)\sqrt{k+2} .
\]
Let $r=\sqrt{k+2}$. Then $\sqrt{k+1}\le r$ and $\sqrt{k+3/2}\le r$, and the
last bound, multiplied by $9/(4k^2)$, is $3r^3/(2k^2)$. Summing
\eqref{eq:propagation-term} over $2\le n\le k+1$, we therefore obtain
\begin{equation}\label{eq:propagation-total}
\sum_{n=2}^{k+1}a_n\,x^n\le c_{\mathrm S}\,\Bigl(\frac{2}{\sqrt{3/2}}
-\frac2r+\frac{3\,(r-1)}{k}+\frac{3\,r^3}{2k^2}\Bigr).
\end{equation}

We claim that the last three terms in the parenthesis add up to at most
$0.87$. After multiplication by $2rk^2$, and with $r^2=k+2$, the claim is
equivalent to
\[
2.5\,k^2+12\,k+6\le r\,k\,(0.87\,k+3).
\]
Since $k\ge12$, we have $r\ge\sqrt{14}\ge3.74$, so that the right-hand side
is at least $3.2538\,k^2+11.22\,k$, and the inequality holds because
$0.7538\,k^2\ge9\,k\ge0.78\,k+6$ for $k\ge12$. Moreover
$2/\sqrt{3/2}\le1.6331$, because $1.5\cdot1.6331^2\ge4$, and
$c_{\mathrm S}\le0.39894229$. The right-hand side of
\eqref{eq:propagation-total} is thus at most
$0.39894229\cdot(1.6331+0.87)<1$, which proves the lemma.
\end{proof}

With Lemmas~\ref{lem:massesdecrease} and~\ref{lem:cayleysum} we can prove
the monotonicity at the critical tilt, which is hypothesis~(1) of
Proposition~\ref{prop:propagation} for $x=x_k$.

\begin{proposition}\label{prop:tilteddecrease}
Let $k\ge2$ be an integer and $x_k$ the critical tilt of column $k$. Then for
every integer $N\ge0$,
\[
z^{(k)}_{N+1}\,x_k^{\,N+1}\le z^{(k)}_N\,x_k^{\,N} .
\]
\end{proposition}

\begin{proof}
We first prove a statement on real sequences. Let $y\ge1$ be an integer,
$x\ge0$ a real number, and $(v_N)_{N\ge0}$ a sequence of real numbers with
$v_0=1$ and
\begin{equation}\label{eq:propagation-recursion}
(N+1)\,v_{N+1}=x\sum_{i=0}^{N}\min(i+1,y)\,v_i\,v_{N-i}\qquad(N\ge0).
\end{equation}
Suppose that
\begin{equation}\label{eq:propagation-conditions}
v_{N+1}\le v_N\quad\text{for }0\le N<y,
\qquad\text{and}\qquad x\sum_{j=1}^{y}v_j\le1 .
\end{equation}
We claim that $v_{N+1}\le v_N$ for every $N\ge0$.

By \eqref{eq:propagation-recursion} and induction, $v_N\ge0$ for every $N$.
Put $\phi_N=v_N-v_{N+1}$. For $N\ge y$ we have
\begin{equation}\label{eq:propagation-difference}
(N+1)\,\phi_N=(1-xy)\,v_N
+x\sum_{i=0}^{N-1}\min(i+1,y)\,v_i\,\phi_{N-1-i} .
\end{equation}
To see this, write $(N+1)\,\phi_N=v_N+N\,v_N-(N+1)\,v_{N+1}$ and subtract
\eqref{eq:propagation-recursion} from the same equation with $N$ replaced by
$N-1$. The term $i=N$ of \eqref{eq:propagation-recursion} is $xy\,v_N$,
because $N+1>y$ and $v_0=1$, and the terms with $i\le N-1$ of the two
equations combine to the sum in \eqref{eq:propagation-difference}.

We prove that $\phi_N\ge0$ by induction on $N$. For $N<y$ this is the first
condition in \eqref{eq:propagation-conditions}. Let $N\ge y$, and suppose
that $\phi_j\ge0$ for all $j<N$. Then $v_N\le v_i$ for every $i\le N$, and
every term of the sum in \eqref{eq:propagation-difference} is nonnegative. We
keep only the terms with $i\ge N-y$ and index them by $j=N-1-i$, so that
$0\le j\le y-1$; these terms occur in the sum because $y\le N$. For them
$\min(N-j,y)\ge y-j$, as $N\ge y$, and $v_{N-1-j}\ge v_N$. Hence
\[
\sum_{i=0}^{N-1}\min(i+1,y)\,v_i\,\phi_{N-1-i}
\ge v_N\sum_{j=0}^{y-1}(y-j)\,\phi_j .
\]
The last sum is evaluated by summation by parts: since
$\phi_j=v_j-v_{j+1}$ and $v_0=1$,
\[
\sum_{j=0}^{y-1}(y-j)\,\phi_j=y\,v_0-\sum_{j=1}^{y}v_j
=y-\sum_{j=1}^{y}v_j .
\]
We insert this into \eqref{eq:propagation-difference}, where the terms
$xy\,v_N$ cancel, and find that
\[
(N+1)\,\phi_N\ge(1-xy)\,v_N+x\,v_N\,\Bigl(y-\sum_{j=1}^{y}v_j\Bigr)
=v_N\,\Bigl(1-x\sum_{j=1}^{y}v_j\Bigr)\ge0
\]
by the second condition in \eqref{eq:propagation-conditions}. This proves the
claim.

We apply the claim with $y=k$, with $x=x_k$ and with
$v_N=u^{(k)}_N=z^{(k)}_N\,x_k^{\,N}$. This sequence has
$v_0=z^{(k)}_0=1$, and it satisfies \eqref{eq:propagation-recursion} by
\eqref{eq:counts-tilted}. For $N\le k$ we have $v_N=c_N\,x_k^{\,N}$ by
Proposition~\ref{prop:freerange}, so that the first condition in
\eqref{eq:propagation-conditions} is Lemma~\ref{lem:massesdecrease} and the
second is the inequality \eqref{eq:propagation-mass}. For $k\ge12$ this
inequality is Lemma~\ref{lem:cayleysum}. For $2\le k\le11$ its left-hand
side, in the form \eqref{eq:propagation-cayley}, is a sum of $k$ rational
numbers, and its exact value is smaller than $1$ in each of the ten cases.
For $k=2$ it is $x_2^{\,2}+(3/2)\,x_2^{\,3}=8/9$, and rounded to three
decimal places the values for $k=2,3,\dots,11$ are $0.889$, $0.850$,
$0.830$, $0.817$, $0.807$, $0.799$, $0.793$, $0.788$, $0.783$ and $0.779$.
Hence $u^{(k)}_{N+1}\le u^{(k)}_N$ for every $N\ge0$, which is the assertion.
\end{proof}

Because of Proposition~\ref{prop:tilteddecrease} the arguments for the rows
beyond the first layer are carried out at the critical tilt. Without the
threshold the
corresponding sequence behaves in the opposite way beyond the
free range: by \eqref{eq:propagation-ratio} and
Lemma~\ref{lem:powermonotone} the numbers $c_N\,x_k^{\,N}$ do not decrease
from index $k-1$ on.

\section{Exact computations}\label{sec:tables}
The signs of the steps in the small columns, which the estimates of the
later sections do not reach, are obtained by exact computation with the
numbers of parking functions. The results are stated in the five lemmas of
this section. They concern the steps $d_{n,k}$ at finitely many positions
of the array, all in rows $n\le300$.

Most of what the lemmas state is the sign of the step at each position of a
finite set, as a strict inequality $d_{n,k}<0$ or $d_{n,k}>0$. The
inequalities are strict because the later
sections need them in this form: the column rule and the diagonal rule
conclude with strict inequalities, and the diagonal rule starts from a step
that may be zero. A step equal to zero therefore has to be excluded at the
positions from which the rules are read, and where no argument excludes
it, the sign is computed. Lemma~\ref{lem:smallcolumns}, on columns
$2\le k\le81$, states facts of other kinds. It says that a positive step
in a row of the first layer is followed by a positive step in the next row
of the layer, and that finitely many sums formed from the cumulative
counts are nonnegative.

\subsection{The integers that are computed}\label{subsec:tables-integers}

All five lemmas concern columns $k\ge2$, and they are proved through the
cumulative counts $A_{n,\ell}$. By \eqref{eq:counts-second} the sign
of a step is decided by a comparison between integers: for $k\ge2$ the step
$d_{n,k}$ is positive if and only if $2A_{n,k-1}<A_{n,k}+A_{n,k-2}$, and it
is negative if and only if $A_{n,k}+A_{n,k-2}<2A_{n,k-1}$.

The numbers $A_{n,\ell}$ are computed from the last-car recursion of
Proposition~\ref{prop:lastcar}, together with the value $A_{0,\ell}=1$. For
a fixed $\ell$ the recursion determines $A_{1,\ell},A_{2,\ell},\dots$ one
after the other, each from all the numbers before it.

The signs and the sums are obtained by computing the integers concerned and
comparing them. We carry out these computations by computer, in exact
integer arithmetic. So that the
reader can repeat the computations, each proof says which integers are
computed, by which recursion, and which comparisons are made. The
numbers $n_0(k)$ and $\zeta(k)$ that enter the statements of four of the
lemmas are listed in full (Definitions~\ref{def:junctionrows}
and~\ref{def:zeta}).

A repetition of the computations can be tested against the following. By
Proposition~\ref{prop:freerange} we have $A_{n,\ell}=n!\,c_n=(n+1)^{n-1}$
whenever $n\le\ell+1$, so that for each $\ell$ the numbers
$A_{0,\ell},\dots,A_{\ell+1,\ell}$, which the recursion produces first, are
known in closed form. The two end values of a row that are recalled in
Section~\ref{sec:intro} say that $A_{n,0}=n!$ for every $n\ge1$ and that
$A_{n,n-1}-A_{n,n-2}=n^{n-2}$ for every $n\ge2$. Rows $n\le7$ of the array
are displayed in Section~\ref{sec:intro}, and the numbers $b_{n,\ell}$ with
$n\le9$ are tabulated in \cite[Table~1]{celano2025}; the cumulative counts
are their partial sums.

The first rows show how the recursion runs. The
recursion gives $A_{1,\ell}=1$ and $A_{2,\ell}=1+\min(2,\ell+1)$, so that
$A_{2,0}=2$ and $A_{2,\ell}=3$ for $\ell\ge1$, and then
\[
A_{3,\ell}=A_{2,\ell}+2\min(2,\ell+1)+\min(3,\ell+1)\,A_{2,\ell}.
\]

Besides the results of the five lemmas, a further fact about small columns
is obtained by evaluation and not by
an estimate; it concerns explicit rational numbers and not the numbers
$A_{n,\ell}$. For each of the ten columns $2\le k\le11$ the proof of
Proposition~\ref{prop:tilteddecrease} evaluates a sum of $k$ rational
numbers and compares it with $1$. The arithmetic with explicit constants that the
estimates of the later sections require is part of their proofs; this
includes the expansion of a polynomial of degree $15$ in the proof of
Proposition~\ref{prop:curvestep}.

\subsection{Negative steps in columns 2 to 9}\label{subsec:tables-negative}

The first lemma gives negative signs, and three statements of the later
sections take such signs from it. Beyond the first layer the step is
positive except at eleven positions, the positions $(2k,k)$ with
$1\le k\le7$ and $(2k+1,k)$ with $2\le k\le5$
(Proposition~\ref{prop:beyondlayer}). We call these eleven positions the
\emph{exceptional positions}. The two rules in this region,
Corollaries~\ref{cor:columnbeyond} and~\ref{cor:diagonalbeyond}, need a
negative step one row above each exceptional position other than $(2,1)$
and at the diagonal neighbour of each exceptional position. In the first
layer of columns $2\le k\le8$ the diagonal rule needs the negative sign
of column $k+1$ in rows $k+2$ to $2k$
(Proposition~\ref{prop:rulesmall}). All these positions are among the
positions $(n,k)$ of the following table, of which there are
$3+4+5+6+6+7+7+7=45$.
\begin{equation}\label{eq:tables-cells}
\begin{array}{l|cccccccc}
\text{column $k$} & 2 & 3 & 4 & 5 & 6 & 7 & 8 & 9\\
\hline
\text{rows $n$} & 3\text{--}5 & 4\text{--}7 & 5\text{--}9 & 6\text{--}11
 & 7\text{--}12 & 8\text{--}14 & 9\text{--}15 & 10\text{--}16
\end{array}
\end{equation}

In this table each column $2\le k\le8$ has every row of its first layer, and
in addition row $2k$ when $k\le7$ and row $2k+1$ when $k\le5$. These
ten positions beyond the first layer are the exceptional positions other
than $(2,1)$. Column $9$ has rows $10\le n\le16$, which are its
first layer without the last row; in row $17$, column $9$ is
positive, by Lemma~\ref{lem:signpair} below and $\zeta(9)=7$. The lemma
describes the positions of the table by four conditions.

\begin{lemma}\label{lem:fortyfive}
Let $k\ge2$ and $n\ge k+1$ be integers such that one of the following four
conditions holds:
\begin{enumerate}
\item $n\le2k+1$ and $k\le5$;
\item $n\le2k$ and $k\le7$;
\item $n\le2k-1$ and $k\le8$;
\item $n\le2k-2$ and $k\le9$.
\end{enumerate}
Then $d_{n,k}<0$.
\end{lemma}

\begin{proof}
By \eqref{eq:counts-second}, the inequality $d_{n,k}<0$ is
equivalent to
\[
A_{n,k}+A_{n,k-2}<2A_{n,k-1}.
\]
We compute the integers $A_{n,\ell}$ with $n\le16$ and $\ell\le9$ from the
last-car recursion; they have at most $19$ decimal digits. We find that
this inequality holds at each of the forty-five positions of
\eqref{eq:tables-cells}; at the thirteen of them that lie in rows $n\le7$
it can also be read from the rows of the array that are displayed in
Section~\ref{sec:intro}. Every pair $(n,k)$ as in
the lemma is one of these positions, because the last row of column $k$ in
\eqref{eq:tables-cells} is $2k+1$ for $k\le5$, at least $2k$ for $k\le7$, at
least $2k-1$ for $k\le8$, and at least $2k-2$ for $k\le9$.
\end{proof}

\subsection{Columns 2 to 81}\label{subsec:tables-small}

For a column $2\le k\le81$ the sign of the step has to be determined in every
row $n\ge2k$ (Proposition~\ref{prop:beyondlayer}), and these are infinitely
many rows. They are divided at a row $n_0(k)$ that is listed for the column.
In the rows below $n_0(k)$ the signs are computed
(Lemma~\ref{lem:junction}). At row $n_0(k)$ a condition is verified
under which Proposition~\ref{prop:propagation} carries the positive step of
row $n_0(k)-1$ to all later rows (Lemma~\ref{lem:smallcolumns}, part 4).
The columns end at $81$ because
from column $82$ on the signs beyond the first layer and the column rule
in the first layer are proved by analysis
(Propositions~\ref{prop:largecolumns} and~\ref{prop:onechange}).

\begin{definition}\label{def:junctionrows}
For $2\le k\le81$ the number $n_0(k)$ is given by the following table.
\[
\begin{array}{r|cccccccccc}
k & 2 & 3 & 4 & 5 & 6 & 7 & 8 & 9 & 10 & 11\\
n_0(k) & 7 & 10 & 13 & 15 & 18 & 20 & 23 & 25 & 28 & 30\\
\hline
k & 12 & 13 & 14 & 15 & 16 & 17 & 18 & 19 & 20 & 21\\
n_0(k) & 32 & 35 & 37 & 39 & 41 & 44 & 46 & 48 & 50 & 52\\
\hline
k & 22 & 23 & 24 & 25 & 26 & 27 & 28 & 29 & 30 & 31\\
n_0(k) & 55 & 57 & 59 & 61 & 63 & 66 & 68 & 70 & 72 & 74\\
\hline
k & 32 & 33 & 34 & 35 & 36 & 37 & 38 & 39 & 40 & 41\\
n_0(k) & 76 & 78 & 80 & 82 & 85 & 87 & 89 & 91 & 93 & 95\\
\hline
k & 42 & 43 & 44 & 45 & 46 & 47 & 48 & 49 & 50 & 51\\
n_0(k) & 97 & 99 & 101 & 103 & 105 & 107 & 109 & 111 & 114 & 116\\
\hline
k & 52 & 53 & 54 & 55 & 56 & 57 & 58 & 59 & 60 & 61\\
n_0(k) & 118 & 120 & 122 & 124 & 126 & 128 & 130 & 132 & 134 & 136\\
\hline
k & 62 & 63 & 64 & 65 & 66 & 67 & 68 & 69 & 70 & 71\\
n_0(k) & 138 & 140 & 142 & 144 & 146 & 148 & 150 & 152 & 154 & 156\\
\hline
k & 72 & 73 & 74 & 75 & 76 & 77 & 78 & 79 & 80 & 81\\
n_0(k) & 158 & 160 & 162 & 164 & 166 & 168 & 170 & 171 & 173 & 175
\end{array}
\]
\end{definition}

Rows $2k\le n<n_0(k)$, which lie between the first layer and row
$n_0(k)$, are called the \emph{junction} of column $k$. Comparing the
entries of the table with $2k$, one finds the following values of the
difference $n_0(k)-2k$.
\[
\begin{array}{r|ccccccc}
n_0(k)-2k & 3 & 4 & 5 & 6 & 7 & 8 & 9\\
\text{columns $k$} & 2 & 3 & 4,\,5 & 6,\,7 & 8,\,9 & 10\text{--}12
 & 13\text{--}16\\
\hline
n_0(k)-2k & 10 & 11 & 12 & 13 & 14 & 13 & \\
\text{columns $k$} & 17\text{--}21 & 22\text{--}26 & 27\text{--}35
 & 36\text{--}49 & 50\text{--}78 & 79\text{--}81 &
\end{array}
\]
In particular $2k+3\le n_0(k)\le2k+14$, so that a junction has between $3$
and $14$ rows.

The numbers $n_0(k)$ are given by the table and not by a formula. All that
the later sections use about $n_0(k)$ is what Lemma~\ref{lem:junction} and
part 4 of Lemma~\ref{lem:smallcolumns} below state, and
any row for which the two statements hold could be listed instead of
$n_0(k)$.

The first lemma on these columns gives the signs in the junctions.

\begin{lemma}\label{lem:junction}
Let $k$ be an integer with $2\le k\le81$, and let $n_0(k)$ be the number of
Definition~\ref{def:junctionrows}. Then the following two statements hold.
\begin{enumerate}
\item One has $2k+3\le n_0(k)$.
\item For every integer $n$ with $2k\le n<n_0(k)$, other than $n=2k$ when
$k\le7$ and $n=2k+1$ when $k\le5$, one has $d_{n,k}>0$.
\end{enumerate}
\end{lemma}

\begin{proof}
Part 1 is read off from the table of Definition~\ref{def:junctionrows}. For
part 2, the inequality $d_{n,k}>0$ is equivalent, by
\eqref{eq:counts-second}, to
\begin{equation}\label{eq:tables-positive}
2A_{n,k-1}<A_{n,k}+A_{n,k-2}.
\end{equation}
The rows of part 2 are rows $2k+2\le n<n_0(k)$ for $2\le k\le5$, rows $2k+1\le
n<n_0(k)$ for $k=6,7$, and rows $2k\le n<n_0(k)$ for
$8\le k\le81$. These are $933$ positions, of which $26$ lie in columns
$2$ to $8$ and $907$ in columns $9$ to $81$, and their rows range from
$6$ to $174$. We compute the integers $A_{n,\ell}$ with $n\le174$ and
$\ell\le81$ from the last-car recursion; the largest of them,
$A_{174,81}$, has $387$ decimal digits. We find that
\eqref{eq:tables-positive} holds at each of the $933$ positions.
\end{proof}

The rows that part 2 of Lemma~\ref{lem:junction} excepts are those of the
ten exceptional positions
$(2k,k)$ with $2\le k\le7$ and $(2k+1,k)$ with $2\le k\le5$, at which the
step is negative by Lemma~\ref{lem:fortyfive}. Part 1 is recorded for the
application of Proposition~\ref{prop:propagation}, which starts from a
positive step in row $n_0(k)-1$: this row has to be a row of the junction
that lies after the excepted rows $2k$ and $2k+1$.

The second lemma on columns $2\le k\le81$ concerns a column in the rows
from its first layer to row $n_0(k)$. The later sections take from it the
two facts about these columns that are not single signs. Part 2 is the
column rule for two consecutive rows of the first layer, which is
Proposition~\ref{prop:layerpersistence} for columns $k\le81$, and part 4 is
the second hypothesis of Proposition~\ref{prop:propagation} with $x=x_k$
and $m=n_0(k)-1$. Parts 1 and 3 are weaker forms of the two parts of
Lemma~\ref{lem:junction} and follow from it.

\begin{lemma}\label{lem:smallcolumns}
Let $k$ be an integer with $2\le k\le81$, let $n_0=n_0(k)$ be the number of
Definition~\ref{def:junctionrows}, let $x_k$ be the critical tilt of column
$k$, and let $\xi_t=\xi^{(k)}_t$ be the tree curvature of column $k$
(Definition~\ref{def:treecurvature}). Then the following four statements
hold.
\begin{enumerate}
\item One has $2k\le n_0$.
\item For every integer $n$ with $k+1\le n\le2k-2$, if $d_{n,k}>0$, then
$d_{n+1,k}>0$.
\item For every integer $n$ with $2k\le n$, $17\le n$ and $n<n_0$, one has
$d_{n,k}\ge0$.
\item For every integer $s$ with $k\le s\le n_0$,
\[
\sum_{t=s}^{n_0}\xi_t\,x_k^{\,t}\ge0.
\]
\end{enumerate}
\end{lemma}

\begin{proof}
Part 1 holds because $2k+3\le n_0$ by part 1 of Lemma~\ref{lem:junction}.
For part 3, let $n$ be an
integer with $2k\le n$, $17\le n$ and $n<n_0$. The rows that part 2 of
Lemma~\ref{lem:junction} excepts are row $2k$ of a column $k\le7$ and row
$2k+1$ of a column $k\le5$, and these are rows $n\le14$. Since $n\ge17$,
row $n$ is not excepted, and that lemma gives $d_{n,k}>0$.

For parts 2 and 4 we compute, for the three values $\ell=k,k-1,k-2$, the
integers $A_{m,\ell}$ with $0\le m\le n_0-1$ from the last-car recursion.
By \eqref{eq:counts-second}, part 2 asserts that for every $n$ with
$k+1\le n\le2k-2$ at least one of the two inequalities
\[
A_{n,k}+A_{n,k-2}\le2A_{n,k-1},\qquad
2A_{n+1,k-1}<A_{n+1,k}+A_{n+1,k-2}
\]
holds; the first of them says that $d_{n,k}\le0$, and the second that
$d_{n+1,k}>0$. Rows $n$ and $n+1$ are among the rows computed, since
$n+1\le2k-1<n_0$ by part 1. For the eighty columns together these are
\[
\sum_{k=2}^{81}(k-2)=3160
\]
pairs $(k,n)$, and we find that for each of them one of the two
inequalities holds.

For part 4 we write the tree curvature with the cumulative counts. By
Definition~\ref{def:treecurvature} and since $z^{(y)}_n=A_{n,y-1}/n!$, we
have, for $t\ge1$,
\begin{multline}\label{eq:tables-curvature}
\xi_t=\frac1{t!}\Bigl(\min(t,k+1)\,A_{t-1,k}-2\min(t,k)\,A_{t-1,k-1}\\
+\min(t,k-1)\,A_{t-1,k-2}\Bigr).
\end{multline}
Let $I_t=t!\,\xi_t$ for $1\le t\le n_0$. This is the integer
in the bracket of \eqref{eq:tables-curvature}. It is formed
from the numbers $A_{t-1,\ell}$ of the three values of $\ell$, which lie in
a row $t-1\le n_0-1$. For $k\le s\le n_0$ let
\[
\Sigma_s=\sum_{t=s}^{n_0}\frac{n_0!}{t!}\;k^{(k-1)(t-s)}\,
(k+1)^{(k-1)(n_0-t)}\;I_t ,
\]
which is an integer as well. Since $x_k=k^{k-1}/(k+1)^{k-1}$, one has
\[
\Sigma_s=n_0!\,(k+1)^{(k-1)n_0}\,k^{-(k-1)s}\sum_{t=s}^{n_0}\xi_t\,x_k^{\,t},
\]
and the factor in front of the sum is positive. Part 4 therefore asserts
that $\Sigma_s\ge0$ for every $k\le s\le n_0$. The integers $\Sigma_s$ are
computed for $s=n_0,n_0-1,\dots,k$ in this order, from $\Sigma_{n_0+1}=0$
and
\[
\Sigma_s=\frac{n_0!}{s!}\,(k+1)^{(k-1)(n_0-s)}\,I_s+k^{k-1}\,\Sigma_{s+1}.
\]

For example, let $k=2$, so that $n_0=7$, $k^{k-1}=2$ and $(k+1)^{k-1}=3$.
The numbers $A_{m,\ell}$ with $1\le m\le6$ and $\ell\le2$ are partial sums
of the rows displayed in Section~\ref{sec:intro}, and with them
\[
(I_2,I_3,\dots,I_7)=(-1,\,-1,\,2,\,51,\,710,\,9819).
\]
The recursion then gives
\[
(\Sigma_7,\Sigma_6,\dots,\Sigma_2)
=(9819,\ 34548,\ 88374,\ 188088,\ 308136,\ 3912),
\]
and these six integers are positive.

For the eighty columns together there are
\[
\sum_{k=2}^{81}\bigl(n_0(k)-k+1\bigr)=4343
\]
integers $\Sigma_s$. They have up to $14737$ decimal digits, a size that is
reached by $\Sigma_{81}$ for $k=81$. We find that $\Sigma_s\ge0$ for every
$2\le k\le81$ and every $k\le s\le n_0(k)$.
\end{proof}

\subsection{Columns 9 to 264}\label{subsec:tables-large}

By Lemma~\ref{lem:fortyfive} each column $2\le k\le8$ is negative in every
row of its first layer, and by Lemma~\ref{lem:junction} columns
$2,3,\dots,8$ are positive in rows $6$, $8$, $10$, $12$, $13$, $15$ and
$16$, respectively. For each column $9\le k\le264$,
Lemma~\ref{lem:signpair} gives a positive step in a row of its first layer,
and Lemma~\ref{lem:nextcolumn} gives the negative signs of column $k+1$ in
the rows from $k+2$ up to that row. These two lemmas end
at column $264$ because the statements that replace them,
Propositions~\ref{prop:curveincrease} and~\ref{prop:diagonalstep}, are
proved from column $265$ on.

The positions that these two lemmas concern lie in the rows up to $300$ and
in the columns up to $265$, but each of them lies in the first layer of its
column, and for such positions the cumulative counts need not be computed
from row $0$ on. Let $\lvert\mathrm{PF}_n\rvert$ denote the number of parking
functions of length $n\ge0$. No displacement of a parking function of length
$n$ exceeds $n-1$, so that $A_{n,\ell}=\lvert\mathrm{PF}_n\rvert$ for
$\ell\ge n-1$, and Proposition~\ref{prop:freerange} gives
$\lvert\mathrm{PF}_n\rvert=n!\,c_n=(n+1)^{n-1}$; in particular
$\lvert\mathrm{PF}_0\rvert=1$. For integers $n\ge0$ and $\ell\ge0$ let
\[
\bar A_{n,\ell}=\lvert\mathrm{PF}_n\rvert-A_{n,\ell}
\]
be the number of parking functions of length $n$ whose maximum displacement
is greater than $\ell$. Then $\bar A_{n,\ell}=0$ for $n\le\ell+1$, and
\eqref{eq:counts-second} becomes
\begin{equation}\label{eq:tables-complement}
d_{n,k}=2\bar A_{n,k-1}-\bar A_{n,k}-\bar A_{n,k-2}
\end{equation}
for all integers $k\ge2$ and $n\ge0$.

The proofs of the two lemmas below compute the numbers $\bar A_{n,\ell}$ by
a recursion in $n$, which the proof of Lemma~\ref{lem:signpair} derives
from Proposition~\ref{prop:chainforward}.

\begin{definition}\label{def:zeta}
For $9\le k\le264$ the number $\zeta(k)$ is given by the following table:
$\zeta(k)$ is the value $\zeta$ in whose range of columns $k$ lies.
\[
\begin{array}{c|c@{\qquad\quad}c|c@{\qquad\quad}c|c}
\zeta & \text{columns $k$} & \zeta & \text{columns $k$} & \zeta
 & \text{columns $k$}\\
\hline
7 & 9 & 17 & 67\text{--}74 & 27 & 163\text{--}174\\
8 & 10\text{--}14 & 18 & 75\text{--}83 & 28 & 175\text{--}185\\
9 & 15\text{--}20 & 19 & 84\text{--}91 & 29 & 186\text{--}197\\
10 & 21\text{--}25 & 20 & 92\text{--}101 & 30 & 198\text{--}210\\
11 & 26\text{--}31 & 21 & 102\text{--}110 & 31 & 211\text{--}222\\
12 & 32\text{--}37 & 22 & 111\text{--}120 & 32 & 223\text{--}235\\
13 & 38\text{--}44 & 23 & 121\text{--}130 & 33 & 236\text{--}248\\
14 & 45\text{--}51 & 24 & 131\text{--}141 & 34 & 249\text{--}261\\
15 & 52\text{--}58 & 25 & 142\text{--}151 & 35 & 262\text{--}264\\
16 & 59\text{--}66 & 26 & 152\text{--}162 & &
\end{array}
\]
\end{definition}

Like the numbers $n_0(k)$, the numbers $\zeta(k)$ are given by the
table and not by a formula. All that the later sections use about
$\zeta(k)$ is what the two lemmas below state and, for $82\le k\le264$, the
bound $100\,(\zeta(k)+1)\le27\,k$, which the proof of
Corollary~\ref{cor:earlypositive} reads from the table; any number for
which these hold could be listed instead of $\zeta(k)$.

The next lemma
says what the numbers mean: the number $\zeta(k)$
locates a positive step of column $k$ and a negative step of column $k+1$
in two consecutive rows. The negative step, which is part 3 of the lemma,
is the case $n=k+\zeta(k)$ of Lemma~\ref{lem:nextcolumn} below; what the
later sections use is the positive step of part 2.

\begin{lemma}\label{lem:signpair}
Let $k$ be an integer with $9\le k\le264$, and let $\zeta=\zeta(k)$ be the
number of Definition~\ref{def:zeta}. Then the following three statements
hold.
\begin{enumerate}
\item One has $\zeta+2\le k$.
\item One has $d_{k+\zeta+1,\,k}>0$.
\item One has $d_{k+\zeta,\,k+1}<0$.
\end{enumerate}
\end{lemma}

\begin{proof}
Part 1 is read off from the table of Definition~\ref{def:zeta}: in each of
its lines the value of $\zeta$ plus two is at most the first column of the
range.

Let $\ell\ge0$ and put $y=\ell+1$, so that $A_{n,\ell}=n!\,z^{(y)}_n$. By
Corollary~\ref{cor:normalised} the sequence $z^{(y)}$ is a solution of the
weighted last-car recursion with the weight $\rho(t)=\min(t,y)$, and this
weight is $t$ for $1\le t\le y$ and $y$ for $t\ge y+1$.
Proposition~\ref{prop:chainforward} therefore applies to $z^{(y)}$ and
gives, for every $n\le2y+1$,
\[
c_n-z^{(y)}_n=\sum_{t=y+1}^{n}c_{n-t}\Bigl(c_{t-1}-\frac yt\,z^{(y)}_{t-1}\Bigr).
\]
We multiply this identity by $n!$. The left-hand side becomes
$\bar A_{n,\ell}$. Since $m!\,c_m=\lvert\mathrm{PF}_m\rvert$ for every
$m\ge0$ and $(t-1)!\,z^{(y)}_{t-1}=A_{t-1,\ell}
=\lvert\mathrm{PF}_{t-1}\rvert-\bar A_{t-1,\ell}$, and since
$n!/((n-t)!\,(t-1)!)=t\binom nt$, a direct computation gives, for all
integers $\ell\ge0$ and $0\le n\le2y+1$, where $y=\ell+1$,
\begin{equation}\label{eq:tables-forward}
\bar A_{n,\ell}=\sum_{t=y+1}^{n}\binom nt\,\lvert\mathrm{PF}_{n-t}\rvert\,
\Bigl((t-y)\,\lvert\mathrm{PF}_{t-1}\rvert+y\,\bar A_{t-1,\ell}\Bigr).
\end{equation}
The right-hand side of \eqref{eq:tables-forward} contains the numbers
$\bar A_{m,\ell}$ with $y\le m\le n-1$ and no others.

By part 1, $k+\zeta+1\le2k-1$, so that rows $n=k+\zeta$ and
$n=k+\zeta+1$ satisfy $n\le2\ell+3$ for each of the four values
$\ell=k-2,k-1,k,k+1$, and \eqref{eq:tables-forward} holds in these rows and
in the rows before them. For these four values of $\ell$ we compute the
integers $\bar A_{m,\ell}$ with $\ell+1\le m\le k+\zeta+1$ from
\eqref{eq:tables-forward}, starting from $\bar A_{\ell+1,\ell}=0$. Their
rows are rows $m\le300$, so that each of them is smaller than
$\lvert\mathrm{PF}_{300}\rvert=301^{299}$, a number of $742$ decimal
digits. We find that for every $9\le k\le264$
\begin{align*}
\bar A_{k+\zeta+1,\,k}+\bar A_{k+\zeta+1,\,k-2}
&<2\,\bar A_{k+\zeta+1,\,k-1},\\
2\,\bar A_{k+\zeta,\,k}
&<\bar A_{k+\zeta,\,k+1}+\bar A_{k+\zeta,\,k-1}.
\end{align*}
These are two comparisons for each of the $256$ columns.
By \eqref{eq:tables-complement} for column $k$ in row $k+\zeta+1$,
the first inequality is part 2. By \eqref{eq:tables-complement} for column
$k+1$ in row $k+\zeta$, the second inequality is part 3.
\end{proof}

Since every value of $\zeta$ is at least $7$, part 1 of
Lemma~\ref{lem:signpair} places row
$k+\zeta(k)+1$ in the first layer of column $k$ and row $k+\zeta(k)$ in the
first layer of column $k+1$.
Proposition~\ref{prop:rulesmall} uses the positive row of part 2 together
with Corollary~\ref{cor:layerpersistenceall}. For
$82\le k\le264$ it is the positive row that
Corollary~\ref{cor:earlypositive} supplies, from which the arguments for
columns $k\ge82$ start.

The last lemma serves the diagonal rule in the first layer of columns
$9\le k\le264$ (Proposition~\ref{prop:rulesmall}): for a step of the first
layer of column $k$ in a row $n\le k+\zeta(k)$ that is not positive, the
rule asks for a negative step of column $k+1$ in row $n+1$, which
is one of the rows $k+2,\dots,k+\zeta(k)+1$. The sign is computed in each
of these rows: from the negative step in row $k+\zeta(k)$,
Corollary~\ref{cor:layerpersistenceall}, applied to column $k+1$, shows
only that the steps in the earlier rows of the first layer are not
positive.

\begin{lemma}\label{lem:nextcolumn}
Let $k$ be an integer with $9\le k\le264$, and let $\zeta(k)$ be the number
of Definition~\ref{def:zeta}. Then $d_{n,k+1}<0$ for every integer $n$ with
$k+2\le n\le k+\zeta(k)+1$.
\end{lemma}

\begin{proof}
By \eqref{eq:tables-complement} for column $k+1$, the inequality
$d_{n,k+1}<0$ is equivalent to
\[
2\,\bar A_{n,k}<\bar A_{n,k+1}+\bar A_{n,k-1}.
\]
The lemma asserts this in the $\zeta(k)$ rows $k+2\le n\le k+\zeta(k)+1$ of
column $k+1$, for every $9\le k\le264$. These are
\[
\sum_{k=9}^{264}\zeta(k)=5934
\]
positions, all of them in rows $n\le300$. Row $n=k+\zeta(k)$ is one of
these rows of column $k+1$, because $\zeta(k)\ge7$, and in this row the
claim is part 3 of Lemma~\ref{lem:signpair}. This settles $256$ of the
positions, one in each column, and $5678$ remain. By part 1 of
Lemma~\ref{lem:signpair}, every row $n\le k+\zeta(k)+1$ satisfies
$n\le2k-1$, hence $n\le2\ell+3$ for the three values $\ell=k-1,k,k+1$, so
that \eqref{eq:tables-forward}, which was shown in the proof of
Lemma~\ref{lem:signpair}, holds in the rows concerned. For these three
values of $\ell$ we compute the integers $\bar A_{m,\ell}$ with
$\ell+1\le m\le k+\zeta(k)+1$ from \eqref{eq:tables-forward}; as in the
proof of Lemma~\ref{lem:signpair}, they are smaller than $301^{299}$. We
find that the inequality holds at each of the $5678$ remaining positions.
\end{proof}

\section{Monotonicity along the diagonal and the diagonal lower bound}\label{sec:diagonal}
The arguments of Sections~\ref{sec:near}, \ref{sec:far}
and~\ref{sec:beyond} compare the entry of the array at a position with the
entry at its diagonal neighbour. The comparison that they use is
Proposition~\ref{prop:diagonallower}: for integers $k\ge2$ and $n\ge1$ the
normalised counts satisfy
$\gamma_k\,\beta^{(k-1)}_{n-1}\le\beta^{(k)}_n$ with the factor
$\gamma_k=(1+1/k)^{k-2}$. We derive it from a monotonicity along the
diagonal for the cumulative counts
(Proposition~\ref{prop:diagonalmonotone}).

\subsection{Monotonicity along the diagonal}\label{subsec:diagonal-monotone}

Let $y\ge1$ and $n\ge0$ be integers. Since there are $(n+1)^{n-1}$ parking
functions of length $n$, the quotient
\[
\pi^{(y)}_n=\frac{z^{(y)}_n}{c_n}=\frac{A_{n,y-1}}{(n+1)^{n-1}}
\]
is the probability that the maximum displacement of a parking
function of length $n$, chosen uniformly at random, is smaller than $y$. By
Section~\ref{sec:counts}, the number $z^{(y)}_n$ is nonnegative and does
not decrease in the threshold $y$, and $z^{(y)}_n=c_n$ in the free range
$n\le y$
(Proposition~\ref{prop:freerange}). Hence $\pi^{(y)}_n$ is nonnegative, it
does not decrease in $y$, and $\pi^{(y)}_n=1$ for $n\le y$.

The last-car recursion expresses these probabilities as averages. For
integers $0\le j\le n$ put
\[
\theta_n(j)=\frac{(j+1)\,c_j\,c_{n-j}}{(n+1)\,c_{n+1}},
\qquad
\psi^{(y)}_n(j)=\frac{\min(j+1,y)}{j+1}\,\pi^{(y)}_j\,\pi^{(y)}_{n-j}.
\]
With these numbers the recursion of Corollary~\ref{cor:normalised}, divided
by $(n+1)\,c_{n+1}$, reads
\begin{equation}\label{eq:diagonal-average}
\pi^{(y)}_{n+1}=\sum_{j=0}^{n}\theta_n(j)\,\psi^{(y)}_n(j).
\end{equation}
The numbers $\theta_n(0),\dots,\theta_n(n)$ are positive, and by
Lemma~\ref{lem:crecursion} they add up to one:
\begin{equation}\label{eq:diagonal-mass}
\sum_{j=0}^{n}\theta_n(j)=1.
\end{equation}
In terms of the parking process, the term of index $j$ of the last-car
recursion (Proposition~\ref{prop:lastcar}) belongs to the parking functions
in which the last car parks in spot $j+1$.

The following proposition states that the probabilities $\pi^{(y)}_n$ do
not decrease along the diagonals of their array,
$\pi^{(y)}_n\le\pi^{(y+1)}_{n+1}$. With $n$ replaced by $n+1$, the two sides
of this inequality are averages of the form \eqref{eq:diagonal-average} in
which both the weights and the values differ.
The proof compares the two averages through their weights and their values
separately: it shows that the weights
$\theta_{n+1}$ arise from the weights $\theta_n$ by moving, for every $j$,
a part of the mass at the point $j$ to the point $j+1$, and that the value
$\psi^{(y)}_n(j)$ is at most each of the two values $\psi^{(y+1)}_{n+1}(j)$
and $\psi^{(y+1)}_{n+1}(j+1)$ once the proposition is known for the
shorter lengths $j$ and $n-j$.

\begin{proposition}\label{prop:diagonalmonotone}
Let $y\ge1$ and $n\ge0$ be integers. Then the normalised counts of the
thresholds $y$ and $y+1$ satisfy
\[
\frac{z^{(y)}_n}{c_n}\le\frac{z^{(y+1)}_{n+1}}{c_{n+1}}.
\]
\end{proposition}

\begin{proof}
In the notation introduced above the claim is
$\pi^{(y)}_n\le\pi^{(y+1)}_{n+1}$, and we prove it for all $y\ge1$ by
induction on $n$. For $n=0$ both sides are equal to $1$, because the lengths
$0$ and $1$ lie in the free ranges of thresholds $y$ and $y+1$. Let
$n\ge0$, suppose that $\pi^{(y)}_i\le\pi^{(y+1)}_{i+1}$ holds for all
integers $0\le i\le n$ and $y\ge1$, and fix $y\ge1$. In the induction step we
write
\[
\theta=\theta_n,\qquad\theta^{+}=\theta_{n+1},\qquad
\psi=\psi^{(y)}_n,\qquad\psi^{+}=\psi^{(y+1)}_{n+1}.
\]
By \eqref{eq:diagonal-average}, the inequality
$\pi^{(y)}_{n+1}\le\pi^{(y+1)}_{n+2}$ that is to be proved reads
\begin{equation}\label{eq:diagonal-goal}
\sum_{j=0}^{n}\theta(j)\,\psi(j)\le\sum_{j=0}^{n+1}\theta^{+}(j)\,\psi^{+}(j).
\end{equation}

We first compare the values: for $0\le j\le n$,
\begin{equation}\label{eq:diagonal-values}
\psi(j)\le\psi^{+}(j)\qquad\text{and}\qquad\psi(j)\le\psi^{+}(j+1).
\end{equation}
By definition,
\begin{align*}
\psi(j)&=\frac{\min(j+1,y)}{j+1}\,\pi^{(y)}_j\,\pi^{(y)}_{n-j},\\
\psi^{+}(j)&=\frac{\min(j+1,y+1)}{j+1}\,\pi^{(y+1)}_j\,\pi^{(y+1)}_{n+1-j},\\
\psi^{+}(j+1)&=\frac{\min(j+2,y+1)}{j+2}\,\pi^{(y+1)}_{j+1}\,\pi^{(y+1)}_{n-j}.
\end{align*}
Each of these values is a product of three nonnegative factors, which we
compare one by one. For the first inequality of
\eqref{eq:diagonal-values}, the first factor does not decrease when $y$ is
replaced by $y+1$; then $\pi^{(y)}_j\le\pi^{(y+1)}_j$ by the monotonicity in
the threshold; and $\pi^{(y)}_{n-j}\le\pi^{(y+1)}_{n-j+1}$ by the induction
hypothesis for the length $n-j$. For the second inequality, the first factor
does not decrease because $\min(j+2,y+1)=\min(j+1,y)+1$ and
$\min(j+1,y)\le j+1$; then $\pi^{(y)}_j\le\pi^{(y+1)}_{j+1}$ by the induction
hypothesis for the length $j$; and $\pi^{(y)}_{n-j}\le\pi^{(y+1)}_{n-j}$ by
the monotonicity in the threshold.

Next we compare the weights: for $0\le j\le n$,
\begin{equation}\label{eq:diagonal-weights}
\theta^{+}(j)\le\theta(j)\qquad\text{and}\qquad\theta^{+}(j+1)\le\theta(j).
\end{equation}
Cancelling the factor $(j+1)\,c_j$ and clearing the denominators, we see
that the first inequality reads
\[
(n+1)\,c_{n+1-j}\,c_{n+1}\le(n+2)\,c_{n-j}\,c_{n+2}.
\]
It holds because $c_{n+1-j}\,c_{n+1}\le c_{n-j}\,c_{n+2}$; since
$c_i=a_{i+1}$, this is Lemma~\ref{lem:cayleyconvex} for the two indices
$n-j+1\le n+2$ and the shift $m=1$. Similarly, the second inequality reads
\[
\frac{(j+2)\,c_{j+1}}{(j+1)\,c_j}\le\frac{(n+2)\,c_{n+2}}{(n+1)\,c_{n+1}}.
\]
By \eqref{eq:prelim-quotient} we have $c_i=(1+1/i)^{i-1}\,c_{i-1}$ for
$i\ge1$. The left-hand side is therefore the first term, and the right-hand
side is the last term, of the chain
\begin{align*}
\Bigl(1+\frac1{j+1}\Bigr)^{j+1}
&\le\Bigl(1+\frac1{n+1}\Bigr)^{n+1}
=\frac{n+2}{n+1}\,\Bigl(1+\frac1{n+1}\Bigr)^{n}\\
&\le\frac{n+2}{n+1}\,\Bigl(1+\frac1{n+2}\Bigr)^{n+1}.
\end{align*}
Here the first step holds because $(1+1/i)^i$ does not decrease in $i\ge1$
(Section~\ref{sec:prelim}) and $j\le n$, and the last step is
Lemma~\ref{lem:powermonotone}.

It remains to derive \eqref{eq:diagonal-goal} from
\eqref{eq:diagonal-mass}, \eqref{eq:diagonal-values}
and~\eqref{eq:diagonal-weights}. Let
\[
\sigma_t=\sum_{j=0}^{t-1}\theta(j)\qquad\text{and}\qquad
\sigma^{+}_t=\sum_{j=0}^{t-1}\theta^{+}(j)
\]
be the partial sums of the two families of weights, the first for
$0\le t\le n+1$ and the second for $0\le t\le n+2$, so that
$\sigma_{n+1}=\sigma^{+}_{n+2}=1$ by \eqref{eq:diagonal-mass}. They interlace:
\begin{equation}\label{eq:diagonal-interlace}
\sigma^{+}_t\le\sigma_t\le\sigma^{+}_{t+1}\qquad(0\le t\le n+1).
\end{equation}
The first inequality follows term by term from the first inequality of
\eqref{eq:diagonal-weights}. The second inequality of
\eqref{eq:diagonal-weights} gives
\[
1-\sigma_t=\sum_{j=t}^{n}\theta(j)\ge\sum_{j=t}^{n}\theta^{+}(j+1)
=1-\sigma^{+}_{t+1},
\]
which is the second inequality of \eqref{eq:diagonal-interlace}. We now show
by induction on $t$ that, for $0\le t\le n+1$,
\begin{equation}\label{eq:diagonal-partial}
\sum_{j=0}^{t-1}\theta(j)\,\psi(j)\le
\sum_{j=0}^{t-1}\theta^{+}(j)\,\psi^{+}(j)
+\bigl(\sigma_t-\sigma^{+}_t\bigr)\,\psi^{+}(t).
\end{equation}
This is trivial for $t=0$. For the step from $t$ to $t+1\le n+1$ we
split the weight $\theta(t)=\sigma_{t+1}-\sigma_t$ into the two parts
$\sigma^{+}_{t+1}-\sigma_t$ and $\sigma_{t+1}-\sigma^{+}_{t+1}$, which are
nonnegative by \eqref{eq:diagonal-interlace}, and we use the first inequality
of \eqref{eq:diagonal-values} on the first part and the second inequality on
the second part:
\[
\theta(t)\,\psi(t)\le\bigl(\sigma^{+}_{t+1}-\sigma_t\bigr)\,\psi^{+}(t)
+\bigl(\sigma_{t+1}-\sigma^{+}_{t+1}\bigr)\,\psi^{+}(t+1).
\]
When this is added to \eqref{eq:diagonal-partial}, the terms with
$\psi^{+}(t)$ combine to
$(\sigma^{+}_{t+1}-\sigma^{+}_t)\,\psi^{+}(t)=\theta^{+}(t)\,\psi^{+}(t)$,
and the result is \eqref{eq:diagonal-partial} for $t+1$. For $t=n+1$ we have
$\sigma_{n+1}=\sigma^{+}_{n+2}$, and the last term of
\eqref{eq:diagonal-partial} is
\[
\bigl(\sigma^{+}_{n+2}-\sigma^{+}_{n+1}\bigr)\,\psi^{+}(n+1)
=\theta^{+}(n+1)\,\psi^{+}(n+1).
\]
So \eqref{eq:diagonal-partial} for $t=n+1$ is \eqref{eq:diagonal-goal},
which completes the induction step.
\end{proof}

\subsection{A lower bound along the diagonal}\label{subsec:diagonal-bound}

For an integer $k\ge2$ the \emph{diagonal factor} is
\[
\gamma_k=\Bigl(1+\frac1k\Bigr)^{k-2}.
\]
We have $\gamma_2=1$ and
\begin{equation}\label{eq:diagonal-factor}
\Bigl(1+\frac1k\Bigr)^{k-1}=\gamma_k\,\Bigl(1+\frac1k\Bigr),
\end{equation}
where the left-hand side is the quotient $a_{k+1}/a_k=c_k/c_{k-1}$ of
consecutive Cayley coefficients, by \eqref{eq:prelim-quotient}. The
diagonal factor does not decrease in $k$: for integers $2\le i\le j$,
Lemma~\ref{lem:powermonotone} and \eqref{eq:diagonal-factor} give
\[
\gamma_i\,\Bigl(1+\frac1i\Bigr)\le\gamma_j\,\Bigl(1+\frac1j\Bigr)
\le\gamma_j\,\Bigl(1+\frac1i\Bigr),
\]
and hence $\gamma_i\le\gamma_j$.

The following proposition is a statement about the
differences $\beta^{(k)}_n=z^{(k+1)}_n-z^{(k)}_n$ and
$\beta^{(k-1)}_{n-1}=z^{(k)}_{n-1}-z^{(k-1)}_{n-1}$ of the cumulative counts.
Its proof is an induction on $n$ that uses
Proposition~\ref{prop:diagonalmonotone} together with the
two comparisons of power series that are established in the proof of
Propositions~\ref{prop:onetree} and~\ref{prop:columninequality}.

\begin{proposition}\label{prop:diagonallower}
Let $k\ge2$ and $n\ge1$ be integers. Then the normalised counts of columns
$k-1$ and $k$ satisfy
\[
\gamma_k\,\beta^{(k-1)}_{n-1}\le\beta^{(k)}_n,
\]
where $\gamma_k$ is the diagonal factor.
\end{proposition}

\begin{proof}
Fix $k\ge2$, and let
\[
\mathfrak D=(Z_{k+1}-Z_k)-\gamma_k\,X\,(Z_k-Z_{k-1}),
\]
where $Z_y$ is the series of threshold $y$ (Section~\ref{sec:series}). We
write $\mathfrak D_n$ for the coefficient of $X^n$ in $\mathfrak D$. Since
$z^{(y)}_0=1$ for every threshold $y\ge1$ and
$\beta^{(\ell)}_n=z^{(\ell+1)}_n-z^{(\ell)}_n$ for $\ell\ge1$
(Section~\ref{sec:counts}), we have $\mathfrak D_0=0$ and
\[
\mathfrak D_n=\beta^{(k)}_n-\gamma_k\,\beta^{(k-1)}_{n-1}\qquad(n\ge1),
\]
so that the proposition asserts that $\mathfrak D_n\ge0$ for every $n\ge1$.

As in Section~\ref{sec:propagation}, for formal power series $S$ and $S^{*}$
with real coefficients we write $S\succeq S^{*}$ if every coefficient of
$S-S^{*}$ is nonnegative, and we put
\begin{gather*}
L_{+}=\Gamma_{k+1}-\Gamma_k,\qquad L_{-}=\Gamma_k-\Gamma_{k-1},\\
Y_{+}=(Z_{k+1}-Z_k)\,(1-K_k),\qquad Y_{-}=(Z_k-Z_{k-1})\,(1-K_k),
\end{gather*}
with the prime kernel series $K_k$ and the series $\Gamma_y$ of
Section~\ref{sec:series}. These are the four series that the proof of
Propositions~\ref{prop:onetree} and~\ref{prop:columninequality} compares.
That proof establishes the two comparisons
\begin{equation}\label{eq:diagonal-comparison}
Y_{+}\succeq L_{+}\qquad\text{and}\qquad L_{-}\succeq Y_{-}
\end{equation}
as relations between power series, with no hypothesis other than $k\ge2$.

By the definitions of $\mathfrak D$, $Y_{+}$ and $Y_{-}$ we have
$\mathfrak D\,(1-K_k)=Y_{+}-\gamma_k\,X\,Y_{-}$. Let
\[
\mathfrak L=L_{+}-\gamma_k\,X\,L_{-},
\]
and write $\mathfrak L_n$ for the coefficient of $X^n$ in $\mathfrak L$. Then
\[
\mathfrak D\,(1-K_k)-\mathfrak L
=(Y_{+}-L_{+})+\gamma_k\,X\,(L_{-}-Y_{-}),
\]
and every coefficient of the right-hand side is nonnegative, by
\eqref{eq:diagonal-comparison} and because $\gamma_k>0$. Hence
\[
\mathfrak D\,(1-K_k)\succeq\mathfrak L .
\]
The coefficients of $1-K_k$ are $1$ at $X^0$ and $-\kappa^{(k)}_m$ at $X^m$
for $m\ge1$ (Section~\ref{sec:series}). Read at the coefficient of $X^n$, the
last relation is therefore
\begin{equation}\label{eq:diagonal-renewal}
\mathfrak D_n\ge\sum_{m=1}^{n}\kappa^{(k)}_m\,\mathfrak D_{n-m}+\mathfrak L_n
\qquad(n\ge0).
\end{equation}

Next we show that
\begin{equation}\label{eq:diagonal-feedback}
(i+1)\,\mathfrak L_{i+1}\ge k\,\mathfrak D_i
\qquad\text{for every integer }i\ge k .
\end{equation}
By the definition of $\Gamma_y$, for $t\ge y$ the coefficient of $X^t$ in
$\Gamma_y$ is $y\,z^{(y)}_{t-1}/t$. Since $i\ge k$, this applies to the four
coefficients that occur in $\mathfrak L_{i+1}$, and
\begin{align*}
(i+1)\,\mathfrak L_{i+1}
&=(k+1)\,z^{(k+1)}_i-k\,z^{(k)}_i
-\gamma_k\,\frac{i+1}{i}\,\bigl(k\,z^{(k)}_{i-1}-(k-1)\,z^{(k-1)}_{i-1}\bigr),\\
k\,\mathfrak D_i
&=k\,\bigl(z^{(k+1)}_i-z^{(k)}_i\bigr)
-k\,\gamma_k\,\bigl(z^{(k)}_{i-1}-z^{(k-1)}_{i-1}\bigr).
\end{align*}
Subtracting and collecting the terms, we find
\begin{align*}
(i+1)\,\mathfrak L_{i+1}-k\,\mathfrak D_i
&=\Bigl(z^{(k+1)}_i-\gamma_k\,\frac{i+1}{i}\,z^{(k)}_{i-1}\Bigr)\\
&\quad+\gamma_k\,\frac{i+1-k}{i}\,\bigl(z^{(k)}_{i-1}-z^{(k-1)}_{i-1}\bigr).
\end{align*}
The second term on the right is nonnegative, because $i+1-k>0$ and
$z^{(k-1)}_{i-1}\le z^{(k)}_{i-1}$. By \eqref{eq:prelim-quotient} and
\eqref{eq:diagonal-factor} we have $c_n=\gamma_n\,(1+1/n)\,c_{n-1}$ for
$n\ge2$, so that Proposition~\ref{prop:diagonalmonotone}, applied with $n$
replaced by $n-1$, says that
\begin{equation}\label{eq:diagonal-cumulative}
\gamma_n\,\Bigl(1+\frac1n\Bigr)\,z^{(y)}_{n-1}\le z^{(y+1)}_n
\qquad(y\ge1,\ n\ge2).
\end{equation}
The first term is nonnegative by
\eqref{eq:diagonal-cumulative} with $y=k$ and $n=i$, because
$\gamma_i\ge\gamma_k$ and $z^{(k)}_{i-1}\ge0$:
\[
z^{(k+1)}_i\ge\gamma_i\,\frac{i+1}{i}\,z^{(k)}_{i-1}
\ge\gamma_k\,\frac{i+1}{i}\,z^{(k)}_{i-1}.
\]
This proves \eqref{eq:diagonal-feedback}.

We can now prove by induction on $n$ that $\mathfrak D_n\ge0$ for every
$n\ge0$. For $n\le k$ we have $\mathfrak D_n=0$: this was noted for $n=0$,
and for $1\le n\le k$ it holds because $z^{(k+1)}_n=z^{(k)}_n$ and
$z^{(k)}_{n-1}=z^{(k-1)}_{n-1}$ by Proposition~\ref{prop:freerange}. Let
$n=i+1$ with $i\ge k$, and suppose that $\mathfrak D_j\ge0$ for all
$j\le i$. On the right-hand side of \eqref{eq:diagonal-renewal} the sum is
nonnegative, because $\kappa^{(k)}_m\ge0$ by Lemma~\ref{lem:kernelsign} with
$x=1$ and $\mathfrak D_{i+1-m}\ge0$ by the induction hypothesis. Moreover
$\mathfrak L_{i+1}\ge k\,\mathfrak D_i/(i+1)\ge0$ by
\eqref{eq:diagonal-feedback} and the induction hypothesis. Hence
$\mathfrak D_{i+1}\ge0$.
\end{proof}

For $n\le k$ both sides of the inequality of
Proposition~\ref{prop:diagonallower} are zero, because
$\beta^{(\ell)}_m=0$ for $m\le\ell$ and $\ell\ge1$
(Section~\ref{sec:counts}).

Section~\ref{sec:beyond} needs Proposition~\ref{prop:diagonallower} at the
critical tilt. Recall from Section~\ref{sec:propagation} that the critical
tilt of a column $k\ge1$ is $x_k=(k/(k+1))^{k-1}$; it is positive and
satisfies $x_k\,(1+1/k)^{k-1}=1$. For a fixed column $k\ge2$, in terms of
the tilted normalised counts $u^{(y)}_N=z^{(y)}_N\,x_k^{\,N}$ of
Section~\ref{sec:propagation}, we define the
\emph{tilted column counts}
\[
U_N=u^{(k)}_N-u^{(k-1)}_N,\qquad U^{+}_N=u^{(k+1)}_N-u^{(k)}_N
\qquad(N\ge0).
\]
Since $\beta^{(\ell)}_N=z^{(\ell+1)}_N-z^{(\ell)}_N$ for $\ell\ge1$, they are
the normalised counts of columns $k-1$ and $k$, both tilted by the
critical tilt of column $k$:
\begin{equation}\label{eq:diagonal-columncounts}
U_N=\beta^{(k-1)}_N\,x_k^{\,N},\qquad U^{+}_N=\beta^{(k)}_N\,x_k^{\,N},
\end{equation}
and their difference $U^{+}_N-U_N=\delta^{(k)}_N\,x_k^{\,N}$ is the
normalised step of the column tilted by $x_k$, for which we write $\eta_N$.
Inserting
\eqref{eq:diagonal-factor} into the equation $x_k\,(1+1/k)^{k-1}=1$ and
multiplying by $k$, we obtain
\begin{equation}\label{eq:diagonal-rational}
(k+1)\,x_k\,\gamma_k=k .
\end{equation}

\begin{corollary}\label{cor:diagonaltilted}
Let $k\ge2$ and $n\ge1$ be integers. Then the tilted column counts $U_N$ and
$U^{+}_N$ of column $k$ satisfy
\[
k\,U_{n-1}\le(k+1)\,U^{+}_n .
\]
\end{corollary}

\begin{proof}
By \eqref{eq:diagonal-columncounts}, \eqref{eq:diagonal-rational} and
Proposition~\ref{prop:diagonallower},
\begin{align*}
k\,U_{n-1}=k\,x_k^{\,n-1}\,\beta^{(k-1)}_{n-1}
&=(k+1)\,x_k^{\,n}\,\gamma_k\,\beta^{(k-1)}_{n-1}\\
&\le(k+1)\,x_k^{\,n}\,\beta^{(k)}_n=(k+1)\,U^{+}_n. \qedhere
\end{align*}
\end{proof}

\section{Persistence on the near rows of the first layer}\label{sec:near}
This section and Section~\ref{sec:far} prove
Proposition~\ref{prop:staypositive}: in a column $k\ge82$, a step that has
turned positive early in the first layer stays positive up to the end of the
layer. The argument for a row $n$ depends on
the distance of $n$ from row $k$, and here we treat the near rows,
$n\le k+0.47k$, for which the hypothesis $k\ge42$ suffices
(Proposition~\ref{prop:nearrows}). The obstacle is what the rows before the
first positive row $q$ contribute to the column equation of row $n$
(Proposition~\ref{prop:columnequation}): the step is not positive in these
rows, and nothing is known about its size. We meet it
by comparing the equation of row $n$ with the equation of row $q$, in which
this contribution is known to be outweighed, because the step in row $q$ is
positive.

\subsection{The diagonal tilt}\label{subsec:near-tilt}

Let $k\ge2$ be an integer. In this subsection column $k$ is fixed, and we
write $\delta_n$ for its normalised step $\delta^{(k)}_n$. The
\emph{diagonal tilt} of column $k$ is the reciprocal of the diagonal factor
$\gamma_k$ of Section~\ref{sec:diagonal},
\[
\hat x_k=\frac{1}{\gamma_k}=\Bigl(\frac{k}{k+1}\Bigr)^{k-2},
\]
so that $\hat x_k>0$ and $\hat x_k\,\gamma_k=1$. The diagonal tilt is to be
distinguished from the critical tilt $x_k=(k/(k+1))^{k-1}$ of
Section~\ref{sec:propagation}. We tilt by $\hat x_k$
the three sequences that occur in the column equation of column $k$, namely
the free prime kernel of Section~\ref{sec:series}, the normalised counts of
column $k-1$ and the normalised step of column $k$: for
integers $m\ge0$, $j\ge0$ and $n\ge0$ let
\[
\hat\varepsilon_m=\hat x_k^{\,m}\,\varepsilon_m,\qquad
\hat u_j=\hat x_k^{\,j}\,\beta^{(k-1)}_j,\qquad
\hat\eta_n=\hat x_k^{\,n}\,\delta_n .
\]
We refer to $\hat\varepsilon$ as the \emph{kernel at the diagonal tilt}, to
$\hat u$ as the \emph{column count at the diagonal tilt} and to $\hat\eta$
as the \emph{step at the diagonal tilt}; the column is not shown in the
notation.

Since $\hat x_k>0$, the number $\hat\eta_n$ has the sign of $\delta_n$
and hence of the step $d_{n,k}$. Moreover, $\hat u_j\ge0$ for every $j\ge0$,
and $\hat\varepsilon_m>0$ for every $m\ge1$.

Multiplying the
equation of Proposition~\ref{prop:columnequation} by $\hat x_k^{\,n}$, we
obtain for $k+1\le n\le2k-1$
\begin{equation}\label{eq:near-tilted-column}
\hat\eta_n=\frac{2\hat x_k}{n}\,\hat u_{n-1}
+\frac{(k+1)\,\hat x_k}{n}\,\hat\eta_{n-1}
+\sum_{m=1}^{n-k}\hat\varepsilon_m\,\hat\eta_{n-m} .
\end{equation}
The first term on the right is the source, which comes from column
$k-1$ and is nonnegative; the second is a multiple of the value
$\hat\eta_{n-1}$ in the row before; and the sum is the feedback of rows
$k,\dots,n-1$, with the kernel $\hat\varepsilon$. In row $k$ itself we have
\begin{equation}\label{eq:near-bottom}
\hat\eta_k=-\hat u_k ,
\end{equation}
because $\delta_k=\beta^{(k)}_k-\beta^{(k-1)}_k$ and $b_{k,k}=0$.

The tilt has the value $1/\gamma_k$ so that the factor $\gamma_k$ disappears
from the lower bound of Proposition~\ref{prop:diagonallower}, by
which $\gamma_k\,\beta^{(k-1)}_{n-1}\le\beta^{(k)}_n$ for every integer
$n\ge1$. Multiplying this inequality by $\hat x_k^{\,n}$ and using
$\hat x_k\,\gamma_k=1$ and $\beta^{(k)}_n=\beta^{(k-1)}_n+\delta_n$, we
obtain
\begin{equation}\label{eq:near-tilted-diagonal}
\hat u_{n-1}-\hat u_n\le\hat\eta_n\qquad\text{for every integer }n\ge1 .
\end{equation}
This inequality bounds the step at the diagonal
tilt from below by the decrease of the column count at this tilt, and
together with \eqref{eq:near-bottom} this is how the size of the steps that
are not positive will be controlled.

\subsection{Three properties of the kernel at the diagonal tilt}\label{subsec:near-kernel}

The comparison of the equation \eqref{eq:near-tilted-column} at two rows
requires that the kernel $\hat\varepsilon$ is log-convex, that it does not
increase on the range of indices that occur, and that it satisfies a condition
which relates its decrease between the two rows to the decrease of the
coefficient $2\hat x_k/n$ of the source. All three properties are read off
from the \emph{step ratio} of the free prime kernel, which is defined for
integers $m\ge1$ by
\[
g_1=\frac12,\qquad
g_m=\frac{m}{m+1}\Bigl(1+\frac{1}{m-1}\Bigr)^{m-1}\quad(m\ge2).
\]
It is the ratio of consecutive values of the kernel,
\begin{equation}\label{eq:near-step-ratio}
\varepsilon_{m+1}=g_m\,\varepsilon_m\qquad\text{for every integer }m\ge1 ,
\end{equation}
as one checks directly from $\varepsilon_1=1$, $\varepsilon_2=1/2$ and
$\varepsilon_m=(m-1)^{m-1}/m!$. The step ratio does not decrease:
\begin{equation}\label{eq:near-ratio-monotone}
g_m\le g_{m+1}\qquad\text{for every integer }m\ge1 .
\end{equation}
Indeed, $g_1=1/2$ and $g_2=4/3$, and for $m\ge2$ neither factor of $g_m$
decreases when $m$ is replaced by $m+1$, the second by the monotonicity of
$(1+1/i)^i$ in $i\ge1$, which is recalled in Section~\ref{sec:prelim}.

The first property is the log-convexity of the kernel on indices $m\ge1$.

\begin{lemma}\label{lem:kernelconvex}
For every integer $m\ge1$ the free prime kernel satisfies
\[
\varepsilon_{m+1}^2\le\varepsilon_m\,\varepsilon_{m+2} .
\]
\end{lemma}

\begin{proof}
By \eqref{eq:near-step-ratio} and \eqref{eq:near-ratio-monotone},
\[
\varepsilon_{m+1}^2=g_m\,\varepsilon_m\,\varepsilon_{m+1}
\le g_{m+1}\,\varepsilon_m\,\varepsilon_{m+1}
=\varepsilon_m\,\varepsilon_{m+2}. \qedhere
\]
\end{proof}

The lemma compares neighbouring indices, and we need to compare two indices
at an arbitrary distance. Lemma~\ref{lem:crossform}, applied to the positive
sequence $(\varepsilon_{n+1})_{n\ge0}$, which satisfies its hypothesis by
Lemma~\ref{lem:kernelconvex}, gives
\begin{equation}\label{eq:near-cross}
\varepsilon_{m+t}\,\varepsilon_r\le\varepsilon_m\,\varepsilon_{r+t}
\qquad\text{for all integers }1\le m\le r\text{ and }t\ge0 .
\end{equation}
A tilt multiplies both sides of \eqref{eq:near-cross} by the same power of
the tilt, so the kernel $\hat\varepsilon$ at the diagonal tilt satisfies the
same inequality, whatever the column.

The other two properties depend on the tilt, and for them we have to know
its size, which the next lemma bounds from both sides.
Here and below we use the Taylor bounds of the logarithm recalled in
Section~\ref{sec:prelim},
\begin{equation}\label{eq:near-taylor}
u-\frac{u^2}{2}\le\log(1+u)\le u-\frac{u^2}{2}+\frac{u^3}{3}
\qquad(u\ge0).
\end{equation}

\begin{lemma}\label{lem:tiltbounds}
For every integer $k\ge2$ the diagonal tilt $\hat x_k$ of column $k$
satisfies
\[
e^{-1}\le\hat x_k\le e^{-1}\,e^{5/(2k)} .
\]
\end{lemma}

\begin{proof}
We have $1+\log\hat x_k=1-(k-2)\log(1+1/k)$, and $k-2\ge0$. The lower bound
of \eqref{eq:near-taylor} with $u=1/k$ gives
\[
1+\log\hat x_k\le1-(k-2)\Bigl(\frac1k-\frac{1}{2k^2}\Bigr)
=\frac{5}{2k}-\frac{1}{k^2}\le\frac{5}{2k} ,
\]
and the upper bound gives
\[
1+\log\hat x_k\ge1-(k-2)\Bigl(\frac1k-\frac{1}{2k^2}+\frac{1}{3k^3}\Bigr)
=\frac{5}{2k}-\frac{4}{3k^2}+\frac{2}{3k^3}\ge\frac{5}{2k}-\frac{4}{3k^2} .
\]
Since $5/(2k)-4/(3k^2)=(15k-8)/(6k^2)\ge0$, it follows that
$0\le1+\log\hat x_k\le5/(2k)$, which is equivalent to the claim.
\end{proof}

Since $\hat x_k\,\gamma_k=1$, the upper bound
of the lemma is equivalent to
\begin{equation}\label{eq:near-factor-lower}
\gamma_k\ge\exp\Bigl(1-\frac{5}{2k}\Bigr)
\qquad\text{for every integer }k\ge2 .
\end{equation}

For $m\ge1$ we have
$\hat\varepsilon_{m+1}=\hat x_k^{\,m+1}\,g_m\,\varepsilon_m$ by
\eqref{eq:near-step-ratio}, and
$\hat\varepsilon_m=\hat x_k^{\,m+1}\,\gamma_k\,\varepsilon_m$. By
Lemma~\ref{lem:kernelgrowth} below,
$\varepsilon_{m+1}\le\gamma_k\,\varepsilon_m$, and with it the second
property, for integers $k\ge2$ and $m\ge1$ with $5\,(m+1)\le3\,k$. Its proof
compares the step ratio $g_m$ with $\gamma_k$ by an upper bound for $g_m$
and by \eqref{eq:near-factor-lower}.

\begin{lemma}\label{lem:kernelgrowth}
Let $k\ge2$ and $m\ge1$ be integers with $5\,(m+1)\le3\,k$. Then the free
prime kernel $\varepsilon$ and the diagonal factor $\gamma_k=(1+1/k)^{k-2}$
satisfy
\[
\varepsilon_{m+1}\le\gamma_k\,\varepsilon_m .
\]
\end{lemma}

\begin{proof}
By \eqref{eq:near-step-ratio} it suffices to show that $g_m\le\gamma_k$.
For $m=1$ we have $g_1=1/2\le1\le\gamma_k$. For
$m\ge2$ we use the bound
\begin{equation}\label{eq:near-ratio-upper}
g_m\le\exp\Bigl(1-\frac{3}{2(m+1)}\Bigr)
\qquad\text{for every integer }m\ge2 ,
\end{equation}
which we prove first. Let $m\ge2$ be an integer and put $i=m-1$, so that
\[
\log g_m=i\,\log\Bigl(1+\frac1i\Bigr)-\log\Bigl(1+\frac{1}{i+1}\Bigr) .
\]
We bound the first logarithm from above and the second from below by
\eqref{eq:near-taylor}, with $u=1/i$ and with $u=1/(i+1)$, and obtain
\[
\log g_m\le1-\frac{1}{2i}+\frac{1}{3i^2}-\frac{1}{i+1}+\frac{1}{2(i+1)^2} .
\]
A straightforward computation shows that the difference between
$1-3/(2(i+2))$ and the right side is
\[
\frac{7i^3+4i^2-4i-4}{6\,i^2(i+1)^2(i+2)} ,
\]
which is positive for $i\ge1$. Since $i+2=m+1$, this proves
\eqref{eq:near-ratio-upper}. Now let $m\ge2$ be as in the lemma. The
hypothesis says that $5/(2k)\le3/(2(m+1))$, and
\eqref{eq:near-ratio-upper} and \eqref{eq:near-factor-lower} give
\[
g_m\le\exp\Bigl(1-\frac{3}{2(m+1)}\Bigr)\le\exp\Bigl(1-\frac{5}{2k}\Bigr)
\le\gamma_k . \qedhere
\]
\end{proof}

The third property concerns the kernel together with the coefficient of the
source. In the equation \eqref{eq:near-tilted-column} of a row $n=k+r$ the
source has the coefficient $2\hat x_k/n$, and row $k$, which lies at distance
$r$, has the weight $\hat\varepsilon_r$. The comparison of this
equation with the equation of an earlier row will require that
between the two rows the weight of row $k$ has fallen at least as much as
the coefficient of the source, and this holds if the product
$(k+m)\,\hat\varepsilon_m$ does not increase with $m$. The next
lemma shows that this is the case for indices $m\le0.47k$.

\begin{lemma}\label{lem:twopointcondition}
Let $k\ge42$ be an integer and let $i\le j$ be integers with $i\ge1$ and
$100\,j\le47\,k$. Then the free prime kernel at the diagonal tilt of column
$k$, $\hat\varepsilon_m=\hat x_k^{\,m}\,\varepsilon_m$, satisfies
\[
(k+j)\,\hat\varepsilon_j\le(k+i)\,\hat\varepsilon_i .
\]
\end{lemma}

\begin{proof}
It suffices to prove that
\begin{equation}\label{eq:near-one-step}
(k+m+1)\,\hat\varepsilon_{m+1}\le(k+m)\,\hat\varepsilon_m
\end{equation}
for every integer $m\ge1$ with $100\,(m+1)\le47\,k$; the lemma follows by
applying this for $m=i,\dots,j-1$, since $m+1\le j$ for these $m$. Since
$\hat\varepsilon_{m+1}=\hat x_k^{\,m+1}\,g_m\,\varepsilon_m$ and
$\hat\varepsilon_m=\hat x_k^{\,m+1}\,\gamma_k\,\varepsilon_m$, the inequality
\eqref{eq:near-one-step} follows from
\begin{equation}\label{eq:near-one-step-ratio}
(k+m+1)\,g_m\le(k+m)\,\gamma_k .
\end{equation}
For $m=1$ this holds because $g_1=1/2$ and $\gamma_k\ge1$.

Let $m\ge2$. By \eqref{eq:near-ratio-upper}, which was shown in the proof of
Lemma~\ref{lem:kernelgrowth}, and by
\eqref{eq:near-factor-lower}, the inequality \eqref{eq:near-one-step-ratio}
holds as soon as
\[
1+\frac{1}{k+m}\le\exp\Bigl(\frac{3}{2(m+1)}-\frac{5}{2k}\Bigr),
\]
and since $1+y\le e^y$ for every real $y$, this follows from
\[
\frac{1}{k+m}\le\frac{3}{2(m+1)}-\frac{5}{2k} .
\]
Bringing the three fractions to a common denominator, we find
\[
\frac{3}{2(m+1)}-\frac{5}{2k}-\frac{1}{k+m}
=\frac{3k^2-5m^2-(4k+5)\,m-7k}{2k\,(m+1)(k+m)} ,
\]
so that it remains to show that
\[
5m^2+(4k+5)\,m+7k\le3k^2 .
\]
The left side increases with $m$, and by hypothesis $m\le0.47k-1$. At
$m=0.47k-1$ the left side equals
\[
5\,(0.47k-1)^2+(4k+5)(0.47k-1)+7k=2.9845\,k^2+0.65\,k ,
\]
and $2.9845\,k^2+0.65\,k\le3k^2$ is equivalent to
$k\ge0.65/0.0155$, which holds because $k\ge42$.
\end{proof}

The rows beyond $k+0.47k$ are treated by a different argument in
Section~\ref{sec:far}.

\subsection{The comparison at two rows}\label{subsec:near-abstract}

We can now say how the equation \eqref{eq:near-tilted-column} is compared at
two rows. Let $q$ be the first positive row of a column $k\ge42$, and let
$n$ be a later row with $n\le k+0.47k$. In the equations of rows $q$ and $n$
the feedback of rows $k,\dots,q-1$ is not positive, and the question is
whether at row $n$ it can outweigh the source although it did not at row
$q$. Three facts work against this. By \eqref{eq:near-bottom} and
\eqref{eq:near-tilted-diagonal}, the proof of Proposition~\ref{prop:tworows}
bounds the sum of the values of $\hat\eta$ in these rows from below in terms
of the column count at the diagonal tilt. By the log-convexity of the kernel, in the form
\eqref{eq:near-cross},
the weight that one of these rows has in the equation of row $n$ is at
most $\hat\varepsilon_{n-k}/\hat\varepsilon_{q-k}$ times its weight in the
equation of row $q$. And by Lemma~\ref{lem:twopointcondition} the
coefficient of the source at row $n$ is at least
$\hat\varepsilon_{n-k}/\hat\varepsilon_{q-k}$ times its coefficient at row
$q$.

This comparison is a statement on arbitrary real sequences, and we prove it
as such. In the application of
Section~\ref{subsec:near-rows} the sequences $v$, $\pi$ and $\phi$ are the
step, the column count and the kernel at the diagonal tilt of a column $k$,
that is, $\hat\eta$, $\hat u$ and $\hat\varepsilon$; the numbers
$\sigma_q$, $\sigma_n$ and $\rho_q$, $\rho_n$ are the coefficients of the
source and of the preceding row in the equation
\eqref{eq:near-tilted-column} of rows $q$ and $n$. We keep the language of
the application and refer to the index $j$ of $v_j$ and of $\pi_j$ as a row.

\begin{proposition}\label{prop:tworows}
Let $(v_j)_{j\ge0}$, $(\pi_j)_{j\ge0}$ and $(\phi_m)_{m\ge0}$ be sequences of
real numbers, let $\sigma_q$, $\sigma_n$, $\rho_q$ and $\rho_n$ be real
numbers with $\sigma_q\ge0$, $\rho_q\ge0$ and $\rho_n>0$, and let $k\ge0$,
$r_q\ge1$ and $t\ge1$ be integers. Put $q=k+r_q$ and $n=q+t$, and assume the
following.
\begin{enumerate}
\item The sequence $v$ satisfies the two equations
\[
v_q=\sigma_q\,\pi_{q-1}+\rho_q\,v_{q-1}+\sum_{m=1}^{r_q}\phi_m\,v_{q-m}
\]
and
\[
v_n=\sigma_n\,\pi_{n-1}+\rho_n\,v_{n-1}+\sum_{m=1}^{r_q+t}\phi_m\,v_{n-m} .
\]
\item One has $v_k=-\pi_k$.
\item For every $j$ with $k+1\le j\le n-1$ one has $\pi_{j-1}-\pi_j\le v_j$.
\item For every $j\ge0$ one has $\pi_j\ge0$.
\item One has $\phi_m>0$ for $1\le m\le r_q+t$, and
$\phi_{m+1}\le\phi_m$ for $1\le m\le r_q+t-1$.
\item For every $m$ with $1\le m\le r_q$ one has
$\phi_{m+t}\,\phi_{r_q}\le\phi_{r_q+t}\,\phi_m$.
\item One has $\phi_{r_q+t}\,\sigma_q\le\phi_{r_q}\,\sigma_n$.
\item One has $v_j\le0$ for $k\le j<q$, and $v_j>0$ for $q\le j<n$.
\end{enumerate}
Then $v_n>0$.
\end{proposition}

\begin{proof}
Suppose, for a contradiction, that $v_n\le0$. We show that then
$\phi_t<\sigma_n<\phi_{t+1}$, which is impossible, because
$\phi_{t+1}\le\phi_t$ by hypothesis~(5).

In the equation of row $n$ we separate rows $q,\dots,n-1$, in which
$v$ is positive, from rows $k,\dots,q-1$, in which it is not. Let
\[
I=\sum_{m=1}^{t}\phi_m\,v_{n-m},\qquad
J_n=\sum_{m=1}^{r_q}\phi_{m+t}\,v_{q-m},\qquad
J_q=\sum_{m=1}^{r_q}\phi_m\,v_{q-m} .
\]
The sum $I$ is the feedback of rows $q,\dots,n-1$ at row $n$, and
$J_n$ and $J_q$ are the feedback of rows $k,\dots,q-1$ at row $n$
and at row $q$. Splitting the sum of the second equation of
hypothesis~(1) at $m=t$ and using $n-(m+t)=q-m$, we see that the two
equations read
\begin{equation}\label{eq:near-split}
\begin{split}
v_n&=\sigma_n\,\pi_{n-1}+\rho_n\,v_{n-1}+I+J_n ,\\
v_q&=\sigma_q\,\pi_{q-1}+\rho_q\,v_{q-1}+J_q .
\end{split}
\end{equation}
Let also
\[
\Sigma^{+}=\sum_{j=q}^{n-1}v_j,\qquad\Sigma^{-}=\sum_{j=k}^{q-1}v_j ,
\]
so that $\Sigma^{+}>0$ and $\Sigma^{-}\le0$ by hypothesis~(8).

By hypothesis~(5) the kernel is positive and does not increase on indices
$1,\dots,r_q+t$. In the sum $I$ we have $\phi_m\ge\phi_t$ and
$v_{n-m}>0$, and in the sum $J_n$ we have $\phi_{m+t}\le\phi_{t+1}$ and
$v_{q-m}\le0$. Hence
\begin{equation}\label{eq:near-feedback}
I\ge\phi_t\,\Sigma^{+},\qquad J_n\ge\phi_{t+1}\,\Sigma^{-} .
\end{equation}
Multiplying the inequality of hypothesis~(6) by $v_{q-m}\le0$ and summing
over $1\le m\le r_q$, we obtain
\begin{equation}\label{eq:near-transfer}
\phi_{r_q+t}\,J_q\le\phi_{r_q}\,J_n .
\end{equation}
Summing the inequality of hypothesis~(3) over rows $q\le j\le n-1$ gives
$\pi_{q-1}-\pi_{n-1}\le\Sigma^{+}$. Summing it over rows
$k+1\le j\le q-1$ gives $\pi_k-\pi_{q-1}\le\Sigma^{-}-v_k$, and by
hypothesis~(2) this is $-\pi_{q-1}\le\Sigma^{-}$. Hence
\begin{equation}\label{eq:near-profile}
\pi_{q-1}\le\pi_{n-1}+\Sigma^{+},\qquad
-(\Sigma^{+}+\Sigma^{-})\le\pi_{n-1} .
\end{equation}

We turn to the upper bound for $\sigma_n$. Since $v_n\le0$ and
$\rho_n\,v_{n-1}>0$, the first equation of \eqref{eq:near-split} gives
\begin{equation}\label{eq:near-debt}
\sigma_n\,\pi_{n-1}+I<-J_n .
\end{equation}
Hence
\[
\sigma_n\,\pi_{n-1}<-J_n-I
\le-\phi_{t+1}\,\Sigma^{-}-\phi_t\,\Sigma^{+}
\le-\phi_{t+1}\,(\Sigma^{+}+\Sigma^{-})
\le\phi_{t+1}\,\pi_{n-1} .
\]
Here the second step is \eqref{eq:near-feedback}, the third uses
$\phi_{t+1}\le\phi_t$ and $\Sigma^{+}>0$, and the fourth is
\eqref{eq:near-profile}, multiplied by $\phi_{t+1}>0$. If $\pi_{n-1}$ were
$0$, this would read $0<0$. Hence $\pi_{n-1}>0$ by hypothesis~(4), and
$\sigma_n<\phi_{t+1}$.

For the lower bound we start at row $q$. There $v_q>0$, $v_{q-1}\le0$
and $\rho_q\ge0$, so that the second equation of \eqref{eq:near-split} gives
$-J_q<\sigma_q\,\pi_{q-1}$. We transfer this inequality to row $n$ by the
chain
\begin{align*}
-\phi_{r_q}\,J_n&\le-\phi_{r_q+t}\,J_q<\phi_{r_q+t}\,\sigma_q\,\pi_{q-1}\\
&\le\phi_{r_q+t}\,\sigma_q\,(\pi_{n-1}+\Sigma^{+})
\le\phi_{r_q}\,\sigma_n\,(\pi_{n-1}+\Sigma^{+}) .
\end{align*}
Its first step is \eqref{eq:near-transfer}. The second is the inequality
just obtained, multiplied by $\phi_{r_q+t}>0$. The third is
\eqref{eq:near-profile}, with $\phi_{r_q+t}\,\sigma_q\ge0$, and the fourth
is hypothesis~(7), with $\pi_{n-1}+\Sigma^{+}\ge0$. Dividing by
$\phi_{r_q}>0$, we get $-J_n<\sigma_n\,(\pi_{n-1}+\Sigma^{+})$, and with
\eqref{eq:near-debt} this gives $I<\sigma_n\,\Sigma^{+}$. By
\eqref{eq:near-feedback} we have
$\phi_t\,\Sigma^{+}\le I<\sigma_n\,\Sigma^{+}$, and since $\Sigma^{+}>0$,
we conclude that $\phi_t<\sigma_n$. This completes the proof.
\end{proof}

\subsection{The near rows}\label{subsec:near-rows}

We now apply Proposition~\ref{prop:tworows} to the step of a column at its
diagonal tilt and obtain the result of this section.

\begin{proposition}\label{prop:nearrows}
Let $k\ge42$, $r_q\ge1$ and $t\ge1$ be integers with
$100\,(r_q+t)\le47\,k$. Assume that $d_{j,k}\le0$ for every $j$ with
$k\le j<k+r_q$, and that $d_{j,k}>0$ for every $j$ with
$k+r_q\le j<k+r_q+t$. Then
\[
d_{k+r_q+t,\,k}>0 .
\]
\end{proposition}

\begin{proof}
Put $q=k+r_q$ and $n=q+t$, and let $\hat\varepsilon$, $\hat u$ and
$\hat\eta$ be the kernel, the column count and the step at the diagonal
tilt $\hat x_k$ of column $k$, as in Section~\ref{subsec:near-tilt}. Since
$\hat\eta_j$ has the sign of $d_{j,k}$, it suffices to show that
$\hat\eta_n>0$. We obtain
this from Proposition~\ref{prop:tworows}, applied to the integers $k$, $r_q$
and $t$, to the sequences
\[
v_j=\hat\eta_j,\qquad\pi_j=\hat u_j,\qquad\phi_m=\hat\varepsilon_m ,
\]
and to the numbers $\sigma_j=2\hat x_k/j$ and $\rho_j=(k+1)\,\hat x_k/j$
for $j=q$ and $j=n$, which are positive. It remains to verify the eight
hypotheses of that proposition.

Hypothesis~(1) holds because $100\,(r_q+t)\le47\,k$ gives
$n\le k+0.47k\le2k-1$. Hence $k+1\le q<n\le2k-1$, and the two equations are
\eqref{eq:near-tilted-column} for rows $q$ and $n$, in which the sum runs up
to $q-k=r_q$ and up to $n-k=r_q+t$. Hypotheses~(2), (3) and~(4) are
\eqref{eq:near-bottom}, the inequality \eqref{eq:near-tilted-diagonal},
which holds for every $j\ge1$, and the nonnegativity of $\hat u_j$.

For hypothesis~(5), the numbers $\hat\varepsilon_m$ with $m\ge1$ are
positive. For $1\le m\le r_q+t-1$ we have $m+1\le0.47k$, hence
$5\,(m+1)\le3\,k$, and the inequality of Lemma~\ref{lem:kernelgrowth},
multiplied by $\hat x_k^{\,m+1}$, is
$\hat\varepsilon_{m+1}\le\hat\varepsilon_m$, because $\hat x_k\,\gamma_k=1$.
Hypothesis~(6) is \eqref{eq:near-cross} with $r=r_q$, multiplied by
$\hat x_k^{\,m+t+r_q}$.

After multiplication by $qn/(2\hat x_k)$, hypothesis~(7) reads
\[
(k+r_q+t)\,\hat\varepsilon_{r_q+t}\le(k+r_q)\,\hat\varepsilon_{r_q} .
\]
This is Lemma~\ref{lem:twopointcondition} with $i=r_q$ and $j=r_q+t$, whose
hypothesis $100\,j\le47\,k$ is the one of the proposition. Finally, the two
conditions of hypothesis~(8) are the two assumptions on the steps $d_{j,k}$,
because $\hat\eta_j$ has the sign of $d_{j,k}$.
\end{proof}

\section{Renewal sequences and the frozen kernel}\label{sec:renewal}
The far rows of the first layer, which are treated in
Section~\ref{sec:far}, require a bound for a weighted sum of terms of a
sequence $v$ against the far part of a kernel, in terms of the last value of
$v$ alone, when all that is known about $v$ is that it dominates its
convolution with the kernel. We obtain such a bound by comparing $v$ with
the renewal sequence of the kernel (Proposition~\ref{prop:extremecolumn});
the comparison rests on the log-convexity of that sequence, a theorem of
de Bruijn and Erd\H{o}s (Proposition~\ref{prop:renewalconvex}). We then
treat the kernel to which the comparison is applied, the frozen kernel, and
its renewal sequence.

\subsection{Renewal sequences and log-convex kernels}\label{subsec:renewal-general}

\begin{definition}\label{def:renewal}
Let $\varphi=(\varphi_i)_{i\ge1}$ and $u=(u_n)_{n\ge0}$ be sequences of real
numbers. The sequence $u$ is \emph{the renewal sequence of the kernel
$\varphi$} if
\[
u_0=1\qquad\text{and}\qquad u_n=\sum_{i=1}^{n}\varphi_i\,u_{n-i}
\quad\text{for every integer } n\ge1 .
\]
\end{definition}

The second condition, for a given $n$, is the \emph{renewal equation} at
$n$. For a given kernel the two conditions determine $u$ term by term:
$u_1=\varphi_1$, $u_2=\varphi_1^2+\varphi_2$, and so on. In terms of formal
power series they say that
\[
\sum_{n\ge0}u_nX^n=\Bigl(1-\sum_{i\ge1}\varphi_iX^i\Bigr)^{-1}.
\]
When the numbers $\varphi_i$ are the probabilities of the lengths of the
steps of a renewal process, $u_n$ is the probability that a renewal occurs at
time $n$. We make no assumption of this
kind, and the sum of the $\varphi_i$ need not be one. Two instances have
occurred already: by Corollary~\ref{cor:crenewal} the sequence $(c_n)_{n\ge0}$ is the
renewal sequence of the kernel $(\varepsilon_m)_{m\ge1}$, and by
Lemma~\ref{lem:nurenewal} the sequence $(\nu_n)_{n\ge0}$ is the renewal
sequence of the kernel $(a_j)_{j\ge1}$; recall that $c_0=\nu_0=1$.

Let $u$ be the renewal sequence of a kernel $\varphi$. If $\varphi_i>0$ for
every $i\ge1$, then $u_n>0$ for every $n\ge0$; this follows from the renewal
equation by induction on $n$.

We think of the index of a sequence as a time. Let $\varphi$ be a kernel and
$v=(v_s)_{s\ge0}$ a sequence of real numbers. The sum
$\varphi_1v_{s-1}+\varphi_2v_{s-2}+\dots+\varphi_sv_0$ is the \emph{feedback}
that $v$ receives at time $s$ from its earlier values through the kernel,
and the difference $\iota_s$ between $v_s$ and this feedback is the
\emph{innovation} of $v$ at time $s$ with respect to $\varphi$. The
renewal sequence of $\varphi$ is the sequence whose innovations are
$\iota_0=1$ and $\iota_s=0$ for $s\ge1$. The following lemma recovers an
arbitrary sequence from its innovations: each innovation is carried to the
later times by the renewal sequence, which is in this sense the resolvent of
the kernel. We need it in the proof of Proposition~\ref{prop:extremecolumn},
where it writes a sequence with nonnegative innovations as a combination of
shifted copies of the renewal sequence with nonnegative coefficients.

\begin{lemma}\label{lem:convolutionform}
Let $\varphi=(\varphi_i)_{i\ge1}$, $u=(u_n)_{n\ge0}$, $v=(v_s)_{s\ge0}$ and
$\iota=(\iota_s)_{s\ge0}$ be sequences of real numbers. Suppose that $u$ is
the renewal sequence of the kernel $\varphi$ in the sense of
Definition~\ref{def:renewal}, and that
\[
v_s=\iota_s+\sum_{i=1}^{s}\varphi_i\,v_{s-i}
\qquad\text{for every integer } s\ge0 .
\]
Then
\[
v_s=\sum_{j=0}^{s}u_{s-j}\,\iota_j
\qquad\text{for every integer } s\ge0 .
\]
\end{lemma}

\begin{proof}
We argue by induction on $s$ and assume the formula for all indices smaller
than $s$. Inserting it for each $v_{s-i}$ with $1\le i\le s$ and exchanging
the order of summation, we obtain
\[
\sum_{i=1}^{s}\varphi_i\,v_{s-i}
=\sum_{j=0}^{s-1}\Bigl(\sum_{i=1}^{s-j}\varphi_i\,u_{s-j-i}\Bigr)\iota_j
=\sum_{j=0}^{s-1}u_{s-j}\,\iota_j ,
\]
where the inner sum is $u_{s-j}$ by the renewal equation at $s-j\ge1$. Adding
$\iota_s=u_0\,\iota_s$ gives the formula at $s$.
\end{proof}

The next statement transfers log-convexity from a kernel to its renewal
sequence. It is a theorem of de Bruijn and Erd\H{o}s
\cite[Theorem~1]{debruijn1953}. There it is stated for a kernel with
nonnegative terms whose sum is one, with strict inequalities in the
hypothesis and in the conclusion, and the remarks that follow the theorem
give the form with non-strict inequalities. We need the theorem for a kernel
with positive terms and with no condition on its sum, in the non-strict
form, and we prove it in this form.

\begin{proposition}\label{prop:renewalconvex}
Let $\varphi=(\varphi_i)_{i\ge1}$ and $u=(u_n)_{n\ge0}$ be sequences of real
numbers such that $u$ is the renewal sequence of the kernel $\varphi$ in the
sense of Definition~\ref{def:renewal}. Suppose that $\varphi_i>0$ for every
$i\ge1$ and that $\varphi_{i+1}^2\le\varphi_i\,\varphi_{i+2}$ for every
$i\ge1$. Then
\[
u_{n+1}^2\le u_n\,u_{n+2}\qquad\text{for every integer } n\ge0 .
\]
\end{proposition}

\begin{proof}
All $u_n$ are positive, as was noted after Definition~\ref{def:renewal}. We
prove the inequality by induction on $n$. For $n=0$ the renewal equations
give $u_1=\varphi_1$ and $u_2=\varphi_1u_1+\varphi_2=\varphi_1^2+\varphi_2$,
so that $u_1^2=\varphi_1^2\le\varphi_1^2+\varphi_2=u_0\,u_2$.

Let $n\ge1$, and suppose that $u_{j+1}^2\le u_j\,u_{j+2}$ for
$0\le j\le n-1$. Then the quotients $u_{j+1}/u_j$ do not decrease as $j$ runs
from $0$ to $n$, and by the hypothesis on the kernel the quotients
$\varphi_{i+2}/\varphi_{i+1}$ do not decrease in $i\ge0$. Consequently, for
all integers $i$ and $j$ with $0\le i\le n$ and $0\le j\le n$,
\begin{equation}\label{eq:renewal-ratios}
\varphi_{i+2}\,\varphi_{n+1}\le\varphi_{i+1}\,\varphi_{n+2}
\qquad\text{and}\qquad
u_{j+1}\,u_n\le u_j\,u_{n+1} .
\end{equation}
Consider the sum over the pairs $(i,j)$ of nonnegative integers with
$i+j=n-1$,
\[
S=\sum_{i+j=n-1}\bigl(\varphi_{n+2}\,\varphi_{i+1}
-\varphi_{n+1}\,\varphi_{i+2}\bigr)
\bigl(u_{n+1}\,u_j-u_n\,u_{j+1}\bigr).
\]
By \eqref{eq:renewal-ratios} both factors of every term are nonnegative, so
that $S\ge0$. On the other hand, $S$ can be evaluated by the renewal
equations at $n$, $n+1$ and $n+2$. The first of them, the second without its
last term $\varphi_{n+1}u_0$, the second without its first term
$\varphi_1u_n$, and the third without its first term $\varphi_1u_{n+1}$ and
its last term $\varphi_{n+2}u_0$ read
\begin{align*}
\sum_{i+j=n-1}\varphi_{i+1}\,u_j&=u_n,\\
\sum_{i+j=n-1}\varphi_{i+1}\,u_{j+1}&=u_{n+1}-\varphi_{n+1},\\
\sum_{i+j=n-1}\varphi_{i+2}\,u_j&=u_{n+1}-\varphi_1\,u_n,\\
\sum_{i+j=n-1}\varphi_{i+2}\,u_{j+1}&=u_{n+2}-\varphi_1\,u_{n+1}-\varphi_{n+2}.
\end{align*}
We expand the product in $S$ and insert these four sums. All terms cancel
except two, and we find
\[
S=\varphi_{n+1}\,\bigl(u_n\,u_{n+2}-u_{n+1}^2\bigr).
\]
Since $S\ge0$ and $\varphi_{n+1}>0$, it follows that
$u_{n+1}^2\le u_n\,u_{n+2}$.
\end{proof}

We come to the comparison with the renewal sequence. Let $u$ be the
renewal sequence of a positive log-convex kernel $\varphi$, and let $v$ be a
sequence whose innovations with respect to $\varphi$ are nonnegative up to a
time $\rho$. Given an integer $M\ge1$ and weights $\theta_m$, the sum
\[
\sum_{m=M}^{\rho}\theta_m\,\varphi_m\,v_{\rho-m}
\]
is the part of the feedback of $v$ at time $\rho$ that passes through the
kernel indices $m\ge M$, taken with the weights $\theta_m$; this is the
\emph{far feedback} of $v$ at time $\rho$.
Proposition~\ref{prop:extremecolumn} below bounds the far feedback of $v$,
relative to the value $v_\rho$, by the far feedback of the renewal sequence
$u$ relative to $u_\rho$, for weights that are nonnegative and nondecreasing
from index $M$ to index $\rho$.

The reason for the bound is the following. By
Lemma~\ref{lem:convolutionform}, $v$ is a combination of shifted copies of
$u$, one for each innovation, with nonnegative coefficients. An innovation
injected at time $\rho-s$ contributes $u_s$ times its size to $v_\rho$ and
\[
\sum_{m=M}^{s}\theta_m\,\varphi_m\,u_{s-m}
\]
times its size to the far feedback. It is therefore enough to see that the
quotient of this sum by $u_s$ is largest for $s=\rho$. By the
renewal equation at $s$, the numbers $\varphi_m\,u_{s-m}/u_s$ with
$1\le m\le s$ are positive and add up to one; in terms of a renewal process,
the $m$-th of them is the share of $u_s$ that comes from the histories whose
last step has length $m$. For a log-convex sequence $u$ every quotient
$u_{s-m}/u_s$ does not increase when $s$ grows. Hence the share of the last
steps shorter than $M$ does not increase with $s$, and the share of
the last steps of length at least $M$ does not decrease. Since the weights
are nondecreasing, a shift of mass towards longer last steps can only
increase the weighted share.

\begin{proposition}\label{prop:extremecolumn}
Let $\varphi=(\varphi_i)_{i\ge1}$, $u=(u_n)_{n\ge0}$, $v=(v_n)_{n\ge0}$ and
$\theta=(\theta_m)_{m\ge1}$ be sequences of real numbers, and let $\rho\ge0$
and $M\ge1$ be integers. Assume the following.
\begin{enumerate}
\item One has $\varphi_i>0$ for every $i\ge1$ and
$\varphi_{i+1}^2\le\varphi_i\,\varphi_{i+2}$ for every $i\ge1$.
\item The sequence $u$ is the renewal sequence of the kernel $\varphi$ in the
sense of Definition~\ref{def:renewal}.
\item One has $\theta_M\ge0$, and $\theta_a\le\theta_b$ for all integers $a$
and $b$ with $M\le a\le b\le\rho$.
\item For every integer $s$ with $0\le s\le\rho$,
\[
\sum_{i=1}^{s}\varphi_i\,v_{s-i}\le v_s .
\]
\end{enumerate}
Then
\[
\Bigl(\sum_{m=M}^{\rho}\theta_m\,\varphi_m\,v_{\rho-m}\Bigr)\,u_\rho
\le
\Bigl(\sum_{m=M}^{\rho}\theta_m\,\varphi_m\,u_{\rho-m}\Bigr)\,v_\rho .
\]
\end{proposition}

\begin{proof}
If $M>\rho$, both sums are empty and there is nothing to prove, so let
$M\le\rho$. By hypotheses~(1) and~(2), $u_n>0$ for every $n\ge0$, as was
noted after Definition~\ref{def:renewal}, and $u_{n+1}^2\le u_n\,u_{n+2}$ for
every $n\ge0$ by Proposition~\ref{prop:renewalconvex}. Hence
Lemma~\ref{lem:crossform}, applied to the sequence $u$ with the indices
$s-m\le\rho-m$ and the shift $m$, gives
\begin{equation}\label{eq:renewal-cross}
u_s\,u_{\rho-m}\le u_{s-m}\,u_\rho\qquad(0\le m\le s\le\rho).
\end{equation}

We first compare the tails of the renewal equations at two times: for every
integer $s$ with $M\le s\le\rho$,
\begin{equation}\label{eq:renewal-tails}
\Bigl(\sum_{m=M}^{s}\varphi_m\,u_{s-m}\Bigr)\,u_\rho
\le\Bigl(\sum_{m=M}^{\rho}\varphi_m\,u_{\rho-m}\Bigr)\,u_s .
\end{equation}
Indeed, by the renewal equations at $s$ and at $\rho$ the two sums in
parentheses are
\[
u_s-\sum_{m=1}^{M-1}\varphi_m\,u_{s-m}
\qquad\text{and}\qquad
u_\rho-\sum_{m=1}^{M-1}\varphi_m\,u_{\rho-m} .
\]
After multiplication by $u_\rho$ and by $u_s$ the leading terms agree,
and \eqref{eq:renewal-tails} becomes
\[
\sum_{m=1}^{M-1}\varphi_m\,u_{\rho-m}\,u_s
\le\sum_{m=1}^{M-1}\varphi_m\,u_{s-m}\,u_\rho ,
\]
which holds term by term by \eqref{eq:renewal-cross}.

Next we bring in the weights. Put
\[
W=\sum_{m=M}^{\rho}\theta_m\,\varphi_m\,u_{\rho-m} .
\]
By hypothesis~(3), $\theta_m\ge\theta_M\ge0$ for $M\le m\le\rho$, so that
$W\ge0$. We claim that for every integer $s$ with $M\le s\le\rho$,
\begin{equation}\label{eq:renewal-weighted}
\Bigl(\sum_{m=M}^{s}\theta_m\,\varphi_m\,u_{s-m}\Bigr)\,u_\rho\le W\,u_s .
\end{equation}
We split $W$ into its terms with $m\le s$ and its terms with $m>s$. Then
\eqref{eq:renewal-weighted} is equivalent to
\begin{equation}\label{eq:renewal-split}
\sum_{m=M}^{s}\theta_m\,\varphi_m\,
\bigl(u_{s-m}\,u_\rho-u_s\,u_{\rho-m}\bigr)
\le u_s\sum_{m=s+1}^{\rho}\theta_m\,\varphi_m\,u_{\rho-m} .
\end{equation}
On the left side every difference in parentheses is nonnegative by
\eqref{eq:renewal-cross}, and $\theta_m\le\theta_s$ for $M\le m\le s$ by
hypothesis~(3). The left side is therefore at most
\[
\theta_s\sum_{m=M}^{s}\varphi_m\,
\bigl(u_{s-m}\,u_\rho-u_s\,u_{\rho-m}\bigr)
\le\theta_s\,u_s\sum_{m=s+1}^{\rho}\varphi_m\,u_{\rho-m} ,
\]
where the inequality is \eqref{eq:renewal-tails}, with the terms $m\le s$ of
its right side moved to the left, multiplied by $\theta_s\ge0$. On the right
side of \eqref{eq:renewal-split} we have $\theta_m\ge\theta_s$ for
$s<m\le\rho$, again by hypothesis~(3), so the last expression is at most the
right side of \eqref{eq:renewal-split}. This proves
\eqref{eq:renewal-weighted}.

Finally we decompose $v$ into its innovations. For $s\ge0$ let
\[
\iota_s=v_s-\sum_{i=1}^{s}\varphi_i\,v_{s-i} ,
\]
so that $\iota_s\ge0$ for $s\le\rho$ by hypothesis~(4). By
Lemma~\ref{lem:convolutionform},
\[
v_s=\sum_{j=0}^{s}u_{s-j}\,\iota_j\qquad\text{for every } s\ge0 .
\]
We insert this for $v_{\rho-m}$ and exchange the two sums:
\[
\sum_{m=M}^{\rho}\theta_m\,\varphi_m\,v_{\rho-m}
=\sum_{j=0}^{\rho-M}\iota_j
\sum_{m=M}^{\rho-j}\theta_m\,\varphi_m\,u_{(\rho-j)-m} .
\]
We multiply by $u_\rho$ and apply \eqref{eq:renewal-weighted} with
$s=\rho-j$ to each inner sum. Since $\iota_j\ge0$, this gives
\[
\Bigl(\sum_{m=M}^{\rho}\theta_m\,\varphi_m\,v_{\rho-m}\Bigr)\,u_\rho
\le W\sum_{j=0}^{\rho-M}u_{\rho-j}\,\iota_j
\le W\sum_{j=0}^{\rho}u_{\rho-j}\,\iota_j
=W\,v_\rho .
\]
The second inequality adds the terms with $\rho-M<j\le\rho$, which are
nonnegative because $W\ge0$, $u_{\rho-j}>0$ and $\iota_j\ge0$, and the last
equality is Lemma~\ref{lem:convolutionform} at $s=\rho$.
\end{proof}

\subsection{The frozen kernel and its renewal sequence}\label{subsec:renewal-frozen}

We now introduce the kernel to which the statements of
Section~\ref{subsec:renewal-general} are applied. Let $\varepsilon$ be the
free prime kernel of Section~\ref{sec:series}. The \emph{frozen kernel} is
the sequence $f=(f_i)_{i\ge1}$ with
\begin{equation}\label{eq:renewal-frozen}
f_1=\frac32\,e^{-1},\qquad f_i=e^{-i}\,\varepsilon_i\quad\text{for } i\ge2 .
\end{equation}
Since $\varepsilon_1=1$, it is the sequence $(e^{-i}\varepsilon_i)_{i\ge1}$
with $1/(2e)$ added to its first term. We also put $f_0=0$.

The kernel is made for one purpose, which explains its form; the argument
itself is part of the proof of Proposition~\ref{prop:farrows}. Let $k\ge2$.
By Lemma~\ref{lem:levelequation} the normalised counts of column $k-1$
satisfy, in rows $k\le n\le2k-1$,
\[
\beta^{(k-1)}_n=\frac{z^{(k)}_{n-1}}{n}
+\frac{k-1}{n}\,\beta^{(k-1)}_{n-1}
+\sum_{m=1}^{n-k}\varepsilon_m\,\beta^{(k-1)}_{n-m} ,
\]
and all terms on the right are nonnegative. The coefficient $(k-1)/n$ of the
preceding row depends on the row, and it is at least $1/2$ as long as
$n\le2k-2$. If this coefficient is replaced by the constant $1/2$ and the
first term is dropped, the equation becomes an inequality.
For the numbers $v_s=e^{-s}\beta^{(k-1)}_{k+s}$ this inequality is
\[
\sum_{i=1}^{s}f_i\,v_{s-i}\le v_s\qquad(1\le s\le k-2),
\]
where the summand $1/(2e)$ of $f_1$ comes from the constant $1/2$. The
sequence $v$ thus has nonnegative innovations with respect to a kernel that
depends neither on the column nor on the row, as hypothesis~(4) of
Proposition~\ref{prop:extremecolumn} requires. The kernel is called frozen
because the coefficient $(k-1)/n$ of the preceding row, which varies with
the row, has been frozen at one value.

The factors $e^{-i}$ tilt the kernel by $1/e$, and this tilt is chosen with
a view to the estimates by which Section~\ref{sec:far} evaluates the
comparison with the renewal sequence. The inequality itself could be written
at any tilt $x>0$, for the numbers $x^s\beta^{(k-1)}_{k+s}$: it suffices to
multiply it by $(e\,x)^s$. The terms of the kernel are then those of $f$
multiplied by the geometric factors $(e\,x)^i$. Section~\ref{sec:far} bounds
the kernel terms with their weights uniformly in $m$
(Lemma~\ref{lem:farweights}) and the terms of the renewal sequence
separately, through a partial sum of the sequence and a lower bound for a
single term.

At the tilt $1/e$ the sum of the kernel $\varepsilon$ is $1-1/e$
(Corollary~\ref{cor:kernelmass}), and $f$ adds $1/(2e)$ to its first term.
The partial sums of the renewal sequence of $f$ are bounded
(Lemma~\ref{lem:resolventmass}), while its terms decay no faster than
$(s+1)^{-3/2}$ (Lemma~\ref{lem:resolventlower}). These are the two bounds
for the renewal sequence that Section~\ref{sec:far} uses.

The frozen kernel is positive: $f_i>0$ for every $i\ge1$, because
$\varepsilon_i>0$ for $i\ge1$. It is also log-convex; the next lemma derives
this from the log-convexity of $\varepsilon$
(Lemma~\ref{lem:kernelconvex}).

\begin{lemma}\label{lem:frozenconvex}
For every integer $i\ge1$, the frozen kernel $f$ of
\eqref{eq:renewal-frozen} satisfies
\[
f_{i+1}^2\le f_i\,f_{i+2} .
\]
\end{lemma}

\begin{proof}
Since $i+1\ge2$, we have $f_{i+1}=e^{-(i+1)}\varepsilon_{i+1}$ and
$f_{i+2}=e^{-(i+2)}\varepsilon_{i+2}$. By Lemma~\ref{lem:kernelconvex},
$\varepsilon_{i+1}^2\le\varepsilon_i\,\varepsilon_{i+2}$, and therefore
\[
f_{i+1}^2=e^{-i}\,e^{-(i+2)}\,\varepsilon_{i+1}^2
\le e^{-i}\varepsilon_i\cdot e^{-(i+2)}\varepsilon_{i+2}
\le f_i\,f_{i+2} ,
\]
because $e^{-i}\varepsilon_i\le f_i$.
\end{proof}

The hypotheses of Propositions~\ref{prop:renewalconvex}
and~\ref{prop:extremecolumn} on the kernel thus hold for $f$. Because $f$
depends neither on the column nor on the row, its renewal sequence is one
sequence, and it can be written down. Let $F$ be the power series with
coefficients $c_n$ and $E=1/F$ its inverse, as in Section~\ref{sec:series}.
For integers $j\ge0$ and $s\ge0$ let $h_{j,s}=[X^s]\,X^jF^{j+1}$ be the
coefficient of $X^s$ in the power series $X^jF^{j+1}$, and put
\begin{equation}\label{eq:renewal-resolvent}
H_s=\sum_{j=0}^{s}2^{-j}\,h_{j,s},\qquad G_s=e^{-s}\,H_s
\qquad(s\ge0).
\end{equation}
The sequence $G$ is the \emph{resolvent sequence} of the frozen kernel.

The formula is found as follows. The tilt by $1/e$ turns the sequence $H$
into $G$, and it turns into $f$ the kernel whose first term is
$\varepsilon_1+1/2$ and whose $m$-th term is $\varepsilon_m$ for $m\ge2$.
Since $E$ is the series with constant term $1$ and coefficient
$-\varepsilon_m$ of $X^m$ for $m\ge1$, the renewal sequence of the latter
kernel has the generating series $1/(E-X/2)$, and
\[
\frac{1}{E-X/2}=\frac{F}{1-XF/2}=\sum_{j\ge0}\Bigl(\frac X2\Bigr)^{j}F^{j+1}.
\]
The coefficient of $X^s$ in this series, to which only the terms with
$j\le s$ contribute, is the number $H_s$ of \eqref{eq:renewal-resolvent}.
The following lemma confirms that $G$ is the renewal sequence of $f$. The
sequence $G$ is defined by its coefficients and not by its renewal equation,
because the coefficients $h_{j,s}$ have a closed form, from which a lower
bound for $G$ can be read off.

\begin{lemma}\label{lem:frozenrenewal}
Let $f$ be the frozen kernel of \eqref{eq:renewal-frozen} and $G$ the
sequence of \eqref{eq:renewal-resolvent}. Then $G$ is the renewal sequence
of the kernel $f$ in the sense of Definition~\ref{def:renewal}.
\end{lemma}

\begin{proof}
Clearly $h_{j,s}=0$ for $s<j$, so that $H_s$ is the sum of
$2^{-j}h_{j,s}$ over $0\le j\le J$, for every $J\ge s$, and $h_{0,s}=c_s$. In
particular $H_0=c_0=1$, and hence $G_0=1$.

Fix integers $j\ge0$ and $s\ge1$. Lemma~\ref{lem:inverseseries} expresses the
coefficients of the product of $E$ with an arbitrary power series; for the
series $X^jF^{j+1}$ it gives
\[
h_{j,s}-\sum_{m=1}^{s}\varepsilon_m\,h_{j,s-m}
=[X^s]\bigl(E\cdot X^jF^{j+1}\bigr)=[X^s]\,X^jF^{j} ,
\]
because $EF=1$. The right side is $0$ for $j=0$, and it is $h_{j-1,s-1}$ for
$j\ge1$. We multiply by $2^{-j}$ and sum over $0\le j\le s$. Since
$h_{j,s-m}=0$ for $j>s-m$, the left side becomes $H_s$ minus the sum of
$\varepsilon_mH_{s-m}$ over $1\le m\le s$, and the right side becomes
\[
\sum_{j=1}^{s}2^{-j}\,h_{j-1,s-1}
=\frac12\sum_{j=0}^{s-1}2^{-j}\,h_{j,s-1}=\frac12\,H_{s-1} .
\]
Hence
\[
H_s=\sum_{m=1}^{s}\varepsilon_m\,H_{s-m}+\frac12\,H_{s-1}
\qquad\text{for every integer } s\ge1 .
\]
Together with $H_0=1$ this says that
$H$ is the renewal sequence of the kernel whose first term is
$\varepsilon_1+1/2=3/2$ and whose $m$-th term is $\varepsilon_m$ for
$m\ge2$. We multiply the last display by $e^{-s}$. Since
$e^{-s}=e^{-m}e^{-(s-m)}$, the term $e^{-s}\varepsilon_mH_{s-m}$ is
$e^{-m}\varepsilon_m\,G_{s-m}$, and $e^{-s}H_{s-1}/2=G_{s-1}/(2e)$. The
contributions at each index $i$ add up to $f_iG_{s-i}$, because $f_i$ is
$e^{-i}\varepsilon_i$, increased by $1/(2e)$ for $i=1$. Hence
$G_s=\sum_{i=1}^{s}f_i\,G_{s-i}$ for every integer $s\ge1$, and together with
$G_0=1$ this says that $G$ is the renewal sequence of $f$.
\end{proof}

Since $f_i>0$ for every $i\ge1$, Lemma~\ref{lem:frozenrenewal} and the
positivity of the renewal sequence of a positive kernel give
\begin{equation}\label{eq:renewal-positive}
G_s>0\qquad\text{for every integer } s\ge0 .
\end{equation}
With Lemmas~\ref{lem:frozenconvex} and~\ref{lem:frozenrenewal},
Proposition~\ref{prop:extremecolumn} can be applied to the frozen kernel
and its renewal sequence $G$, and it bounds a far feedback by an expression
in the values of $G$. We now prove the two bounds for $G$ by which
Section~\ref{sec:far} evaluates this expression. The lower bound is the next
lemma.

\begin{lemma}\label{lem:resolventlower}
Let $G$ be the sequence of \eqref{eq:renewal-resolvent}. For every integer
$s\ge6$,
\[
G_s\ge \frac{e}{s+1}\cdot
\frac{1-\dfrac{1}{12(s+1)}}{\sqrt{2\pi(s+1)}}\cdot
\Bigl(4-\frac{21}{s+1}\Bigr).
\]
\end{lemma}

\begin{proof}
For $0\le j\le s$ the number $h_{j,s}$ is the coefficient of $X^{s-j}$ in
$F^{j+1}$. Since $F=e^{T}$ and $T=Xe^{T}$ (Section~\ref{sec:series}), the
Lagrange inversion formula \cite[Theorem~A.2]{flajoletsedgewick2009},
applied to the function $u\mapsto e^{(j+1)u}$, gives for every integer
$n\ge1$
\[
[X^n]\,F^{j+1}=\frac1n\,[u^{n-1}]\bigl((j+1)\,e^{(j+1+n)u}\bigr)
=\frac{(j+1)\,(j+1+n)^{n-1}}{n!} ,
\]
and the constant term of $F^{j+1}$ is $1$. For $n=s-j$ this shows that
$h_{j,s}=(j+1)\,(s+1)^{s-j-1}/(s-j)!$, also for $j=s$. Dividing by
$h_{0,s}=c_s=(s+1)^{s-1}/s!$ we obtain, for $0\le j\le s$,
\[
h_{j,s}=c_s\,(j+1)\,\frac{s!}{(s-j)!\,(s+1)^{j}}
=c_s\,(j+1)\prod_{t=1}^{j}\Bigl(1-\frac{t}{s+1}\Bigr).
\]
A product of factors $1-x_t$ with numbers $x_t$ in the interval $[0,1]$ is
at least one minus the sum of the $x_t$, and the numbers $t/(s+1)$ with
$1\le t\le j$ add up to $j(j+1)/(2(s+1))$. Therefore
\[
h_{j,s}\ge c_s\,(j+1)\Bigl(1-\frac{j(j+1)}{2(s+1)}\Bigr)
\qquad(0\le j\le s).
\]
We multiply by $2^{-j}$ and sum over $0\le j\le s$:
\[
H_s\ge c_s\sum_{j=0}^{s}\frac{j+1}{2^{j}}
-\frac{c_s}{2(s+1)}\sum_{j=0}^{s}\frac{j\,(j+1)^2}{2^{j}} .
\]
It is easy to check by induction on $s$ that
\begin{gather*}
\sum_{j=0}^{s}\frac{j+1}{2^{j}}=4-\frac{s+3}{2^{s}},\\
\sum_{j=0}^{s}\frac{j\,(j+1)^2}{2^{j}}
=40-\frac{s^3+8s^2+27s+40}{2^{s}} .
\end{gather*}
In particular the second sum is at most $40$. Moreover
$(s+1)(s+3)\le2^{s}$ for $s\ge6$, as follows by induction from the case
$s=6$, which is $63\le64$. Hence $(s+3)/2^{s}\le1/(s+1)$ for $s\ge6$, and
\[
H_s\ge c_s\Bigl(4-\frac{1}{s+1}-\frac{20}{s+1}\Bigr)
=c_s\Bigl(4-\frac{21}{s+1}\Bigr)\qquad(s\ge6).
\]

It remains to bound $e^{-s}c_s$. Since $c_s=(s+1)^{s}/(s+1)!$, we have
$e^{-s}c_s=e\,P(s+1)/(s+1)$ with the Poisson mass $P$, and
the lower bound of \eqref{eq:prelim-stirling} gives
\[
P(s+1)\ge \Bigl(1-\frac{1}{12(s+1)}\Bigr)\frac{1}{\sqrt{2\pi(s+1)}} .
\]
For $s\ge6$ the number $4-21/(s+1)$ is nonnegative, so that
\begin{align*}
G_s=e^{-s}H_s&\ge e^{-s}c_s\Bigl(4-\frac{21}{s+1}\Bigr)\\
&\ge \frac{e}{s+1}\cdot
\frac{1-\dfrac{1}{12(s+1)}}{\sqrt{2\pi(s+1)}}\cdot
\Bigl(4-\frac{21}{s+1}\Bigr).\qedhere
\end{align*}
\end{proof}

\subsection{Total masses}\label{subsec:renewal-mass}

It remains to bound the partial sums of $G$. If one sets $X=1$ in the
identity between generating series of Section~\ref{subsec:renewal-general},
the sum of a renewal sequence appears as $1/(1-\sigma)$, where $\sigma$ is
the sum of its kernel. A bound for the sum of $G$ therefore amounts to a
bound below one for the sum of the frozen kernel, and by
\eqref{eq:renewal-frozen} this requires the sum of
$e^{-m}\varepsilon_m$ over $m\ge1$. The proof of
Corollary~\ref{cor:kernelmass} obtains this sum from the sum of the sequence
$(e^{-n}c_n)_{n\ge0}$, which is given by the total mass of the Borel
distribution (Lemma~\ref{lem:borelmass}), and from the fact that the
sequence $(c_n)_{n\ge0}$ is the renewal sequence of the kernel
$(\varepsilon_m)_{m\ge1}$ (Corollary~\ref{cor:crenewal}). It does not
substitute into an identity between generating series, and it argues with
partial sums instead.

Let $u$ be the renewal sequence of a kernel $\varphi$. For integers $N\ge0$
we write
\[
\sigma_N=\sum_{i=1}^{N}\varphi_i,\qquad\Sigma_N=\sum_{n=0}^{N}u_n
\]
for the partial sums of the kernel and of its renewal sequence.

The numbers $e^{-j}a_j=e^{-j}j^{j-1}/j!$, $j\ge1$, are the probabilities of
the Borel distribution with parameter $1$, and the next lemma is the
statement that this distribution has total mass one.

\begin{lemma}\label{lem:borelmass}
The series of the nonnegative numbers $e^{-j}a_j$, $j\ge1$, converges, and
\[
\sum_{j=1}^{\infty}e^{-j}\,\frac{j^{j-1}}{j!}=1 .
\]
\end{lemma}

\begin{proof}
Aldous and Pitman state two facts about the total progeny of a
Galton--Watson process \cite[Section~2.2 and~(10)]{aldouspitman1998}, which
we use for the Poisson offspring law with mean $1$: the total progeny is
then finite with probability one, and it equals $j$ with probability
$e^{-j}j^{j-1}/j!$, for every integer $j\ge1$. From these two facts we
deduce the lemma. When the total progeny is finite, it is one of the
integers $j\ge1$, so that the probability that it is finite is the sum of
the numbers $e^{-j}j^{j-1}/j!$ over $j\ge1$. Hence this series converges,
and its sum is $1$.
\end{proof}

As a consequence we obtain the sum of the kernel $\varepsilon$ at the tilt
$1/e$.

\begin{corollary}\label{cor:kernelmass}
The series of the nonnegative numbers $e^{-m}\varepsilon_m$, $m\ge1$,
converges, and
\[
\sum_{m=1}^{\infty}e^{-m}\,\frac{(m-1)^{m-1}}{m!}=1-\frac1e .
\]
\end{corollary}

\begin{proof}
Let $u$ be the renewal sequence of a kernel $\varphi$ with $\varphi_i>0$
for every $i\ge1$, so that $u_n>0$ for every $n\ge0$, as was noted after
Definition~\ref{def:renewal}. Adding the renewal equations at
$n=1,\dots,N$ and collecting the terms with the same $i$, we obtain
\begin{equation}\label{eq:renewal-summed}
\Sigma_N-1=\sum_{n=1}^{N}\sum_{i=1}^{n}\varphi_i\,u_{n-i}
=\sum_{i=1}^{N}\varphi_i\,\Sigma_{N-i}\qquad(N\ge0).
\end{equation}
The numbers $\Sigma_N$ do not decrease in $N$, and $\Sigma_0=1$. If every
$\Sigma_{N-i}$ in \eqref{eq:renewal-summed} is replaced by the larger number
$\Sigma_N$, we get $\Sigma_N-1\le\sigma_N\,\Sigma_N$, and dividing by
$\Sigma_N\ge1$ gives
\begin{equation}\label{eq:renewal-lower}
1-\frac{1}{\Sigma_N}\le\sigma_N\qquad(N\ge0).
\end{equation}
If, in \eqref{eq:renewal-summed} with $N$ replaced by $N+L$, the terms with
$i>N$ are discarded and $\Sigma_{N+L-i}$ is replaced by the smaller number
$\Sigma_L$ in the terms with $i\le N$, we get
\begin{equation}\label{eq:renewal-upper}
\sigma_N\,\Sigma_L\le\Sigma_{N+L}-1\qquad(N\ge0,\ L\ge0).
\end{equation}

We now take $\varphi_m=e^{-m}\varepsilon_m$ for $m\ge1$ and
$u_n=e^{-n}c_n$ for $n\ge0$. By Corollary~\ref{cor:crenewal} and $c_0=1$,
the sequence $(c_n)$ is the renewal sequence of the kernel
$(\varepsilon_m)_{m\ge1}$. We multiply its renewal equation at $n\ge1$ by
$e^{-n}$; since $e^{-n}=e^{-m}e^{-(n-m)}$, the term
$e^{-n}\varepsilon_mc_{n-m}$ is $\varphi_m\,u_{n-m}$, so that
$u_n=\sum_{m=1}^{n}\varphi_m\,u_{n-m}$. Together with $u_0=c_0=1$ this says
that $u$ is the renewal sequence of $\varphi$. Since $\varepsilon_m>0$ for
$m\ge1$, the inequalities \eqref{eq:renewal-lower} and
\eqref{eq:renewal-upper} hold for the partial sums $\sigma_N$ of $\varphi$
and the partial sums $\Sigma_N$ of $u$.

Since $c_n=a_{n+1}$, the number $e^{-n}c_n$ is $e$ times the term
$e^{-(n+1)}a_{n+1}$ of the series of Lemma~\ref{lem:borelmass}. Hence
$\Sigma_N\le e$ for every $N$, and $\Sigma_N$ tends to $e$. By
\eqref{eq:renewal-upper}, $\sigma_N\,\Sigma_L\le e-1$ for all $N$ and $L$,
and letting $L$ tend to infinity we get $\sigma_N\le1-1/e$. By
\eqref{eq:renewal-lower},
$\sigma_N\ge1-1/\Sigma_N$, and the right side tends to $1-1/e$. Hence
$\sigma_N$ tends to $1-1/e$.
\end{proof}

We can now bound the partial sums of the resolvent sequence. The bound of
the next lemma is the value of $1/(1-\sigma)$ at
$\sigma=(1-1/e)+1/(2e)=1-1/(2e)$, the sum of the kernel $\varepsilon$ at the
tilt $1/e$ (Corollary~\ref{cor:kernelmass}) increased by the summand
$1/(2e)$ of $f_1$.

\begin{lemma}\label{lem:resolventmass}
Let $G$ be the sequence of \eqref{eq:renewal-resolvent}. For every integer
$N\ge0$,
\[
\sum_{s=0}^{N-1}G_s\le2e .
\]
\end{lemma}

\begin{proof}
By Lemma~\ref{lem:frozenrenewal}, $G$ is the renewal sequence of the frozen
kernel $f$, and $f_i>0$ for every $i\ge1$. Hence \eqref{eq:renewal-lower},
which was shown in the proof of Corollary~\ref{cor:kernelmass} for the
renewal sequence of every kernel with positive terms, holds for the partial
sums $\sigma_N$ of $f$ and the partial sums $\Sigma_N$ of $G$. One has
$\sigma_0=0$, and by \eqref{eq:renewal-frozen}, $\sigma_N$ is for $N\ge1$ the
sum of $e^{-i}\varepsilon_i$ over $1\le i\le N$ plus $1/(2e)$,
where this sum is at most $1-1/e$ by Corollary~\ref{cor:kernelmass}. Hence
\[
1-\frac{1}{\Sigma_N}\le\sigma_N\le1-\frac1e+\frac{1}{2e}=1-\frac{1}{2e}
\qquad(N\ge0),
\]
and therefore $\Sigma_N\le2e$ for every $N\ge0$. The sum of the lemma is
empty for $N=0$, and it is $\Sigma_{N-1}$ for $N\ge1$.
\end{proof}

\section{Persistence on the far rows of the first layer}\label{sec:far}
We return to the first layer of a column $k\ge82$.
Proposition~\ref{prop:nearrows} shows that a column which has turned
positive stays positive on the near rows $n\le k+0.47k$. Here we prove that
the step in a far row $n$, with $k+0.47k<n\le2k-1$, is positive if the steps
in rows $k+r_q,\dots,n-1$ are positive for an integer $r_q\ge0$ with
$3\,r_q\le k$ (Proposition~\ref{prop:farrows}). We then combine the two
statements. Call a row $k+r_q$ early if $r_q\ge1$ and $3\,r_q\le k$. A
column $k\ge82$ whose steps in rows $k,\dots,k+r_q-1$ are not positive and
whose step in an early row $k+r_q$ is positive has a positive step in every
row from $k+r_q$ to $2k-1$, the end of its first layer
(Proposition~\ref{prop:staypositive}). In
Section~\ref{sec:curve} this proposition is applied together with
Corollary~\ref{cor:earlypositive}, which provides in every column $k\ge82$
a positive row $j$ with $k+1\le j\le k+27k/100$, so that the first positive
row of the column is early in this sense.

The comparison of the kernel on which the argument of
Section~\ref{sec:near} rests is established only up to row
$k+0.47k$. On the far rows we use instead that the first positive row is
early, so that a long run of rows directly before row $n$ is known to be
positive; with such a run the step in row $n$ can be bounded from below by
the entries of column $k-1$ alone.

\subsection{A lower bound for the step after a run of nonnegative
steps}\label{subsec:far-inequality}

Let $k\ge2$. Throughout the section, $\delta_n=\delta^{(k)}_n$ is the
normalised step of column $k$, which has the sign of
$d_{n,k}$, and $\beta^{(k-1)}_n$ is the normalised count of column $k-1$
(Section~\ref{sec:counts}). For a row $k+1\le n\le2k-1$ of the first layer
of column $k$,
the column equation of Proposition~\ref{prop:columnequation} reads
\[
\delta_n=\frac{k+1}{n}\,\delta_{n-1}+\frac2n\,\beta^{(k-1)}_{n-1}
+\sum_{m=1}^{n-k}\varepsilon_m\,\delta_{n-m}.
\]
It expresses $\delta_n$ in terms of the source
$(2/n)\,\beta^{(k-1)}_{n-1}$ and of the
steps $\delta_{n-1},\dots,\delta_k$ of the earlier rows, and nothing is known
about the size of the steps before the first positive row. Every step,
however, is bounded from below in terms of column $k-1$. Subtracting
$\beta^{(k-1)}_j$ from both sides of the inequality
$\gamma_k\,\beta^{(k-1)}_{j-1}\le\beta^{(k)}_j$ of
Proposition~\ref{prop:diagonallower}, in which $\gamma_k=(1+1/k)^{k-2}$ is
the diagonal factor of Section~\ref{sec:diagonal}, we get
\begin{equation}\label{eq:far-diagonal}
\gamma_k\,\beta^{(k-1)}_{j-1}-\beta^{(k-1)}_j\le\delta_j\qquad(j\ge1).
\end{equation}
The kernel of the column equation is the free prime kernel $\varepsilon_m$,
which is positive. So \eqref{eq:far-diagonal} may be inserted for every step
$\delta_{n-m}$ in the sum of the column equation, the feedback of the
earlier rows, whatever the sign of that step. The entry
$\beta^{(k-1)}_{n-1-m}$ then occurs twice, with the
coefficient $\gamma_k\varepsilon_m$ from row $n-m$ and with the
coefficient $-\varepsilon_{m+1}$ from row $n-m-1$, so that the insertion
leaves $(\gamma_k\varepsilon_m-\varepsilon_{m+1})\,\beta^{(k-1)}_{n-1-m}$ at
distance $m$. This is a summation by parts, and by itself it is of no
use: the entry $\beta^{(k-1)}_{n-1}$ receives the coefficient
$-\varepsilon_1=-1$ from row $n-1$, against the coefficient $2/n$ that it
has in the source.

The rows directly before $n$ can be treated better when their steps are known
to be nonnegative. For such a row the weight $\varepsilon_m$ may be replaced
by any smaller nonnegative weight before \eqref{eq:far-diagonal} is inserted.
Suppose that the steps in the $\tau-1$ rows before $n$ are nonnegative. On
these rows we replace $\varepsilon_m$ by the geometric weight
$\psi\,\gamma_k^{\,m-1}$, where $\psi=\varepsilon_\tau/\gamma_k^{\,\tau-1}$.
These weights are at most $\varepsilon_m$ as soon as
$\varepsilon_{m+1}\le\gamma_k\,\varepsilon_m$ for $m<\tau$; the insertions of
\eqref{eq:far-diagonal} telescope exactly for them; and the term that the
telescoping leaves at the far end of the run cancels against the first term
of the summation by parts over the remaining rows. What the run costs is the
single term $-\psi\,\beta^{(k-1)}_{n-1}$, and $\psi$ is small when the run is
long. The following lemma carries this out for arbitrary sequences, with
$\phi$, $u$, $\Delta$ and $\chi$ in the roles of $\varepsilon$,
$\beta^{(k-1)}$, $\delta$ and $\gamma_k$.

\begin{lemma}\label{lem:runbound}
Let $(\phi_m)_{m\ge0}$, $(u_j)_{j\ge0}$ and $(\Delta_j)_{j\ge0}$ be sequences
of real numbers, let $\chi>0$ be a real number, and let $k\ge1$, $n$ and
$\tau$ be integers with $k+1\le n$ and $2\le\tau\le n-k$. Suppose that the
following six conditions hold.
\begin{enumerate}
\item $\phi_m>0$ for every $m\ge1$.
\item $u_{k-1}=0$.
\item The sequences satisfy
\[
\frac2n\,u_{n-1}+\frac{k+1}{n}\,\Delta_{n-1}
+\sum_{m=1}^{n-k}\phi_m\,\Delta_{n-m}\le\Delta_n .
\]
\item $\chi\,u_{j-1}-u_j\le\Delta_j$ for every $j$ with $k\le j\le n-1$.
\item $\phi_{m+1}\le\chi\,\phi_m$ for every $m$ with $1\le m<\tau$.
\item $\Delta_j\ge0$ for every $j$ with $n-\tau+1\le j\le n-1$.
\end{enumerate}
Then
\[
\Bigl(\frac2n-\frac{\phi_\tau}{\chi^{\tau-1}}\Bigr)\,u_{n-1}
+\sum_{m=\tau}^{n-k-1}\bigl(\chi\,\phi_m-\phi_{m+1}\bigr)\,u_{n-1-m}
\le\Delta_n .
\]
\end{lemma}

\begin{proof}
Put $\psi=\phi_\tau/\chi^{\tau-1}$, which is positive by~(1). We have
\begin{equation}\label{eq:far-weights}
\psi\,\chi^{m-1}\le\phi_m\qquad(1\le m\le\tau).
\end{equation}
Indeed, for $m=\tau$ this is an equality, and the other cases follow
from~(5) by downward induction on $m$.

Since $\tau\ge2$, we have $\Delta_{n-1}\ge0$ by~(6), and (3) gives
\begin{equation}\label{eq:far-start}
\frac2n\,u_{n-1}+\sum_{m=1}^{n-k}\phi_m\,\Delta_{n-m}\le\Delta_n .
\end{equation}
For $1\le m\le n-k$, condition~(4) with $j=n-m$ reads
$\chi\,u_{n-1-m}-u_{n-m}\le\Delta_{n-m}$.

For $1\le m\le\tau-1$ we have $\Delta_{n-m}\ge0$ by~(6), so that by
\eqref{eq:far-weights} and~(4)
\[
\phi_m\,\Delta_{n-m}\ge\psi\,\chi^{m-1}\,\Delta_{n-m}
\ge\psi\,\bigl(\chi^{m}\,u_{n-1-m}-\chi^{m-1}\,u_{n-m}\bigr).
\]
The right-hand sides telescope, and their sum over these $m$ is
$\psi\,\chi^{\tau-1}u_{n-\tau}-\psi\,u_{n-1}$. As
$\psi\,\chi^{\tau-1}=\phi_\tau$, this gives
\begin{equation}\label{eq:far-telescope}
\sum_{m=1}^{\tau-1}\phi_m\,\Delta_{n-m}\ge\phi_\tau\,u_{n-\tau}-\psi\,u_{n-1}.
\end{equation}
For $\tau\le m\le n-k$ we use only that $\phi_m>0$, which gives
$\phi_m\,\Delta_{n-m}\ge\phi_m\,(\chi\,u_{n-1-m}-u_{n-m})$. We sum the
right-hand sides over these $m$ and collect the terms that contain the same
entry of $u$. The entry $u_{k-1}$, which occurs with the coefficient
$\chi\,\phi_{n-k}$, vanishes by~(2), and we obtain
\begin{equation}\label{eq:far-parts}
\sum_{m=\tau}^{n-k}\phi_m\,\Delta_{n-m}\ge-\phi_\tau\,u_{n-\tau}
+\sum_{m=\tau}^{n-k-1}\bigl(\chi\,\phi_m-\phi_{m+1}\bigr)\,u_{n-1-m}.
\end{equation}
When \eqref{eq:far-telescope} and \eqref{eq:far-parts} are added, the terms
$\phi_\tau\,u_{n-\tau}$ and $-\phi_\tau\,u_{n-\tau}$ cancel. Inserting the
result into \eqref{eq:far-start}, we obtain the assertion.
\end{proof}

\subsection{The cost of the run and the far coefficients}\label{subsec:far-estimates}

When Lemma~\ref{lem:runbound} is applied to a column $k$, we have
$\phi=\varepsilon$ and $\chi=\gamma_k$, and we write $t=\tau-1$ for the
length of the run. The run then costs the term
$-\psi\,\beta^{(k-1)}_{n-1}$ with
$\psi=\varepsilon_{t+1}/\gamma_k^{\,t}=\varepsilon_{t+1}\,\hat x_k^{\,t}$,
where $\hat x_k=1/\gamma_k$ is the diagonal tilt of column $k$
(Section~\ref{sec:near}). The proof of Proposition~\ref{prop:farrows}
compares $n\psi/2$ with $1$. The next lemma bounds $n\psi/2$, for a run of about $\sigma k$
rows and a row $n\le Lk$, by a quantity that decreases like $k^{-1/2}$. In
it $Y$ stands for $n$, and the factor $C$ accounts for the facts that $t$ is
an integer between $\sigma k$ and $\sigma k+1$ and that $\hat x_k$ exceeds
$e^{-1}$ by a factor of at most $e^{5/(2k)}$.

\begin{lemma}\label{lem:lossterm}
Let $k\ge15$ and $t\ge1$ be integers, and let $\sigma$, $C$, $L$ and $Y$ be
real numbers such that $\sigma>0$, $C>0$ and
\[
\sigma\,(k+3)\le C\,(\sigma k+1),\qquad
\sigma k\le t\le\sigma k+1,\qquad
0\le Y\le L\,k .
\]
Then the free prime kernel $\varepsilon$ and the diagonal tilt $\hat x_k$ of
column $k$ satisfy
\[
\frac Y2\;\varepsilon_{t+1}\,\hat x_k^{\,t}
\le\frac{C\,L\,e^{5\sigma/2}}{2\sigma\sqrt{2\pi\sigma k}}.
\]
\end{lemma}

\begin{proof}
Put $\psi=\varepsilon_{t+1}\,\hat x_k^{\,t}$.
Since $(t+1)\,\varepsilon_{t+1}=\nu_t$ and $e^{-t}\nu_t=P(t)$, where
$t\ge1$, the upper bound of \eqref{eq:prelim-stirling} gives
$(t+1)\,\varepsilon_{t+1}\le e^{t}/\sqrt{2\pi t}$, and
Lemma~\ref{lem:tiltbounds} gives $\hat x_k^{\,t}\le e^{-t}e^{5t/(2k)}$. As
$t\le\sigma k+1$, we obtain
\[
\psi\,(t+1)\sqrt{2\pi t}\le e^{5t/(2k)}\le e^{5\sigma/2}\,e^{5/(2k)} .
\]
It is easy to check that $e^{h}\le1/(1-h)\le1+6h/5$ for every real number
$h$ with $0\le h\le1/6$. As $5/(2k)\le1/6$ for $k\ge15$, this gives
$e^{5/(2k)}\le1+3/k$. By the hypothesis on $C$ and by $\sigma k\le t$,
\[
\sigma k\,\Bigl(1+\frac3k\Bigr)=\sigma\,(k+3)\le C\,(\sigma k+1)\le C\,(t+1).
\]
We multiply the first display by $\sigma k$, insert these two bounds, and
divide by $t+1$. This gives
$\psi\,\sigma k\sqrt{2\pi t}\le C\,e^{5\sigma/2}$, and here
$\sqrt{2\pi t}\ge\sqrt{2\pi\sigma k}$. Finally $0\le Y\le L\,k$ and
$\psi\ge0$, so that
\[
\frac Y2\,\psi\le\frac{L\,k}{2}\cdot
\frac{C\,e^{5\sigma/2}}{\sigma k\sqrt{2\pi\sigma k}}
=\frac{C\,L\,e^{5\sigma/2}}{2\sigma\sqrt{2\pi\sigma k}} .\qedhere
\]
\end{proof}

We turn to the sum in the conclusion of Lemma~\ref{lem:runbound}.
Recall the step ratio $g_m$ of the free prime kernel, for which
$\varepsilon_{m+1}=g_m\,\varepsilon_m$ (Section~\ref{sec:near}). The
coefficient that belongs to distance $m$ is
$\gamma_k\,\varepsilon_m-\varepsilon_{m+1}=-(g_m-\gamma_k)\,\varepsilon_m$.
It is nonnegative as long as $m+1\le3k/5$, by
Lemma~\ref{lem:kernelgrowth}, and it may be negative beyond. In the
proof of Proposition~\ref{prop:farrows} the coefficients with $m+1>3k/5$
occur with $\varepsilon_m$ replaced by $e^{-m}\varepsilon_m$. For every
$m\ge2$ the number $e^{-m}\varepsilon_m$ is the term $f_m$ of the frozen
kernel of Section~\ref{sec:renewal}. What has to be bounded
is therefore $(g_m-\gamma_k)\,f_m$, and the next lemma bounds it, uniformly
in $m$, by a quantity of the order $k^{-5/2}$.

\begin{lemma}\label{lem:farweights}
Let $k$ and $m\ge49$ be integers with $m+1\le k$ and $3k\le5(m+1)$. Then the
step ratio $g_m$ of the free prime kernel, the diagonal factor $\gamma_k$
and the frozen kernel $f$ satisfy
\[
(g_m-\gamma_k)\,f_m\le\frac{105}{100}\cdot\frac{c_{\mathrm S}}{k^2\sqrt k}.
\]
\end{lemma}

\begin{proof}
Put $j=m+1$, so that $50\le j\le k$ and $3k\le5j$. We first show that
\begin{equation}\label{eq:far-excess}
g_m-\gamma_k\le e\cdot\frac{5j-3k}{2kj}.
\end{equation}
Since $m\ge2$ and $m+1=j$, the bound \eqref{eq:near-ratio-upper} for the
step ratio reads $g_m\le e^{1-3/(2j)}$. On the other hand
\eqref{eq:near-factor-lower} gives $\gamma_k\ge e^{1-5/(2k)}$, because
$k\ge2$. As $5/(2k)-3/(2j)=(5j-3k)/(2kj)$ and
$e^{-h}\ge1-h$ for every real $h$,
\[
\gamma_k\ge e^{1-3/(2j)}\exp\Bigl(-\frac{5j-3k}{2kj}\Bigr)
\ge e^{1-3/(2j)}\Bigl(1-\frac{5j-3k}{2kj}\Bigr).
\]
Together with the bound for $g_m$ this yields \eqref{eq:far-excess}, because
$5j-3k\ge0$ and $e^{1-3/(2j)}\le e$.

Next, $m\,\varepsilon_m=\nu_{m-1}$ and $e^{-(m-1)}\nu_{m-1}=P(m-1)$, so that
by the upper bound of \eqref{eq:prelim-stirling}
\[
f_m=e^{-m}\varepsilon_m=\frac{e^{-1}P(m-1)}{m}
\le\frac{e^{-1}c_{\mathrm S}}{m\sqrt{m-1}}
=\frac{e^{-1}c_{\mathrm S}}{(j-1)\sqrt{j-2}} .
\]
Since $f_m\ge0$ and the right-hand side of \eqref{eq:far-excess} is
nonnegative, multiplying the two bounds gives
\[
(g_m-\gamma_k)\,f_m\le\frac{(5j-3k)\,c_{\mathrm S}}{2kj\,(j-1)\sqrt{j-2}} .
\]
It remains to show that
$(5j-3k)\,k^{3/2}\le1.05\cdot2j\,(j-1)\sqrt{j-2}$. For every real
$\rho\ge0$,
\[
2\rho^5-5\rho^2+3=(\rho-1)^2\,(2\rho^3+4\rho^2+6\rho+3)\ge0 .
\]
For $\rho=\sqrt{j/k}$ and after multiplication by $k^{5/2}$ this is
$(5j-3k)\,k^{3/2}\le2j^{5/2}$. Moreover
$j^{3/2}\le1.05\,(j-1)\sqrt{j-2}$: after squaring and division by
$j^3$ this is the inequality $1\le1.1025\,(1-1/j)^2(1-2/j)$, and for
$j\ge50$ its right-hand side is at least $1.1025\cdot0.98^2\cdot0.96>1$.
Multiplying the last inequality by $2j$ and combining it with the one before,
we obtain the claim.
\end{proof}

\subsection{The far rows}\label{subsec:far-rows}

The following proposition treats a far row $k+r$ of the first layer of a
column $k\ge82$. The proof applies Lemma~\ref{lem:runbound} to a run
of positive rows before row $k+r$ whose length is proportional to $k$, and
it bounds the cost of the run by Lemma~\ref{lem:lossterm}. For $r\le3k/5$
this suffices, because the sum in the conclusion of
Lemma~\ref{lem:runbound} then has no negative term. For $r>3k/5$ the terms
of the sum that may be negative are bounded by
Proposition~\ref{prop:extremecolumn}, applied to column $k-1$ and to the
renewal sequence $G$ of the frozen kernel, and by
Lemma~\ref{lem:farweights} and the two bounds of Section~\ref{sec:renewal}
for $G$; as the bound that results grows with $r$, these rows are divided
once more at $4k/5$, beyond which the run is longer and the bound for its
cost is smaller.

\begin{proposition}\label{prop:farrows}
Let $k\ge82$, $r\ge0$ and $r_q\ge0$ be integers such that
\[
47\,k<100\,r,\qquad r+1\le k,\qquad 3\,r_q\le k,
\]
and suppose that $d_{j,k}>0$ for every $j$ with $k+r_q\le j<k+r$. Then
$d_{k+r,k}>0$.
\end{proposition}

\begin{proof}
Let $n=k+r$; the hypotheses on $r$ say that $k+0.47k<n\le2k-1$. We write
$\delta_j=\delta^{(k)}_j=d_{j,k}/j!$, and we have to show that $\delta_n>0$.

Let $t_0\ge1$ be an
integer such that $\delta_j\ge0$ for every $j$ with $n-t_0\le j\le n-1$ and
\begin{equation}\label{eq:far-cut}
t_0+1\le r,\qquad 5\,(t_0+1)\le3k,
\end{equation}
and put $\psi=\varepsilon_{t_0+1}\,\hat x_k^{\,t_0}$, where
$\hat x_k=1/\gamma_k$. Then
\begin{equation}\label{eq:far-run}
\Bigl(\frac2n-\psi\Bigr)\,\beta^{(k-1)}_{n-1}
+\sum_{m=t_0+1}^{r-1}\bigl(\gamma_k\,\varepsilon_m-\varepsilon_{m+1}\bigr)\,
\beta^{(k-1)}_{n-1-m}\le\delta_n .
\end{equation}
This is the conclusion of Lemma~\ref{lem:runbound} for the sequences
$\phi_m=\varepsilon_m$, $u_j=\beta^{(k-1)}_j$ and $\Delta_j=\delta_j$, with
$\chi=\gamma_k$ and $\tau=t_0+1$. Here $n-k=r$, so that the requirement
$2\le\tau\le n-k$ of the lemma is the first part of \eqref{eq:far-cut}, and
$\phi_\tau/\chi^{\tau-1}=\psi$. The six conditions of the lemma hold.
Conditions~(1) and~(2) are clear, since $\varepsilon_m>0$ for $m\ge1$ and
$\beta^{(k-1)}_{k-1}=b_{k-1,k-1}/(k-1)!=0$.
Condition~(3) holds with equality: it is the column equation of
Proposition~\ref{prop:columnequation}, which is available because
$k+1\le n\le2k-1$. Condition~(4) is \eqref{eq:far-diagonal}. Condition~(5)
follows from Lemma~\ref{lem:kernelgrowth}, since
$5\,(m+1)\le5\,(t_0+1)\le3k$ for $m\le t_0$ by the second part of
\eqref{eq:far-cut}, and condition~(6) is the assumption on rows
$n-t_0,\dots,n-1$.

Next we fix the length $t_0$ of the run. We put
\[
\sigma=
\begin{cases}
\dfrac{47}{100}-\dfrac13=\dfrac{41}{300}&\text{if }r\le\dfrac{3k}5,\\[2ex]
\dfrac35-\dfrac13=\dfrac4{15}&\text{if }\dfrac{3k}5<r\le\dfrac{4k}5,\\[2ex]
\dfrac45-\dfrac13=\dfrac7{15}&\text{if }r>\dfrac{4k}5,
\end{cases}
\qquad\qquad t_0=\lfloor\sigma k\rfloor+1 .
\]
Then $\sigma k\le t_0\le\sigma k+1$, and in each of the three cases
$r>(\sigma+1/3)\,k$. Together with $3r_q\le k$ this gives $r-r_q>\sigma k$,
and hence $r-r_q\ge t_0$, because $r-r_q$ is an integer. Rows
$n-t_0,\dots,n-1$ are therefore among rows $k+r_q,\dots,n-1$, in which
the step is positive by hypothesis. Moreover
$t_0+1\le\sigma k+2<(\sigma+1/3)\,k<r$ and
$5\,(t_0+1)\le5\sigma k+10\le7k/3+10\le3k$, both because $k\ge15$. Thus
$t_0$ satisfies \eqref{eq:far-cut}, and \eqref{eq:far-run} holds with this
$t_0$.

Let $L$ and $C>0$ be real numbers with
$n\le L\,k$ and
$\sigma\,(k+3)\le C\,(\sigma k+1)$. Lemma~\ref{lem:lossterm}, applied with
$t=t_0$ and $Y=n$, bounds the nonnegative number $n\psi/2$ by
$C\,L\,e^{5\sigma/2}/\bigl(2\sigma\sqrt{2\pi\sigma k}\bigr)$. Squaring, we
get
\begin{equation}\label{eq:far-loss}
\Bigl(\frac{n\psi}{2}\Bigr)^2\le\frac{C^2L^2\,e^{5\sigma}}{8\pi\,\sigma^3\,k}.
\end{equation}
We shall evaluate the right-hand side with $\pi>3.14$ and with the bounds
\[
e^{41/60}\le2,\qquad e^{4/3}\le3.8,\qquad e^{7/3}\le10.32
\]
for the three values of $e^{5\sigma}$. These bounds hold because $e<2.7183$
and $2.7183^{41}<2^{60}$, $2.7183^{4}<3.8^{3}$, $2.7183^{7}<10.32^{3}$.

We also note that $\beta^{(k-1)}_{n-1}>0$,
because the preference sequence $(1,2,\dots,n-2,\,n-k)$ of length $n-1$ is
clearly a parking function with maximum displacement $k-1$: its last car
prefers spot $n-k$ and parks in spot $n-1$.

It remains to derive $\delta_n>0$ from \eqref{eq:far-run} in each of the
three ranges of $r$ that occur in the definition of $\sigma$. Suppose first
that $r\le3k/5$. Then every index $m$ of the sum in
\eqref{eq:far-run} satisfies $5\,(m+1)\le5r\le3k$, so that its coefficient
is nonnegative by Lemma~\ref{lem:kernelgrowth}; the entries
$\beta^{(k-1)}_{n-1-m}$ are nonnegative as well. Dropping the sum, we get
\begin{equation}\label{eq:far-first}
\frac2n\,\Bigl(1-\frac{n\psi}{2}\Bigr)\,\beta^{(k-1)}_{n-1}\le\delta_n .
\end{equation}
Here $n\le8k/5$ and $3\sigma\le1$, so that \eqref{eq:far-loss} holds with
$L=8/5$ and $C=1$:
\[
\Bigl(\frac{n\psi}{2}\Bigr)^2
\le\frac{(8/5)^2\cdot2}{8\cdot3.14\cdot(41/300)^3\,k}
\le\frac{79.85}{k}<1,
\]
where the last step uses $k\ge82$. Hence $n\psi/2<1$, and
\eqref{eq:far-first} gives $\delta_n>0$, because $\beta^{(k-1)}_{n-1}>0$.

Now let $r>3k/5$. In this case the sum in \eqref{eq:far-run} may have
negative terms, and we bound it from below. Let $m_0=\lfloor3k/5\rfloor$. Then
$49\le m_0\le r-1\le k-2$, the first inequality because $k\ge82$. For
$m\ge1$ put $\theta_m=\max(g_m-\gamma_k,\,0)$. As
$\varepsilon_{m+1}=g_m\,\varepsilon_m$, we have
\[
\gamma_k\,\varepsilon_m-\varepsilon_{m+1}
=-(g_m-\gamma_k)\,\varepsilon_m\ge-\theta_m\,\varepsilon_m\qquad(m\ge1).
\]
Moreover $\theta_m=0$ for $m<m_0$, because then
$5\,(m+1)\le5m_0\le3k$ and Lemma~\ref{lem:kernelgrowth} gives
$g_m\le\gamma_k$. As the entries $\beta^{(k-1)}_{n-1-m}$ are nonnegative,
the sum in \eqref{eq:far-run} is therefore at least $-\Sigma$, where
\[
\Sigma=\sum_{m=m_0}^{r-1}\theta_m\,\varepsilon_m\,\beta^{(k-1)}_{n-1-m},
\]
and \eqref{eq:far-run} implies
\begin{equation}\label{eq:far-split}
\Bigl(\frac2n-\psi\Bigr)\,\beta^{(k-1)}_{n-1}-\Sigma\le\delta_n .
\end{equation}
To bound $\Sigma$ we compare column $k-1$ with the renewal sequence of
the frozen kernel. Let $v_i=e^{-i}\beta^{(k-1)}_{k+i}$ for $i\ge0$. We claim
that
\begin{equation}\label{eq:far-innovation}
\sum_{m=1}^{i}f_m\,v_{i-m}\le v_i\qquad(0\le i\le k-2).
\end{equation}
For $i=0$ this says that $v_0\ge0$. For $1\le i\le k-2$ we apply
Lemma~\ref{lem:levelequation} to column $k-1$ and row $k+i$, which is
at most $2k-1$:
\[
\beta^{(k-1)}_{k+i}=\frac{z^{(k)}_{k+i-1}}{k+i}
+\frac{k-1}{k+i}\,\beta^{(k-1)}_{k+i-1}
+\sum_{m=1}^{i}\varepsilon_m\,\beta^{(k-1)}_{k+i-m}.
\]
The first term on the right is nonnegative, and $(k-1)/(k+i)\ge1/2$ because
$i\le k-2$. Dropping the first term and multiplying by $e^{-i}$, we get
\[
v_i\ge\frac{e^{-1}}{2}\,v_{i-1}+\sum_{m=1}^{i}e^{-m}\varepsilon_m\,v_{i-m}.
\]
The right-hand side is the left-hand side of \eqref{eq:far-innovation},
since $f_1=e^{-1}\varepsilon_1+e^{-1}/2$ and $f_m=e^{-m}\varepsilon_m$ for
$m\ge2$. This proves \eqref{eq:far-innovation}. We need it for
$i\le r-1$, and these $i$ are at most $k-2$ by the hypothesis $r\le k-1$.

The frozen kernel is positive, and it is log-convex by
Lemma~\ref{lem:frozenconvex}. The sequence $G$ is its renewal sequence by
Lemma~\ref{lem:frozenrenewal}. The weights $\theta_m$ are nonnegative, and
they do not decrease when $m$ grows, because $g_m$ does not
(Section~\ref{sec:near}). With \eqref{eq:far-innovation} for $i\le r-1$,
these are the hypotheses of Proposition~\ref{prop:extremecolumn} for the
kernel $f$, its renewal sequence $G$, the sequence $v$ and the weights
$\theta_m$ with $m_0\le m\le r-1$. The proposition gives
\[
\Bigl(\sum_{m=m_0}^{r-1}\theta_m\,f_m\,v_{r-1-m}\Bigr)\,G_{r-1}
\le\Bigl(\sum_{m=m_0}^{r-1}\theta_m\,f_m\,G_{r-1-m}\Bigr)\,v_{r-1}.
\]
For $m\ge m_0$ we have $f_m=e^{-m}\varepsilon_m$, so that
$f_m\,v_{r-1-m}=e^{-(r-1)}\,\varepsilon_m\,\beta^{(k-1)}_{n-1-m}$ by the
definition of $v$, and $v_{r-1}=e^{-(r-1)}\beta^{(k-1)}_{n-1}$. After the
factor $e^{-(r-1)}$ has been cancelled, the inequality reads
\begin{equation}\label{eq:far-extreme}
\Sigma\,G_{r-1}\le\Bigl(\sum_{m=m_0}^{r-1}\theta_m\,f_m\,G_{r-1-m}\Bigr)\,
\beta^{(k-1)}_{n-1}.
\end{equation}
For $m_0\le m\le r-1$ we have $m\ge49$, $m+1\le k$ and
$3k<5\,(m_0+1)\le5\,(m+1)$, so that Lemma~\ref{lem:farweights} gives
$\theta_m\,f_m\le1.05\,c_{\mathrm S}\,k^{-5/2}$. The numbers $G_i$ are
positive by \eqref{eq:renewal-positive}. Hence the sum in
\eqref{eq:far-extreme} is at most $1.05\,c_{\mathrm S}\,k^{-5/2}$ times
$G_0+G_1+\dots+G_{r-1-m_0}$. We multiply
\eqref{eq:far-split} by $G_{r-1}$, insert \eqref{eq:far-extreme} with this
bound, and divide by $G_{r-1}$ again. With
\[
\Omega=\frac n2\cdot\frac{1.05\,c_{\mathrm S}}{k^{5/2}}\cdot
\frac{1}{G_{r-1}}\sum_{i=0}^{r-1-m_0}G_i
\]
the result is
\begin{equation}\label{eq:far-second}
\frac2n\,\Bigl(1-\frac{n\psi}{2}-\Omega\Bigr)\,\beta^{(k-1)}_{n-1}
\le\delta_n .
\end{equation}
Since $\beta^{(k-1)}_{n-1}>0$, it remains to show that $n\psi/2+\Omega<1$.

For this we bound $G_{r-1}$ from below. For an integer $j\ge6$ let
$\upsilon(j)=\bigl(1-1/(12j)\bigr)\bigl(4-21/j\bigr)$; both factors are positive
and increase with $j$. As $r-1\ge m_0\ge6$, Lemma~\ref{lem:resolventlower}
gives $G_{r-1}\ge e\,c_{\mathrm S}\,r^{-3/2}\,\upsilon(r)$. As $n=k+r$, it follows
that, if $r_0\ge6$ is an integer with $r\ge r_0$, so that
$\upsilon(r)\ge\upsilon(r_0)$, then
\begin{equation}\label{eq:far-feedback}
\Omega\le\frac{1.05}{2}\cdot\frac{(k+r)\,r^{3/2}}{k^{5/2}}\cdot
\frac{1}{e\,\upsilon(r_0)}\sum_{i=0}^{r-1-m_0}G_i .
\end{equation}

In the range $3k/5<r\le4k/5$ we have $n\le9k/5$ and $3\sigma\le1$, so
that \eqref{eq:far-loss} holds with $L=9/5$ and $C=1$:
\[
\Bigl(\frac{n\psi}{2}\Bigr)^2
\le\frac{(9/5)^2\cdot3.8}{8\cdot3.14\cdot(4/15)^3\,k}
\le\frac{25.85}{k},
\]
and hence $n\psi/2\le0.57$, because $k\ge82$. In \eqref{eq:far-feedback} we
may take $r_0=50$,
because $r>3k/5\ge49.2$, and $\upsilon(50)=(599/600)\cdot(179/50)>3.574$. The sum
of the $G_i$ is at most $2e$ by Lemma~\ref{lem:resolventmass}, and
\[
(k+r)\,r^{3/2}\le\frac95\,\Bigl(\frac45\Bigr)^{3/2}k^{5/2}\le1.2881\,k^{5/2},
\]
since $(4/5)^3\le0.7156^2$. Therefore
$\Omega\le1.05\cdot1.2881/3.574<0.385$, and
$n\psi/2+\Omega<0.57+0.385<1$.

Finally, in the range $r>4k/5$ we have $n\le2k$, and \eqref{eq:far-loss}
holds with $L=2$ and $C=1.011$: for $\sigma=7/15$ the condition
$\sigma\,(k+3)\le C\,(\sigma k+1)$ is equivalent to
$7/5-1.011\le0.011\cdot7k/15$, which holds for $k\ge76$. Thus
\[
\Bigl(\frac{n\psi}{2}\Bigr)^2
\le\frac{1.011^2\cdot4\cdot10.32}{8\cdot3.14\cdot(7/15)^3\,k}
\le\frac{16.53}{k}.
\]
In \eqref{eq:far-feedback} we have $(k+r)\,r^{3/2}\le2k^{5/2}$, since
$r\le k$.

Suppose first that $k\ge110$. Then $(n\psi/2)^2\le16.53/110$, hence
$n\psi/2\le0.395$. We may take $r_0=89$, because $r>4k/5\ge88$, and
$\upsilon(89)=(1067/1068)\cdot(335/89)>3.76$. With the bound $2e$ of
Lemma~\ref{lem:resolventmass} for the sum of the $G_i$ we get
$\Omega\le2.1/3.76<0.56$, and $n\psi/2+\Omega<0.395+0.56<1$.

Suppose now that $82\le k\le109$. Then $(n\psi/2)^2\le16.53/82$,
hence $n\psi/2\le0.455$. We may take $r_0=66$, because $r>4k/5\ge65.6$, and
$\upsilon(66)=(791/792)\cdot(81/22)>3.677$. The bound $2e$ for the sum of the
$G_i$ is too weak here, and we bound the sum more
carefully. Let $J=r-m_0+1$, so that the sum in question is
$G_0+G_1+\dots+G_{J-2}$. From $4k/5<r\le k-1$ and $3k/5-1<m_0\le3k/5$ we get
$k/5<r-m_0<2k/5$, and hence $18\le J\le44$, because $82\le k\le109$. For
$i\ge17$ we have $1-1/(12(i+1))\ge215/216$ and
$4-21/(i+1)\ge17/6$, so that Lemma~\ref{lem:resolventlower} and the bound
$c_{\mathrm S}\ge0.39894$ give
\[
G_i\ge\frac{e\,c_{\mathrm S}}{(i+1)^{3/2}}\cdot\frac{215}{216}\cdot\frac{17}{6}
\ge\frac{1.12\,e}{(i+1)^{3/2}}\qquad(i\ge17).
\]
Inserting here the bound
$1/(j\sqrt j)\ge2\,\bigl(1/\sqrt j-1/\sqrt{j+1}\bigr)$ of
Section~\ref{sec:prelim} with $j=i+1$ and summing over
$J-1\le i\le16J-2$, where $i\ge17$, we get
\[
\sum_{i=J-1}^{16J-2}G_i
\ge2\cdot1.12\,e\,\Bigl(\frac{1}{\sqrt J}-\frac{1}{\sqrt{16J}}\Bigr)
=\frac{1.68\,e}{\sqrt J}.
\]
The right-hand side is at least $1.68\,e/6.64\ge0.253\,e$, because
$J\le44\le6.64^2$. On the other hand $G_0+G_1+\dots+G_{16J-2}\le2e$ by
Lemma~\ref{lem:resolventmass}. Subtracting, we obtain
$G_0+G_1+\dots+G_{J-2}\le(2-0.253)\,e=1.747\,e$. Now
\eqref{eq:far-feedback} gives $\Omega\le1.05\cdot1.747/3.677<0.5$, and
$n\psi/2+\Omega<0.455+0.5<1$.

In each of the last two ranges of $r$ we have shown that
$n\psi/2+\Omega<1$, so that \eqref{eq:far-second} gives $\delta_n>0$.
\end{proof}

\subsection{Columns with an early first positive
row}\label{subsec:far-persistence}

We now combine Propositions~\ref{prop:nearrows} and~\ref{prop:farrows} by an
induction along the column.

\begin{proposition}\label{prop:staypositive}
Let $k\ge82$ and $r_q\ge1$ be integers with $3\,r_q\le k$. Suppose that
$d_{j,k}\le0$ for every $j$ with $k\le j<k+r_q$, and that $d_{k+r_q,k}>0$.
Then $d_{n,k}>0$ for every $n$ with $k+r_q\le n\le2k-1$.
\end{proposition}

\begin{proof}
We prove $d_{n,k}>0$ by induction on $n$. For $n=k+r_q$ this is a hypothesis.
Let $k+r_q<n\le2k-1$, and suppose that $d_{j,k}>0$ for every $j$ with
$k+r_q\le j<n$. Write $n=k+r$, so that $r_q<r\le k-1$. If $100\,r>47\,k$,
then Proposition~\ref{prop:farrows}, applied with these $r$ and $r_q$, gives
$d_{n,k}>0$. If $100\,r\le47\,k$, we apply Proposition~\ref{prop:nearrows}
with $t=r-r_q\ge1$. Its hypotheses hold: we have $100\,(r_q+t)\le47\,k$
and $d_{j,k}\le0$ for $k\le j<k+r_q$, and the positivity of $d_{j,k}$ for
$k+r_q\le j<k+r_q+t=n$ is the induction hypothesis. It gives
$d_{k+r_q+t,k}=d_{n,k}>0$.
\end{proof}

\section{Weighted sums over compositions}\label{sec:compositions}
In Section~\ref{sec:frozen} the normalised step at a position of the first
layer is expressed through the chain sums $\Lambda_m(y;N)$ of the
introduction. A chain sum is a sum over the compositions of the integer $m$
in which a part of size $p$ has weight $a_p$ and a cut at position $s$ has
weight $y/(N-s)$. Here we define such sums for an arbitrary weight $f(s)$
of the cuts, since the later sections compare the weight $y/(N-s)$ with
constant weights, for which the sums can be evaluated, and we prove the
recursions and identities for them on which the proofs of
Propositions~\ref{prop:curveincrease} and~\ref{prop:diagonalstep} rest.

\subsection{Compositions and cut profiles}\label{subsec:compositions-weights}

Let $m\ge1$ be an integer. A \emph{composition} of $m$ is a finite sequence
$\mathbf p=(p_1,\dots,p_q)$ of positive integers with $p_1+\dots+p_q=m$. The numbers
$p_i$ are its \emph{parts}, and $q\ge1$ is the number of its parts. The
\emph{cut points} of $\mathbf p$ are the $q-1$ partial sums $s_i=p_1+\dots+p_i$ with
$1\le i\le q-1$. With $s_0=0$ and $s_q=m$ we have $0=s_0<s_1<\dots<s_q=m$ and
$p_i=s_i-s_{i-1}$. We think of $\mathbf p$ as a way of cutting a segment of length
$m$ at some of its interior integer points $1,\dots,m-1$ into pieces of the
lengths $p_1,\dots,p_q$. A composition is determined by the set of its cut
points, and every subset of $\{1,\dots,m-1\}$ is the set of cut points of
exactly one composition of $m$. The integer $0$ has no composition.

A \emph{cut profile} is a function $f\colon\mathbb N\to\mathbb R$, where
$\mathbb N=\{0,1,2,\dots\}$; the number $f(s)$ is the weight of a cut at
position $s$. The \emph{weight} of the composition $\mathbf p$ for the cut
profile
$f$ is
\[
\operatorname{wt}_f(\mathbf p)=\prod_{i=1}^{q}a_{p_i}\ \prod_{i=1}^{q-1}f(s_i),
\]
the product of the Cayley coefficients of its parts and of the values of $f$
at its cut points. These values, one for each cut point of $\mathbf p$, are
called the \emph{cut factors} of $\mathbf p$. For integers $q\ge0$ and
$m\ge0$ we write $V_q(m;f)$ for
the sum of the weights $\operatorname{wt}_f(\mathbf p)$ of the compositions $\mathbf p$ of
$m$ with exactly $q$ parts, so that
\[
V_q(m;f)=\sum_{0=s_0<s_1<\dots<s_q=m}\ \prod_{i=1}^{q}a_{s_i-s_{i-1}}\
\prod_{i=1}^{q-1}f(s_i)\qquad(q\ge1),
\]
and we write
\[
\Lambda_m[f]=\sum_{q=1}^{m}V_q(m;f)
\]
for the total weight of all compositions of $m$. We have
$V_q(m;f)=0$ for $q=0$, for $m=0$ and for $q>m$,
since a composition of $m$ has at least one part and at most $m$ parts, and
$\Lambda_0[f]=0$. The only composition with one part is $(m)$, which has no
cut point, so that $V_1(m;f)=a_m$ for every $m\ge0$ and every profile; for
$m=0$ this holds because $a_0=0$. For example, $a_1=a_2=1$ and $a_3=3/2$, and
the compositions $(3)$, $(1,2)$, $(2,1)$ and $(1,1,1)$ of $3$ have the sets of
cut points $\emptyset$, $\{1\}$, $\{2\}$ and $\{1,2\}$. Hence
$V_1(3;f)=3/2$, $V_2(3;f)=f(1)+f(2)$ and $V_3(3;f)=f(1)\,f(2)$, and
\[
\Lambda_3[f]=\tfrac32+f(1)+f(2)+f(1)\,f(2).
\]

The following properties are immediate from the definition. Let $f$ and $g$
be cut profiles and let $m\ge0$. The weight of a composition of $m$ involves
the values of the profile at the points $1,\dots,m-1$ only. Hence, if
$f(s)=g(s)$ for $0<s<m$, then $V_q(m;f)=V_q(m;g)$ for every $q$, and
$\Lambda_m[f]=\Lambda_m[g]$. If $f(s)\ge0$ for $0<s<m$, then
$V_q(m;f)\ge0$ for every $q$. If $f(s)>0$ for $0<s<m$, then $V_q(m;f)>0$
for $1\le q\le m$. If $0\le f(s)\le g(s)$ for $0<s<m$, then
$\Lambda_m[f]\le\Lambda_m[g]$.

When $m\ge1$, $f$ is nonnegative at the points $1,\dots,m-1$ and
$\Lambda_m[f]>0$, the numbers
$V_q(m;f)/\Lambda_m[f]$ with $1\le q\le m$ are nonnegative and add up to
one. They are the probabilities with which a composition of $m$ that is
drawn with probability proportional to its weight has $q$ parts. Several
statements of this section describe this law of the number of parts; we use
the probabilistic language only to say what a sum means, and every such
statement is an identity between finite sums.

The identities of this section rest on one observation: a composition can be
split at each of its cut points. For a cut profile $f$ and an integer
$b\ge0$ we write $f^{+b}$ for the \emph{shifted profile} $s\mapsto f(b+s)$.
Let $\mathbf p=(p_1,\dots,p_q)$ be a composition of $m$ and let $b=s_j$ be one of its
cut points. Then $\mathbf p'=(p_1,\dots,p_j)$ is a composition of $b$, and
$\mathbf p''=(p_{j+1},\dots,p_q)$ is a composition of $m-b$. The cut points of $\mathbf p$
are the cut points of $\mathbf p'$, the point $b$, and the points $b+s$ for the cut
points $s$ of $\mathbf p''$.
Conversely, for $0<b<m$ every pair $(\mathbf p',\mathbf p'')$ of a composition of $b$ and a
composition of $m-b$ arises in this way from exactly one composition $\mathbf p$ of
$m$ that has a cut point at $b$, namely from the sequence $\mathbf p'$ followed by
the sequence $\mathbf p''$. If $\mathbf p'$ has $j$ parts and $\mathbf p''$ has $r$ parts, then $\mathbf p$
has $j+r$ parts, and $b$ is its $j$-th cut point.

Let $q\ge1$. A composition of $m$ with $q+1$ parts has a last cut point
$s=s_q$. Split at this point, it consists of a composition $\mathbf p'$ of $s$
with $q$ parts and of the composition $\mathbf p''=(m-s)$ with one part, which
has no cut point. The cut factors of the composition are therefore those of
$\mathbf p'$ together with $f(s)$, and its last part contributes $a_{m-s}$.
Sorting these compositions by their last cut point gives, for all $q\ge1$
and $m\ge0$,
\begin{equation}\label{eq:compositions-lastpart}
V_{q+1}(m;f)=\sum_{0<s<m}V_q(s;f)\,f(s)\,a_{m-s}.
\end{equation}
Together with the values $V_0(m;f)=0$ and
$V_1(m;f)=a_m$, the recursion \eqref{eq:compositions-lastpart} determines
all the numbers $V_q(m;f)$. The following two lemmas are recursions for the
total weights, by the last cut and by the first part.

\begin{lemma}\label{lem:lastcut}
For every function $f\colon\mathbb N\to\mathbb R$ and every integer $m\ge0$,
\[
\Lambda_m[f]=a_m+\sum_{0<s<m}\Lambda_s[f]\,f(s)\,a_{m-s}.
\]
\end{lemma}

\begin{proof}
For $m=0$ both sides are zero. Let $m\ge1$. Since $V_1(m;f)=a_m$, the
recursion \eqref{eq:compositions-lastpart} gives
\begin{align*}
\Lambda_m[f]&=a_m+\sum_{q=1}^{m-1}V_{q+1}(m;f)\\
&=a_m+\sum_{0<s<m}\Bigl(\sum_{q=1}^{m-1}V_q(s;f)\Bigr)f(s)\,a_{m-s}.
\end{align*}
For $0<s<m$ the inner sum is $\Lambda_s[f]$, because $s\le m-1$ and
$V_q(s;f)=0$ for $q>s$.
\end{proof}

\begin{lemma}\label{lem:firstpart}
For every function $f\colon\mathbb N\to\mathbb R$ and every integer $m\ge0$,
\[
\Lambda_m[f]=a_m+\sum_{0<p<m}a_p\,f(p)\,\Lambda_{m-p}[f^{+p}].
\]
\end{lemma}

\begin{proof}
Let $\mathbf p$ be a composition of $m$, let $b$ be one of its cut points, and
let $\mathbf p'$ and $\mathbf p''$ be the compositions of $b$ and of $m-b$ into
which $b$ splits it. The cut points of $\mathbf p$ are the cut points of
$\mathbf p'$, the point $b$, and the points $b+s$ for the cut points $s$ of
$\mathbf p''$. Consequently
\begin{equation}\label{eq:compositions-split}
\operatorname{wt}_f(\mathbf p)=\operatorname{wt}_f(\mathbf p')\,f(b)\,
\operatorname{wt}_{f^{+b}}(\mathbf p'').
\end{equation}
By the correspondence between the pairs $(\mathbf p',\mathbf p'')$ and the
compositions of $m$ with a cut point at $b$ described above, we may sum
\eqref{eq:compositions-split} over the pairs $(\mathbf p',\mathbf p'')$ with given numbers
$j\ge1$ and $r\ge1$ of parts and obtain, for $0<b<m$,
\begin{equation}\label{eq:compositions-parts}
\sum_{\mathbf p}\operatorname{wt}_f(\mathbf p)=V_j(b;f)\,f(b)\,V_r(m-b;f^{+b}),
\end{equation}
where $\mathbf p$ runs through the compositions of $m$ with $j+r$ parts whose $j$-th
cut point is $b$. Let $q\ge1$. A composition of $m$ with $q+1$ parts has a
first cut point $p=s_1$. We sort these compositions by their first cut
point, for which $j=1$, $r=q$ and $V_1(p;f)=a_p$. This gives, for every
$m\ge0$,
\begin{equation}\label{eq:compositions-firstpart}
V_{q+1}(m;f)=\sum_{0<p<m}a_p\,f(p)\,V_q(m-p;f^{+p}).
\end{equation}

For $m=0$ both sides of the assertion are zero. For $m\ge1$ we argue as in
the proof of
Lemma~\ref{lem:lastcut}, with \eqref{eq:compositions-firstpart} instead of
\eqref{eq:compositions-lastpart}:
\[
\Lambda_m[f]=a_m+\sum_{0<p<m}a_p\,f(p)\sum_{q=1}^{m-1}V_q(m-p;f^{+p}),
\]
and the inner sum is $\Lambda_{m-p}[f^{+p}]$, because $m-p\le m-1$.
\end{proof}

The recursion by the last cut keeps the
profile, whereas the recursion by the first part changes it: the cut points
of the composition that remains lie at positions $p+s$ of the original
segment.

The cut profile that the parking problem produces is the following. For real
numbers $y$ and $N$ the \emph{threshold profile} $\lambda^{y,N}$ is the cut
profile
\[
\lambda^{y,N}(s)=\frac{y}{N-s}\qquad(s\ne N).
\]
If $N$ is a nonnegative integer, the quotient is not defined at $s=N$, and
we put $\lambda^{y,N}(N)=0$, so that the profile is defined at every point
of $\mathbb N$ for all real numbers $y$ and $N$. The total weight for the
threshold profile is written
\[
\Lambda_m(y;N)=\Lambda_m[\lambda^{y,N}]
\]
and is called a \emph{chain sum}, and $y$ is called its \emph{threshold}.
By the cut factors of the chain sum $\Lambda_m(y;N)$ we mean those of the
compositions of $m$ for the threshold profile, the numbers
$\lambda^{y,N}(s)$ with $0<s<m$.

A chain sum is thus defined for all real numbers $y$ and $N$ and every
integer $m\ge0$. Suppose first that $m\le N$. Then every cut point of a
composition of $m$ is an integer $s$ with $0<s<m\le N$.
The cut factors of the chain sum are therefore the numbers
$y/(N-s)$ with $0<s<m$, their denominators are positive, and the value that
we gave to the profile at the point $N$ does not occur. For $y>0$ the weight
$y/(N-s)$ of a cut increases with $s<N$. For $m\ge1$ the chain sum is then
the sum that was displayed in the introduction.
If $m>N$, on the other hand, a composition of $m$ can have cut
points $s\ge N$, at which the denominator $N-s$ is zero or negative. The
identities below that are stated for all real $N$ and every $m$ include this
case, and in them the profile has the value $0$ at the point $N$.

Chain sums arise in Section~\ref{sec:frozen}, and the way in which they
arise there is the reason for the form of the threshold profile. For an
integer $k\ge0$ and a threshold $j\ge1$ that is one of the integers $k$ and
$k+1$, the linear relation of Proposition~\ref{prop:chainforward} is applied
there repeatedly to the normalised counts $z^{(j)}_n$ with $0\le n\le2k+1$,
and with $N=n+1$ this expresses the count, apart from one term, as a sum
over the decreasing chains $N>t_1>\dots>t_i>k$ of integers. Such a chain
corresponds to a composition of the integer $m=N-t_i$, in which the integers
$t_l$ of the chain other than the last give the cut points $s_l=N-t_l$, and
the factor $j/t_l$ with which such an integer occurs is a cut factor of the
form $y/(N-s)$, the value of the threshold profile $\lambda^{j,N}$ at the
cut point $s_l$. Since $t_i>0$, the chain sums that arise in this way have
$m<N$.

The properties listed above take the following form for the threshold
profile. If $y\ge0$, then $\lambda^{y,N}(s)\ge0$ for every integer $s$ with
$0\le s\le N$. Hence
for $y\ge0$ and an integer $m$ with $1\le m\le N$ the numbers
$V_q(m;\lambda^{y,N})$ are nonnegative and $\Lambda_m(y;N)>0$.
A composition with $q$ parts has $q-1$ cut points, and
$\lambda^{y,N}(s)=y\,\lambda^{1,N}(s)$ at every point $s$ of
$\mathbb N$, the point $s=N$ included. Hence
\begin{align}
V_q(m;\lambda^{y,N})&=y^{q-1}\,V_q(m;\lambda^{1,N})\qquad(q\ge1),
\label{eq:compositions-homogeneity}\\
\Lambda_m(y;N)&=\sum_{q=1}^{m}y^{q-1}\,V_q(m;\lambda^{1,N})\notag
\end{align}
for all real $y$ and $N$ and every $m\ge0$. The chain sum is thus a
polynomial in the threshold, in which the compositions with $q$ parts
contribute to the coefficient of $y^{q-1}$; the coefficients depend on $m$
and $N$ only, and they are nonnegative when $m\le N$. Finally, the shift of
a threshold profile is a threshold profile,
$(\lambda^{y,N})^{+p}=\lambda^{y,N-p}$. Lemma~\ref{lem:firstpart} therefore
reads
\[
\Lambda_m(y;N)=a_m+\sum_{0<p<m}a_p\,\frac{y}{N-p}\,\Lambda_{m-p}(y;N-p)
\qquad(m\le N).
\]
Removing a first part of size $p$ lowers the integer $N$ to $N-p$, and in
this form the recursion by the first part is one step of the iteration
described above.

\subsection{Marked cut points}\label{subsec:compositions-marked}

In the estimates that lead to Proposition~\ref{prop:curveincrease} a chain
is separated at a cut point of its composition, that is, at one of its
integers other than the last, and the two pieces are estimated separately.
The sums that arise in this way are the following.
For a cut profile $f$ and an integer $m\ge0$ we put
\begin{align*}
\Lambda^{(1)}_m[f]&=\sum_{0<b<m}\Lambda_b[f]\,f(b)\,\Lambda_{m-b}[f^{+b}],\\
\Lambda^{(2)}_m[f]&=\sum_{0<b<m}\Lambda^{(1)}_b[f]\,f(b)\,\Lambda_{m-b}[f^{+b}].
\end{align*}
A composition together with one of its cut points is a composition with one
\emph{marked} cut point, and $\Lambda^{(1)}_m[f]$ is the total weight of
the compositions of $m$ with one marked cut point. The second definition
repeats the first one level higher, and $\Lambda^{(2)}_m[f]$ is the total
weight of the compositions of $m$ with two marked cut points. Lemmas~\ref{lem:onemarked}
and~\ref{lem:twomarked} below make these descriptions precise. The first sum
is empty for $m\le1$, so that
$\Lambda^{(1)}_0[f]=\Lambda^{(1)}_1[f]=0$, and therefore
$\Lambda^{(2)}_m[f]=0$ for $m\le2$. For the threshold profile we write
$\Lambda^{(1)}_m(y;N)=\Lambda^{(1)}_m[\lambda^{y,N}]$.

For the proof of Proposition~\ref{prop:curvaturepair} we need a recursion by
the last cut for the sums with marked cut points as well. When the last
part of a composition with two marked cut points is removed, its last cut
point $s$ was either the larger of the two marked points, and a composition
of $s$ with one marked cut point remains, or it was not marked, and a
composition of $s$ with two marked cut points remains.

\begin{lemma}\label{lem:markedrecursion}
For every function $f\colon\mathbb N\to\mathbb R$ and every integer $m\ge0$,
\[
\Lambda^{(2)}_m[f]=\sum_{0<s<m}\bigl(\Lambda^{(1)}_s[f]+\Lambda^{(2)}_s[f]\bigr)
\,f(s)\,a_{m-s}.
\]
\end{lemma}

\begin{proof}
Let $0<b<m$. Lemma~\ref{lem:lastcut}, applied to the profile $f^{+b}$ and to
the integer $m-b$, gives, after a shift of the index of summation by $b$,
\[
\Lambda_{m-b}[f^{+b}]
=a_{m-b}+\sum_{b<s<m}\Lambda_{s-b}[f^{+b}]\,f(s)\,a_{m-s}.
\]
We multiply by $\Lambda^{(1)}_b[f]\,f(b)$ and sum over $0<b<m$. The left
side becomes $\Lambda^{(2)}_m[f]$. On the right side the terms $a_{m-b}$
contribute the terms of the asserted sum that contain $\Lambda^{(1)}$, and
the other terms form a double sum over $0<b<s<m$. Summed first over $b$, it
is
\[
\sum_{0<s<m}\Bigl(\sum_{0<b<s}\Lambda^{(1)}_b[f]\,f(b)\,\Lambda_{s-b}[f^{+b}]
\Bigr)f(s)\,a_{m-s},
\]
and the sum in brackets is $\Lambda^{(2)}_s[f]$ by definition.
\end{proof}

\subsection{Compositions without cut factors}
\label{subsec:compositions-constant}

We now turn to the constant cut profile $\mathbf 1$ with the value $1$, for
which the cuts carry no factor. For
integers $m\ge0$ and $q\ge0$ let
$T_{m,q}=V_q(m;\mathbf 1)$. Thus $T_{m,0}=0$ and
\[
T_{m,q}=\sum_{\substack{p_1+\dots+p_q=m\\ p_i\ge1}}a_{p_1}\cdots a_{p_q}
\qquad(q\ge1).
\]
By Section~\ref{subsec:compositions-weights} we have $T_{m,1}=a_m$,
$T_{m,q}\ge0$ for all $m$ and $q$, $T_{m,q}>0$ for $1\le q\le m$, and
$T_{m,q}=0$ for $q>m$. For a real number $h$, regarded as a constant cut
profile, and every integer $m\ge0$, the total weight $\Lambda_m[h]$ is the
sum of $T_{m,q}\,h^{q-1}$ over $1\le q\le m$. The numbers $T_{m,q}$ are in this
sense the weights of the model in which a composition is weighted by its
parts alone, and Sections~\ref{sec:gain} and~\ref{sec:growth} compare the
weights of the threshold profile with them.

The numbers $T_{m,q}$ can be evaluated in closed form by the Lagrange
inversion formula. We state the result, which
Sections~\ref{sec:birthday}, \ref{sec:gain} and~\ref{sec:growth} use, with
the probabilities of the birthday problem. For
integers $m\ge1$ and $j\ge0$ let
\[
B_{m,j}=\prod_{i=0}^{j-1}\Bigl(1-\frac im\Bigr)
=\frac{m(m-1)\cdots(m-j+1)}{m^{j}}.
\]
The number $B_{m,j}$, the \emph{birthday tail}, is the probability that $j$
independent draws, each uniformly distributed over $m$ values, are pairwise
distinct. Clearly $B_{m,0}=B_{m,1}=1$ and
\begin{equation}\label{eq:compositions-tail}
B_{m,j+1}=B_{m,j}\Bigl(1-\frac jm\Bigr)\qquad(j\ge0).
\end{equation}
For $j\le m$ we have $B_{m,j}=m!/((m-j)!\,m^{j})>0$, and for $j\ge m+1$ the
product contains the factor with $i=m$, which is zero, so that $B_{m,j}=0$;
in particular $B_{m,m+1}=0$, and $B_{m,j}\ge0$ for all $j$.

\begin{proposition}\label{prop:partslaw}
Let $m\ge0$ and $q\ge1$ be integers. If $m\ge1$, then
\[
T_{m,q}=a_m\,q\,B_{m,q}.
\]
If $m=0$, then $T_{0,q}=0$.
\end{proposition}

\begin{proof}
The integer $0$ has no composition, so that $T_{0,q}=0$. Let $m\ge1$. If
$q>m$, then $T_{m,q}=0$ and $B_{m,q}=0$. Let $1\le q\le m$. The tree series
$T=\sum_{j\ge1}a_j\,X^j$ of Section~\ref{sec:series} (the letter $T$ without
indices denotes this series) has no constant term. The coefficient of $X^m$
in its $q$-th power is therefore the sum of the products
$a_{p_1}\cdots a_{p_q}$ over the sequences $(p_1,\dots,p_q)$ of positive
integers with sum $m$, which is $T_{m,q}$. Since the tree series satisfies
the equation $T=Xe^{T}$, the Lagrange inversion formula
\cite[Theorem~A.2]{flajoletsedgewick2009} applies to it with the series
$\phi(u)=e^{u}$, whose constant term is not zero. In its form for the powers
of the solution of such an equation, the formula gives
\[
[X^m]\,T^q=\frac qm\,[u^{m-q}]\,\phi(u)^m=\frac qm\,[u^{m-q}]\,e^{mu},
\]
where $[u^{m-q}]$ denotes the coefficient of $u^{m-q}$ in a power series in
$u$. This coefficient of $e^{mu}$ is $m^{m-q}/(m-q)!$. Since
$a_m=m^{m-1}/m!$ and $B_{m,q}=m!/((m-q)!\,m^{q})$, it follows that
\[
T_{m,q}=\frac qm\cdot\frac{m^{m-q}}{(m-q)!}=\frac{q\,m^{m-q-1}}{(m-q)!}
=a_m\,q\,B_{m,q}.\qedhere
\]
\end{proof}

For instance, the proposition gives $T_{m,m}=1$, the weight of the
composition of $m$ into $m$ parts of size $1$, and for $m=3$ it gives
$3/2$, $2$ and $1$, the values of the example of
Section~\ref{subsec:compositions-weights} for the profile $\mathbf 1$.

Let $m\ge1$. By
\eqref{eq:compositions-tail}, $(q/m)\,B_{m,q}=B_{m,q}-B_{m,q+1}$, and
therefore
\begin{align}
T_{m,q}&=m\,a_m\,\bigl(B_{m,q}-B_{m,q+1}\bigr)\qquad(q\ge1),\notag\\
\sum_{q=1}^{m}T_{m,q}&=m\,a_m\,\bigl(B_{m,1}-B_{m,m+1}\bigr)=m\,a_m,
\label{eq:compositions-law}
\end{align}
the sum being a telescoping one. The \emph{birthday law of $m$} is the law
of a random variable $Q$ with $\Pr(Q\ge j)=B_{m,j}$ for $j\ge1$. It is the
law of the number of
draws that precede the first repetition when one draws independently and
uniformly from $m$ values, since that number is at least $j$ exactly when
the first $j$ draws are pairwise distinct. It is carried by the integers
$1,\dots,m$, and its mean, which is the sum of $\Pr(Q\ge j)$
over $j\ge1$, is $B_{m,1}+\dots+B_{m,m}$. Section~\ref{sec:birthday}
studies this mean.

\subsection{The number of cut points}\label{subsec:compositions-counting}

We return to an arbitrary cut profile and evaluate the sums with marked cut
points of Section~\ref{subsec:compositions-marked} by the number of parts. A
composition with $q$ parts has $q-1$ cut points, so that there are $q-1$
ways to mark one of them.

\begin{lemma}\label{lem:onemarked}
For every function $f\colon\mathbb N\to\mathbb R$ and every integer $m\ge0$,
\[
\Lambda^{(1)}_m[f]=\sum_{q=1}^{m}(q-1)\,V_q(m;f).
\]
\end{lemma}

\begin{proof}
Let $0<b<m$. We sum \eqref{eq:compositions-parts}, which was shown in the
proof of Lemma~\ref{lem:firstpart}, over $j\ge1$ and $r\ge1$ and obtain
\[
\sum_{\mathbf p}\operatorname{wt}_f(\mathbf p)=\Lambda_b[f]\,f(b)\,\Lambda_{m-b}[f^{+b}],
\]
where $\mathbf p$ runs through all compositions of $m$ that have a cut point at
$b$. So the term of index $b$ in the definition of $\Lambda^{(1)}_m[f]$
is the sum of the weights
$\operatorname{wt}_f(\mathbf p)$ of the compositions $\mathbf p$ of $m$ that have a cut
point at $b$. In the sum over $b$ every composition of $m$ therefore occurs
once for each of its cut points, that is, $q-1$ times if it has $q$ parts.
\end{proof}

If $m\ge1$, $f$ is nonnegative at the points $1,\dots,m-1$ and
$\Lambda_m[f]>0$, the lemma
says that $\Lambda^{(1)}_m[f]/\Lambda_m[f]$ is the mean number of cut points
of a composition of $m$ that is drawn with probability proportional to its
weight.

A composition with $q$ parts has $(q-1)(q-2)/2$ pairs of cut points, and the
sum with two marked cut points counts them.

\begin{lemma}\label{lem:twomarked}
For every function $f\colon\mathbb N\to\mathbb R$ and every integer $m\ge0$,
\[
\Lambda^{(2)}_m[f]=\sum_{q=1}^{m}\frac{(q-1)(q-2)}{2}\,V_q(m;f).
\]
\end{lemma}

\begin{proof}
Let $0<b<m$. By Lemma~\ref{lem:onemarked}, applied to the integer $b$, the
number $\Lambda^{(1)}_b[f]$ is the sum of the weights
$\operatorname{wt}_f(\mathbf p')$ of the compositions $\mathbf p'$ of $b$, each multiplied
by the number of cut points of $\mathbf p'$. The number $\Lambda_{m-b}[f^{+b}]$ is
the sum of the weights $\operatorname{wt}_{f^{+b}}(\mathbf p'')$ of the compositions
$\mathbf p''$ of $m-b$. The pairs $(\mathbf p',\mathbf p'')$ correspond to the compositions $\mathbf p$ of
$m$ that have a cut point at $b$, the weights are related by
\eqref{eq:compositions-split}, which was shown in the proof of
Lemma~\ref{lem:firstpart}, and the cut points of $\mathbf p'$ are the cut points
of $\mathbf p$ that are smaller than $b$. Hence the term of index $b$ in the
definition of $\Lambda^{(2)}_m[f]$ is the sum of the weights
$\operatorname{wt}_f(\mathbf p)$ of the compositions $\mathbf p$ of $m$ that have a cut
point at $b$, each multiplied by the number of its cut points that are
smaller than $b$. In the sum over $b$ every composition $\mathbf p$ of $m$ therefore
occurs with a multiplicity which is the number of pairs $b'<b$ of cut points
of $\mathbf p$, and a composition with $q$ parts has $(q-1)(q-2)/2$ such pairs.
\end{proof}

The two lemmas together give the second factorial moment of the number of
parts. It is needed for the lower bound of the curvature term in
Proposition~\ref{prop:curvaturepair}: for the threshold profile the
compositions with $q$ parts contribute to the coefficient of $y^{q}$ in
$y\,\Lambda_m(y;N)$, and a second difference of the power $y^{q}$ brings
the factor $q(q-1)$.

\begin{corollary}\label{cor:partsmoment}
For every function $f\colon\mathbb N\to\mathbb R$ and every integer $m\ge0$,
\[
\sum_{q=1}^{m}q\,(q-1)\,V_q(m;f)=2\,\Lambda^{(1)}_m[f]+2\,\Lambda^{(2)}_m[f].
\]
\end{corollary}

\begin{proof}
By Lemmas~\ref{lem:onemarked} and~\ref{lem:twomarked} the right side is the
sum of $\bigl(2(q-1)+(q-1)(q-2)\bigr)\,V_q(m;f)$ over
$1\le q\le m$, and $2(q-1)+(q-1)(q-2)=q\,(q-1)$.
\end{proof}

The identity $q(q-1)/2=(q-1)+(q-1)(q-2)/2$ behind the corollary has a direct
meaning. The $q-1$ cut points of a composition of $m$ with $q$ parts and the
end point $m$ of the segment are $q$ points, and a pair of these points
consists either of the end point and a cut point or of two cut points. So
$\Lambda^{(1)}_m[f]+\Lambda^{(2)}_m[f]$ counts every composition with its
weight times the number of pairs among its $q$ points.

\subsection{The difference $\nu_m-\Lambda_m[f]$}
\label{subsec:compositions-violation}

The inequalities of Section~\ref{sec:pair} bound a total weight
$\Lambda_m[f]$ from below. The proof of
Proposition~\ref{prop:curveincrease} also needs an upper bound, for
the chain sum on the left side of the criterion of
Proposition~\ref{prop:criterion}. The identity of this subsection makes such
a bound possible. It expresses the amount by which $\Lambda_m[f]$ falls
short of the number $\nu_m=m^m/m!$ as a combination of the totals
$\Lambda_s[f]$ of the smaller integers $s<m$, and for a profile with values
in the interval $[0,1]$ the coefficients of this combination are
nonnegative. Lower bounds for the $\Lambda_s[f]$ then give an upper bound
for $\Lambda_m[f]$, and this is how Proposition~\ref{prop:boundaryupper} is
proved.

\begin{lemma}\label{lem:firstviolation}
For every function $f\colon\mathbb N\to\mathbb R$ and every integer $m\ge1$,
\[
\nu_m-\Lambda_m[f]=\sum_{0<s<m}\nu_{m-s}\,\bigl(1-f(s)\bigr)\,\Lambda_s[f].
\]
\end{lemma}

\begin{proof}
Let $n\ge1$. By Lemma~\ref{lem:nurenewal},
\[
\sum_{j=0}^{n}a_j\,\nu_{n-j}=\nu_n.
\]
In this sum the term $j=0$ vanishes, the term $j=n$ is $a_n\,\nu_0=a_n$, and
with $s=n-j$ the other terms give
\begin{equation}\label{eq:compositions-renewal}
\nu_n=a_n+\sum_{0<s<n}a_{n-s}\,\nu_s\qquad(n\ge1).
\end{equation}
Put $\varphi_s=(1-f(s))\,\Lambda_s[f]$ and $Y_m=\nu_m-\Lambda_m[f]$. We
prove
\[
Y_m=\sum_{0<s<m}\nu_{m-s}\,\varphi_s
\]
by induction on $m\ge1$, assuming it for all smaller positive integers.
Subtracting the recursion of
Lemma~\ref{lem:lastcut} from \eqref{eq:compositions-renewal} with $n=m$, we
get
\[
Y_m=\sum_{0<s<m}a_{m-s}\bigl(\nu_s-f(s)\,\Lambda_s[f]\bigr)
=\sum_{0<s<m}a_{m-s}\bigl(\varphi_s+Y_s\bigr).
\]
By the induction hypothesis, $Y_s$ is the sum of
$\nu_{s-t}\,\varphi_t$ over $0<t<s$, for every $s$ with $0<s<m$. We insert
this and collect the coefficient of each $\varphi_t$:
\[
Y_m=\sum_{0<t<m}\varphi_t\,\Bigl(a_{m-t}+\sum_{t<s<m}a_{m-s}\,\nu_{s-t}
\Bigr).
\]
With $u=s-t$ the inner sum is the sum of $a_{m-t-u}\,\nu_u$ over
$0<u<m-t$, so that the expression in brackets is $\nu_{m-t}$ by
\eqref{eq:compositions-renewal} with $n=m-t\ge1$. This is the assertion.
\end{proof}

The identity of Lemma~\ref{lem:firstviolation} can be understood as follows.
For a composition with the cut
points $s_1<\dots<s_j$,
\[
1-\prod_{i=1}^{j}f(s_i)=\sum_{i=1}^{j}\bigl(1-f(s_i)\bigr)\prod_{l<i}f(s_l),
\]
a telescoping sum. In the term of index $i$ the cut points before $s_i$ keep
their factors, the point $s_i$ carries the factor $1-f(s_i)$, and the cut
points after $s_i$ carry the factor $1$. When $0\le f\le1$, one
may read $f(s)$ as the probability that the cut point $s$ passes a test,
independently of the other cut points. The identity then sorts the
compositions with a failing cut point by the first such point.

For the profile $\mathbf 1$ the lemma gives the total weight of all
compositions of $m$ when the cuts carry no factor.

\begin{corollary}\label{cor:unitprofile}
For every integer $m\ge1$,
\[
\Lambda_m[\mathbf 1]=\nu_m=\frac{m^m}{m!}.
\]
\end{corollary}

\begin{proof}
For the profile $f=\mathbf 1$ every term on the right side of
Lemma~\ref{lem:firstviolation} vanishes.
\end{proof}

Since $\Lambda_m[\mathbf 1]$ is the sum of $T_{m,q}$ over
$1\le q\le m$ and $\nu_m=m\,a_m$, the corollary says the same as
\eqref{eq:compositions-law}. With the
corollary the sums with marked cut points become convolutions for the
profile $\mathbf 1$. A shift does not change a constant profile, so that
\[
\Lambda^{(1)}_m[\mathbf 1]=\sum_{0<b<m}\nu_b\,\nu_{m-b},\qquad
\Lambda^{(2)}_m[\mathbf 1]=\sum_{0<b<m}\Lambda^{(1)}_b[\mathbf 1]\,\nu_{m-b}.
\]
Lemmas~\ref{lem:onemarked} and~\ref{lem:twomarked} and
Proposition~\ref{prop:partslaw} express the same two sums through the
numbers $B_{m,q}$, and Section~\ref{sec:birthday} draws the consequences.

\section{Convolutions of the Poisson mass and the birthday law}\label{sec:birthday}
The estimates that lead to Proposition~\ref{prop:curveincrease} are
expressed through two convolutions of the Poisson mass $P(m)$. We evaluate
them here in closed form through the mean $s_m$ of the birthday law of $m$
(Lemmas~\ref{lem:twofold} and~\ref{lem:threefold}), compute the mean
positions in them, and bound $s_m$ from both sides
(Lemmas~\ref{lem:birthdaylower} and~\ref{lem:birthdayupper}). The
estimates for Proposition~\ref{prop:curveincrease} use the lower bound,
through Corollary~\ref{cor:twofoldlower}; the upper bound is used only in
Section~\ref{sec:growth}, for
Proposition~\ref{prop:diagonalstep}. The two bounds are deduced from a
theorem of Flajolet, Grabner, Kirschenhofer and Prodinger
\cite[Theorem~7]{flajolet1995}.

For an integer $m\ge0$ the \emph{twofold convolution} and the
\emph{threefold convolution} of the Poisson mass are
\[
P_2(m)=\sum_{0<g<m}P(g)\,P(m-g),\qquad
P_3(m)=\sum_{0<g<m}P_2(g)\,P(m-g),
\]
where $g$ runs over the integers with $0<g<m$. The proof of
Proposition~\ref{prop:curveincrease} rests on the criterion of
Proposition~\ref{prop:criterion}, and Section~\ref{sec:pair} bounds the two
sides of that criterion by expressions in $P$, $P_2$ and $P_3$. The reason
is the following. By
Corollary~\ref{cor:unitprofile}, for $m\ge1$ the total weight
$\Lambda_m[\mathbf 1]$ of the compositions of $m$ for the cut profile
$\mathbf 1$ is $\nu_m=e^{m}P(m)$. The chain sums on the two sides of the
criterion have cut factors that are a little smaller than $1$, and they are
estimated by comparison with the profile $\mathbf 1$. Where a chain of
length $m$ is divided at one marked point or at two, its pieces contribute a
product of two or three of the numbers $\nu$ whose indices add up to $m$,
which is $e^{m}$ times a product of Poisson masses, and the sum over the
positions of the marked points is $e^{m}P_2(m)$ or $e^{m}P_3(m)$.

The sum
that defines $P_2(m)$ runs over the splittings $m=g+(m-g)$ of $m$ into two
positive parts, with the splitting point $g$; the terms $g=0$ and
$g=m$ of the full convolution, which are both equal to $P(m)$ because
$P(0)=1$, are left out. Clearly $P_2(0)=P_2(1)=0$ and $P_3(m)=0$ for
$m\le2$, and $P_2(m)\ge0$ and $P_3(m)\ge0$ for every $m$. Inserting the
definition of $P_2(g)$, as a sum over the integers $h$ with $0<h<g$, into
that of $P_3(m)$, and writing $g_1=h$, $g_2=g-h$ and $g_3=m-g$, we obtain
\begin{equation}\label{eq:birthday-triples}
P_3(m)=\sum_{g_1+g_2+g_3=m}P(g_1)\,P(g_2)\,P(g_3).
\end{equation}
Here and below, a sum over $g_1+\dots+g_q=m$ runs over the compositions
$(g_1,\dots,g_q)$ of $m$ into $q$ positive parts. The right-hand side of
\eqref{eq:birthday-triples} does not change when the three parts are
permuted. In particular, the threefold convolution can also be formed with
the single factor first:
\[
P_3(m)=\sum_{0<g<m}P(g)\,P_2(m-g).
\]

\subsection{Moments of the birthday law}\label{subsec:birthday-moments}

For an integer $m\ge1$ we put
\[
s_m=\sum_{j=1}^{m}B_{m,j},
\]
and we put $s_0=0$. Here $B_{m,j}$ is the birthday tail of
Section~\ref{sec:compositions}. A random variable $Q$ with the birthday law
of $m$, which is defined there, takes the values $1,\dots,m$, with
$\Pr(Q=q)=(q/m)\,B_{m,q}$ for $1\le q\le m$ and $\Pr(Q\ge j)=B_{m,j}$ for
$j\ge1$. The mean of a random variable with values in the
positive integers is the sum of its tail probabilities, so that $s_m$ is
the mean of $Q$, the \emph{mean of the birthday law} of $m$. The symbol
$s_m$ is not related to the cut points $s_0<s_1<\dots<s_q$ of a composition.

By Proposition~\ref{prop:partslaw} the compositions of $m$ with $q$ parts
have the total weight $T_{m,q}=a_m\,q\,B_{m,q}$. Sums over compositions
therefore lead to sums in which a function of $q$ is multiplied by
$q\,B_{m,q}$, which is $m$ times a point mass of the birthday law, whereas
$s_m$ is a sum of the tails $B_{m,q}$. The following summation by parts
passes from the first form to the second.

\begin{lemma}\label{lem:abelsummation}
Let $m\ge1$ be an integer and let $\psi$ be a real-valued function on the
nonnegative integers. Then the birthday tail $B_{m,q}$ satisfies
\[
\sum_{q=1}^{m}\psi(q)\,q\,B_{m,q}
=m\,\Bigl(\psi(0)+\sum_{q=1}^{m}\bigl(\psi(q)-\psi(q-1)\bigr)\,B_{m,q}\Bigr).
\]
\end{lemma}

\begin{proof}
For every integer $q\ge0$ the recursion $B_{m,q+1}=B_{m,q}\,(1-q/m)$ gives
$q\,B_{m,q}=m\,(B_{m,q}-B_{m,q+1})$. Hence the left-hand side of the claim
is $m$ times
\begin{align*}
\sum_{q=1}^{m}\psi(q)\,\bigl(B_{m,q}-B_{m,q+1}\bigr)
&=\sum_{q=1}^{m}\psi(q)\,B_{m,q}-\sum_{q=2}^{m+1}\psi(q-1)\,B_{m,q}\\
&=\sum_{q=1}^{m}\bigl(\psi(q)-\psi(q-1)\bigr)\,B_{m,q}\\
&\qquad+\psi(0)\,B_{m,1}-\psi(m)\,B_{m,m+1},
\end{align*}
and we have $B_{m,1}=1$ and $B_{m,m+1}=0$.
\end{proof}

For the constant function $\psi=1$, for
$\psi(j)=j$ and for $\psi(j)=j-1$ the lemma gives, for every integer
$m\ge1$,
\begin{equation}\label{eq:birthday-first}
\begin{gathered}
\sum_{q=1}^{m}q\,B_{m,q}=m,\qquad
\sum_{q=1}^{m}q^{2}\,B_{m,q}=m\,s_m,\\
\sum_{q=1}^{m}(q-1)\,q\,B_{m,q}=m\,(s_m-1).
\end{gathered}
\end{equation}
The second and the third factorial moment of the birthday law,
$\mathbb{E}\,Q(Q-1)$ and $\mathbb{E}\,Q(Q-1)(Q-2)$, are $1/m$ times the
sums of the next two
corollaries. The first of them enters the closed form of $P_3$ in
Lemma~\ref{lem:threefold}, and both are used in
Section~\ref{sec:growth}.

\begin{corollary}\label{cor:secondfactorial}
For every integer $m\ge1$,
\[
\sum_{q=1}^{m}q^{2}\,(q-1)\,B_{m,q}=2m\,(m-s_m).
\]
\end{corollary}

\begin{proof}
We apply Lemma~\ref{lem:abelsummation} to $\psi(j)=j(j-1)$, for which
$\psi(0)=0$ and $\psi(q)-\psi(q-1)=2(q-1)$. The sum of the claim is
therefore
\[
2m\sum_{q=1}^{m}(q-1)\,B_{m,q}=2m\,(m-s_m),
\]
by the first identity of \eqref{eq:birthday-first} and the definition
of $s_m$.
\end{proof}

\begin{corollary}\label{cor:thirdfactorial}
For every integer $m\ge1$,
\[
\sum_{q=1}^{m}q^{2}\,(q-1)(q-2)\,B_{m,q}=3m\,\bigl((m+2)\,s_m-3m\bigr).
\]
\end{corollary}

\begin{proof}
We apply Lemma~\ref{lem:abelsummation} to $\psi(j)=j(j-1)(j-2)$, for which
$\psi(0)=0$ and $\psi(q)-\psi(q-1)=3(q-1)(q-2)=3q^{2}-9q+6$. The sum of the
claim is therefore
\[
m\sum_{q=1}^{m}\bigl(3q^{2}-9q+6\bigr)\,B_{m,q}
=m\,\bigl(3m\,s_m-9m+6\,s_m\bigr),
\]
by the first two identities of \eqref{eq:birthday-first} and the definition
of $s_m$. The right-hand side is $3m\,((m+2)\,s_m-3m)$.
\end{proof}

\subsection{Closed forms of the two convolutions}
\label{subsec:birthday-closed}

We now evaluate the twofold convolution. The proof reads $e^{m}P_2(m)$
as the total weight, for the cut profile $\mathbf 1$, of the compositions of
$m$ with one marked cut
point, in which a composition with $q$ parts is counted $q-1$ times, and it
evaluates this weight by Proposition~\ref{prop:partslaw}.

\begin{lemma}\label{lem:twofold}
For every integer $m\ge1$, the twofold convolution of the Poisson mass and
the mean $s_m$ of the birthday law of $m$ satisfy
\[
P_2(m)=P(m)\,(s_m-1).
\]
\end{lemma}

\begin{proof}
Let $g$ and $m$ be integers with $0<g<m$. By the definition of the Poisson
mass, $P(g)\,P(m-g)=e^{-m}\,\nu_g\,\nu_{m-g}$, and by
Corollary~\ref{cor:unitprofile}, $\nu_g=\Lambda_g[\mathbf 1]$ and
$\nu_{m-g}=\Lambda_{m-g}[\mathbf 1]$. The constant profile $\mathbf 1$ is
equal to its shift $\mathbf 1^{+g}$, so that the sum of these products over
$g$ is the sum that defines $\Lambda^{(1)}_m[\mathbf 1]$ in
Section~\ref{sec:compositions}:
\begin{equation}\label{eq:birthday-marked}
e^{m}P_2(m)=\sum_{0<g<m}\Lambda_g[\mathbf 1]\,\Lambda_{m-g}[\mathbf 1]
=\Lambda^{(1)}_m[\mathbf 1].
\end{equation}
This holds for every integer $m\ge0$; for $m\le1$ the sums are empty. Let
now $m\ge1$. By Lemma~\ref{lem:onemarked} for the profile $\mathbf 1$, for
which $V_q(m;\mathbf 1)=T_{m,q}$, and by Proposition~\ref{prop:partslaw},
\[
\Lambda^{(1)}_m[\mathbf 1]=\sum_{q=1}^{m}(q-1)\,T_{m,q}
=a_m\sum_{q=1}^{m}(q-1)\,q\,B_{m,q}=m\,a_m\,(s_m-1),
\]
where the last step is the third identity of \eqref{eq:birthday-first}.
Since $m\,a_m=m^m/m!=\nu_m$ for $m\ge1$, we conclude with
\eqref{eq:birthday-marked} that
$P_2(m)=e^{-m}\,\nu_m\,(s_m-1)=P(m)\,(s_m-1)$.
\end{proof}

The threefold convolution is treated in the same way, with two marked cut
points instead of one. A composition with $q$ parts has
$(q-1)(q-2)/2$ pairs of cut points, which brings in the second factorial
moment of the birthday law. The result expresses $P_3$ in terms of $P$ and
$P_2$, and Sections~\ref{sec:sums} and~\ref{sec:curve} use it to remove
$P_3$ from their estimates.

\begin{lemma}\label{lem:threefold}
For every integer $m\ge1$, the threefold convolution of the Poisson mass
satisfies
\[
P_3(m)=(m-1)\,P(m)-2\,P_2(m).
\]
\end{lemma}

\begin{proof}
Let $g$ be an integer with $0<g<m$. For $0<h<g$ the definition of the
Poisson mass and Corollary~\ref{cor:unitprofile} give
$P(h)\,P(g-h)=e^{-g}\,\Lambda_h[\mathbf 1]\,\Lambda_{g-h}[\mathbf 1]$, and
since $\mathbf 1^{+h}=\mathbf 1$, summing over $h$ gives
$e^{g}P_2(g)=\Lambda^{(1)}_g[\mathbf 1]$. With this and with
Corollary~\ref{cor:unitprofile} at $m-g$,
\[
P_2(g)\,P(m-g)=e^{-m}\,\Lambda^{(1)}_g[\mathbf 1]\,
\Lambda_{m-g}[\mathbf 1],
\]
and, as in the proof of Lemma~\ref{lem:twofold}, the sum of the right-hand
sides over $g$ is $e^{-m}$ times the sum that defines
$\Lambda^{(2)}_m[\mathbf 1]$ in Section~\ref{sec:compositions}. Hence, by
Lemma~\ref{lem:twomarked} and Proposition~\ref{prop:partslaw},
\[
e^{m}P_3(m)=\Lambda^{(2)}_m[\mathbf 1]
=\sum_{q=1}^{m}\frac{(q-1)(q-2)}{2}\,T_{m,q}
=a_m\sum_{q=1}^{m}\frac{(q-1)(q-2)}{2}\,q\,B_{m,q}.
\]
Since $(q-1)(q-2)\,q/2=q^{2}(q-1)/2-(q-1)\,q$,
Corollary~\ref{cor:secondfactorial} and the third identity of
\eqref{eq:birthday-first} give
\[
\sum_{q=1}^{m}\frac{(q-1)(q-2)}{2}\,q\,B_{m,q}
=m\,(m-s_m)-m\,(s_m-1)=m\,(m+1-2s_m).
\]
With $m\,a_m=\nu_m$ we obtain
\[
P_3(m)=P(m)\,(m+1-2s_m)=(m-1)\,P(m)-2\,P(m)\,(s_m-1),
\]
and $P(m)\,(s_m-1)=P_2(m)$ by Lemma~\ref{lem:twofold}.
\end{proof}

\subsection{Mean positions}\label{subsec:birthday-means}

In the estimates of Section~\ref{sec:pair} the terms of the convolutions
$P_2$ and $P_3$ are multiplied by a function of the splitting point that is
linear, or that is bounded by a linear function. The sums that result are
evaluated by the next three lemmas, which give the mean position of a
splitting point. Their proofs use the symmetry of a product of Poisson
masses in its factors.

\begin{lemma}\label{lem:twofoldmean}
For every integer $m\ge0$,
\[
\sum_{0<g<m}g\,P(g)\,P(m-g)=\frac{m}{2}\,P_2(m).
\]
\end{lemma}

\begin{proof}
With $g_1=g$ and $g_2=m-g$, the left-hand side of the claim is the sum of
$g_1\,P(g_1)\,P(g_2)$ over the compositions $(g_1,g_2)$ of $m$ into two
parts. The product $P(g_1)\,P(g_2)$ does not change when the two parts are
exchanged, so that this sum equals the sum of $g_2\,P(g_1)\,P(g_2)$. The
two sums add up to $m$ times the sum of $P(g_1)\,P(g_2)$, which is
$m\,P_2(m)$, and hence each of them is $(m/2)\,P_2(m)$.
\end{proof}

For the threefold convolution there are two sums of this kind, according to
whether the splitting point separates the first two of the three factors
from the third, or the first factor from the other two.

\begin{lemma}\label{lem:threefoldleft}
For every integer $m\ge0$,
\[
\sum_{0<g<m}g\,P_2(g)\,P(m-g)=\frac{2m}{3}\,P_3(m).
\]
\end{lemma}

\begin{proof}
As in the derivation of \eqref{eq:birthday-triples}, in which
$g=g_1+g_2$, the left-hand side of the claim is
\[
\sum_{g_1+g_2+g_3=m}(g_1+g_2)\,P(g_1)\,P(g_2)\,P(g_3).
\]
The product $P(g_1)\,P(g_2)\,P(g_3)$ does not change when the three parts
are permuted, so that the sums of $g_1$, of $g_2$ and of $g_3$ times this
product over the compositions of $m$ into three parts are equal. They add
up to $m$ times the sum of the product, which is $m\,P_3(m)$ by
\eqref{eq:birthday-triples}, and hence each of them is $(m/3)\,P_3(m)$.
The sum displayed above is the sum of the first two, that is,
$2\cdot(m/3)\,P_3(m)$.
\end{proof}

\begin{lemma}\label{lem:threefoldright}
For every integer $m\ge0$,
\[
\sum_{0<g<m}g\,P(g)\,P_2(m-g)=\frac{m}{3}\,P_3(m).
\]
\end{lemma}

\begin{proof}
Inserting the definition of $P_2(m-g)$, we see that the left-hand side of
the claim is the sum of $g_1\,P(g_1)\,P(g_2)\,P(g_3)$ over the compositions
of $m$ into three parts, with $g_1=g$. The product $P(g_1)\,P(g_2)\,P(g_3)$
does not change when the three parts are permuted, so that the sums of
$g_1$, of $g_2$ and of $g_3$ times this product are equal. They add up to
$m$ times the sum of the product, which is $m\,P_3(m)$ by
\eqref{eq:birthday-triples}, and hence the first of them is
$(m/3)\,P_3(m)$.
\end{proof}

For
an integer $R\ge0$ and the integers $s$ with $0<s<R$ let
\[
\vartheta_s=P(R-s)\,(R-1-s)\,P_2(s),
\]
where the dependence on $R$ is not shown. The numbers $\vartheta_s$ occur
in the proof of Proposition~\ref{prop:boundaryupper}, in a sum
over $s$ in which $\vartheta_s$ is multiplied by the value at $s$ of a
concave function. That proof needs the total of the weights and a lower
bound for their mean position, the sum of the numbers $s\,\vartheta_s$
divided by the sum of the weights. The total follows from
Lemma~\ref{lem:threefoldleft}.

\begin{corollary}\label{cor:weightsum}
For every integer $R\ge0$, the weights $\vartheta_s$ satisfy
\[
\sum_{0<s<R}\vartheta_s=\frac{R-3}{3}\,P_3(R).
\]
\end{corollary}

\begin{proof}
Writing $R-1-s=(R-1)-s$, we have
\[
\sum_{0<s<R}\vartheta_s
=(R-1)\sum_{0<s<R}P_2(s)\,P(R-s)-\sum_{0<s<R}s\,P_2(s)\,P(R-s).
\]
The first sum on the right is $P_3(R)$, by the definition of the threefold
convolution, and the second is $(2R/3)\,P_3(R)$, by
Lemma~\ref{lem:threefoldleft}. The claim follows.
\end{proof}

The lower bound for the mean position rests on Lemma~\ref{lem:threefold},
which turns the sums over $s$ into sums over the compositions of $R$ into
five parts and into four parts.

\begin{lemma}\label{lem:weightmean}
For every integer $R\ge0$, the weights $\vartheta_s$ satisfy
\[
\frac{2R}{5}\sum_{0<s<R}\vartheta_s\ \le\ \sum_{0<s<R}s\,\vartheta_s.
\]
\end{lemma}

\begin{proof}
Let $s$ be an integer with $0<s<R$. Lemma~\ref{lem:threefold} at
$R-s\ge1$ gives $(R-s-1)\,P(R-s)=P_3(R-s)+2\,P_2(R-s)$, and hence
\[
\vartheta_s=P_2(s)\,P_3(R-s)+2\,P_2(s)\,P_2(R-s).
\]
For $q=4$ and $q=5$ let $\Sigma_q$ denote the sum of the products
$P(g_1)\cdots P(g_q)$ over the compositions $R=g_1+\dots+g_q$ of $R$ into
$q$ parts. A composition of $s$ into two parts, followed by a composition of
$R-s$ into three parts, is a composition of $R$ into five parts with
$g_1+g_2=s$, and every composition of $R$ into five parts arises in this
way for exactly one $s$. By the definition of $P_2$ and by
\eqref{eq:birthday-triples}, the sum of the products $P_2(s)\,P_3(R-s)$
over $0<s<R$ is therefore $\Sigma_5$, and
\[
\sum_{0<s<R}s\,P_2(s)\,P_3(R-s)
=\sum_{g_1+\dots+g_5=R}(g_1+g_2)\,P(g_1)\cdots P(g_5)
=\frac{2R}{5}\,\Sigma_5,
\]
where the last step holds because the product $P(g_1)\cdots P(g_5)$ does
not change when the five parts are permuted: the sums of $g_1,\dots,g_5$
times this product over the compositions of $R$ into five parts are
therefore equal, they add up to $R\,\Sigma_5$, and each of them is
$(R/5)\,\Sigma_5$. In the same way, with compositions into four parts instead
of five, the sum of the products $P_2(s)\,P_2(R-s)$ over $0<s<R$ is
$\Sigma_4$, and
\[
\sum_{0<s<R}s\,P_2(s)\,P_2(R-s)=\frac{2R}{4}\,\Sigma_4.
\]
Consequently
\[
\sum_{0<s<R}\vartheta_s=\Sigma_5+2\,\Sigma_4,\qquad
\sum_{0<s<R}s\,\vartheta_s=\frac{2R}{5}\,\Sigma_5+R\,\Sigma_4.
\]
The right-hand side of the claim minus its left-hand side is thus
$R\,\Sigma_4-(4R/5)\,\Sigma_4=(R/5)\,\Sigma_4$, which is nonnegative.
\end{proof}

\subsection{Bounds for the mean of the birthday law}
\label{subsec:birthday-bounds}

By Lemmas~\ref{lem:twofold} and~\ref{lem:threefold}, bounds for $P_2(m)$
and $P_3(m)$ come down to bounds for the Poisson mass $P(m)$ and for
$s_m$. For $P(m)$ these are the bounds \eqref{eq:prelim-stirling}. In this
subsection we bound $s_m$ from both sides.

The sum $s_m=1+(m-1)/m+(m-1)(m-2)/m^{2}+\cdots$ is known as Ramanujan's
$Q$-function \cite[(1.3)]{flajolet1995}, and we deduce the two bounds from
a theorem of Flajolet, Grabner, Kirschenhofer and Prodinger, which is stated
in terms of the partial sums of the exponential series. These partial sums
are related to $s_m$ as follows. For integers $m\ge1$ and $0\le j\le m$ we
have $B_{m,j}=m!/((m-j)!\,m^{j})$, so
that $(m^{m}/m!)\,B_{m,j}=m^{m-j}/(m-j)!$. Summing over $0\le j\le m$ and
using $B_{m,0}=1$, we get
\begin{equation}\label{eq:birthday-partial}
\frac{m^{m}}{m!}\,(1+s_m)=\sum_{j=0}^{m}\frac{m^{m-j}}{(m-j)!}
=\sum_{i=0}^{m}\frac{m^{i}}{i!}.
\end{equation}
In the classical notation, a number $\theta_m$ is
defined for every integer $m\ge0$ by
\[
\frac12\,e^{m}=\sum_{i=0}^{m-1}\frac{m^{i}}{i!}+\theta_m\,\frac{m^{m}}{m!}.
\]
For $m\ge1$ the sum on the right-hand side is the last sum of
\eqref{eq:birthday-partial} without its term $m^{m}/m!$ of index $i=m$, and
hence equal to $(m^{m}/m!)\,s_m$. The definition of $\theta_m$ therefore
reads $e^{m}/2=(m^{m}/m!)\,(s_m+\theta_m)$, and since
$P(m)=e^{-m}\,m^{m}/m!$, it gives
\begin{equation}\label{eq:birthday-theta}
s_m=\frac{1}{2\,P(m)}-\theta_m\qquad(m\ge1).
\end{equation}
The theorem \cite[Theorem~7]{flajolet1995} states that for every integer
$m\ge0$ there is a real number $k_m$ with
\[
\theta_m=\frac13+\frac{4}{135\,(m+k_m)}\qquad\text{and}\qquad
\frac{2}{21}\le k_m\le\frac{8}{45}.
\]
Since $k_m>0$, this gives
\begin{equation}\label{eq:birthday-thetabounds}
\frac13<\theta_m<\frac13+\frac{4}{135\,m}\qquad(m\ge1).
\end{equation}
The two lemmas of this subsection are stated with the constants that the
later sections use.

\begin{lemma}\label{lem:birthdaylower}
For every integer $m\ge2$, the mean $s_m$ of the birthday law of $m$
satisfies
\[
\frac{1}{2\,P(m)}-\frac13-\frac{1}{15\,m}\ \le\ s_m.
\]
\end{lemma}

\begin{proof}
Since $4/135<1/15$, the relation \eqref{eq:birthday-theta} and the
right-hand inequality of \eqref{eq:birthday-thetabounds} give
\[
s_m=\frac{1}{2\,P(m)}-\theta_m
>\frac{1}{2\,P(m)}-\frac13-\frac{4}{135\,m}
>\frac{1}{2\,P(m)}-\frac13-\frac{1}{15\,m}.\qedhere
\]
\end{proof}

With the closed form of Lemma~\ref{lem:twofold}, the lower bound of $s_m$
becomes a lower bound of the twofold convolution.

\begin{corollary}\label{cor:twofoldlower}
For every integer $m\ge2$, the twofold convolution of the Poisson mass
satisfies
\[
\frac12-\Bigl(\frac43+\frac{1}{15\,m}\Bigr)\,P(m)\ \le\ P_2(m).
\]
\end{corollary}

\begin{proof}
Multiplying the inequality of Lemma~\ref{lem:birthdaylower} by $P(m)$, we
obtain $1/2-P(m)/3-P(m)/(15m)\le P(m)\,s_m$, and
$P(m)\,s_m=P_2(m)+P(m)$ by Lemma~\ref{lem:twofold}.
\end{proof}

We turn to the upper bound. Here the left-hand inequality of
\eqref{eq:birthday-thetabounds} is combined with Robbins' bound for the
factorial, which turns $1/(2P(m))$ into an explicit function of $m$.

\begin{lemma}\label{lem:birthdayupper}
For every integer $m\ge1$, the mean $s_m$ of the birthday law of $m$
satisfies
\[
s_m\le\sqrt{\frac{\pi m}{2}}\;e^{1/(12m)}-\frac13.
\]
\end{lemma}

\begin{proof}
By \eqref{eq:birthday-theta} and the left-hand inequality of
\eqref{eq:birthday-thetabounds}, we have $s_m<1/(2P(m))-1/3$. Moreover,
$1/(2P(m))=m!/(2\,(m/e)^{m})$, and the right-hand inequality of
\eqref{eq:prelim-robbins} shows
that
\[
\frac{1}{2\,P(m)}<\frac{\sqrt{2\pi m}}{2}\;e^{1/(12m)}
=\sqrt{\frac{\pi m}{2}}\;e^{1/(12m)}.\qedhere
\]
\end{proof}

\section{The frozen polynomial and a criterion for a positive step}\label{sec:frozen}
By \eqref{eq:counts-step}, the normalised step $\delta^{(k)}_n$ of a column
$k\ge2$ is a second difference of the normalised counts in the threshold.
In this section we write this step as a second difference of a polynomial
in a real variable $y$, the frozen polynomial, minus a multiple of a chain
sum. Proposition~\ref{prop:frozenstep} does this for the steps of two
neighbouring columns in one row, with the same polynomial, and its two
formulas are used
separately. The first gives the expression of a step of the first layer in
Corollary~\ref{cor:curvatureboundary}, from which
Proposition~\ref{prop:diagonalstep} is proved in Section~\ref{sec:step}.
The second gives a sufficient condition for a positive step in which the
step itself no longer occurs (Proposition~\ref{prop:criterion}), from which
Sections~\ref{sec:pair} to~\ref{sec:curve} derive
Proposition~\ref{prop:curveincrease}.

\subsection{The kernel deficit}\label{subsec:frozen-deficit}

For integers $k\ge0$ and $t\ge0$ the \emph{kernel deficit} is the number
\[
\varkappa_k(t)=a_t-\sum_{k<i<t}a_i\,a_{t-i} .
\]
The sum is empty for $t\le k+1$, so that
$\varkappa_k(t)=a_t$ in this range.
The kernel deficits arise as the terms without the threshold in a linear
equation for the normalised counts of two consecutive thresholds, which is
derived in the proof of Proposition~\ref{prop:frozenstep}, and they are the
coefficients of the frozen polynomial. The next
lemma shows that they are positive in a range that contains all their
arguments in the first layer of column $k$.

\begin{lemma}\label{lem:deficitpositive}
Let $k\ge0$ and $t$ be integers with $1\le t\le 2k+1$. Then
$\varkappa_k(t)>0$.
\end{lemma}

\begin{proof}
The reflection $i\mapsto t-i$ maps the set of integers $i$ with $k<i<t$
onto the set of integers $j$ with $0<j<t-k$, and it leaves the product
$a_i\,a_{t-i}$ unchanged. An index $j<t-k$ satisfies $j\le k$, because
$t\le 2k+1$; so the two sets are disjoint. Since the numbers $a_i$ are
nonnegative, it follows that
\[
2\sum_{k<i<t}a_i\,a_{t-i}\ \le\ \sum_{i=0}^{t}a_i\,a_{t-i} .
\]
By Lemma~\ref{lem:cayleyconvolution},
\[
\sum_{i=0}^{t-1}i\,a_i\,a_{t-i}=(t-1)\,a_t .
\]
Since $a_0=0$, the sum may be extended to $i=t$. Adding to the identity so
obtained the same identity with $i$ replaced by $t-i$, we get
\[
t\sum_{i=0}^{t}a_i\,a_{t-i}=2\,(t-1)\,a_t .
\]
Therefore $t$ times the sum over $k<i<t$ is at most $(t-1)\,a_t$, and
$t\,\varkappa_k(t)\ge a_t>0$.
\end{proof}

\subsection{The frozen polynomial and the curvature
sums}\label{subsec:frozen-polynomial}

From here on we use the sums over compositions of
Section~\ref{sec:compositions}: the weights $V_q(m;f)$ and $\Lambda_m[f]$
for a cut profile $f$, the weights $\Lambda^{(1)}_m[f]$ and
$\Lambda^{(2)}_m[f]$ with one and with two marked cut points, the threshold
profile $\lambda^{y,N}$ and the chain sums $\Lambda_m(y;N)$. The homogeneity
\eqref{eq:compositions-homogeneity} of the weights $V_q(m;\lambda^{y,N})$ in
the threshold $y$, summed over the number $q$ of parts, gives
\begin{equation}\label{eq:frozen-homogeneous}
\Lambda_m(y;N)=\sum_{q=1}^{m}y^{\,q-1}\,V_q\bigl(m;\lambda^{1,N}\bigr).
\end{equation}
We also use two facts that are noted in Section~\ref{sec:compositions}:
the weights $V_q(m;f)$ of a cut profile $f$ with $f(s)\ge0$ for $0<s<m$ are
nonnegative, and a threshold profile shifted by an integer $g\ge0$ is again
a threshold profile, namely $(\lambda^{y,N})^{+g}=\lambda^{y,N-g}$.

The frozen polynomial is modelled on the repeated substitution into itself
of the linear equation of the counts that the proof of
Proposition~\ref{prop:frozenstep} derives. Let $k\ge0$ be an integer, let
$j\ge1$ be one of the integers $k$ and $k+1$, and let $N\ge1$ be an
integer. In this substitution for the count $z^{(j)}_{N-1}$, each count on
the right-hand side is replaced by the right-hand side of the same equation
at a smaller argument; a term that arises in this way belongs to a
decreasing chain of integers $N>t_1>\dots>t_i>k$, and every integer $t$
that the chain visits contributes a factor $j/t$.
Put $m=N-t_i$. The differences $N-t_1$, $t_1-t_2$, \dots, $t_{i-1}-t_i$ of
a chain are the parts of a composition of $m$, whose cut points are
$s_l=N-t_l$ for $1\le l<i$, and the factor $j/t_l$ of the integer that the
chain visits at the cut $s_l$ is the value $\lambda^{j,N}(s_l)$ of the
threshold profile. The last integer of the chain contributes
$(j/t_i)\,\varkappa_k(t_i)$. This is the origin of the following
definition, in which the threshold has become a real variable; the frozen
polynomial that it introduces is not related to the frozen kernel of
Section~\ref{sec:renewal}.

\begin{definition}\label{def:frozenpolynomial}
Let $k\ge0$ and $n\ge0$ be integers, and put $N=n+1$. The \emph{frozen
polynomial} of column $k$ and row $n$
is the function of the real variable $y$ given by
\[
\Phi_{n,k}(y)=\varkappa_k(N)+\sum_{m=1}^{n-k}\frac{\varkappa_k(N-m)}{N-m}\;y\,\Lambda_m(y;N).
\]
The sum is empty when $n\le k$.
\end{definition}

By \eqref{eq:frozen-homogeneous} the function $y\,\Lambda_m(y;N)$ is a
polynomial in $y$ of degree at most $m$ without constant term. So
$\Phi_{n,k}$ is a polynomial, of degree at most $n-k$ when $n\ge k$, with
constant term $\varkappa_k(N)$; for $n\le k$ it is the constant
$a_{n+1}$.

The proof of Proposition~\ref{prop:frozenstep} identifies the counts of the
thresholds $k$ and $k+1$ with values of the frozen polynomial of column
$k$. The name of the frozen polynomial refers to the weight of the last-car
recursion that equals $t$ for $t\le k$ and a free value for $t>k$: the point
$k$ at which the weight stops growing is fixed, while the value at which it
stops is left free.

The coefficients of the frozen polynomial are built from kernel deficits
and from composition weights whose signs are known in the first layer.
Indeed, for
$k\le n\le 2k-1$ the arguments $N$ and $N-m$ of the kernel deficits in
Definition~\ref{def:frozenpolynomial} lie between $k+1$ and $2k$, so that
these kernel deficits are positive by Lemma~\ref{lem:deficitpositive}, and
the coefficients $V_q(m;\lambda^{1,N})$ in \eqref{eq:frozen-homogeneous} are
nonnegative.

The frozen polynomial depends on $y$ only through the functions
$y\,\Lambda_m(y;N)$. Their second differences with step one are the
quantities in which the steps will be expressed.

\begin{definition}\label{def:curvaturesum}
For real numbers $y$ and $N$ and an integer $m\ge0$, the \emph{curvature
sum} is
\[
\mathsf C_m(y;N)=(y+1)\,\Lambda_m(y+1;N)-2\,y\,\Lambda_m(y;N)+(y-1)\,\Lambda_m(y-1;N).
\]
\end{definition}

In this section the curvature sum is used with positive integers $y$ and
$N$ and with $m<N$, so that the denominators $N-s$ of the threshold profile
are positive at the cut points $0<s<m$.

\begin{lemma}\label{lem:frozendifference}
Let $k\ge0$ and $n\ge0$ be integers, let $y$ be a real number, and put
$N=n+1$. Then
\[
\Phi_{n,k}(y+1)-2\,\Phi_{n,k}(y)+\Phi_{n,k}(y-1)
=\sum_{m=1}^{n-k}\frac{\varkappa_k(N-m)}{N-m}\;\mathsf C_m(y;N).
\]
The sum is empty when $n\le k$.
\end{lemma}

\begin{proof}
This is immediate from Definitions~\ref{def:frozenpolynomial}
and~\ref{def:curvaturesum}, since the constant term $\varkappa_k(N)$ and the
coefficients $\varkappa_k(N-m)/(N-m)$ do not depend on $y$.
\end{proof}

A curvature sum can in turn be taken power by power in the threshold. The
later statements about curvature sums start from the resulting expansion by
the number of parts.

\begin{lemma}\label{lem:curvatureparts}
For all real numbers $y$ and $N$ and every integer $m\ge0$,
\[
\mathsf C_m(y;N)=\sum_{q=1}^{m}\bigl((y+1)^q-2\,y^q+(y-1)^q\bigr)\,
V_q\bigl(m;\lambda^{1,N}\bigr).
\]
\end{lemma}

\begin{proof}
Multiplying \eqref{eq:frozen-homogeneous} by the threshold, we have for
every real number $x$
\[
x\,\Lambda_m(x;N)=\sum_{q=1}^{m}x^{q}\,V_q\bigl(m;\lambda^{1,N}\bigr).
\]
Inserting this with $x=y+1$, $x=y$ and $x=y-1$ into
Definition~\ref{def:curvaturesum} gives the statement.
\end{proof}

The weights $V_q(m;\lambda^{1,N})$ of this expansion do not contain $y$,
and they are nonnegative when
$N\ge m$. By the binomial theorem,
\begin{equation}\label{eq:frozen-binomial}
(y+1)^q-2\,y^q+(y-1)^q=2\sum_{i\ge1}\binom{q}{2i}\,y^{\,q-2i}
\end{equation}
for every real $y$ and every integer $q\ge0$. This is zero for $q=1$, and
for $y>0$ and $q\ge2$ all its terms are nonnegative and the first of them is
$q(q-1)\,y^{q-2}$. The curvature sums satisfy $\mathsf C_1(y;N)=0$ and
$\mathsf C_m(y;N)>0$ for $y>0$, $m\ge2$ and $N\ge m$.

\subsection{The step as a curvature term minus a boundary
term}\label{subsec:frozen-step}

By \eqref{eq:counts-step} the step of column $k$ in row $n$
is the second difference of the counts of thresholds $k+1$, $k$ and
$k-1$. The next proposition writes the steps of columns $k$ and $k+1$ in
row $n$ as a second difference of the frozen polynomial $\Phi_{n,k}$ minus a
multiple of a chain sum.

\begin{proposition}\label{prop:frozenstep}
Let $k\ge2$ and $r$ be integers with $1\le r\le k-1$, and put $n=k+r$
and $N=n+1$. Then the normalised steps of columns $k$ and $k+1$ in row
$n$ are
\begin{align}
\delta^{(k)}_n&=\bigl(\Phi_{n,k}(k+1)-2\,\Phi_{n,k}(k)+\Phi_{n,k}(k-1)\bigr)
-\frac{a_k}{k}\,\Lambda_{r+1}(k-1;N),\label{eq:frozen-stepk}\\
\delta^{(k+1)}_n&=\bigl(\Phi_{n,k}(k+2)-2\,\Phi_{n,k}(k+1)+\Phi_{n,k}(k)\bigr)
-\frac{a_{k+1}}{k+1}\,\Lambda_{r}(k+2;N).\label{eq:frozen-stepnext}
\end{align}
\end{proposition}

\begin{proof}
We first compare the frozen polynomials of two consecutive columns, for
arbitrary integers $k$ and $n$ with $1\le k\le n$: for every real number
$y$, with $N=n+1$,
\begin{equation}\label{eq:frozen-knot}
\Phi_{n,k-1}(y)-\Phi_{n,k}(y)=(y-k)\,\frac{a_k}{k}\,\Lambda_{N-k}(y;N).
\end{equation}
To see this, we compare the two definitions term by term. Lowering $k$ by
one adds the term of index $i=k$ to the sum that defines the kernel
deficit: for every integer $t>k$,
\begin{equation}\label{eq:frozen-lowering}
\varkappa_{k-1}(t)=\varkappa_k(t)-a_k\,a_{t-k},
\qquad\text{and}\qquad
\varkappa_{k-1}(k)=a_k .
\end{equation}
The sum of
$\Phi_{n,k-1}$ has one term more than that of $\Phi_{n,k}$, the term
$m=N-k$, which is $(a_k/k)\,y\,\Lambda_{N-k}(y;N)$ because
$\varkappa_{k-1}(k)=a_k$. In the constant term and in the terms with
$1\le m\le n-k$ the arguments of the kernel deficits are larger than $k$,
so that by \eqref{eq:frozen-lowering} the kernel deficits of the two
polynomials differ by $-a_k\,a_{N-k}$ and by $-a_k\,a_{N-m-k}$. These terms
therefore contribute
\[
-a_k\Bigl(a_{N-k}+\sum_{m=1}^{n-k}\Lambda_m(y;N)\,\frac{y}{N-m}\,a_{N-k-m}\Bigr)
=-a_k\,\Lambda_{N-k}(y;N)
\]
to the difference, by the recursion of Lemma~\ref{lem:lastcut} for
the profile $\lambda^{y,N}$ and the compositions of $N-k$. The two
contributions add up to the right-hand side of \eqref{eq:frozen-knot}.

Next we show that the counts of two consecutive thresholds are values of one
frozen polynomial: if $k\ge0$ is an integer and $j\ge1$ is one of the
integers $k$ and $k+1$, then
\begin{equation}\label{eq:frozen-values}
z^{(j)}_n=\Phi_{n,k}(j)\qquad\text{for }0\le n\le 2k+1 .
\end{equation}
By Corollary~\ref{cor:normalised} the sequence $z^{(j)}$ is the solution of the
weighted last-car recursion with weight $\rho(t)=\min(t,j)$. This
weight satisfies $\rho(t)=t$ for $1\le t\le k$ and $\rho(t)=j$ for
$t\ge k+1$. Proposition~\ref{prop:chainforward}, applied with this $k$,
therefore gives for every integer $n$ with $0\le n\le 2k+1$
\[
c_n-z^{(j)}_n=\sum_{t=k+1}^{n}c_{n-t}\Bigl(c_{t-1}-\frac{j}{t}\,z^{(j)}_{t-1}\Bigr).
\]
We solve this for $z^{(j)}_n$. Since $c_i=a_{i+1}$, the terms that contain
no count are $a_{n+1}$ and the products $-a_{n+1-t}\,a_t$ with $k<t\le n$,
and they add up to $\varkappa_k(n+1)$. With $N=n+1$ the equation becomes
\begin{equation}\label{eq:frozen-volterra}
z^{(j)}_{N-1}=\varkappa_k(N)+\sum_{k<t<N}a_{N-t}\,\frac{j}{t}\,z^{(j)}_{t-1}
\qquad(1\le N\le 2k+2).
\end{equation}

With $\tau=N-m$ as the index of summation,
Definition~\ref{def:frozenpolynomial} reads, for every integer $N\ge1$ and
every real number $y$,
\begin{equation}\label{eq:frozen-sizeform}
\Phi_{N-1,k}(y)=\varkappa_k(N)+\sum_{k<\tau<N}\varkappa_k(\tau)\,
\frac{y}{\tau}\,\Lambda_{N-\tau}(y;N).
\end{equation}
Let $k$ and $j$ be as above. We show that $z^{(j)}_{N-1}=\Phi_{N-1,k}(j)$
for $1\le N\le 2k+2$ by induction on $N$. Let $N$ be an integer of this
range, and suppose that the identity holds for all smaller values of $N$.
By this hypothesis we may replace every $z^{(j)}_{t-1}$, $k<t<N$, in
\eqref{eq:frozen-volterra} by the right-hand side of
\eqref{eq:frozen-sizeform}, taken with $t$ instead of $N$ and with $y=j$.
We then collect the terms that contain $\varkappa_k(\tau)$ for a given
$\tau$. The result is
\[
z^{(j)}_{N-1}=\varkappa_k(N)+\sum_{k<\tau<N}\varkappa_k(\tau)\,\frac{j}{\tau}\,
\Bigl(a_{N-\tau}+\sum_{\tau<t<N}a_{N-t}\,\frac{j}{t}\,\Lambda_{t-\tau}(j;t)\Bigr),
\]
where $a_{N-\tau}$ comes from the term $t=\tau$ of
\eqref{eq:frozen-volterra}. The
expression in parentheses is $\Lambda_{N-\tau}(j;N)$. Indeed, the profile
$\lambda^{j,N}$ shifted by $g$ is $\lambda^{j,N-g}$, so that the recursion
of Lemma~\ref{lem:firstpart}, for the profile $\lambda^{j,N}$ and the
compositions of $N-\tau$, is
\[
\Lambda_{N-\tau}(j;N)=a_{N-\tau}+\sum_{0<g<N-\tau}a_g\,\frac{j}{N-g}\,
\Lambda_{N-\tau-g}(j;N-g),
\]
and with $t=N-g$ the right-hand side is the expression in parentheses.
Comparing with \eqref{eq:frozen-sizeform} for $y=j$, we obtain
$z^{(j)}_{N-1}=\Phi_{N-1,k}(j)$. This completes the induction, and with
$n=N-1$ it proves \eqref{eq:frozen-values}.

Now let $k$, $r$, $n$ and $N$ be as in the statement. By
\eqref{eq:frozen-values} for column $k$ we have $z^{(k)}_n=\Phi_{n,k}(k)$
and $z^{(k+1)}_n=\Phi_{n,k}(k+1)$, since $n\le 2k+1$. For column $k-1$ and
its threshold $k-1\ge1$, \eqref{eq:frozen-values} gives
$z^{(k-1)}_n=\Phi_{n,k-1}(k-1)$; this requires $n\le 2(k-1)+1$, and it is
here that the hypothesis $r\le k-1$ is used. For column $k+1$ and its
threshold $k+2$ it gives $z^{(k+2)}_n=\Phi_{n,k+1}(k+2)$. We pass from the
polynomials of columns $k-1$ and $k+1$ to $\Phi_{n,k}$ by
\eqref{eq:frozen-knot}. At $y=k-1$, where $N-k=r+1$, this identity gives
\[
\Phi_{n,k-1}(k-1)=\Phi_{n,k}(k-1)-\frac{a_k}{k}\,\Lambda_{r+1}(k-1;N).
\]
With $k$ replaced by $k+1$, which is allowed because $k+1\le n$, and at
$y=k+2$, where $N-(k+1)=r$, it gives
\[
\Phi_{n,k+1}(k+2)=\Phi_{n,k}(k+2)-\frac{a_{k+1}}{k+1}\,\Lambda_{r}(k+2;N).
\]
Inserting the four values into \eqref{eq:counts-step}, for column
$k$ and for column $k+1$, we obtain \eqref{eq:frozen-stepk} and
\eqref{eq:frozen-stepnext}.
\end{proof}

The rows $n=k+r$ of Proposition~\ref{prop:frozenstep} are the rows
$k+1\le n\le 2k-1$, which form the first layer of column $k$; so
\eqref{eq:frozen-stepk} holds at every position of this layer. In each of the
formulas \eqref{eq:frozen-stepk} and \eqref{eq:frozen-stepnext} the second
difference of the frozen polynomial is the \emph{curvature term} of the
step, and the multiple of a chain sum is its \emph{boundary term}.

Writing out the curvature term of the first formula,
\eqref{eq:frozen-stepk}, by Lemma~\ref{lem:frozendifference}, we obtain the
expression of a first-layer step from which
Proposition~\ref{prop:diagonalstep} is proved in Section~\ref{sec:step}.

\begin{corollary}\label{cor:curvatureboundary}
Let $k\ge2$ and $r$ be integers with $1\le r\le k-1$, and put $N=k+r+1$.
Then
\[
\frac{d_{k+r,k}}{(k+r)!}=\sum_{m=1}^{r}\frac{\varkappa_k(N-m)}{N-m}\;
\mathsf C_m(k;N)-\frac{a_k}{k}\,\Lambda_{r+1}(k-1;N).
\]
\end{corollary}

\begin{proof}
The left-hand side is $\delta^{(k)}_{k+r}$. By
Lemma~\ref{lem:frozendifference} for column $k$, row $n=k+r$ and
$y=k$, where $n+1=N$ and $n-k=r$, the second difference in
\eqref{eq:frozen-stepk} is the sum on the right.
\end{proof}

The boundary term is positive: the
chain sum $\Lambda_{r+1}(k-1;N)$ is a sum of nonnegative weights
$V_q(r+1;\lambda^{k-1,N})$, of which the one with $q=1$ is $a_{r+1}>0$. In
the curvature term the coefficients are positive by
Lemma~\ref{lem:deficitpositive}, since $k+1\le N-m\le 2k-1$; the curvature
sum $\mathsf C_1(k;N)$ vanishes, and those with $m\ge2$ are positive, as
noted after Lemma~\ref{lem:curvatureparts}. The sign of the step is decided
by the comparison of the two terms.

\subsection{The pair sum and the criterion}\label{subsec:frozen-criterion}

For the criterion we read the second formula of
Proposition~\ref{prop:frozenstep} for column $k-1$. Let $k\ge3$ and $R$
be integers with $1\le R\le k-2$, and put $n=k+R-1$ and $N=n+1=k+R$. The
proposition applies to column $k-1\ge2$ with $r=R\le(k-1)-1$, for which
its row is $(k-1)+R=n$ and its frozen polynomial is $\Phi_{n,k-1}$. Formula
\eqref{eq:frozen-stepnext} then expresses the step of column $k$:
\begin{multline}\label{eq:frozen-otherform}
\delta^{(k)}_n=\bigl(\Phi_{n,k-1}(k+1)-2\,\Phi_{n,k-1}(k)+\Phi_{n,k-1}(k-1)\bigr)\\
-\frac{a_k}{k}\,\Lambda_R(k+1;N).
\end{multline}
The rows are $k\le n\le 2k-3$, and $n-k=R-1$. For $R\ge2$, position
$(n,k)$ lies in the first layer of column $k$ and has the deficit $R-1$,
so that \eqref{eq:frozen-stepk} holds at this position as well, with
$r=R-1$. The number $\delta^{(k)}_n$ then has two expressions as a curvature
term minus a boundary term, \eqref{eq:frozen-stepk} and
\eqref{eq:frozen-otherform}, and these are different decompositions of the
step. The curvature term of \eqref{eq:frozen-stepk} is a second difference
of $\Phi_{n,k}$, the frozen polynomial of column $k$, whereas that of
\eqref{eq:frozen-otherform} is a second difference of $\Phi_{n,k-1}$, the
frozen polynomial of column $k-1$ in the same row; both are taken at the
points $k+1$, $k$ and $k-1$. The chain sums of the two boundary terms run
over the compositions of the same number $R$, with threshold $k-1$ in
\eqref{eq:frozen-stepk} and with threshold $k+1$ in
\eqref{eq:frozen-otherform}.

We bound the curvature term of \eqref{eq:frozen-otherform} from below by a
quantity in which the frozen polynomial no longer occurs. The quantity is
suggested by the chains on which Definition~\ref{def:frozenpolynomial} is
modelled, taken for column $k-1$: the decreasing chains of integers
$N>t_1>\dots>t_i\ge k$, in which every integer $t_l$ that a chain visits
contributes one factor $y/t_l$, so that a chain that visits $i$ integers
belongs to the power $y^i$. In the proof of
Proposition~\ref{prop:curvaturepair}, $k^2$ times the second
difference of a power $y^i$ at $y=k$ is bounded from below by
$i(i-1)\,k^i$, and $i(i-1)$ is twice the number of pairs of visited
integers. This leads to the chains in which two of the visited integers are
marked. We sort these chains by
the lower of their two marked integers, which we write $N-b$. This integer
contributes the factor $k/(N-b)$. The part of the chain from $N$
down to $N-b$ is a composition of $b$, and the upper marked integer is one
of the integers that the chain visits between $N$ and $N-b$, that is, one
of the cut points of this composition. In the following quantity the
compositions of $b$ with their marked cut point are collected in the weight
$\Lambda^{(1)}_b$, and the part of the chain below $N-b$ is replaced by the
count $z^{(k)}_{N-b-1}$.

\begin{definition}\label{def:markedpairs}
Let $k\ge1$ and $R\ge0$ be integers, and put $N=k+R$. The \emph{pair sum} is
\[
\mathcal P(k,R)=\sum_{b=1}^{R}\Lambda^{(1)}_b(k;N)\cdot\frac{k}{N-b}\cdot
z^{(k)}_{N-b-1}.
\]
\end{definition}

Here $\Lambda^{(1)}_b(k;N)=\Lambda^{(1)}_b[\lambda^{k,N}]$ is the weight of
the compositions of $b$ with one marked cut point for the threshold profile.
We write $a$ for the marked cut point, so that $0<a<b$. From here to the
end of the section the letters $a$ and $b$ without an index denote integers
of this kind, the positions of the two marked points; the Cayley
coefficients $a_j$ and the numbers $b_{n,\ell}$ always carry their indices.
By the definition of the weight $\Lambda^{(1)}_b$, and since the profile
$\lambda^{k,N}$ shifted by $a$ is $\lambda^{k,N-a}$,
\[
\mathcal P(k,R)=\sum_{1\le a<b\le R}\Lambda_a(k;N)\,\frac{k}{N-a}\;
\Lambda_{b-a}(k;N-a)\,\frac{k}{N-b}\;z^{(k)}_{N-b-1},
\]
a sum over the pairs $a<b$ of marked points. In the chains from which we
started, the two marked integers are $N-a>N-b$. The factor $\Lambda_a(k;N)$
is the sum over the part of the chain above $N-a$, which is a composition
of $a$, and the factor $\Lambda_{b-a}(k;N-a)$ is the sum over the part
between the two marked integers, which is a composition of $b-a$; the
factors $k/(N-a)$ and $k/(N-b)$ belong to the marked integers themselves.

The next proposition bounds $k^2$ times the curvature term of
\eqref{eq:frozen-otherform} from below by twice the pair sum. In its proof,
which does not expand into chains, the second differences of the powers produce a
sum that
is organised by the last point $m$ of a composition, with the coefficient
$\varkappa_{k-1}(N-m)$, whereas the pair sum is organised by the lower marked
point $b$, with everything below it collected in a count. The proof
identifies the two sums by a duality between the linear equation of the
counts $z^{(k)}_{N-b-1}$, which expresses the count at a point $b$ by the
counts at larger points, and the recursion of
Lemma~\ref{lem:markedrecursion} for the weights with one or two marked cut
points, which expresses the weight at a point by the weights at smaller
points; the two equations have the same kernel.

\begin{proposition}\label{prop:curvaturepair}
Let $k\ge1$ and $R\ge0$ be integers such that either $R=0$, or $R\ge1$ and
$k\ge R+2$, and put $n=k+R-1$. Then
\[
2\,\mathcal P(k,R)\ \le\ k^2\,\bigl(\Phi_{n,k-1}(k+1)-2\,\Phi_{n,k-1}(k)
+\Phi_{n,k-1}(k-1)\bigr).
\]
\end{proposition}

\begin{proof}
For $R=0$ the pair sum is an empty sum and the polynomial
$\Phi_{n,k-1}=\Phi_{k-1,k-1}$ is constant, so that both sides are $0$. Let
$R\ge1$; then $k\ge R+2\ge3$. We put $N=k+R=n+1$, and we write
$\varkappa$ for $\varkappa_{k-1}$, $\lambda$
for the threshold profile $\lambda^{k,N}$, so that $\lambda(s)=k/(N-s)$,
$\Lambda^{(1)}_m$ and $\Lambda^{(2)}_m$ for $\Lambda^{(1)}_m[\lambda]$ and
$\Lambda^{(2)}_m[\lambda]$, and $\sigma_m$ for the sum
$\Lambda^{(1)}_m+\Lambda^{(2)}_m$.

By Lemma~\ref{lem:frozendifference} for column $k-1$, row $n$ and
$y=k$, where $n-(k-1)=R$,
\begin{multline}\label{eq:frozen-pairstart}
\Phi_{n,k-1}(k+1)-2\,\Phi_{n,k-1}(k)+\Phi_{n,k-1}(k-1)\\
=\sum_{m=1}^{R}\frac{\varkappa(N-m)}{N-m}\;\mathsf C_m(k;N).
\end{multline}
Fix $m$ with $1\le m\le R$, and
write $V_q$ for $V_q(m;\lambda^{1,N})$; these numbers are nonnegative,
because $m<N$. By Lemma~\ref{lem:curvatureparts}, $\mathsf C_m(k;N)$ is the
sum over $1\le q\le m$ of
$\bigl((k+1)^q-2k^q+(k-1)^q\bigr)V_q$. For $y=k$ all terms on the right-hand
side of \eqref{eq:frozen-binomial} are nonnegative, and the first of them is
$q(q-1)\,k^{q-2}$; hence
\[
k^2\bigl((k+1)^q-2\,k^q+(k-1)^q\bigr)\ \ge\ q\,(q-1)\,k^q
\]
for every $q\ge1$. On the other hand, $V_q(m;\lambda)=k^{q-1}\,V_q$, and
Corollary~\ref{cor:partsmoment} for the profile $\lambda$ states that the
numbers $q(q-1)\,V_q(m;\lambda)$, $1\le q\le m$, add up to $2\,\sigma_m$.
Together, for every $m$ with $1\le m\le R$,
\begin{equation}\label{eq:frozen-oneterm}
k^2\,\mathsf C_m(k;N)\ \ge\ \sum_{q=1}^{m}q\,(q-1)\,k^q\,V_q
=2\,k\,\sigma_m .
\end{equation}
For these $m$ we have $1\le N-m\le 2k-3$, so that $\varkappa(N-m)>0$ by
Lemma~\ref{lem:deficitpositive}. We multiply \eqref{eq:frozen-oneterm} by
$\varkappa(N-m)/(N-m)$ and sum over $m$. Since $k/(N-m)=\lambda(m)$, this
gives
\begin{equation}\label{eq:frozen-lowerbound}
k^2\sum_{m=1}^{R}\frac{\varkappa(N-m)}{N-m}\;\mathsf C_m(k;N)\ \ge\
2\sum_{m=1}^{R}\varkappa(N-m)\,\lambda(m)\,\sigma_m ,
\end{equation}
and by \eqref{eq:frozen-pairstart} the left-hand side is the right-hand side
of the proposition.

It remains to show that the sum on the right is the pair sum:
\begin{equation}\label{eq:frozen-duality}
\sum_{m=1}^{R}\varkappa(N-m)\,\lambda(m)\,\sigma_m=\mathcal P(k,R).
\end{equation}
For $1\le b\le R$ put $v_b=z^{(k)}_{N-b-1}$.
Equation~\eqref{eq:frozen-volterra}, which was shown in the proof of
Proposition~\ref{prop:frozenstep}, holds for column $k-1$ and its
threshold $k$ when the $N$ of that equation is replaced by any integer from
$1$ to $2k$. We use it for the integer $N-b$, which is allowed because
$1\le N-b\le 2k-3$, and with $m=N-t$ as the index of summation. Since
$k/t=\lambda(m)$ and $z^{(k)}_{t-1}=v_m$, it reads
\begin{equation}\label{eq:frozen-forward}
v_b=\varkappa(N-b)+\sum_{b<m\le R}a_{m-b}\,\lambda(m)\,v_m
\qquad(1\le b\le R).
\end{equation}
On the other side, Lemma~\ref{lem:markedrecursion} for the profile $\lambda$
states that $\Lambda^{(2)}_m$ is the sum of
$\sigma_s\,\lambda(s)\,a_{m-s}$ over $0<s<m$. Since
$\sigma_m=\Lambda^{(1)}_m+\Lambda^{(2)}_m$, this gives
\begin{equation}\label{eq:frozen-backward}
\sigma_m=\Lambda^{(1)}_m+\sum_{0<s<m}a_{m-s}\,\lambda(s)\,\sigma_s
\qquad(1\le m\le R).
\end{equation}
We solve \eqref{eq:frozen-forward} for $\varkappa(N-s)$,
multiply by $\lambda(s)\,\sigma_s$ and sum over $1\le s\le R$:
\[
\sum_{s=1}^{R}\varkappa(N-s)\,\lambda(s)\,\sigma_s
=\sum_{s=1}^{R}v_s\,\lambda(s)\,\sigma_s
-\sum_{1\le s<m\le R}\lambda(m)\,v_m\;a_{m-s}\,\lambda(s)\,\sigma_s .
\]
Summing over $s$ first in the double sum and using
\eqref{eq:frozen-backward}, we see that the right-hand side is
\[
\sum_{m=1}^{R}v_m\,\lambda(m)\,\sigma_m
-\sum_{m=1}^{R}\lambda(m)\,v_m\,\bigl(\sigma_m-\Lambda^{(1)}_m\bigr)
=\sum_{m=1}^{R}\Lambda^{(1)}_m\,\lambda(m)\,v_m ,
\]
which is $\mathcal P(k,R)$ by Definition~\ref{def:markedpairs}. This proves
\eqref{eq:frozen-duality}, and the proposition follows from
\eqref{eq:frozen-lowerbound}.
\end{proof}

\begin{remark}
The one inequality in the proof keeps, for each $q$, the first term of the
sum \eqref{eq:frozen-binomial} at $y=k$, and the second term is smaller than
the first by the factor $(q-2)(q-3)/(12k^2)$.
\end{remark}

Inserted into \eqref{eq:frozen-otherform}, the proposition gives the
criterion.

\begin{proposition}\label{prop:criterion}
Let $k$ and $R$ be integers with $1\le R\le k-2$, so that $k\ge3$, and put
$N=k+R$ and $n=k+R-1$. If
\[
k\,a_k\,\Lambda_R(k+1;N)<2\,\mathcal P(k,R),
\]
then the normalised step $\delta^{(k)}_n$ of column $k$ in row $n$
is positive.
\end{proposition}

\begin{proof}
We multiply \eqref{eq:frozen-otherform} by $k^2$ and apply
Proposition~\ref{prop:curvaturepair}, whose hypothesis holds because
$R\ge1$ and $k\ge R+2$. The hypothesis of the criterion then gives
\[
k^2\,\delta^{(k)}_n\ \ge\ 2\,\mathcal P(k,R)-k\,a_k\,\Lambda_R(k+1;N)\ >\ 0 .
\qedhere
\]
\end{proof}

\section{The two sides of the criterion}\label{sec:pair}
Proposition~\ref{prop:criterion} reduces the positivity of the step of
column $k$ in row $k+R-1$, for integers $k$ and $R$ with $1\le R\le k-2$, to
an inequality in which the step no longer occurs: with $N=k+R$, the step is
positive as soon as
\[
k\,a_k\,\Lambda_R(k+1;N)<2\,\mathcal P(k,R).
\]
We now estimate the two sides. Proposition~\ref{prop:pairlower}
bounds the pair sum $\mathcal P(k,R)$ from below, and
Proposition~\ref{prop:boundaryupper} bounds the chain sum $\Lambda_R(k+1;N)$
of the boundary term from above. Both bounds are expressed in terms of the
Poisson mass $P$ and of its convolutions $P_2$ and $P_3$ of
Section~\ref{sec:birthday}. Section~\ref{sec:sums} bounds the three sums of
these convolutions that occur in the first bound, and
Section~\ref{sec:curve} then inserts both bounds into the criterion.

\subsection{A tangent inequality for a Volterra equation}
\label{subsec:pair-volterra}

The left-hand
side of the criterion contains a chain sum over the compositions of $R$,
and the pair sum is built from chain sums and from normalised counts
$z^{(k)}_n$. In a chain sum $\Lambda_m(y;N)$ a composition of $m$ carries,
besides the Cayley coefficients of its parts, the value $y/(N-s)$ of the
threshold profile at each of its cut points $s$; these values are its cut
factors (Section~\ref{sec:compositions}). For the counts, the proof of
Lemma~\ref{lem:countjensen} uses a linear equation in which an integer
$t>k$ carries the factor $k/t$; in this section we call these
factors cut factors as well.

If every cut factor is replaced by $1$, the chain
sum $\Lambda_R(k+1;N)$ becomes
$\nu_R=e^{R}P(R)$ (Corollary~\ref{cor:unitprofile}). The two bounds of this
section take the cut factors into account instead of replacing them by $1$.

The common tool is a lower bound for sums of products of cut factors. Write
a cut factor as $e^{-x}$ with a real number $x$. A product of cut factors is
then $e^{-X}$, where $X$ is the sum of the exponents, and the exponential
function lies above each of its tangents: for all real numbers $X$ and $\mu$,
\begin{equation}\label{eq:pair-tangent}
e^{-X}\ge e^{-\mu}\,(1+\mu-X).
\end{equation}
The right-hand side is linear in $X$, so that its sum over the compositions
is known as soon as the sum of the exponents
is known. The sums in question are solutions of linear equations of one
type, and Lemma~\ref{lem:volterra} carries out this summation for all of
them at once. The inequalities that follow from it apply directly to the
pair sum, for which a lower bound is wanted. For the chain sum of the
boundary term, which has to be bounded from above, they are applied to the
amount by which the chain sum falls short of its value without cut factors,
which Lemma~\ref{lem:firstviolation} expresses through the chain sums over
the compositions of the integers smaller than $R$.

We consider real-valued functions
on the nonnegative integers and equations of the form
\begin{equation}\label{eq:pair-volterra}
u(M)=h(M)+\sum_{t_0<t<M}a_{M-t}\,w(t)\,u(t),
\end{equation}
in which the value of a function $u$ at $M$ is expressed in terms of its
values at smaller arguments, of a function $h$ and of a weight $w$; here
$t_0\ge0$ is
an integer. Such an equation is a Volterra equation of the second kind in
discrete form, whose kernel $a_{M-t}$ is formed from the Cayley
coefficients. Two sequences of this paper satisfy an equation of
this form. The
total weights $\Lambda_M[f]$ of the compositions of $M$ for a cut profile
$f$ satisfy it with $h(M)=a_M$, $w=f$ and $t_0=0$; this is the last-cut
recursion of Lemma~\ref{lem:lastcut}. The numbers $\nu_M$ satisfy it with
the same $h$ and $t_0$ and with the weight $1$: in the identity of
Lemma~\ref{lem:nurenewal} the term $j=0$ vanishes and the term $j=m$ is
$a_m\,\nu_0=a_m$, so that
\begin{equation}\label{eq:pair-nurenewal}
\nu_m=a_m+\sum_{0<s<m}a_{m-s}\,\nu_s\qquad(m\ge1).
\end{equation}
The proof of Lemma~\ref{lem:countjensen} applies Lemma~\ref{lem:volterra}
also to the normalised counts $z^{(k)}_{M-1}$ of a threshold $k\ge1$.

The following lemma compares the solution $u$ of \eqref{eq:pair-volterra}
for a weight of the form $w(t)=e^{-x(t)}$ with the solution $v$ of the same
equation without weights.

\begin{lemma}\label{lem:volterra}
Let $t_0\ge0$ and $B\ge0$ be integers, and let $h$, $v$, $u$ and $x$ be
real-valued functions on the nonnegative integers such that for every integer
$M$ with $1\le M\le B$ the following three conditions hold.
\begin{enumerate}
\item $h(M)\ge0$.
\item $v(M)=h(M)+\sum_{t_0<t<M}a_{M-t}\,v(t)$.
\item $u(M)=h(M)+\sum_{t_0<t<M}a_{M-t}\,e^{-x(t)}\,u(t)$.
\end{enumerate}
Then for every real number $\mu$ and every integer $M$ with $1\le M\le B$,
\[
e^{-\mu}\Bigl((1+\mu)\,v(M)-\sum_{t_0<t<M}\nu_{M-t}\,v(t)\,x(t)\Bigr)\le u(M).
\]
\end{lemma}

\begin{proof}
For an integer $M\ge0$ let $S(M)=\sum_{t_0<t<M}\nu_{M-t}\,v(t)\,x(t)$. We
first show that for every $M\ge0$
\begin{equation}\label{eq:pair-kernelidentity}
\sum_{t_0<t<M}a_{M-t}\,\bigl(v(t)\,x(t)+S(t)\bigr)=S(M).
\end{equation}
When the definition of $S(t)$ is inserted on the left-hand side and the two
summations are exchanged, a number $v(c)\,x(c)$ with $t_0<c<M$ occurs with
the total coefficient $a_{M-c}+\sum_{0<s<M-c}a_{M-c-s}\,\nu_s$, and this
equals $\nu_{M-c}$ by \eqref{eq:pair-nurenewal} for $m=M-c\ge1$.

We now prove the inequality of the lemma by strong induction on $M$, for all
real numbers $\mu$ at once. Let $1\le M\le B$, let $\mu$ be a real number,
and suppose that the inequality holds at every integer $t$ with $1\le t<M$
and for every real value of the parameter. Let $t$ be an integer with
$t_0<t<M$. Then $1\le t\le B$, and the inequality at $t$ with the parameter
$\mu-x(t)$ reads
\[
e^{-(\mu-x(t))}\Bigl(\bigl(1+\mu-x(t)\bigr)\,v(t)-S(t)\Bigr)\le u(t).
\]
Multiplying it by $a_{M-t}\,e^{-x(t)}\ge0$, we obtain
\begin{equation}\label{eq:pair-inductionterm}
e^{-\mu}\,a_{M-t}\bigl((1+\mu)\,v(t)-v(t)\,x(t)-S(t)\bigr)
\le a_{M-t}\,e^{-x(t)}\,u(t).
\end{equation}
The shift of the parameter by $x(t)$ is what turns the weight $e^{-x(t)}$ in
the equation of $u$ into the common factor $e^{-\mu}$. Moreover
$1+\mu\le e^{\mu}$ and $h(M)\ge0$, so that
\[
e^{-\mu}\,(1+\mu)\,h(M)\le h(M).
\]
We add this inequality and the inequalities \eqref{eq:pair-inductionterm}
for all $t$ with $t_0<t<M$. On the right we obtain $u(M)$, by the equation of
$u$. On the left we obtain $e^{-\mu}$ times
\begin{multline*}
(1+\mu)\Bigl(h(M)+\sum_{t_0<t<M}a_{M-t}\,v(t)\Bigr)
-\sum_{t_0<t<M}a_{M-t}\,\bigl(v(t)\,x(t)+S(t)\bigr)\\
=(1+\mu)\,v(M)-S(M),
\end{multline*}
by the equation of $v$ and by \eqref{eq:pair-kernelidentity}. This is the
inequality at $M$.
\end{proof}

Lemma~\ref{lem:volterra} can be read as follows, although its proof does not
proceed in this way. Iterating the equation of $u$ expresses $u(M)$ as a
sum over the decreasing chains $M>t_1>\dots>t_j>t_0$ of integers, $j\ge0$.
Such a chain has the weight
\[
a_{M-t_1}\,a_{t_1-t_2}\cdots a_{t_{j-1}-t_j}\,h(t_j),
\]
the empty chain having the weight $h(M)$, and it carries in addition the
factor $e^{-X}$, where $X=x(t_1)+\dots+x(t_j)$ is the sum of the exponents
at the points that the chain visits. The function $v$ is the same sum
without the factors $e^{-X}$.
By the first condition the weights are nonnegative, and
\eqref{eq:pair-tangent}, summed over the chains, bounds $u(M)$ from below by
$e^{-\mu}$ times the difference of $(1+\mu)\,v(M)$ and the sum of $X$ over
the chains, each taken with its weight. To this last sum a point $t$
contributes $x(t)$ once for every chain that visits it. Such a chain
consists of a chain from $M$ down to $t$ and of a chain that starts at $t$.
The chains of the first kind are the compositions of $M-t$ without cut
factors, whose total weight is $\nu_{M-t}$
(Corollary~\ref{cor:unitprofile}), and the total weight of the chains of the
second kind is $v(t)$. So the sum of $X$ over the chains is
$\sum_{t_0<t<M}\nu_{M-t}\,v(t)\,x(t)$, which is the sum in the lemma.

If $v(M)>0$, one can take $\mu=S(M)/v(M)$ in the lemma, in the notation of
its proof. The expression in parentheses then equals $v(M)$, and the lemma
reads $v(M)\,e^{-S(M)/v(M)}\le u(M)$. In the reading by chains this is
Jensen's inequality for the convex function $X\mapsto e^{-X}$: the mean of
$e^{-X}$ over the chains is at least $e^{-\bar X}$, where $\bar X$ is the
mean of $X$. Corollary~\ref{cor:compositionjensen} and
Lemma~\ref{lem:countjensen} are of this form. The proof of
Proposition~\ref{prop:boundaryupper} uses the lemma for a general $\mu$,
with one and the same value for several chain sums.

For the total weight $\Lambda_m[f]$ of the compositions of $m$ for a cut
profile $f$ (Section~\ref{sec:compositions}) the lemma takes the following
form. For a real-valued function $x$ on the nonnegative integers we write
$e^{-x}$ for the cut profile $s\mapsto e^{-x(s)}$; every positive cut
profile is of this form.

\begin{lemma}\label{lem:compositiontangent}
Let $x$ be a real-valued function on the nonnegative integers, let $\mu$ be a
real number and let $m\ge1$ be an integer. Then
\[
e^{-\mu}\Bigl((1+\mu)\,\nu_m-\sum_{0<c<m}\nu_{m-c}\,\nu_c\,x(c)\Bigr)
\le\Lambda_m[e^{-x}].
\]
\end{lemma}

\begin{proof}
We apply Lemma~\ref{lem:volterra} with $t_0=0$, $B=m$, $h(M)=a_M$,
$v(M)=\nu_M$ and $u(M)=\Lambda_M[e^{-x}]$. Its first condition is clear, the
equation of $v$ is \eqref{eq:pair-nurenewal}, and the equation of $u$ is the
last-cut recursion of Lemma~\ref{lem:lastcut} for the profile $e^{-x}$. The
conclusion of Lemma~\ref{lem:volterra} at $M=m$ is the claim.
\end{proof}

In the reading of Lemma~\ref{lem:volterra} by chains, the chains
are now the compositions of $m$, and
the sum in the lemma, divided by $\nu_m$, is the mean of the sum
of the exponents at the cut points, for a composition that is drawn with a
probability proportional to the product of the Cayley coefficients of its
parts. Taking this mean for $\mu$, we obtain Jensen's inequality.

\begin{corollary}\label{cor:compositionjensen}
Let $x$ be a real-valued function on the nonnegative integers and let
$m\ge1$ be an integer. Then
\[
\nu_m\,\exp\Bigl(-\frac{1}{\nu_m}\sum_{0<c<m}\nu_{m-c}\,\nu_c\,x(c)\Bigr)
\le\Lambda_m[e^{-x}].
\]
\end{corollary}

\begin{proof}
Let $S=\sum_{0<c<m}\nu_{m-c}\,\nu_c\,x(c)$. We apply
Lemma~\ref{lem:compositiontangent} with $\mu=S/\nu_m$, for which
$(1+\mu)\,\nu_m-S=\nu_m$.
\end{proof}

A term of the pair sum contains, besides two chain sums, a normalised count
$z^{(k)}_n$ of threshold $k$ (Section~\ref{sec:counts}), and the next
lemma is the inequality of the same kind for these numbers.

\begin{lemma}\label{lem:countjensen}
Let $k\ge1$ and $M$ be integers with $1\le M\le2k+1$. Then the normalised
count of threshold $k$ satisfies
\[
a_M\,\exp\Bigl(-\frac{1}{a_M}\sum_{k<t<M}\nu_{M-t}\,a_t\,\log\frac tk\Bigr)
\le z^{(k)}_{M-1}.
\]
\end{lemma}

\begin{proof}
By \eqref{eq:frozen-volterra}, which was shown in the proof of
Proposition~\ref{prop:frozenstep}, taken with $j=k$, which is admissible
there because $k\ge1$, and with $M$ in the place of $N$, the normalised counts
$z^{(k)}_{M-1}$ satisfy for every integer $M$ with $1\le M\le2k+2$ an
equation of the form \eqref{eq:pair-volterra}, with $t_0=k$, with the
weight $w(t)=k/t$ and with the kernel deficit $\varkappa_k$ of
Section~\ref{sec:frozen} as the function $h$.

We now apply Lemma~\ref{lem:volterra} with $t_0=k$, $B=2k+1$,
$h=\varkappa_k$ and $v(M)=a_M$, with $u(M)=z^{(k)}_{M-1}$ for $M\ge1$ and
with $x(t)=\log(t/k)$ for $t\ge1$; the values of $u$ and $x$ at $0$ do not
enter. Its three conditions hold for $1\le M\le2k+1$: in this range
$\varkappa_k(M)>0$ by Lemma~\ref{lem:deficitpositive}; the equation of $v$
is the definition of $\varkappa_k(M)$; and the equation of $u$ is
\eqref{eq:frozen-volterra} with $j=k$, in which $k/t=e^{-x(t)}$. Hence,
for every real number $\mu$,
\[
e^{-\mu}\bigl((1+\mu)\,a_M-S\bigr)\le z^{(k)}_{M-1},\qquad
S=\sum_{k<t<M}\nu_{M-t}\,a_t\,\log\frac tk.
\]
For $\mu=S/a_M$ we have $(1+\mu)\,a_M-S=a_M$, and the inequality becomes the
claim.
\end{proof}

The number $a_M$ on the left-hand side of Lemma~\ref{lem:countjensen} is
$c_{M-1}$, the number of all parking functions of length $M-1$ divided by
$(M-1)!$, so that the lemma bounds from below the proportion of the parking
functions of length $M-1$ all of whose displacements are smaller than $k$.

\subsection{The Borel distribution}\label{subsec:pair-borel}

The numbers $e^{-m}a_m$, $m\ge1$, are the probabilities of the Borel
distribution with parameter $1$; their sum is $1$ by
Lemma~\ref{lem:borelmass}. The lower bound of the pair sum contains
quotients $e^{-i}a_i/(e^{-j}a_j)$ of two of these probabilities with
$i\le j$, which come from the Cayley coefficients in
Lemma~\ref{lem:countjensen}. The next lemma shows that $(j/i)^{3/2}$ is an
upper bound for such a quotient.

\begin{lemma}\label{lem:borelmonotone}
For all integers $i$ and $j$ with $1\le i\le j$, the Cayley coefficients
satisfy
\[
\log a_i-i+\tfrac32\log i\le\log a_j-j+\tfrac32\log j.
\]
\end{lemma}

\begin{proof}
For an integer $m\ge1$ let $\phi(m)=\log a_m-m+\frac32\log m$. It suffices
to show that $\phi(j)\le\phi(j+1)$ for every $j\ge1$. Since
$a_{j+1}=(1+1/j)^{j-1}a_j$ by \eqref{eq:prelim-quotient}, we have
\begin{align*}
\phi(j+1)-\phi(j)&=(j-1)\log\Bigl(1+\frac1j\Bigr)-1
+\frac32\log\Bigl(1+\frac1j\Bigr)\\
&=\Bigl(j+\frac12\Bigr)\log\Bigl(1+\frac1j\Bigr)-1.
\end{align*}
For every real number $u\ge0$ we have $\log(1+u)\ge2u/(u+2)$, as one checks
by differentiating the difference of the two sides, which vanishes at $u=0$.
At $u=1/j$ this inequality reads $\log(1+1/j)\ge2/(2j+1)$, so that
$(j+\frac12)\log(1+1/j)\ge1$ and $\phi(j+1)-\phi(j)\ge0$.
\end{proof}

In the proof
of Proposition~\ref{prop:pairlower} a bound for the quotient
$e^{-i}a_i/(e^{-j}a_j)$ is multiplied
by nonnegative weights and summed over one of the two indices, and the sums
of those weights against the index and against its square are known or
easily bounded. A power with the exponent $3/2$ cannot be summed in this
way, whereas a polynomial can. We therefore replace the bound
$(j/i)^{3/2}=(1+u)^{3/2}$, with
$u=(j-i)/i$, by its Taylor polynomial of degree two at $u=0$, which lies
above it for $u\ge0$.

\begin{lemma}\label{lem:borelquotient}
For all integers $i$ and $j$ with $1\le i\le j$, the Cayley coefficients
satisfy
\[
\frac{e^{\,j-i}\,a_i}{a_j}\le1+\frac32\cdot\frac{j-i}{i}
+\frac38\Bigl(\frac{j-i}{i}\Bigr)^2.
\]
\end{lemma}

\begin{proof}
Let $u=(j-i)/i$, so that $u\ge0$ and $1+u=j/i$, and let
$q=e^{\,j-i}a_i/a_j$. By Lemma~\ref{lem:borelmonotone},
\[
\log q=(j-i)+\log a_i-\log a_j\le\tfrac32\,(\log j-\log i)
=\tfrac32\log(1+u),
\]
and hence $q^2\le(1+u)^3$. On the other hand, a direct computation gives
\[
\Bigl(1+\frac32\,u+\frac38\,u^2\Bigr)^2
=(1+u)^3+\frac18\,u^3+\frac9{64}\,u^4\ge(1+u)^3.
\]
Thus $q^2\le(1+\frac32u+\frac38u^2)^2$, and the claim follows on taking
square roots.
\end{proof}

\subsection{The lower bound of the pair sum}\label{subsec:pair-lower}

Let $k\ge1$ and $R\ge0$ be integers and $N=k+R$.
By Definition~\ref{def:markedpairs} and by the definition of the weight
$\Lambda^{(1)}_b[f]$ of the compositions of $b$ with one marked cut point
(Section~\ref{sec:compositions}), the pair sum is a sum over the pairs of
marked cut points $a<b$,
\begin{equation}\label{eq:pair-terms}
\begin{split}
\mathcal P(k,R)&=\sum_{b=1}^{R}\ \sum_{0<a<b}\mathcal P_{ab},\\
\mathcal P_{ab}&=\Lambda_a(k;N)\,\frac{k}{N-a}\,\Lambda_{b-a}(k;N-a)\,
\frac{k}{N-b}\,z^{(k)}_{N-b-1}.
\end{split}
\end{equation}
Here we have used that the threshold profile $\lambda^{k,N}$, shifted by $a$,
is the threshold profile $\lambda^{k,N-a}$. In this subsection $a$ and $b$
without an index denote cut points; the Cayley coefficients always carry an
index.

The cut factor at a cut point
$s\le R$ is $k/(N-s)$, and $N-s=k+(R-s)$. We write it as
\[
\frac{k}{N-s}=e^{-L(R-s)},\qquad\text{where}\quad
L(y)=\log\Bigl(1+\frac yk\Bigr).
\]
Each of the three sums in $\mathcal P_{ab}$, the two chain sums and the normalised
count, is thus a sum of products of exponentials, and
Corollary~\ref{cor:compositionjensen} and Lemma~\ref{lem:countjensen} bound
it from below by a single exponential. For the pair $a<b$ this gives a lower
bound of the form $a_k\,e^{R}\,P(a)\,P(b-a)\,e^{-X_{ab}}$. Here
$P(a)\,P(b-a)$ is, up to the factor $e^{b}$, the value $\nu_a\,\nu_{b-a}$ of
the two chain sums without cut factors, and the exponent $X_{ab}$ consists
of means of sums of values of $L$. Two tangents then make the sum over the
pairs computable. The inequality \eqref{eq:pair-tangent} replaces
$e^{-X_{ab}}$ by an expression that is linear in $X_{ab}$. And the concave
function $L$ lies below each of its tangents: for a real number
$\theta\ge0$ the tangent of $L$ at the point $y=k\theta$ is a line
$\alpha+\beta y$, whose coefficients are written out in
Proposition~\ref{prop:pairlower}. When $L$ is replaced by a line, the sum
of the exponents over the pairs is evaluated by the mean positions in the
convolutions of the Poisson mass that were computed in
Section~\ref{sec:birthday}. For the sum of the logarithms themselves no such
evaluation is available, and this is the reason for the second tangent.

The result is expressed in terms of three sums. For a real number $N$ and an
integer
$c\ne N$ let
\[
\varrho(N,c)=\frac{c}{2\,(N-c)}+\frac{c^2}{8\,(N-c)^2}.
\]
Clearly $\varrho(k+R,c)\ge0$ for integers $k\ge0$, $R\ge0$ and $c$ with
$0\le c\le R$ and $c<k+R$. For integers $k\ge0$ and $R\ge0$
let
\begin{equation}\label{eq:pair-weight}
\mathcal W(R)=\sum_{b=1}^{R}P_2(b),
\end{equation}
\begin{equation}\label{eq:pair-logweight}
\begin{split}
\mathcal G(k,R)={}&\sum_{b=1}^{R}\Bigl(\frac72\,P_2(b)+2\,P_3(b)\Bigr)\\
&+\sum_{0<c<R}P_3(c)\,\bigl(1+\varrho(k+R,c)\bigr),
\end{split}
\end{equation}
\begin{equation}\label{eq:pair-excess}
\begin{split}
\mathcal A(k,R)={}&\sum_{b=1}^{R}\Bigl(\bigl(R-\tfrac b2\bigr)
\bigl(P_2(b)+2\,P_3(b)\bigr)+\tfrac52\,(R-b)\,P_2(b)\Bigr)\\
&+\sum_{0<c<R}(R-c)\,P_3(c)\,\bigl(1+\varrho(k+R,c)\bigr),
\end{split}
\end{equation}
where $P_2$ and $P_3$ are the twofold and the threefold convolution of the
Poisson mass (Section~\ref{sec:birthday}).

By the definition of $P_2$, the sum $\mathcal W(R)$ is the sum of the
weights $P(a)\,P(b-a)$ over the pairs $0<a<b\le R$; this is the total
weight of the pairs. When every value $L(y)$ in the exponents $X_{ab}$ is
replaced
by the line $\alpha+\beta y$, the sum of the exponents over the pairs
becomes $\alpha$ times the total weight of these values plus $\beta$ times
the total of the arguments $y$ at which they are taken, and
$\mathcal G(k,R)$ and $\mathcal A(k,R)$ are bounds for these two totals.
The summand of the first sum of \eqref{eq:pair-logweight} collects, for a
fixed larger marked point $b$, the two marked points with the weight
$P_2(b)$ each, the cut points of the two chain sums with the weight
$P_3(b)$ each, and a term with the weight $\frac32P_2(b)$ that comes from
the comparison of the Cayley coefficients $a_k$ and $a_{N-b}$; the second
sum belongs to the normalised counts. Column $k$ enters $\mathcal G$ and
$\mathcal A$ only through the correction $\varrho(k+R,c)$, which comes from
Lemma~\ref{lem:borelquotient}. The proof of the proposition first
proves the lower bound described above for the term $\mathcal P_{ab}$ of
one pair. It then shows that the sum of the
exponents $X_{ab}$ over the pairs, taken with the weights $P(a)\,P(b-a)$,
is at most $\alpha\,\mathcal G(k,R)+\beta\,\mathcal A(k,R)$, which is its
longest part, and it concludes by \eqref{eq:pair-tangent}.

\begin{proposition}\label{prop:pairlower}
Let $k\ge1$ and $R\ge0$ be integers with $R\le k$, and let $\theta\ge0$ and
$\mu$ be real numbers. Put
\[
\alpha=\log(1+\theta)-\frac{\theta}{1+\theta},\qquad
\beta=\frac{1}{k\,(1+\theta)}.
\]
Then the pair sum $\mathcal P(k,R)$ of Definition~\ref{def:markedpairs}
satisfies
\[
a_k\,e^{R}\,e^{-\mu}\Bigl((1+\mu)\,\mathcal W(R)
-\bigl(\alpha\,\mathcal G(k,R)+\beta\,\mathcal A(k,R)\bigr)\Bigr)
\le\mathcal P(k,R),
\]
where $\mathcal W(R)$, $\mathcal G(k,R)$ and $\mathcal A(k,R)$ are the sums
\eqref{eq:pair-weight}, \eqref{eq:pair-logweight} and
\eqref{eq:pair-excess}.
\end{proposition}

\begin{proof}
Throughout the proof $N=k+R$, and $L(y)=\log(1+y/k)$ for real numbers
$y\ge0$. We use the following properties of $L$. It is nonnegative, and for
all integers $j\ge k$ and $0\le s\le R$,
\begin{equation}\label{eq:pair-logexcess}
L(j-k)=\log\frac jk,\qquad\frac{k}{N-s}=e^{-L(R-s)}.
\end{equation}
Moreover, for every real number $y\ge0$,
\begin{equation}\label{eq:pair-line}
L(y)\le\alpha+\beta\,y.
\end{equation}
Indeed, the inequality $\log t\le t-1$ at the point $t=(1+y/k)/(1+\theta)$
gives
\[
L(y)-\log(1+\theta)\le\frac{1+y/k}{1+\theta}-1
=-\frac{\theta}{1+\theta}+\frac{y}{k\,(1+\theta)}.
\]

We first bound the term of one pair from below. Let $a$ and $b$ be integers
with $0<a<b\le R$. We define six numbers, whose indices show on which of the
two marked points they depend:
\begin{gather*}
X^{(1)}_{a}=\frac{1}{\nu_a}\sum_{0<c<a}\nu_{a-c}\,\nu_c\,L(R-c),\qquad
X^{(2)}_{a}=L(R-a),\\
X^{(3)}_{ab}=\frac{1}{\nu_{b-a}}\sum_{0<c<b-a}\nu_{b-a-c}\,\nu_c\,L(R-a-c),
\qquad X^{(4)}_{b}=L(R-b),\\
X^{(5)}_{b}=\log\frac{a_k\,e^{R-b}}{a_{N-b}},\qquad
X^{(6)}_{b}=\frac{1}{a_{N-b}}\sum_{k<t<N-b}\nu_{N-b-t}\,a_t\,\log\frac tk,
\end{gather*}
and we let $X_{ab}$ be their sum. We claim that
\begin{equation}\label{eq:pair-oneterm}
a_k\,e^{R}\,P(a)\,P(b-a)\,e^{-X_{ab}}\le\mathcal P_{ab}.
\end{equation}
The five factors of $\mathcal P_{ab}$ are bounded from below one by one. By
\eqref{eq:pair-logexcess} the threshold profile $\lambda^{k,N}$ agrees at
the cut points $0<s<a$ with the profile $e^{-x}$ for a function $x$ with
$x(s)=L(R-s)$ at these points, and the total weight $\Lambda_a[f]$ of the
compositions of $a$ involves the values of $f$ at the cut points $0<s<a$
only (Section~\ref{sec:compositions}). Corollary~\ref{cor:compositionjensen}
therefore gives
\[
\nu_a\,e^{-X^{(1)}_{a}}\le\Lambda_a(k;N).
\]
In the same way the profile $\lambda^{k,N-a}$ agrees at the cut points
$0<s<b-a$ with $e^{-x}$ for $x(s)=L(R-a-s)$, and
\[
\nu_{b-a}\,e^{-X^{(3)}_{ab}}\le\Lambda_{b-a}(k;N-a).
\]
The two cut factors are $k/(N-a)=e^{-X^{(2)}_{a}}$ and
$k/(N-b)=e^{-X^{(4)}_{b}}$, again by \eqref{eq:pair-logexcess}.
Lemma~\ref{lem:countjensen} applies with $M=N-b$, since $1\le b\le R\le k$
gives $k\le N-b\le2k-1$, and it states that
\[
a_{N-b}\,e^{-X^{(6)}_{b}}\le z^{(k)}_{N-b-1}.
\]
Finally
$a_{N-b}=a_k\,e^{R-b}\,e^{-X^{(5)}_{b}}$ by the definition of $X^{(5)}_{b}$.
Multiplying the bounds for the five factors, all of which are positive, we
obtain
\[
\nu_a\,\nu_{b-a}\,a_k\,e^{R-b}\,e^{-X_{ab}}\le\mathcal P_{ab},
\]
and this is \eqref{eq:pair-oneterm}, because $\nu_a=e^{a}P(a)$ and
$\nu_{b-a}=e^{b-a}P(b-a)$.

Next we bound the sum of the exponents over the pairs. We claim that
\begin{equation}\label{eq:pair-exponents}
\sum_{b=1}^{R}\ \sum_{0<a<b}P(a)\,P(b-a)\,X_{ab}
\le\alpha\,\mathcal G(k,R)+\beta\,\mathcal A(k,R).
\end{equation}
To prove this we use that the Poisson mass is positive, that $P_2$ and $P_3$
are nonnegative, and that $e^{-m}\,\nu_g\,\nu_{m-g}=P(g)\,P(m-g)$ for
$0\le g\le m$. We also use repeatedly the following consequence of
\eqref{eq:pair-line}: for an integer $m\ge0$ and a real number $y\ge m$,
\begin{equation}\label{eq:pair-row}
\sum_{0<g<m}P(g)\,P(m-g)\,L(y-g)
\le\Bigl(\alpha+\beta\,\bigl(y-\tfrac m2\bigr)\Bigr)\,P_2(m).
\end{equation}
Indeed, $y-g\ge0$ and hence $L(y-g)\le\alpha+\beta y-\beta g$ by
\eqref{eq:pair-line}, the weights $P(g)\,P(m-g)$ add up to $P_2(m)$ by
definition, and
$\sum_{0<g<m}g\,P(g)\,P(m-g)=\frac m2\,P_2(m)$ by
Lemma~\ref{lem:twofoldmean}.

Fix $b$ with $1\le b\le R$. We sum the parts of $X_{ab}$ over $a$ one after
the other. For the cut factor at $a$, \eqref{eq:pair-row} with $m=b$ and
$y=R$ gives directly
\[
\sum_{0<a<b}P(a)\,P(b-a)\,X^{(2)}_{a}
\le\Bigl(\alpha+\beta\,\bigl(R-\tfrac b2\bigr)\Bigr)\,P_2(b).
\]
For the first chain sum we have
$P(a)\,X^{(1)}_{a}=\sum_{0<c<a}P(c)\,P(a-c)\,L(R-c)$, because
$P(a)/\nu_a=e^{-a}$, and by \eqref{eq:pair-row} with $m=a$ and $y=R$ this is
at most $(\alpha+\beta(R-a/2))\,P_2(a)$. We multiply by $P(b-a)$ and sum
over $a$. Since $\sum_{0<a<b}P_2(a)\,P(b-a)=P_3(b)$ by definition and
$\sum_{0<a<b}a\,P_2(a)\,P(b-a)=\frac{2b}3\,P_3(b)$ by
Lemma~\ref{lem:threefoldleft}, we obtain
\[
\sum_{0<a<b}P(a)\,P(b-a)\,X^{(1)}_{a}
\le\Bigl(\alpha+\beta\,\bigl(R-\tfrac b3\bigr)\Bigr)\,P_3(b).
\]
For the second chain sum we have in the same way
\[
P(b-a)\,X^{(3)}_{ab}=\sum_{0<c<b-a}P(c)\,P(b-a-c)\,L(R-a-c),
\]
and \eqref{eq:pair-row} with $m=b-a$ and $y=R-a$ bounds this by
$(\alpha+\beta(R-b/2-a/2))\,P_2(b-a)$. We multiply by $P(a)$ and sum over
$a$. The substitution $a\mapsto b-a$ shows that
$\sum_{0<a<b}P(a)\,P_2(b-a)=P_3(b)$, and
$\sum_{0<a<b}a\,P(a)\,P_2(b-a)=\frac b3\,P_3(b)$ by
Lemma~\ref{lem:threefoldright}. Hence
\[
\sum_{0<a<b}P(a)\,P(b-a)\,X^{(3)}_{ab}
\le\Bigl(\alpha+\beta\,\bigl(R-\tfrac{2b}3\bigr)\Bigr)\,P_3(b).
\]
The parts $X^{(4)}_{b}$ and $X^{(5)}_{b}$ do not depend on $a$. By
Lemma~\ref{lem:borelmonotone} with $i=k$ and $j=N-b$, and by
\eqref{eq:pair-logexcess},
\[
X^{(5)}_{b}=\log a_k+(R-b)-\log a_{N-b}
\le\tfrac32\,\bigl(\log(N-b)-\log k\bigr)=\tfrac32\,L(R-b).
\]
Hence
$X^{(4)}_{b}+X^{(5)}_{b}\le\frac52L(R-b)\le\frac52(\alpha+\beta(R-b))$ by
\eqref{eq:pair-line}, and since the weights add up to $P_2(b)$,
\[
\sum_{0<a<b}P(a)\,P(b-a)\,\bigl(X^{(4)}_{b}+X^{(5)}_{b}\bigr)
\le\tfrac52\,\bigl(\alpha+\beta\,(R-b)\bigr)\,P_2(b).
\]
The sum of these four bounds is
\[
\alpha\Bigl(\tfrac72\,P_2(b)+2\,P_3(b)\Bigr)
+\beta\Bigl(\bigl(R-\tfrac b2\bigr)\bigl(P_2(b)+2\,P_3(b)\bigr)
+\tfrac52\,(R-b)\,P_2(b)\Bigr),
\]
where we have used that $(R-b/3)+(R-2b/3)=2(R-b/2)$. Summed over
$1\le b\le R$, this is $\alpha$ times the first sum of
\eqref{eq:pair-logweight} plus $\beta$ times the first sum of
\eqref{eq:pair-excess}.

It remains to treat the parts $X^{(6)}_{b}$. They do not depend on $a$
either, so that their contribution to the left-hand side of
\eqref{eq:pair-exponents} is $\sum_{b=1}^{R}P_2(b)\,X^{(6)}_{b}$. We show
that
\begin{equation}\label{eq:pair-countrow}
\sum_{b=1}^{R}P_2(b)\,X^{(6)}_{b}
\le\sum_{0<c<R}P_3(c)\,\bigl(1+\varrho(N,c)\bigr)\,
\bigl(\alpha+\beta\,(R-c)\bigr).
\end{equation}
The right-hand side is $\alpha$ times the second sum of
\eqref{eq:pair-logweight} plus $\beta$ times the second sum of
\eqref{eq:pair-excess}, so that \eqref{eq:pair-exponents} follows from
\eqref{eq:pair-countrow}. In the
sum that defines $X^{(6)}_{b}$ we substitute $t=N-c$. The range $k<t<N-b$
becomes $b<c<R$, and $\log(t/k)=L(R-c)$ by \eqref{eq:pair-logexcess}, so
that
\[
X^{(6)}_{b}=\sum_{b<c<R}\frac{\nu_{c-b}\,a_{N-c}}{a_{N-b}}\,L(R-c).
\]
By Lemma~\ref{lem:borelquotient} with $i=N-c$ and $j=N-b$, and since
$\nu_{c-b}=e^{c-b}P(c-b)$,
\[
\frac{\nu_{c-b}\,a_{N-c}}{a_{N-b}}
\le P(c-b)\,\Bigl(1+\frac32\cdot\frac{c-b}{N-c}
+\frac38\Bigl(\frac{c-b}{N-c}\Bigr)^2\Bigr).
\]
We insert this bound, using that $L(R-c)\ge0$ and $P_2(b)\ge0$, and
exchange the two summations, so that $c$ runs over $0<c<R$ and $b$ over
$0<b<c$. For a fixed $c$ let $D=N-c>0$ and $u_b=(c-b)/D$ for $0<b<c$.
Since $0\le c-b\le c$, we have $u_b^2\le(c/D)\,u_b$, and therefore
\[
1+\frac32\,u_b+\frac38\,u_b^2
\le1+(c-b)\Bigl(\frac{3}{2D}+\frac{3c}{8D^2}\Bigr).
\]
We multiply by $P_2(b)\,P(c-b)\ge0$ and sum over $0<b<c$. By the definition
of $P_3$ and by Lemma~\ref{lem:threefoldleft}, these weights add up to
$P_3(c)$ and $\sum_{0<b<c}(c-b)\,P_2(b)\,P(c-b)=\frac c3\,P_3(c)$. Hence
\begin{align*}
\sum_{0<b<c}P_2(b)\,P(c-b)\,\Bigl(1+\frac32\,u_b+\frac38\,u_b^2\Bigr)
&\le P_3(c)\,\Bigl(1+\frac c3\Bigl(\frac{3}{2D}+\frac{3c}{8D^2}\Bigr)\Bigr)\\
&=P_3(c)\,\bigl(1+\varrho(N,c)\bigr).
\end{align*}
Finally $L(R-c)\le\alpha+\beta(R-c)$ by \eqref{eq:pair-line}, and the factor
$P_3(c)\,(1+\varrho(N,c))$ is nonnegative. This proves
\eqref{eq:pair-countrow}.

By \eqref{eq:pair-tangent} with $X=X_{ab}$ and by
\eqref{eq:pair-oneterm}, for all integers $a$ and $b$ with $0<a<b\le R$,
\[
a_k\,e^{R}\,e^{-\mu}\,P(a)\,P(b-a)\,\bigl(1+\mu-X_{ab}\bigr)\le\mathcal P_{ab}.
\]
We sum over the pairs. The right-hand sides add up to $\mathcal P(k,R)$, by
\eqref{eq:pair-terms}. On the left the weights $P(a)\,P(b-a)$ add up to
$\sum_{b=1}^{R}P_2(b)=\mathcal W(R)$, so that the sum of the left-hand sides
is
\[
a_k\,e^{R}\,e^{-\mu}\Bigl((1+\mu)\,\mathcal W(R)
-\sum_{b=1}^{R}\ \sum_{0<a<b}P(a)\,P(b-a)\,X_{ab}\Bigr).
\]
By \eqref{eq:pair-exponents} the double sum is at most
$\alpha\,\mathcal G(k,R)+\beta\,\mathcal A(k,R)$, which gives the claim.
\end{proof}

The two parameters
are left free because the proposition is used
with bounds for the three sums instead of the sums themselves.
Section~\ref{sec:sums} proves a lower bound for $\mathcal W(R)$ and upper
bounds for $\mathcal G(k,R)$ and $\mathcal A(k,R)$, and
Section~\ref{sec:curve} inserts them and chooses $\theta$ and $\mu$ in terms
of them.

To see what the proposition gives, let $R\ge2$, write $\mathcal W$,
$\mathcal G$ and $\mathcal A$ for the three sums, suppose that they are
positive, and take $\theta=\mathcal A/(k\,\mathcal G)$ and
$\mu=(\mathcal G/\mathcal W)\log(1+\theta)$. For this $\theta$ we have
$\beta\,\mathcal A=\theta\,\mathcal G/(1+\theta)$, so that
$\alpha\,\mathcal G+\beta\,\mathcal A=\mathcal G\log(1+\theta)
=\mu\,\mathcal W$, and the proposition reads
\[
\mathcal P(k,R)\ge a_k\,e^{R}\;\mathcal W\,
\Bigl(1+\frac{\mathcal A}{k\,\mathcal G}\Bigr)^{-\mathcal G/\mathcal W}.
\]
In this form the bound can be read as follows: there are
$\mathcal W$ pairs, counted with their weights, and each of them carries in
the mean $\mathcal G/\mathcal W$ cut factors $k/(k+y)$, with
$y=\mathcal A/\mathcal G$ in the mean.

\subsection{The upper bound of the boundary chain sum}
\label{subsec:pair-boundary}

The left-hand side of the criterion of Proposition~\ref{prop:criterion}
contains the chain sum
$\Lambda_R(k+1;k+R)$, whose cut factors are $(k+1)/(k+R-s)$ for $0<s<R$.
The
following proposition bounds $e^{-R}$ times this chain sum from above by
$P(R)$ minus an explicit expression. It concerns the chain sum
$\Lambda_R(y+1;y+R)$, in which
column $k$ is replaced by a real number $y\ge0$, and it contains two
parameters. The number $K$ is an upper bound for
$y$: the bound holds for every $y\le K$, and Section~\ref{sec:curve} uses
this to apply
the proposition with a number $K$ that depends on $R$ alone. The number
$\mu$ is the point at which the tangent \eqref{eq:pair-tangent} is taken;
Section~\ref{sec:curve} uses one fixed numerical value.

\begin{proposition}\label{prop:boundaryupper}
Let $R\ge1$ be an integer, let $y$ and $K$ be real numbers with
$0\le y\le K$ and $R\le K$, and let $\mu$ be a real number with
$1+\mu\ge0$. Then the chain sum $\Lambda_R(y+1;y+R)$ satisfies
\begin{multline*}
e^{-R}\,\Lambda_R(y+1;y+R)\le P(R)\\
-e^{-\mu}\Bigl((1+\mu)\,\frac{R-2}{2\,(K+R-1)}\,P_2(R)
-\frac{1}{K+1}\cdot\frac{R-3}{3}\,P_3(R)\cdot
\frac{\frac45R-1}{K+\frac35R}\Bigr).
\end{multline*}
\end{proposition}

\begin{proof}
For an integer $s$ with $0<s<R$ let
\[
\varphi(s)=1-\frac{K+1}{K+R-s}=\frac{R-1-s}{K+R-s}.
\]
Since $K\ge R\ge1$, the denominator is positive and $0\le\varphi(s)<1$. For
$1\le s\le R$ we abbreviate
$\widetilde\Lambda_s=e^{-s}\,\Lambda_s(K+1;K+R)$.

We first replace $y$ by $K$ and show that
\begin{equation}\label{eq:pair-raise}
e^{-R}\,\Lambda_R(y+1;y+R)\le\widetilde\Lambda_R.
\end{equation}
At a cut point $0<s<R$ the cut factor $(y+1)/(y+R-s)$ is positive, and it
is at most the cut factor $(K+1)/(K+R-s)$, since cross-multiplication turns
this inequality into $(K-y)(R-s-1)\ge0$. The total weight $\Lambda_R[f]$
of the compositions of $R$ does not decrease when a nonnegative cut profile
$f$ is increased at the cut points, because the weight of every composition
is a product of nonnegative Cayley coefficients and of values of the
profile at its cut points (Section~\ref{sec:compositions}). This gives
\eqref{eq:pair-raise}.

Next we show that, for every real number $\mu$,
\begin{equation}\label{eq:pair-twosums}
\widetilde\Lambda_R\le P(R)-e^{-\mu}\Bigl((1+\mu)\,I_1
-\frac{1}{K+1}\,I_2\Bigr),
\end{equation}
where
\begin{align*}
I_1&=\sum_{0<s<R}\varphi(s)\,P(s)\,P(R-s),\\
I_2&=\sum_{0<s<R}\varphi(s)\,\bigl(R-1-\tfrac s2\bigr)\,P(R-s)\,P_2(s).
\end{align*}
By Lemma~\ref{lem:firstviolation} for the threshold profile
$f=\lambda^{K+1,K+R}$, for which $1-f(s)=\varphi(s)$,
\[
\nu_R-\Lambda_R(K+1;K+R)
=\sum_{0<s<R}\nu_{R-s}\,\varphi(s)\,\Lambda_s(K+1;K+R).
\]
We multiply by $e^{-R}=e^{-s}\,e^{-(R-s)}$ and obtain
\begin{equation}\label{eq:pair-shortfall}
P(R)-\widetilde\Lambda_R=\sum_{0<s<R}\varphi(s)\,P(R-s)\,\widetilde\Lambda_s.
\end{equation}
The chain sums on the right are bounded from below by
Lemma~\ref{lem:compositiontangent}. Let $0<s<R$. At a cut point $0<g<s$ the
cut factor is $(K+1)/(K+R-g)=e^{-x(g)}$ with
$x(g)=\log\bigl((K+R-g)/(K+1)\bigr)$. The lemma, with $m=s$ and the same $\mu$
and multiplied by $e^{-s}$, therefore gives
\[
e^{-\mu}\Bigl((1+\mu)\,P(s)-\sum_{0<g<s}P(g)\,P(s-g)\,x(g)\Bigr)
\le\widetilde\Lambda_s;
\]
here we have used that $e^{-s}\,\nu_g\,\nu_{s-g}=P(g)\,P(s-g)$. Since
$\log t\le t-1$, we have $x(g)\le(K+R-g)/(K+1)-1=(R-1-g)/(K+1)$. The weights
$P(g)\,P(s-g)$ are positive, their sum is $P_2(s)$, and
$\sum_{0<g<s}g\,P(g)\,P(s-g)=\frac s2\,P_2(s)$ by
Lemma~\ref{lem:twofoldmean}. Hence
\begin{equation}\label{eq:pair-before}
e^{-\mu}\Bigl((1+\mu)\,P(s)-\frac{1}{K+1}\,\bigl(R-1-\tfrac s2\bigr)\,
P_2(s)\Bigr)\le\widetilde\Lambda_s.
\end{equation}
In \eqref{eq:pair-shortfall} every term has the nonnegative factors
$\varphi(s)$ and $P(R-s)$, so that \eqref{eq:pair-before} can be inserted,
and the result is \eqref{eq:pair-twosums}.

It remains to bound $I_1$ from below and $I_2$ from above. For
$0<s<R$ we have $R-1-s\ge0$ and $K+R-s\le K+R-1$, hence
$\varphi(s)\ge(R-1-s)/(K+R-1)$. By the definition of $P_2$ and by
Lemma~\ref{lem:twofoldmean},
$\sum_{0<s<R}(R-1-s)\,P(s)\,P(R-s)=(R-1-R/2)\,P_2(R)$, and therefore
\begin{equation}\label{eq:pair-firstsum}
I_1\ge\frac{R-2}{2\,(K+R-1)}\,P_2(R).
\end{equation}
For the second sum we write $I_2=\sum_{0<s<R}\vartheta_s\,\chi(s)$ with
the weights $\vartheta_s=P(R-s)\,(R-1-s)\,P_2(s)$ of
Section~\ref{sec:birthday}, which are nonnegative, and with
\[
\chi(s)=\frac{R-1-\frac s2}{K+R-s}
=\frac12-\frac{K-R+2}{2\,(K+R-s)}.
\]
Since $K-R+2>0$, the second expression shows that $\chi$ is a concave
function of the real variable $s<K+R$. We know the sum of the weights and
a bound for their mean position:
$\sum_{0<s<R}\vartheta_s=\frac{R-3}{3}\,P_3(R)$ by
Corollary~\ref{cor:weightsum}, and
$\frac{2R}{5}\sum_{0<s<R}\vartheta_s\le\sum_{0<s<R}s\,\vartheta_s$ by
Lemma~\ref{lem:weightmean}. Accordingly we compare $\chi$ with its tangent
at the point $s^*=2R/5$. Let $D=K+R-s^*=K+\frac35R$. The derivative of
$\chi$ at $s^*$ is $-(K-R+2)/(2D^2)$, and a concave function lies below its
tangents, so that
\[
\chi(s)\le\chi(s^*)-\frac{K-R+2}{2D^2}\,(s-s^*)\qquad(0<s<R).
\]
We multiply by $\vartheta_s\ge0$ and sum over $0<s<R$:
\begin{align*}
I_2&\le\chi(s^*)\sum_{0<s<R}\vartheta_s-\frac{K-R+2}{2D^2}
\Bigl(\sum_{0<s<R}s\,\vartheta_s-\frac{2R}{5}\sum_{0<s<R}\vartheta_s\Bigr)\\
&\le\chi(s^*)\sum_{0<s<R}\vartheta_s,
\end{align*}
because $K-R+2>0$ and the expression in parentheses is nonnegative. Since
$\chi(s^*)=(\frac45R-1)/(K+\frac35R)$, this is
\begin{equation}\label{eq:pair-secondsum}
I_2\le\frac{R-3}{3}\,P_3(R)\cdot\frac{\frac45R-1}{K+\frac35R}.
\end{equation}
The proposition now follows from \eqref{eq:pair-raise} and
\eqref{eq:pair-twosums}, when $I_1$ and $I_2$ are replaced by the
bounds \eqref{eq:pair-firstsum} and \eqref{eq:pair-secondsum}; for $I_1$
we use that $1+\mu\ge0$.
\end{proof}

The expression in parentheses on the right-hand side of
Proposition~\ref{prop:boundaryupper} has a main term, with
$P_2(R)$, and a correction of the opposite sign, with $P_3(R)$. The main
term is the effect of the cut factors to first order: in the chain sum
$\Lambda_R(K+1;K+R)$ the cut factor at $s$ is $(K+1)/(K+R-s)$, so that a cut
at $s$ loses the proportion $(R-1-s)/(K+R-s)$ of the weight, and summing
this over the position of one cut gives a multiple of $P_2(R)$. The
correction accounts for the cut factors at the cut points before the one
considered, which have reduced the weight already.

\section{Three sums of convolutions}\label{sec:sums}
Proposition~\ref{prop:pairlower} bounds the pair sum $\mathcal P(k,R)$ from
below by an expression that contains three sums of the convolutions $P_2$
and $P_3$ of the Poisson mass: the total weight of the pairs,
$\mathcal W(R)$, and the sums $\mathcal G(k,R)$ and $\mathcal A(k,R)$ of
Section~\ref{sec:pair}. In that bound $\mathcal W$ carries a factor that is
positive for the parameters that will be chosen, and $\mathcal G$ and
$\mathcal A$ are subtracted, each with a nonnegative factor. We therefore
bound $\mathcal W$ from below and $\mathcal G$ and $\mathcal A$ from above,
by explicit polynomials in $s=\sqrt R$ (Lemmas~\ref{lem:weightlower},
\ref{lem:logupper} and~\ref{lem:excessupper}), in the range in which the
proof of Proposition~\ref{prop:curvestep} needs them: for $R\ge42$ and,
where the sum depends on $k$, for $R^{3/2}\le1.07593\,k$.

\subsection{The three sums and the scheme of the estimates}
\label{subsec:sums-scheme}

The three sums are defined in Section~\ref{sec:pair}, for integers $k\ge0$
and $R\ge0$: the total weight of the pairs $\mathcal W(R)$ in
\eqref{eq:pair-weight}, the sum $\mathcal G(k,R)$ in
\eqref{eq:pair-logweight} and the sum $\mathcal A(k,R)$ in
\eqref{eq:pair-excess}. The definitions of $\mathcal G$ and $\mathcal A$
contain the numbers $\varrho(k+R,c)$, which we call corrections; the
function $\varrho$ is defined there as well, immediately before
\eqref{eq:pair-weight}.

The three proofs follow one scheme. Each of the sums $\mathcal G$ and
$\mathcal A$ consists of a \emph{main part}, a sum of multiples of $P_2(b)$
and $P_3(b)$ over $1\le b\le R$, and of a \emph{correction part}, in which
the numbers $P_3(c)$ with $0<c<R$ are multiplied by the corrections
$\varrho(k+R,c)$. In the main part the closed form
$P_3(b)=(b-1)\,P(b)-2\,P_2(b)$ of Lemma~\ref{lem:threefold} removes $P_3$.
What is left of the term of index $b$ is a nonnegative multiple of the
Poisson mass $P(b)$ minus a positive multiple of $P_2(b)$. An upper bound of
the term therefore needs an upper bound of $P(b)$, which is that of
\eqref{eq:prelim-stirling}, and a lower bound of $P_2(b)$, which
is Corollary~\ref{cor:twofoldlower}; the same lower bound of $P_2$ gives the
lower bound of $\mathcal W$ directly. With these inserted, the term is a
combination of the powers $b^{3/2}$, $b^{1/2}$, $b^{-1/2}$ and $b^{-3/2}$
and of a polynomial in $b$, and the sum over $b$ is bounded by an explicit
function of $s$. The correction part is bounded more roughly.

For $\mathcal W$ and for $\mathcal A$ we state the bounds that the estimates
themselves produce: the two polynomials consist of the leading terms that result from
the scheme above, with decimal coefficients that are rounded in the
direction that keeps the inequality.

In the numerical comparisons we use that, for every integer $R\ge42$,
\begin{equation}\label{eq:sums-halfline}
s=\sqrt R\ge6.48,
\end{equation}
because $6.48^2\le42$.

\subsection{Sums of half-integer powers}\label{subsec:sums-powers}

For an integer $R\ge2$ we write
\[
\Sigma_j=\sum_{b=2}^{R}b^{\,j/2}\qquad\text{for }j\in\{-3,-1,1,3\},
\]
without showing the dependence on $R$. These are the sums to which
$\mathcal W$ and the main parts of $\mathcal G$ and $\mathcal A$ lead. They
begin at $b=2$ because the terms of index $1$ will be seen
to vanish, and because the lower bound of $P_2(b)$ in
Corollary~\ref{cor:twofoldlower} holds from $b=2$ on. The correction parts
lead to two sums of the powers $c^{3/2}$ over the integers $c$ with
$0<c<R$. We use the following bounds. For every integer $R\ge2$,
\begin{align}
\Sigma_{-1}=\sum_{b=2}^{R}\frac1{\sqrt b}
&\le2\sqrt R+\frac1{2\sqrt R}-2.44948,\label{eq:sums-minusone}\\
\Sigma_{-3}=\sum_{b=2}^{R}\frac1{b\sqrt b}
&\le1.63300,\label{eq:sums-minusthree}\\
\Sigma_{1}=\sum_{b=2}^{R}\sqrt b
&\le\frac23\,R\sqrt R+\frac12\,\sqrt R-1.17851,\label{eq:sums-one}\\
\Sigma_{3}=\sum_{b=2}^{R}b\sqrt b
&\ge\frac25\,R^2\sqrt R+\frac12\,R\sqrt R-0.84853,\label{eq:sums-three}
\end{align}
and for every integer $R\ge0$,
\begin{align}
\sum_{0<c<R}c\sqrt c&\le\frac25\,R^2\sqrt R,\label{eq:sums-open}\\
\sum_{0<c<R}c\sqrt c\,(R-c)&\le\frac4{35}\,(R+1)^3\sqrt{R+1}.
\label{eq:sums-weighted}
\end{align}
The bound \eqref{eq:sums-three} is a lower bound because $\Sigma_3$ will
occur with a negative coefficient.

Each of the six bounds compares a sum with an integral. For the first four
we use the midpoint rule and the trapezoid rule, which for a convex or a
concave function err in a known direction. Let $\varphi$ be a
differentiable convex function on an interval $[x,x+1]$. Then
\begin{equation}\label{eq:sums-convex}
\varphi\Bigl(x+\frac12\Bigr)\le\int_x^{x+1}\varphi(t)\,dt
\le\frac{\varphi(x)+\varphi(x+1)}2 .
\end{equation}
Indeed, $\varphi$ lies above its tangent at the point $x+1/2$ and below its
chord. For a concave function both inequalities are reversed.

The functions $t^{-1/2}$ and $t^{-3/2}$ are convex on the positive
half-line. We apply the first inequality of \eqref{eq:sums-convex} to them
on the intervals $[b-1/2,\,b+1/2]$ with $2\le b\le R$, which together make
up the interval $[3/2,\,R+1/2]$, and obtain
\begin{align*}
\Sigma_{-1}&\le\int_{3/2}^{R+1/2}\frac{dt}{\sqrt t}
=2\sqrt{R+\frac12}-\sqrt6,\\
\Sigma_{-3}&\le\int_{3/2}^{R+1/2}\frac{dt}{t\sqrt t}
=\sqrt{\frac83}-\frac2{\sqrt{R+1/2}}.
\end{align*}
Since $\sqrt{R+1/2}\le\sqrt R+1/(4\sqrt R)$ and $\sqrt6\ge2.44948$, the
first of these gives \eqref{eq:sums-minusone}. In the second we leave out
the negative term and use $\sqrt{8/3}\le1.63300$; this gives
\eqref{eq:sums-minusthree}.

The function $\sqrt t$ is concave and the function $t\sqrt t$ is convex on
the half-line $t\ge0$. We apply the second inequality of
\eqref{eq:sums-convex}, reversed for $\sqrt t$, on the intervals $[b,b+1]$
with $2\le b\le R-1$, and add. Since
\[
\sum_{b=2}^{R-1}\frac{\varphi(b)+\varphi(b+1)}2
=\sum_{b=2}^{R}\varphi(b)-\frac{\varphi(2)+\varphi(R)}2
\]
for every function $\varphi$ and every integer $R\ge2$, this gives
\begin{align*}
\Sigma_1-\frac{\sqrt2+\sqrt R}2&\le\int_2^R\sqrt t\,dt
=\frac23\bigl(R\sqrt R-2\sqrt2\bigr),\\
\Sigma_3-\frac{2\sqrt2+R\sqrt R}2&\ge\int_2^Rt\sqrt t\,dt
=\frac25\bigl(R^2\sqrt R-4\sqrt2\bigr),
\end{align*}
and therefore
\[
\Sigma_1\le\frac23\,R\sqrt R+\frac12\,\sqrt R-\frac56\,\sqrt2,\qquad
\Sigma_3\ge\frac25\,R^2\sqrt R+\frac12\,R\sqrt R-\frac35\,\sqrt2 .
\]
Since $(5/6)\sqrt2\ge1.17851$ and $(3/5)\sqrt2\le0.84853$, the
bounds \eqref{eq:sums-one} and \eqref{eq:sums-three} follow.

For the last two bounds monotonicity is enough, since they serve for the
correction parts only. The functions $t\sqrt t$ and $t^2\sqrt t$ increase on
the half-line $t\ge0$, so that the value of each at an integer $c\ge0$ is at
most its integral over $[c,c+1]$. For $t\sqrt t$ this gives
\eqref{eq:sums-open}:
\[
\sum_{0<c<R}c\sqrt c\le\int_0^{R}t\sqrt t\,dt=\frac25\,R^2\sqrt R .
\]
In \eqref{eq:sums-weighted} the summand $c\sqrt c\,(R-c)$ is not monotone in
$c$. For $R\le1$ the sum is empty, so let $R\ge2$. We write the factor $R-c$
of the summand as the number of integers $n$ with $c\le n<R$, which turns the
sum into a double sum of an increasing function:
\begin{align*}
\sum_{0<c<R}c\sqrt c\,(R-c)=\sum_{n=1}^{R-1}\,\sum_{c=1}^{n}c\sqrt c
&\le\frac25\sum_{n=1}^{R-1}(n+1)^2\sqrt{n+1}\\
&\le\frac25\int_2^{R+1}t^2\sqrt t\,dt
\le\frac25\cdot\frac27\,(R+1)^3\sqrt{R+1}.
\end{align*}
Here the first inequality is \eqref{eq:sums-open} with $R$ replaced by
$n+1$, and the second is the same comparison for $t^2\sqrt t$ at the
integers $m=n+1$ from $2$ to $R$, whose intervals $[m,m+1]$ make up the
interval $[2,R+1]$.

\subsection{The single terms}\label{subsec:sums-terms}

The terms of the three sums are bounded through the upper bound
$P(b)\le c_{\mathrm S}/\sqrt b$ of \eqref{eq:prelim-stirling}, through two
inequalities for the convolutions $P_2$ and $P_3$ at one point, and through
a bound of the correction. For the constant $c_{\mathrm S}$ we use the bound
$c_{\mathrm S}\le0.39894229$ of Section~\ref{sec:prelim}.

The proofs below derive the two inequalities for $P_2$ and $P_3$ from
Corollary~\ref{cor:twofoldlower} and Lemma~\ref{lem:threefold}. The numbers
$P_2(m)$ and $P_3(m)$ are sums of products of values of the Poisson mass and
hence nonnegative for every integer $m\ge0$. We also use that $P_2(1)=0$.

The correction is bounded as follows. Let $k\ge1$ and $R$ be integers such
that $s=\sqrt R$ satisfies $s\ge6.48$, and let $R^{3/2}\le1.07593\,k$. Then
for every integer $c$ with $0<c<R$,
\begin{equation}\label{eq:sums-correction}
0\le\varrho(k+R,c)\le0.560302\,\frac{c}{s^3}.
\end{equation}
To see this, let $N=k+R$, so that $N-c>k$ and
\[
\varrho(N,c)=\frac{c}{2\,(N-c)}+\frac{c^2}{8\,(N-c)^2}
=\frac{c}{2\,(N-c)}\Bigl(1+\frac{c}{4\,(N-c)}\Bigr),
\]
which is nonnegative. The hypothesis on $k$ gives
$1/(N-c)\le1/k\le1.07593/s^3$. With $c\le R=s^2$ and
$s\ge6.48$ it follows that
$c/(N-c)\le1.07593/s\le1.07593/6.48$. Hence
\[
\varrho(N,c)\le\frac{1.07593\,c}{2s^3}
\Bigl(1+\frac{1.07593}{4\cdot6.48}\Bigr)\le0.560302\,\frac{c}{s^3}.
\]

\subsection{The total weight of the pairs}\label{subsec:sums-weight}

We begin with the sum $\mathcal W$, which has no correction part and does
not contain $P_3$.

\begin{lemma}\label{lem:weightlower}
Let $R\ge42$ be an integer. Then the total weight of the pairs satisfies
\[
\frac R2-1.06385\,\sqrt R+0.7184\le\mathcal W(R).
\]
\end{lemma}

\begin{proof}
Let $s=\sqrt R$. For every integer $b\ge2$,
\[
P_2(b)\ge\frac12-\frac43\cdot\frac{c_{\mathrm S}}{\sqrt b}
-\frac1{15}\cdot\frac{c_{\mathrm S}}{b\sqrt b},
\]
by Corollary~\ref{cor:twofoldlower} and the upper bound of
\eqref{eq:prelim-stirling}. The term $b=1$ of $\mathcal W(R)$ is $P_2(1)=0$, and the
sum of this bound over the $R-1$ integers $b$ with
$2\le b\le R$ gives
\[
\mathcal W(R)\ge\frac{R-1}2-\frac43\,c_{\mathrm S}\,\Sigma_{-1}
-\frac1{15}\,c_{\mathrm S}\,\Sigma_{-3}.
\]
Here $c_{\mathrm S}$ may be replaced by its upper bound $0.39894229$. By
\eqref{eq:sums-minusthree} we have $\Sigma_{-3}\le1.633$. In
\eqref{eq:sums-minusone} we have $\sqrt R=s$ and, by
\eqref{eq:sums-halfline}, $1/(2s)\le1/12.96<0.07717$, so that
$\Sigma_{-1}\le2s-2.37231$. With $R=s^2$ we obtain
\begin{align*}
\mathcal W(R)\ge{}&\frac{s^2}2-\frac12
-\frac43\cdot0.39894229\cdot(2s-2.37231)\\
&-\frac1{15}\cdot0.39894229\cdot1.633 .
\end{align*}
On the right-hand side the coefficient of $s$ is
$-(8/3)\cdot0.39894229\ge-1.06385$, and the constant term is
\begin{multline*}
-\frac12+\frac43\cdot0.39894229\cdot2.37231\\
-\frac1{15}\cdot0.39894229\cdot1.633\ge0.7184 .
\end{multline*}
Hence the right-hand side is at least $s^2/2-1.06385\,s+0.7184$, which is
the claim.
\end{proof}

\subsection{The sum $\mathcal G$}\label{subsec:sums-log}

We turn to the two sums that are subtracted in
Proposition~\ref{prop:pairlower}, and begin with $\mathcal G$. The
comparison of Section~\ref{sec:curve} goes through with a weaker bound of
$\mathcal G$ than the estimate produces, and the following lemma states such
a bound: its polynomial is the polynomial that bounds $\mathcal A(k,R)$ in
Lemma~\ref{lem:excessupper}, multiplied by $2/s^2$. With this ratio, one of
the two free parameters of Proposition~\ref{prop:pairlower} can be taken
there so that, after the sums have been replaced by their bounds, the term
with $\mathcal A$ cancels against a part of the term with $\mathcal G$.

\begin{lemma}\label{lem:logupper}
Let $k\ge1$ and $R\ge42$ be integers with $R^{3/2}\le1.07593\,k$, and let
$s=\sqrt R$. Then
\[
\mathcal G(k,R)\le2\,\bigl(0.478731\,s^3-0.974453\,s^2+0.642741\,s\bigr).
\]
\end{lemma}

\begin{proof}
We split the factor $1+\varrho(k+R,c)$ of the second sum in the definition
\eqref{eq:pair-logweight} of $\mathcal G(k,R)$ into its two terms. The sum
of $P_3(c)$ over $0<c<R$ is at most the sum over $1\le c\le R$, because
$P_3(R)\ge0$. Hence
$\mathcal G(k,R)\le\mathcal G'+\mathcal G''$, where
\[
\mathcal G'=\sum_{b=1}^{R}\Bigl(\frac72\,P_2(b)+3\,P_3(b)\Bigr),\qquad
\mathcal G''=\sum_{0<c<R}P_3(c)\,\varrho(k+R,c).
\]
By Lemma~\ref{lem:threefold} we have $P_3(b)=(b-1)\,P(b)-2\,P_2(b)$ for
every $b\ge1$, so that the term of index $b$ of $\mathcal G'$ is
\[
\frac72\,P_2(b)+3\,P_3(b)=3\,(b-1)\,P(b)-\frac52\,P_2(b).
\]
For $b=1$ it is zero, because $P_2(1)=0$. For $2\le b\le R$,
\[
P_2(b)\ge\frac12-\frac43\cdot\frac{c_{\mathrm S}}{\sqrt b}
-\frac1{15}\cdot\frac{c_{\mathrm S}}{b\sqrt b}
\]
by Corollary~\ref{cor:twofoldlower} and the upper bound of
\eqref{eq:prelim-stirling}, and we insert this lower bound and the upper
bound of \eqref{eq:prelim-stirling}:
\begin{align*}
3\,(b-1)\,P(b)-\frac52\,P_2(b)
&\le3\,c_{\mathrm S}\Bigl(\sqrt b-\frac1{\sqrt b}\Bigr)-\frac54
+\frac{10}3\cdot\frac{c_{\mathrm S}}{\sqrt b}
+\frac16\cdot\frac{c_{\mathrm S}}{b\sqrt b}\\
&=c_{\mathrm S}\Bigl(3\sqrt b+\frac13\cdot\frac1{\sqrt b}
+\frac16\cdot\frac1{b\sqrt b}\Bigr)-\frac54 .
\end{align*}
Summing over the $R-1$ integers $b$ with $2\le b\le R$, we obtain
\[
\mathcal G'\le c_{\mathrm S}\Bigl(3\,\Sigma_1+\frac13\,\Sigma_{-1}
+\frac16\,\Sigma_{-3}\Bigr)-\frac54\,(R-1).
\]
For every integer $c\ge1$, Lemma~\ref{lem:threefold} gives
$P_3(c)=(c-1)\,P(c)-2\,P_2(c)\le c\,P(c)$, and
$c\,P(c)\le c_{\mathrm S}\sqrt c$ by the upper bound of
\eqref{eq:prelim-stirling}, so that
\[
P_3(c)\le c_{\mathrm S}\,\sqrt c .
\]
The two factors of a term of $\mathcal G''$ are bounded by this inequality
and by \eqref{eq:sums-correction}, which applies because $s\ge6.48$ by
\eqref{eq:sums-halfline}; with
\eqref{eq:sums-open} and $R^2\sqrt R=s^5$ this gives
\[
\mathcal G''\le\frac{0.560302\,c_{\mathrm S}}{s^3}\sum_{0<c<R}c\sqrt c
\le0.560302\cdot\frac25\,c_{\mathrm S}\,s^2 .
\]

It remains to insert the bounds of the power sums and of $c_{\mathrm S}$.
By \eqref{eq:sums-one} we have $3\,\Sigma_1\le2s^3+3s/2-3.53553$. By
\eqref{eq:sums-minusone}, in which $1/(2s)\le1/12.96<0.07717$, we have
$\Sigma_{-1}\le2s-2.37231$ and hence $\Sigma_{-1}/3\le2s/3-0.79077$. By
\eqref{eq:sums-minusthree} we have $\Sigma_{-3}/6\le0.27217$. Adding these
three bounds, we find
\[
3\,\Sigma_1+\frac13\,\Sigma_{-1}+\frac16\,\Sigma_{-3}
\le2s^3+\frac{13}6\,s-4.0541 .
\]
Since $c_{\mathrm S}\le0.39894229$, it follows that
\begin{align*}
c_{\mathrm S}\Bigl(3\,\Sigma_1+\frac13\,\Sigma_{-1}
+\frac16\,\Sigma_{-3}\Bigr)
&\le0.39894229\,\Bigl(2s^3+\frac{13}6\,s-4.0541\Bigr)\\
&\le0.79788458\,s^3+0.86438\,s-1.61735 .
\end{align*}
Likewise $0.560302\cdot(2/5)\,c_{\mathrm S}\le0.08942$. With $R=s^2$ we
arrive at
\begin{align*}
\mathcal G(k,R)
&\le0.79788458\,s^3+0.86438\,s-1.61735-\frac54\,(s^2-1)+0.08942\,s^2\\
&\le0.79788458\,s^3-1.16058\,s^2+0.86438\,s ,
\end{align*}
where the constant $5/4-1.61735$, which is negative, has been left out. The
right-hand side of the claim minus this bound is
\[
0.15957742\,s^3-0.788326\,s^2+0.421102\,s ,
\]
and this is positive, because by \eqref{eq:sums-halfline}
\[
0.15957742\,s\ge0.15957742\cdot6.48>0.788326 .\qedhere
\]
\end{proof}

\subsection{The sum $\mathcal A$}\label{subsec:sums-excess}

The bound of $\mathcal A$ is obtained in the same way as that of
$\mathcal G$, with the weights $R-b/2$, $R-b$ and $R-c$ in the sums and with
a longer computation.

\begin{lemma}\label{lem:excessupper}
Let $k\ge1$ and $R\ge42$ be integers with $R^{3/2}\le1.07593\,k$, and let
$s=\sqrt R$. Then
\[
\mathcal A(k,R)\le0.478731\,s^5-0.974453\,s^4+0.642741\,s^3 .
\]
\end{lemma}

\begin{proof}
We split the factor $1+\varrho(k+R,c)$ of the second sum in the definition
\eqref{eq:pair-excess} of $\mathcal A(k,R)$ into its two terms. The sum of
$(R-c)\,P_3(c)$ over $0<c<R$ does not change when the term $c=R$, which is
zero, is added to it. Hence
$\mathcal A(k,R)=\mathcal A'+\mathcal A''$, where the correction part is
\[
\mathcal A''=\sum_{0<c<R}(R-c)\,P_3(c)\,\varrho(k+R,c)
\]
and the main part $\mathcal A'$ is the sum over $1\le b\le R$ of the terms
\[
\Bigl(R-\frac b2\Bigr)\bigl(P_2(b)+2\,P_3(b)\bigr)
+\frac52\,(R-b)\,P_2(b)+(R-b)\,P_3(b).
\]

\emph{Step 1: the terms of the main part.}
Inserting $P_3(b)=(b-1)\,P(b)-2\,P_2(b)$, which holds for $b\ge1$ by
Lemma~\ref{lem:threefold}, and collecting terms, we find that the term of
index $b$ of $\mathcal A'$ is
\[
(3R-2b)\,(b-1)\,P(b)-\Bigl(\frac52\,R-b\Bigr)P_2(b).
\]
For $b=1$ it is zero, because $P_2(1)=0$. Let $2\le b\le R$. Then
$(3R-2b)(b-1)\ge0$ and $5R/2-b>0$, and
\[
P_2(b)\ge\frac12-\frac43\cdot\frac{c_{\mathrm S}}{\sqrt b}
-\frac1{15}\cdot\frac{c_{\mathrm S}}{b\sqrt b}
\]
by Corollary~\ref{cor:twofoldlower} and the upper bound of
\eqref{eq:prelim-stirling}, so that the upper bound of
\eqref{eq:prelim-stirling} and this lower bound may be inserted:
\begin{align*}
(3R-2b)\,(b-1)\,P(b)
&\le c_{\mathrm S}\,(3R-2b)\Bigl(\sqrt b-\frac1{\sqrt b}\Bigr),\\
-\Bigl(\frac52\,R-b\Bigr)P_2(b)
&\le-\frac12\Bigl(\frac52\,R-b\Bigr)
+c_{\mathrm S}\Bigl(\frac52\,R-b\Bigr)
\Bigl(\frac4{3\sqrt b}+\frac1{15\,b\sqrt b}\Bigr).
\end{align*}
Multiplying out the two products on the right, we see that the term of
index $b$ of $\mathcal A'$ is at most
\begin{multline*}
c_{\mathrm S}\Bigl(-2\,b\sqrt b+\Bigl(3R+\frac23\Bigr)\sqrt b
+\Bigl(\frac R3-\frac1{15}\Bigr)\frac1{\sqrt b}
+\frac R6\cdot\frac1{b\sqrt b}\Bigr)\\
-\frac12\Bigl(\frac52\,R-b\Bigr).
\end{multline*}

\emph{Step 2: the sum of the main part.}
We sum the bound of Step~1 over $2\le b\le R$. Its last terms add up to
\[
\sum_{b=2}^{R}\frac12\Bigl(\frac52\,R-b\Bigr)=R^2-\frac32\,R+\frac12 .
\]
The factors of $c_{\mathrm S}$ add up to the number
\[
Y=-2\,\Sigma_3+\Bigl(3R+\frac23\Bigr)\Sigma_1
+\Bigl(\frac R3-\frac1{15}\Bigr)\Sigma_{-1}+\frac R6\,\Sigma_{-3} ,
\]
and therefore
\[
\mathcal A'\le c_{\mathrm S}\,Y-\Bigl(R^2-\frac32\,R+\frac12\Bigr).
\]
In $Y$ only $\Sigma_3$ has a negative coefficient. We replace $\Sigma_1$,
$\Sigma_{-1}$ and $\Sigma_{-3}$ by their upper bounds \eqref{eq:sums-one},
\eqref{eq:sums-minusone} and \eqref{eq:sums-minusthree} and $\Sigma_3$ by
its lower bound \eqref{eq:sums-three}. With $\sqrt R=s$ and $R=s^2$ this
gives
\begin{multline*}
Y\le-2\,\Bigl(\frac25\,s^5+\frac12\,s^3-0.84853\Bigr)
+\Bigl(3s^2+\frac23\Bigr)\Bigl(\frac23\,s^3+\frac12\,s-1.17851\Bigr)\\
+\Bigl(\frac{s^2}3-\frac1{15}\Bigr)\Bigl(2s+\frac1{2s}-2.44948\Bigr)
+\frac{s^2}6\cdot1.633 .
\end{multline*}
Expanding the products, we find that the right-hand side is a polynomial in
$s$ plus the term $-1/(30s)$, which is negative. The polynomial has the
coefficients $-4/5+2=6/5$ of $s^5$, $-1+3/2+4/9+2/3=29/18$ of $s^3$ and
$1/3+1/6-2/15=11/30$ of $s$, and for
its coefficient of $s^2$ and its constant term we have
\begin{align*}
-3\cdot1.17851-\frac{2.44948}3+\frac{1.633}6&\le-4.07985,\\
2\cdot0.84853-\frac23\cdot1.17851+\frac{2.44948}{15}&\le1.07469 .
\end{align*}
Hence $Y\le Y_0$, where
\[
Y_0=\frac65\,s^5+\frac{29}{18}\,s^3-4.07985\,s^2+\frac{11}{30}\,s
+1.07469 .
\]
The number $Y_0$ is positive, because
$(29/18)\,s^3\ge(29/18)\cdot6.48\,s^2>4.07985\,s^2$ by
\eqref{eq:sums-halfline} and the other terms of $Y_0$ are positive. Since
$c_{\mathrm S}>0$, we have $c_{\mathrm S}\,Y\le c_{\mathrm S}\,Y_0$, and
since $Y_0>0$ and $c_{\mathrm S}\le0.39894229$, we have
$c_{\mathrm S}\,Y_0\le0.39894229\,Y_0$. Multiplying the coefficients of
$Y_0$ by $0.39894229$ and rounding upwards, we obtain
\[
c_{\mathrm S}\,Y
\le0.478731\,s^5+0.642741\,s^3-1.62762\,s^2+0.14628\,s+0.42874 .
\]

\emph{Step 3: the correction part.}
For every integer $c\ge1$, Lemma~\ref{lem:threefold} gives
$P_3(c)=(c-1)\,P(c)-2\,P_2(c)\le c\,P(c)$, and
$c\,P(c)\le c_{\mathrm S}\sqrt c$ by the upper bound of
\eqref{eq:prelim-stirling}, so that $P_3(c)\le c_{\mathrm S}\sqrt c$. By this
inequality and by \eqref{eq:sums-correction}, which applies because
$s\ge6.48$ by \eqref{eq:sums-halfline},
\[
\mathcal A''\le\frac{0.560302\,c_{\mathrm S}}{s^3}
\sum_{0<c<R}c\sqrt c\,(R-c),
\]
and by \eqref{eq:sums-weighted}, in which $R+1=s^2+1$, the sum on the right
is at most $4/35$ times $(s^2+1)^3\sqrt{s^2+1}$. We have
\[
\sqrt{s^2+1}\le s+\frac1{2s}=\frac{2s^2+1}{2s}
\]
and
\[
(s^2+1)^3\,(2s^2+1)=2s^8+7s^6+9s^4+5s^2+1\le2s^8+7.22\,s^6 .
\]
For the last inequality we use twice that $s^2\ge6.48^2\ge41.99$ by
\eqref{eq:sums-halfline}: this gives $0.22\,s^6\ge9.2378\,s^4$, and then
$0.2378\,s^4\ge9.98\,s^2\ge5s^2+1$. Therefore
\[
\sum_{0<c<R}c\sqrt c\,(R-c)\le\frac4{35}\cdot\frac{2s^8+7.22\,s^6}{2s}
=\frac4{35}\,\bigl(s^7+3.61\,s^5\bigr).
\]
Multiplying by $0.560302\,c_{\mathrm S}/s^3$ and using
$c_{\mathrm S}\le0.39894229$, we obtain
\[
\mathcal A''
\le0.560302\cdot\frac4{35}\cdot0.39894229\,\bigl(s^4+3.61\,s^2\bigr).
\]
Since $0.560302\cdot(4/35)\cdot0.39894229\le0.025547$ and
$0.025547\cdot3.61\le0.09223$, it follows that
\[
\mathcal A''\le0.025547\,s^4+0.09223\,s^2 .
\]

\emph{Step 4: the comparison.}
By the bounds of $\mathcal A'$ and of $c_{\mathrm S}\,Y$ in Step~2, in
which $R=s^2$, and by the bound of $\mathcal A''$ in Step~3, the sum
$\mathcal A(k,R)=\mathcal A'+\mathcal A''$ is at most
\begin{multline*}
0.478731\,s^5+0.642741\,s^3-1.62762\,s^2+0.14628\,s+0.42874\\
-\Bigl(s^4-\frac32\,s^2+\frac12\Bigr)+0.025547\,s^4+0.09223\,s^2 .
\end{multline*}
Collecting terms, we see that this polynomial is the right-hand side of the
claim plus
\[
-0.03539\,s^2+0.14628\,s-0.07126 ,
\]
and this is negative, because by \eqref{eq:sums-halfline}
\[
0.03539\,s\ge0.03539\cdot6.48>0.14628 .\qedhere
\]
\end{proof}

\section{A positive step next to the curve and the change of sign in the first layer}\label{sec:curve}
In this section we prove that every column $k\ge265$ has a positive step
at most one row beyond the curve $k^2=(n-k)^3$: there is a row $j>k$ with
$(j-k-1)^3\le k^2$ and $d_{j,k}>0$
(Proposition~\ref{prop:curveincrease}). For
its proof it remains to insert the bounds of Sections~\ref{sec:pair}
and~\ref{sec:sums} into the criterion of Proposition~\ref{prop:criterion}.
We do this at the positions $(k+r,k)$ with $r\ge41$ and
$(r-1)^3\le k^2\le r^3$, and prove that the step is positive there
(Proposition~\ref{prop:curvestep}).

Combined with Lemma~\ref{lem:signpair} and with
Proposition~\ref{prop:staypositive},
Proposition~\ref{prop:curveincrease} gives the single
change of sign in the first layer of every column $k\ge82$
(Proposition~\ref{prop:onechange}). From this and
Lemma~\ref{lem:smallcolumns} we obtain the column rule in the first layer
of every column $k\ge2$: a positive step in a row of the layer is followed
by a positive step in the next row of the layer
(Proposition~\ref{prop:layerpersistence}).

\subsection{A positive step next to the curve}\label{subsec:curve-step}

As in Sections~\ref{sec:frozen} to~\ref{sec:sums}, the position under study
is written $(k+R-1,k)$, so that $R$ is one more than its deficit, and
$s=\sqrt R$. By Proposition~\ref{prop:criterion}, the step at such a
position with $1\le R\le k-2$ is positive as soon as
\begin{equation}\label{eq:curve-criterion}
k\,a_k\,\Lambda_R(k+1;k+R)<2\,\mathcal P(k,R),
\end{equation}
and Sections~\ref{sec:pair} and~\ref{sec:sums} bound the two sides of this
inequality. We prove \eqref{eq:curve-criterion} at the positions $(k+r,k)$
with $r\ge41$ and $(r-1)^3\le k^2\le r^3$, for which $R=r+1$. The first of
the two inequalities
says that the position one row above, $(k+r-1,k)$, lies behind the curve,
and the second that position $(k+r,k)$ itself lies on the curve, when
$k^2=r^3$, or beyond it, when $k^2<r^3$.

The proof divides
\eqref{eq:curve-criterion} by $k\,a_k\,e^{R}$ and separates its two sides by
one explicit function of $s$. For the right side it inserts the bounds of
Section~\ref{sec:sums} into Proposition~\ref{prop:pairlower} and chooses
the two parameters of that proposition; here the hypothesis on $k$ enters
in the form $k\le s^3\le1.07593\,k$. For the left side it applies
Proposition~\ref{prop:boundaryupper} with a number $K$ that depends on $R$
alone, which the inequality $k^2\le r^3$ permits, and the comparison of the
resulting bound with the function of $s$ comes down to the positivity of
one polynomial of degree $15$ on the half-line $s\ge6.48$.

\begin{proposition}\label{prop:curvestep}
Let $k$ and $r$ be nonnegative
integers with
\[
r\ge41\qquad\text{and}\qquad(r-1)^3\le k^2\le r^3 .
\]
Then $d_{k+r,k}>0$.
\end{proposition}

\begin{proof}
\emph{Step 1: reduction to two inequalities.}
Let $R=r+1$ and $s=\sqrt R$. Then $R\ge42$, row $k+r$ is row
$k+R-1$, and the hypothesis reads
\begin{equation}\label{eq:curve-window}
(R-2)^3\le k^2\le(R-1)^3 .
\end{equation}
The normalised step $\delta^{(k)}_n=d_{n,k}/n!$ has the sign of $d_{n,k}$,
so that we have to show that $\delta^{(k)}_{k+R-1}>0$. Since $R\ge42$,
we have $s\ge6.48$ by \eqref{eq:sums-halfline}.

We first show that
\begin{equation}\label{eq:curve-order}
k\le s^3\le1.07593\,k .
\end{equation}
By \eqref{eq:curve-window}, $k^2\le(R-1)^3\le R^3=(s^3)^2$, and hence
$k\le s^3$. For the other inequality we use that
$R\le\frac{21}{20}\,(R-2)$, which is equivalent to $R\ge42$. Taking third
powers and using \eqref{eq:curve-window} again, we get
\[
(s^3)^2=R^3\le\Bigl(\frac{21}{20}\Bigr)^3(R-2)^3
\le\Bigl(\frac{21}{20}\Bigr)^3k^2\le(1.07593\,k)^2,
\]
where the last step holds because $(21/20)^3\le1.07593^2$. This proves
\eqref{eq:curve-order}.

From \eqref{eq:curve-order} we get $k\ge R+2$: since
$1.07593\,k\ge s^3=R\,s\ge6.48\,R$ and $6.48\ge6\cdot1.07593$, we have
$k\ge6R\ge R+2$. Thus $1\le R\le k-2$, which is the hypothesis of
Proposition~\ref{prop:criterion}, so that it suffices to prove
\eqref{eq:curve-criterion}. Moreover $k\ge1$, so that $a_k>0$.

For a real number $t$ let
\begin{equation}\label{eq:curve-polynomials}
\begin{aligned}
\mathcal W_0(t)&=\tfrac12\,t^2-1.06385\,t+0.7184,\\
\mathcal Q_0(t)&=0.478731\,t^2-0.974453\,t+0.642741 .
\end{aligned}
\end{equation}
We shall prove the two inequalities
\begin{align}
e^{-R}\,\Lambda_R(k+1;k+R)&<
\frac{2\cdot0.38385\,\mathcal W_0(s)\,\bigl(1+\frac{1}{10s}\bigr)}{s^3},
\label{eq:curve-upper}\\
\mathcal P(k,R)&\ge a_k\,e^{R}\cdot0.38385\,\Bigl(1+\frac{1}{10s}\Bigr)
\frac{k}{s^3}\,\mathcal W_0(s).
\label{eq:curve-lower}
\end{align}
They give \eqref{eq:curve-criterion}: multiplying \eqref{eq:curve-upper}
by the positive number
$k\,a_k\,e^{R}$ and using \eqref{eq:curve-lower}, we get
\begin{align*}
k\,a_k\,\Lambda_R(k+1;k+R)
&<2\,a_k\,e^{R}\cdot0.38385\,\Bigl(1+\frac{1}{10s}\Bigr)\frac{k}{s^3}\,
\mathcal W_0(s)\\
&\le2\,\mathcal P(k,R).
\end{align*}

\emph{Step 2: the lower bound of the pair sum.}
We prove \eqref{eq:curve-lower}. Since $R\ge42$ and
$s^3\le1.07593\,k$, Lemmas~\ref{lem:weightlower}, \ref{lem:logupper}
and~\ref{lem:excessupper} give
\[
\mathcal W(R)\ge\mathcal W_0(s),\qquad
\mathcal G(k,R)\le2s\,\mathcal Q_0(s),\qquad
\mathcal A(k,R)\le s^3\,\mathcal Q_0(s).
\]
Both $\mathcal W_0(s)$ and $\mathcal Q_0(s)$ are positive, since
$s/2>1.06385$ and $0.478731\,s>0.974453$. Let
\[
\theta=\frac{s^2}{2k},\qquad x=\frac{s^3}{k},\qquad
\mu=\frac{2s\,\mathcal Q_0(s)\,\log(1+\theta)}{\mathcal W_0(s)},\qquad
h=0.9575-\frac{1}{10s},
\]
and, as in Proposition~\ref{prop:pairlower},
\[
\alpha=\log(1+\theta)-\frac{\theta}{1+\theta},\qquad
\beta=\frac{1}{k\,(1+\theta)}.
\]
Then $\theta>0$ and $\mu>0$, we have $2s\,\theta=x$, and
\eqref{eq:curve-order} says that $1\le x\le1.07593$. Since $k\ge1$ and
$R\le k$, Proposition~\ref{prop:pairlower} applies to these values of
$\theta$ and $\mu$, and it gives
\begin{equation}\label{eq:curve-pairbound}
a_k\,e^{R}\,e^{-\mu}\,\bigl((1+\mu)\,\mathcal W(R)
-\alpha\,\mathcal G(k,R)-\beta\,\mathcal A(k,R)\bigr)\le\mathcal P(k,R).
\end{equation}

We first show that
\begin{equation}\label{eq:curve-bracket}
(1+\mu)\,\mathcal W(R)-\alpha\,\mathcal G(k,R)-\beta\,\mathcal A(k,R)
\ge\mathcal W_0(s) .
\end{equation}
The three coefficients $1+\mu$, $\alpha$ and $\beta$ are nonnegative; for
$\alpha$ this is the inequality
$\log u\ge1-1/u$ at $u=1+\theta$. By the bounds of the three sums, the
left side of
\eqref{eq:curve-bracket} is therefore at least
\[
(1+\mu)\,\mathcal W_0(s)
-\Bigl(\log(1+\theta)-\frac{\theta}{1+\theta}\Bigr)\,2s\,\mathcal Q_0(s)
-\frac{s^3\mathcal Q_0(s)}{k\,(1+\theta)} .
\]
Here
$2s\,\mathcal Q_0(s)\,\theta/(1+\theta)=s^3\mathcal Q_0(s)/(k\,(1+\theta))$,
because $2s\,\theta=s^3/k$, so that these two terms cancel; and
$\mu\,\mathcal W_0(s)=2s\,\mathcal Q_0(s)\log(1+\theta)$ by the choice of
$\mu$,
so that these two terms cancel as well. What remains is $\mathcal W_0(s)$.

Next we show that $\mu\le x\,h$. By the Taylor bound
$\log(1+u)\le u-u^2/2+u^3/3$ for $u\ge0$, which is recalled in
Section~\ref{sec:prelim}, and by $\theta=x/(2s)$,
\begin{align*}
\mu\,\mathcal W_0(s)=2s\,\mathcal Q_0(s)\log(1+\theta)
&\le2s\,\mathcal Q_0(s)
\Bigl(\frac{x}{2s}-\frac{x^2}{8s^2}+\frac{x^3}{24s^3}\Bigr)\\
&=\frac{x\,\mathcal Q_0(s)}{s^2}
\Bigl(s^2-\frac{x\,s}{4}+\frac{x^2}{12}\Bigr).
\end{align*}
The last bracket is largest at the smallest value $x=1$, since
\[
\Bigl(s^2-\frac s4+\frac{1}{12}\Bigr)
-\Bigl(s^2-\frac{x\,s}{4}+\frac{x^2}{12}\Bigr)
=(x-1)\Bigl(\frac s4-\frac{x+1}{12}\Bigr)\ge0 ;
\]
here $x\ge1$ and $s/4\ge1.62>(x+1)/12$. Hence
\begin{equation}\label{eq:curve-exponent}
\mu\,\mathcal W_0(s)\,s^2\le x\,\mathcal Q_0(s)\,
\Bigl(s^2-\frac s4+\frac{1}{12}\Bigr).
\end{equation}
We compare the right side with $x\,h\,\mathcal W_0(s)\,s^2$. Since
$h\,s^2=0.9575\,s^2-s/10$, multiplying out the polynomials
\eqref{eq:curve-polynomials} gives
\begin{multline*}
\mathcal W_0(s)\,\Bigl(0.9575\,s^2-\frac{s}{10}\Bigr)
-\mathcal Q_0(s)\,\Bigl(s^2-\frac s4+\frac{1}{12}\Bigr)\\
=\frac{19}{10^{6}}\,s^4
+s^2\,\Bigl(\frac{40799}{1600000}\,s-\frac{263991}{2000000}\Bigr)
+\Bigl(\frac{510149}{3000000}\,s-\frac{214247}{4000000}\Bigr).
\end{multline*}
For $s\ge6.48$ the three terms on the right are positive, since the two
brackets increase with $s$ and are positive at $s=6.48$. Together with
\eqref{eq:curve-exponent} this gives
$\mu\,\mathcal W_0(s)\,s^2\le x\,h\,\mathcal W_0(s)\,s^2$, and hence
$\mu\le x\,h$.

For decimal values of the exponential function
we use that, for every real number $y$ with $0\le y\le1$,
\begin{equation}\label{eq:curve-taylor}
\sum_{j=0}^{7}\frac{y^j}{j!}\le e^{y}\le
\sum_{j=0}^{7}\frac{y^j}{j!}+\frac{9\,y^8}{8\cdot8!}.
\end{equation}
The lower bound is clear. For the upper bound, $(8+i)!\ge8!\cdot9^{i}$ for
every $i\ge0$, so that the terms of the exponential series from index $8$ on
add up to at most $(y^8/8!)\sum_{i\ge0}9^{-i}=9y^8/(8\cdot8!)$.

We now show that
\begin{equation}\label{eq:curve-exponential}
e^{-xh}\ge0.38385\,\Bigl(1+\frac{1}{10s}\Bigr)\,\frac1x .
\end{equation}
Since $-x\,h=-0.9575+1/(10s)-(x-1)\,h$, the left side is the product
of the three positive numbers $e^{-0.9575}$, $e^{1/(10s)}$ and
$e^{-(x-1)h}$, and we bound each of them from below. By
\eqref{eq:curve-taylor}, applied to the number $0.9575$, we have
$e^{0.9575}\le2.60518$, and $0.38385\cdot2.60518\le1$; hence
$e^{-0.9575}\ge0.38385$. Next, $e^{1/(10s)}\ge1+1/(10s)$. For the third
factor, the Taylor bound $\log(1+u)\ge u-u^2/2$ for $u\ge0$, also recalled
in Section~\ref{sec:prelim}, gives at $u=x-1\le0.07593$
\[
\log x\ge(x-1)\Bigl(1-\frac{x-1}{2}\Bigr)\ge0.962\,(x-1)\ge(x-1)\,h,
\]
because $h\le0.9575$; hence $e^{-(x-1)h}\ge1/x$. The product of the
three lower bounds is the right side of \eqref{eq:curve-exponential}.

Finally, since $\mu\le x\,h$, we have $e^{-\mu}\ge e^{-xh}$. With
\eqref{eq:curve-exponential} and $1/x=k/s^3$, and since $a_k\,e^{R}>0$ and
$\mathcal W_0(s)>0$,
\[
a_k\,e^{R}\cdot0.38385\,\Bigl(1+\frac{1}{10s}\Bigr)\,\frac{k}{s^3}\,
\mathcal W_0(s)
\le a_k\,e^{R}\,e^{-\mu}\,\mathcal W_0(s) .
\]
By \eqref{eq:curve-bracket}, the right side is at most the left side of
\eqref{eq:curve-pairbound}, and hence at most $\mathcal P(k,R)$. This
proves \eqref{eq:curve-lower}.

\emph{Step 3: the upper bound of the chain sum.}
It remains to prove \eqref{eq:curve-upper}.
Proposition~\ref{prop:boundaryupper} bounds the chain sum
$\Lambda_R(k+1;k+R)$ for every real number $K\ge\max(k,R)$, and its bound
contains a second real parameter, which may take any value not smaller than
$-1$. We take $K=(R-1)(2R-1)/(2s)$, and for the second parameter the value
$0.35$. The number $K$ depends on $R$ alone and is a rational function of
$s$; in this way the bound of the chain sum no longer depends on $k$.
We have $R\le K$, because
\[
2s\,(K-R)=(s^2-1)(2s^2-1)-2s^3=s^2\,(2s^2-2s-3)+1>0 ,
\]
the bracket being positive for $s\ge6.48$. Also $k\le K$: we have
$K=(R-1)\bigl(s-1/(2s)\bigr)$ and
$\bigl(s-1/(2s)\bigr)^2=R-1+1/(4R)\ge R-1$, so that
$K^2\ge(R-1)^3\ge k^2$ by \eqref{eq:curve-window}.
Proposition~\ref{prop:boundaryupper} therefore applies to the integer
$R\ge1$, to the real numbers $k$ and $K$, which satisfy $0\le k\le K$ and
$R\le K$, and to the value $0.35$ of its second parameter. With the
abbreviations
\[
\chi_1=\frac{R-2}{2\,(K+R-1)},\qquad
\chi_2=\frac{1}{K+1}\cdot\frac{R-3}{3}\cdot\frac{\frac45R-1}{K+\frac35R},
\]
which are positive numbers, it gives
\begin{equation}\label{eq:curve-chain}
e^{-R}\,\Lambda_R(k+1;k+R)\le P(R)-e^{-0.35}\,
\bigl(1.35\,\chi_1\,P_2(R)-\chi_2\,P_3(R)\bigr).
\end{equation}

Let
\[
\phi=\frac43+\frac{1}{15R},\qquad\bar v=\frac{c_{\mathrm S}}{s}.
\]
By the upper bound of \eqref{eq:prelim-stirling}, $P(R)\le\bar v$, and
$0<c_{\mathrm S}\le0.39894229$ by the enclosure of $c_{\mathrm S}$ in
Section~\ref{sec:prelim}. By Corollary~\ref{cor:twofoldlower} and
Lemma~\ref{lem:threefold},
\[
P_2(R)\ge\frac12-\phi\,P(R),\qquad
P_3(R)=(R-1)\,P(R)-2\,P_2(R),
\]
and hence
\[
P_3(R)\le(R-1+2\phi)\,P(R)-1 .
\]
Finally, \eqref{eq:curve-taylor}, applied to the number $0.35$, shows that
$e^{0.35}$ lies between $1.4190675$ and $1.4190676$. Since
\[
0.704688\cdot1.4190676\le1\le0.704689\cdot1.4190675,
\]
it follows that
\[
0.704688\le e^{-0.35}\le0.704689 .
\]

Since $e^{-0.35}$, $\chi_1$ and $\chi_2$ are positive, the bounds for
$P_2(R)$ and $P_3(R)$ can be inserted into \eqref{eq:curve-chain}: its
right side is at most
\[
P(R)-e^{-0.35}\,\Bigl(1.35\,\chi_1\bigl(\tfrac12-\phi\,P(R)\bigr)
-\chi_2\bigl((R-1+2\phi)\,P(R)-1\bigr)\Bigr).
\]
In this expression $P(R)$ has the coefficient
\[
1+e^{-0.35}\,\bigl(1.35\,\chi_1\phi+\chi_2\,(R-1+2\phi)\bigr),
\]
which is positive, so that
the expression does not decrease when $P(R)$ is replaced by the larger
number $\bar v$:
\begin{equation}\label{eq:curve-affine}
\begin{split}
e^{-R}\,\Lambda_R(k+1;k+R)\le\bar v-e^{-0.35}\,
\Bigl(&1.35\,\chi_1\bigl(\tfrac12-\phi\,\bar v\bigr)\\
&-\chi_2\bigl((R-1+2\phi)\,\bar v-1\bigr)\Bigr).
\end{split}
\end{equation}

Next we bound the right side of \eqref{eq:curve-affine} by a multiple of
$\bar v$. Let
\[
J_1=1.253314\,s-\phi,\qquad J_2=s^2-2.5066\,s+1.67 .
\]
Both numbers are positive: we have $\phi\le4/3+1/600=1.335$, because
$15R\ge600$, so that $J_1\ge1.253314\cdot6.48-1.335>0$, and
$J_2=s\,(s-2.5066)+1.67$. We compare the two inner brackets of
\eqref{eq:curve-affine} with $\bar v\,J_1$ and $\bar v\,J_2$, using that
$s\,\bar v=c_{\mathrm S}\le0.39894229$. First,
\[
\bar v\,J_1=1.253314\,c_{\mathrm S}-\phi\,\bar v\le\tfrac12-\phi\,\bar v,
\]
because $1.253314\cdot0.39894229\le1/2$. Secondly,
\[
\bar v\,J_2-\bigl((R-1+2\phi)\,\bar v-1\bigr)
=1-2.5066\,c_{\mathrm S}+(2.67-2\phi)\,\bar v\ge0,
\]
because $2.5066\cdot0.39894229\le1$ and $2\phi\le2.67$.
Inserting the two comparisons into \eqref{eq:curve-affine}, we see that
its right side is at most
\[
\bar v-e^{-0.35}\,\bigl(1.35\,\chi_1J_1-\chi_2J_2\bigr)\,\bar v .
\]
Here the term
$\chi_1J_1\bar v$, which is subtracted, and the term $\chi_2J_2\bar v$,
which is added, are both positive. We use the lower bound of $e^{-0.35}$
in the
first and its upper bound in the second, and arrive at
\begin{equation}\label{eq:curve-omega}
e^{-R}\,\Lambda_R(k+1;k+R)\le\bar v\,\Omega_1,
\end{equation}
where
\[
\Omega_1=1-1.35\cdot0.704688\,\chi_1J_1+0.704689\,\chi_2J_2 .
\]

\emph{Step 4: the comparison of the two bounds.}
Let
\[
\Omega_2=\frac{1.9243\,\mathcal W_0(s)\,(10s+1)}{10\,s^3}.
\]
We claim
that $\Omega_1<\Omega_2$. To see this we write $\Omega_1$ as a rational
function of $s$. By the choice of $K$, $2s\,K=(s^2-1)(2s^2-1)$. Let
\begin{align*}
M_1&=(s^2-1)(2s^2+2s-1),\\
M_2&=(s^2-1)(2s^2-1)+\tfrac65\,s^3,\\
M_3&=(s^2-1)(2s^2-1)+2s .
\end{align*}
Then $2s\,(K+R-1)=M_1$, $2s\,(K+\frac35R)=M_2$ and
$2s\,(K+1)=M_3$, and the three numbers are positive. Consequently
\begin{gather*}
\chi_1=\frac{s\,(s^2-2)}{M_1},\qquad
\chi_2=\frac{4s^2\,(s^2-3)\bigl(\frac45s^2-1\bigr)}{3\,M_2M_3},\\
J_1=\frac{18.79971\,s^3-20\,s^2-1}{15\,s^2},
\end{gather*}
the last because $15s^2\phi=20s^2+1$ and $15\cdot1.253314=18.79971$. We
multiply $\Omega_2-\Omega_1$ by the positive number
$30\,s^3M_1M_2M_3$. The result is a polynomial in $s$:
\[
30\,s^3\,M_1M_2M_3\,(\Omega_2-\Omega_1)=\mathcal X(s),
\]
where
\begin{align*}
\mathcal X(s)={}&5.7729\,\mathcal W_0(s)\,(10s+1)\,M_1M_2M_3
-30\,s^3\,M_1M_2M_3\\
&+1.35\cdot0.704688\cdot2s^2\,(s^2-2)\,
\bigl(18.79971\,s^3-20\,s^2-1\bigr)\,M_2M_3\\
&-0.704689\cdot40\,s^5\,(s^2-3)\bigl(\tfrac45s^2-1\bigr)
\bigl(s^2-2.5066\,s+1.67\bigr)\,M_1 .
\end{align*}
The four terms come from $\Omega_2$, from the term $1$ of $\Omega_1$, and
from its terms with $\chi_1J_1$ and with $\chi_2J_2$; the factor $5.7729$ is
$3\cdot1.9243$. Each of them is a polynomial of degree $15$ in $s$ whose
coefficients are finite decimal fractions. Multiplying out the four
products and collecting the powers
of $s$, we obtain
\begin{align*}
\mathcal X(s)={}&88.893548437184\,s^{15}-481.1809899040896\,s^{14}\\
&-940.74918031232\,s^{13}+2462.6087604838976\,s^{12}\\
&+2054.9679357758624\,s^{11}-5874.3007648560576\,s^{10}\\
&-1064.5268494392928\,s^{9}+7091.8315979794336\,s^{8}\\
&-1196.745291572152\,s^{7}-4127.091835775184\,s^{6}\\
&+1376.539082916408\,s^{5}+1128.517422882\,s^{4}\\
&-439.556846403\,s^{3}-108.63749913\,s^{2}\\
&+35.331013935\,s+4.14725136 .
\end{align*}
These decimal coefficients are exact; none of them is rounded.

The coefficients are of both signs, so that the sign of $\mathcal X(s)$
cannot be read off from them. We separate the four terms of
highest degree from the others. For a real number $t$ let
\begin{align*}
\mathcal X_1(t)={}&88.893548437184\,t^{3}-481.1809899040896\,t^{2}\\
&-940.74918031232\,t+2462.6087604838976 ,
\end{align*}
and let $\mathcal X_0(s)$ be the sum of the twelve terms of $\mathcal X(s)$
of degree at most $11$, so that
\[
\mathcal X(s)=s^{12}\,\mathcal X_1(s)+\mathcal X_0(s).
\]
Four terms are taken together because the sum of the three terms of highest
degree is negative at $s=6.48$. It suffices to show that
$\mathcal X_1(s)$ and $\mathcal X_0(s)$ are positive.

The polynomial $\mathcal X_1$ increases on the half-line $t\ge6.48$. Indeed,
for $t\ge0$ its derivative satisfies
\[
\mathcal X_1'(t)\ge t\,(3\cdot88.89\,t-2\cdot481.19)-940.75 ,
\]
because $88.89$ is smaller than the coefficient of $t^3$ in
$\mathcal X_1(t)$, and $481.19$ and $940.75$ are larger than the absolute
values of its coefficients of $t^2$ and of $t$. For $t\ge6.48$ the bracket
is at least $765$, because $3\cdot88.89\cdot6.48>1728$ and
$2\cdot481.19<963$, and therefore
\[
\mathcal X_1'(t)\ge765\cdot6.48-940.75>0 .
\]
Moreover $\mathcal X_1(6.48)>0$. We have $6.48^2=41.9904$ and
$6.48^3=272.097792$. Replacing the first of these powers by $42$ and the
second by $272.09$, and rounding the coefficients in the same way, the
constant term to $2462.6$, we only make the value smaller, so that
\begin{align*}
\mathcal X_1(6.48)
&\ge88.89\cdot272.09-481.19\cdot42-940.75\cdot6.48+2462.6\\
&\ge24186.08-20209.98-6096.06+2462.6=342.64 .
\end{align*}
Hence $\mathcal X_1(s)\ge\mathcal X_1(6.48)>0$, because $s\ge6.48$.

We turn to $\mathcal X_0(s)$. Its two terms of highest degree
are $2054.96\ldots\,s^{11}$ and $-5874.30\ldots\,s^{10}$, and since
$s\ge6.48$ their sum is at least
\[
(2054\cdot6.48-5875)\,s^{10}\ge7000\,s^{10} .
\]
Of the coefficients of $s^{9},s^{8},\dots,s^{0}$ none is smaller than
$-4128$. We add the terms of $\mathcal X_0(s)$ one at a time, in the order
of decreasing degree, and show for $j=10,9,\dots,0$ that the sum of the
terms of degree at least $j$ is at least $7000\,s^{j}$.
For $j=10$ this has just been shown. If it holds for $j+1$, where
$0\le j\le9$, then the sum of the terms of degree at least $j$ is at least
\[
7000\,s^{j+1}-4128\,s^{j}=(7000\,s-4128)\,s^{j}\ge7000\,s^{j} ,
\]
because $7000\cdot6.48-4128>7000$. For $j=0$ this says that
$\mathcal X_0(s)\ge7000$.

Hence $\mathcal X(s)>0$, and therefore
$\Omega_1<\Omega_2$.

Since $\bar v$ is positive, \eqref{eq:curve-omega} and the
inequality $\Omega_1<\Omega_2$ show that the left side of
\eqref{eq:curve-upper} is smaller than $\bar v\,\Omega_2$. By the
definitions of $\bar v$ and
of $\Omega_2$,
\[
\bar v\,\Omega_2=1.9243\,c_{\mathrm S}\cdot
\frac{\mathcal W_0(s)\,\bigl(1+\frac{1}{10s}\bigr)}{s^3},
\]
and the fraction on the right is positive, because
$\mathcal W_0(s)>0$. Moreover
\[
1.9243\,c_{\mathrm S}\le1.9243\cdot0.39894229\le2\cdot0.38385 .
\]
This proves \eqref{eq:curve-upper}.
\end{proof}

We apply Proposition~\ref{prop:curvestep} in the columns $k\ge265$, for
which Lemma~\ref{lem:signpair}
supplies no positive row, in the first row of the column beyond the
curve. This gives the following bound for the first
positive row of the column. The bound is used in
Section~\ref{sec:staircase}, together
with Corollary~\ref{cor:layerpersistenceall} and
Proposition~\ref{prop:diagonalstep}, for the diagonal rule in the first
layer of such a column.

\begin{proposition}\label{prop:curveincrease}
For every integer $k\ge265$ there is
an integer $j$ with
\[
k<j,\qquad(j-k-1)^3\le k^2,\qquad d_{j,k}>0 .
\]
\end{proposition}

\begin{proof}
Let $r$ be the least nonnegative integer with $k^2<r^3$. Then
$(r-1)^3\le k^2<r^3$, and $r\ge42$: for $r\le41$ we would have
$r^3\le41^3<265^2\le k^2$. Proposition~\ref{prop:curvestep}
therefore gives $d_{k+r,k}>0$, and the integer $j=k+r$ has the required
properties, since $(j-k-1)^3=(r-1)^3\le k^2$.
\end{proof}

\subsection{The change of sign in the first layer}\label{subsec:curve-layer}

Proposition~\ref{prop:staypositive} shows that a column
$k\ge82$ which is not positive in the rows $k$ to $q-1$ and is positive in
a row $q\ge k+1$ stays positive from row $q$ to the end of the
first layer, under the assumption that $q$ lies in the first third of the
layer, $3\,(q-k)\le k$, and the argument of Section~\ref{sec:beyond} for
the rows beyond the first layer needs the bound $100\,(q-k)\le27\,k$. The
following corollary supplies a positive row within this distance of row
$k$.

\begin{corollary}\label{cor:earlypositive}
For every integer $k\ge82$ there is
an integer $r$ with
\[
r\ge1,\qquad100\,r\le27\,k,\qquad d_{k+r,k}>0 .
\]
\end{corollary}

\begin{proof}
First let $82\le k\le264$, and let $\zeta(k)$ be the number that
Definition~\ref{def:zeta} lists for column $k$. By
Lemma~\ref{lem:signpair}, which holds for $9\le k\le264$, column $k$ is
positive in row $k+\zeta(k)+1$. We take $r=\zeta(k)+1$, so that $r\ge1$,
and have to verify that $100\,(\zeta(k)+1)\le27\,k$. This is read from the
table of Definition~\ref{def:zeta}: we have $\zeta(k)\le20$ for $k\le101$,
$\zeta(k)\le24$ for $k\le141$, and $\zeta(k)\le35$ for every $k\le264$.
Hence $100\,(\zeta(k)+1)$ is at most $2100\le27\cdot82$ for
$82\le k\le101$, at most $2500\le27\cdot102$ for $102\le k\le141$, and at
most $3600\le27\cdot142$ for $142\le k\le264$.

Now let $k\ge265$. By Proposition~\ref{prop:curveincrease} there is an
integer $j>k$ with $(j-k-1)^3\le k^2$ and $d_{j,k}>0$. Let $t=j-k-1$ and
$r=t+1$. Then $t\ge0$ and $t^3\le k^2$, we have $r\ge1$, and
$d_{k+r,k}=d_{j,k}>0$. It remains to show that $100\,(t+1)\le27\,k$. Suppose
that $100\,(t+1)>27\,k$. Then $400\,t>108\,k-400\ge100\,k$, because
$8k\ge400$, and so $k<4t$. Hence $k^3<64\,t^3\le64\,k^2$, so that $k<64$,
which contradicts $k\ge265$.
\end{proof}

The proof of the next proposition applies
Proposition~\ref{prop:staypositive} to the first positive row $q$ of a
column $k\ge82$ instead of the row of the corollary.
Proposition~\ref{prop:onechange} describes the signs of the column from
row $k$ to the end of the first layer: read down the column, the step
changes its sign exactly once. The signs in the rows
from $k$ to $2k-1$ are also
the start of the induction on the row in Section~\ref{sec:beyond}.

\begin{proposition}\label{prop:onechange}
For every integer $k\ge82$ there is
an integer $q$ with the following three properties.
\begin{enumerate}
\item One has $k+1\le q$ and $100\,(q-k)\le27\,k$.
\item One has $d_{j,k}\le0$ for every integer $j$ with $k\le j<q$.
\item One has $d_{n,k}>0$ for every integer $n$ with $q\le n\le2k-1$.
\end{enumerate}
\end{proposition}

\begin{proof}
By Corollary~\ref{cor:earlypositive} there is an integer $r\ge1$ with
$100\,r\le27\,k$ and $d_{k+r,k}>0$. Let $q$ be the first positive row of
column $k$, that is, the least row $j\ge k+1$ with $d_{j,k}>0$. Then
$k+1\le q\le k+r$, hence $100\,(q-k)\le100\,r\le27\,k$, and $d_{q,k}>0$.

For the second property let $k\le j<q$. If $j=k$, then
$d_{k,k}/k!=\delta^{(k)}_k=-a_k/k<0$ by Lemma~\ref{lem:initialstep}, and if
$k+1\le j<q$, then $d_{j,k}\le0$ by the choice of $q$.

For the third property let $r_q=q-k$. Then $r_q\ge1$, and $3\,r_q\le k$,
because $100\,r_q\le27\,k$. We have just seen that $d_{j,k}\le0$ for every
$j$ with $k\le j<k+r_q$, and that $d_{k+r_q,k}>0$. These are the hypotheses
of Proposition~\ref{prop:staypositive} for column $k\ge82$, and its
conclusion is that $d_{n,k}>0$ for every $n$ with $k+r_q\le n\le2k-1$.
\end{proof}

We can now prove the column rule in the case in which its conclusion lies
in the first layer, for all columns $k\ge2$. The case in which the
conclusion lies beyond the first layer is
treated in Section~\ref{sec:beyond} (Corollary~\ref{cor:columnbeyond}).

\begin{proposition}\label{prop:layerpersistence}
Let $k\ge2$ and $n$ be integers with
$k+1\le n\le2k-2$. If $d_{n,k}>0$, then $d_{n+1,k}>0$.
\end{proposition}

\begin{proof}
For $2\le k\le81$ the claim is part 2 of
Lemma~\ref{lem:smallcolumns}.

Now let $k\ge82$, and let $q$ be a row with the properties of
Proposition~\ref{prop:onechange}. Since $d_{n,k}>0$ and $n\ge k$, the
second property of $q$ gives $q\le n$. Then $q\le n+1\le2k-1$, and the
third property gives $d_{n+1,k}>0$.
\end{proof}

Applied repeatedly, Proposition~\ref{prop:layerpersistence} carries a
positive step to every later row of the first layer.
Section~\ref{sec:staircase} uses this for the diagonal rule, where a
positive row of the first layer is known and a step of the first layer
that is not positive must then lie in an earlier row.

\begin{corollary}\label{cor:layerpersistenceall}
Let $k\ge2$, $m$ and $n$ be integers with
$k+1\le m\le n\le2k-1$. If $d_{m,k}>0$, then $d_{n,k}>0$.
\end{corollary}

\begin{proof}
This follows by induction on $n$ from
Proposition~\ref{prop:layerpersistence}, which applies to every row $n$
with $m\le n\le2k-2$.
\end{proof}

\section{The prime kernel at the critical tilt}\label{sec:kernel}
Beyond the first layer the column equation of
Proposition~\ref{prop:columnequation} is no longer available. What remains
is the
inequality of Proposition~\ref{prop:columninequality}, in which the
feedback of the earlier rows is formed with the prime kernel
$\kappa^{(k)}$ of threshold $k$. From this inequality
Section~\ref{sec:beyond} derives that the step of
a column $k\ge82$ is positive in every row $n\ge2k$
(Proposition~\ref{prop:largecolumns}), by an argument at the critical tilt
$x_k$ of Section~\ref{sec:propagation}. Here we prove what that argument
uses about the tilted kernel $\bar\kappa_m=\kappa^{(k)}_m\,x_k^{\,m}$: its
monotonicity and decay, its size for $k\ge82$, and the order of the weights
that it gives to the rows before and after the first positive row of the
column (Lemmas~\ref{lem:orderbefore} and~\ref{lem:orderafter}).

\subsection{Monotonicity and harmonic decay}\label{subsec:kernel-monotone}

Recall from Section~\ref{sec:propagation} that the critical tilt of a column
$k\ge1$ is $x_k=(k/(k+1))^{k-1}$, that $0<x_k\le1$, and that at this tilt
the tilted Cayley masses $\omega_n=c_n\,x_k^{\,n}$ do not increase up to
index $k$ (Lemma~\ref{lem:massesdecrease}). We first show that at the same
tilt the tilted free prime kernel $\varepsilon_m\,x_k^{\,m}$ does not
increase up to index $k$ either; here $\varepsilon_m=(m-1)^{m-1}/m!$, with
$\varepsilon_1=1$, is the free prime kernel of Section~\ref{sec:series}.
The proof derives this monotonicity from that of the tilted Cayley masses.

\begin{lemma}\label{lem:freekerneldecrease}
Let $k$ and $m$ be integers with $1\le m<k$. Then the free prime kernel
$\varepsilon$ and the critical tilt $x_k$ of column $k$ satisfy
\[
\varepsilon_{m+1}\,x_k^{\,m+1}\le\varepsilon_m\,x_k^{\,m} .
\]
\end{lemma}

\begin{proof}
For an integer $i\ge1$ let $\varpi_i=(1+1/i)^i$, and let $\varpi_0=1$. Then
\begin{equation}\label{eq:kernel-shifted}
a_{i+1}=\varpi_i\,\varepsilon_{i+1}\qquad\text{for every integer }i\ge0 ,
\end{equation}
as one checks directly from the definitions. The factor $\varpi_i$ does not
decrease with $i$: we have $\varpi_0=1\le2=\varpi_1$, and
$\varpi_{i-1}\le\varpi_i$ for $i\ge2$ by the monotonicity of $(1+1/i)^i$
that is recalled in Section~\ref{sec:prelim}.

Write $x=x_k$. Since $a_i=c_{i-1}$, we have $a_i\,x^i=x\,\omega_{i-1}$ for
every integer $i\ge1$. Since $m\le k$, Lemma~\ref{lem:massesdecrease} gives
$\omega_m\le\omega_{m-1}$, and hence $a_{m+1}\,x^{m+1}\le a_m\,x^m$. With
\eqref{eq:kernel-shifted} for indices $m$ and $m-1$ it follows that
\[
\varpi_m\,\varepsilon_{m+1}\,x^{m+1}=a_{m+1}\,x^{m+1}\le a_m\,x^m
=\varpi_{m-1}\,\varepsilon_m\,x^m\le\varpi_m\,\varepsilon_m\,x^m ,
\]
where the last step uses $\varpi_{m-1}\le\varpi_m$. Dividing by $\varpi_m$
gives the claim.
\end{proof}

We now consider an arbitrary threshold $y\ge2$ and an arbitrary tilt
$x\ge0$. Recall from Section~\ref{sec:series} that
$(\kappa^{(y)}_m\,x^m)_{m\ge0}$ is the tilted prime kernel of threshold
$y$, and that in the free range $1\le m\le y$ it agrees with the tilted free
prime kernel, $\kappa^{(y)}_m\,x^m=\varepsilon_m\,x^m$
(Lemma~\ref{lem:kernelfree}). Beyond the free range the tilted prime kernel
is given by the recursion \eqref{eq:series-tilted-recursion}, which holds
for every integer $N\ge y$ and in which the tilted weights
$(y-\gamma)\,a_\gamma\,x^{\gamma}$ are nonnegative and do not depend
on $N$.

The next two lemmas concern a tilt at which the tilted prime kernel does
not increase on the free range. They show that the numbers
$m\,\kappa^{(y)}_m\,x^m$ then do not increase beyond the free range, so
that the tilted prime kernel decays at least like $1/m$, which is its
\emph{harmonic decay}, and that the tilted prime kernel itself does not
increase from index $1$ on. The two lemmas are proved by one induction.

\begin{lemma}\label{lem:harmonic}
Let $y\ge2$ be an integer and $x\ge0$ a real number, and suppose that the
free prime kernel $\varepsilon$ satisfies
\[
\varepsilon_{m+1}\,x^{m+1}\le\varepsilon_m\,x^m
\qquad\text{for every integer $m$ with }1\le m<y .
\]
Then the prime kernel $\kappa^{(y)}$ of threshold $y$ satisfies
\[
(m+1)\,\kappa^{(y)}_{m+1}\,x^{m+1}\le m\,\kappa^{(y)}_m\,x^m
\qquad\text{for every integer }m\ge y .
\]
\end{lemma}

\begin{proof}
We prove by induction on $m\ge1$ that
\begin{equation}\label{eq:kernel-induction}
\kappa^{(y)}_{m+1}\,x^{m+1}\le\kappa^{(y)}_m\,x^m ,
\end{equation}
and that the inequality of the lemma holds at $m$ if $m\ge y$. Let $m\ge1$,
and suppose that \eqref{eq:kernel-induction} holds with $m$ replaced by any
index $i$ with $1\le i<m$.

If $m+1\le y$, the indices $m$ and $m+1$ lie in the free range, where
$\kappa^{(y)}_m=\varepsilon_m$ and $\kappa^{(y)}_{m+1}=\varepsilon_{m+1}$ by
Lemma~\ref{lem:kernelfree}, and \eqref{eq:kernel-induction} is the
hypothesis of the lemma.

Let $m\ge y$. By \eqref{eq:series-tilted-recursion} for $N=m+1$ and for
$N=m$,
\begin{align*}
(m+1)\,\kappa^{(y)}_{m+1}\,x^{m+1}
&=\sum_{\gamma=1}^{y-1}(y-\gamma)\,a_\gamma\,x^{\gamma}
\cdot\kappa^{(y)}_{m+1-\gamma}\,x^{m+1-\gamma},\\
m\,\kappa^{(y)}_m\,x^m
&=\sum_{\gamma=1}^{y-1}(y-\gamma)\,a_\gamma\,x^{\gamma}
\cdot\kappa^{(y)}_{m-\gamma}\,x^{m-\gamma} .
\end{align*}
For $1\le\gamma\le y-1$ the index $i=m-\gamma$ satisfies $1\le i<m$,
because $m\ge y$, so that the induction hypothesis gives
\[
\kappa^{(y)}_{m+1-\gamma}\,x^{m+1-\gamma}
\le\kappa^{(y)}_{m-\gamma}\,x^{m-\gamma} .
\]
Multiplying by the weight $(y-\gamma)\,a_\gamma\,x^{\gamma}\ge0$ and
summing over $\gamma$, we obtain the inequality of the lemma at $m$.
Moreover $\kappa^{(y)}_{m+1}\,x^{m+1}\ge0$ by Lemma~\ref{lem:kernelsign},
and hence
\[
m\,\kappa^{(y)}_{m+1}\,x^{m+1}\le(m+1)\,\kappa^{(y)}_{m+1}\,x^{m+1}
\le m\,\kappa^{(y)}_m\,x^m .
\]
Dividing by $m$, we obtain \eqref{eq:kernel-induction} at $m$. This
completes the induction.
\end{proof}

\begin{lemma}\label{lem:kerneldecrease}
Let $y\ge2$ be an integer and $x\ge0$ a real number, and suppose that the
free prime kernel $\varepsilon$ satisfies
\[
\varepsilon_{m+1}\,x^{m+1}\le\varepsilon_m\,x^m
\qquad\text{for every integer $m$ with }1\le m<y .
\]
Then the prime kernel $\kappa^{(y)}$ of threshold $y$ satisfies
\[
\kappa^{(y)}_{m+1}\,x^{m+1}\le\kappa^{(y)}_m\,x^m
\qquad\text{for every integer }m\ge1 .
\]
\end{lemma}

\begin{proof}
The hypotheses are those of Lemma~\ref{lem:harmonic}, and the claim is the
inequality \eqref{eq:kernel-induction}, which the proof of that lemma
establishes for every integer $m\ge1$.
\end{proof}

We return to the critical tilt. Let $k\ge2$ be an integer. By
Lemma~\ref{lem:freekerneldecrease}, the hypothesis of
Lemmas~\ref{lem:harmonic} and~\ref{lem:kerneldecrease} holds for threshold
$y=k$ and tilt $x=x_k$. The tilted prime kernel of threshold $k$ at this
tilt, $\bar\kappa_m=\kappa^{(k)}_m\,x_k^{\,m}$,
therefore has the following five properties. The first is
Lemma~\ref{lem:kernelsign}, the second is Lemma~\ref{lem:kernelfree}, the
third and the fourth are the conclusions of
Lemmas~\ref{lem:kerneldecrease} and~\ref{lem:harmonic}, and the last is
\eqref{eq:series-tilted-recursion} for threshold $y=k$ and tilt $x=x_k$:
\begin{align}
\bar\kappa_m&\ge0&&(m\ge0),\label{eq:kernel-sign}\\
\bar\kappa_m&=\varepsilon_m\,x_k^{\,m}&&(1\le m\le k),
\label{eq:kernel-free}\\
\bar\kappa_{m+1}&\le\bar\kappa_m&&(m\ge1),\label{eq:kernel-monotone}\\
(m+1)\,\bar\kappa_{m+1}&\le m\,\bar\kappa_m&&(m\ge k),
\label{eq:kernel-harmonicstep}\\
N\,\bar\kappa_N&=\sum_{\gamma=1}^{k-1}(k-\gamma)\,a_\gamma\,x_k^{\,\gamma}\,
\bar\kappa_{N-\gamma}&&(N\ge k).\label{eq:kernel-recursion}
\end{align}
By \eqref{eq:kernel-monotone} we have $\bar\kappa_i\le\bar\kappa_m$ whenever
$1\le m\le i$. These properties are used throughout the rest of the section,
in which $k\ge2$ is a column and $\bar\kappa$ always denotes the prime
kernel of threshold $k$ at the critical tilt $x_k$; the notation does not
show the column, and we refer to $\bar\kappa$ as the kernel.
Section~\ref{sec:beyond} uses the harmonic decay
\eqref{eq:kernel-harmonicstep} in the following iterated form.

\begin{corollary}\label{cor:harmoniccritical}
Let $k\ge2$ and $L\ge k$ be integers. Then the prime kernel $\bar\kappa$ of
threshold $k$ at the critical tilt satisfies
\[
L\,\bar\kappa_L\le k\,\bar\kappa_k .
\]
\end{corollary}

\begin{proof}
This follows from \eqref{eq:kernel-harmonicstep} by induction on $L$,
starting from the equality at $L=k$.
\end{proof}

\subsection{Faster decay beyond the free range}\label{subsec:kernel-decay}

The next lemma bounds the kernel from index $2k-1$ on in terms of its value
$\bar\kappa_k$. Its proof carries a bound for the ratio of consecutive values
of the kernel from a window of $k-1$ indices to the next index by the
recursion \eqref{eq:kernel-recursion}. Section~\ref{sec:beyond} uses the
lemma for the rows that are far from the first positive row of the column.

\begin{lemma}\label{lem:quadratic}
Let $k\ge2$ and $L\ge2k-1$ be integers. Then the prime kernel $\bar\kappa$
of threshold $k$ at the critical tilt satisfies
\[
(L-1)\,L\,\bar\kappa_L\le2k\,(k-1)\,\bar\kappa_k .
\]
\end{lemma}

\begin{proof}
We claim that for all integers $t\ge1$ and $m$ with
$m\ge k+(t-1)(k-1)$,
\begin{equation}\label{eq:kernel-rate}
(m+1)\,\bar\kappa_{m+1}\le(m+1-t)\,\bar\kappa_m .
\end{equation}
For $t=1$ this is \eqref{eq:kernel-harmonicstep}. Let
$m\ge k$ be an integer and $\rho$ a real number such that
$\bar\kappa_{i+1}\le\rho\,\bar\kappa_i$ for every integer $i$ with
$m-k+1\le i\le m-1$. Then
\begin{equation}\label{eq:kernel-onestep}
(m+1)\,\bar\kappa_{m+1}\le\rho\,m\,\bar\kappa_m .
\end{equation}
Indeed, by \eqref{eq:kernel-recursion} for $N=m+1$ and for $N=m$,
\begin{align*}
(m+1)\,\bar\kappa_{m+1}
&=\sum_{\gamma=1}^{k-1}(k-\gamma)\,a_\gamma\,x_k^{\,\gamma}\,
\bar\kappa_{m+1-\gamma}\\
&\le\rho\sum_{\gamma=1}^{k-1}(k-\gamma)\,a_\gamma\,x_k^{\,\gamma}\,
\bar\kappa_{m-\gamma}=\rho\,m\,\bar\kappa_m ,
\end{align*}
because the weights $(k-\gamma)\,a_\gamma\,x_k^{\,\gamma}$ are nonnegative
and the indices $i=m-\gamma$ with $1\le\gamma\le k-1$ are those of the
window.

We prove \eqref{eq:kernel-rate} by induction on $t$; the case $t=1$ has
just been noted. Suppose that \eqref{eq:kernel-rate} holds for some
$t\ge1$, and let $m\ge k+t\,(k-1)$. Every index $i$ with
$m-k+1\le i\le m-1$ satisfies $i\ge k+(t-1)(k-1)$, so that
$(i+1)\,\bar\kappa_{i+1}\le(i+1-t)\,\bar\kappa_i$ by the induction
hypothesis, and $(m-i-1)\,\bar\kappa_{i+1}\le(m-i-1)\,\bar\kappa_i$ by
\eqref{eq:kernel-monotone}. Adding these two inequalities, we get
$m\,\bar\kappa_{i+1}\le(m-t)\,\bar\kappa_i$. The hypothesis of
\eqref{eq:kernel-onestep} thus holds with $\rho=(m-t)/m$, and
\eqref{eq:kernel-onestep} gives
$(m+1)\,\bar\kappa_{m+1}\le(m-t)\,\bar\kappa_m$, which is
\eqref{eq:kernel-rate} with $t$ replaced by $t+1$.

We now prove the claim of the lemma by induction on $L$. For $L=2k-1$,
Corollary~\ref{cor:harmoniccritical} gives
$(2k-1)\,\bar\kappa_{2k-1}\le k\,\bar\kappa_k$, and multiplying by $2k-2$
we obtain the claim. Let $L\ge2k-1$, and suppose that the claim holds for
$L$. Since $L\ge k+(k-1)$, the inequality \eqref{eq:kernel-rate} with $t=2$
and $m=L$ gives $(L+1)\,\bar\kappa_{L+1}\le(L-1)\,\bar\kappa_L$. Multiplying
by $L$ and using the induction hypothesis, we get
\[
L\,(L+1)\,\bar\kappa_{L+1}\le(L-1)\,L\,\bar\kappa_L
\le2k\,(k-1)\,\bar\kappa_k ,
\]
which is the claim for $L+1$.
\end{proof}

\subsection{The size of the kernel in the free range}
\label{subsec:kernel-size}

Corollary~\ref{cor:harmoniccritical} and Lemma~\ref{lem:quadratic} bound the
kernel beyond the free range in terms of its value $\bar\kappa_k$. We now
bound the kernel in the free range, where it is explicit, in absolute terms.
The proof of the next lemma writes the kernel in the free range with the
Poisson mass $P$ of Section~\ref{sec:prelim}, and the lemma is the first of
the two bounds that Section~\ref{sec:beyond} uses.

\begin{lemma}\label{lem:kernelatk}
For every integer $k\ge82$, the prime kernel $\bar\kappa$ of threshold $k$
at the critical tilt satisfies
\[
\bar\kappa_k\le\frac{0.662}{k\,\sqrt k} .
\]
\end{lemma}

\begin{proof}
Let $1\le m\le k$ be an integer. By \eqref{eq:kernel-free} we have
$\bar\kappa_m=\varepsilon_m\,x_k^{\,m}$, and
\[
m\,\varepsilon_m=\frac{(m-1)^{m-1}}{(m-1)!}=\nu_{m-1}=e^{m-1}\,P(m-1),
\]
where $\nu$ and the Poisson mass $P$ are as in Section~\ref{sec:prelim}.
Hence
\[
\bar\kappa_m=\frac{P(m-1)\,(e\,x_k)^m}{e\,m}\qquad(1\le m\le k).
\]
For $m\ge2$
we have $P(m-1)\le1/\sqrt{2\pi(m-1)}$ by the upper bound of
\eqref{eq:prelim-stirling}, and
$(e\,x_k)^m\le e^{3m/(2k)}$ by Lemma~\ref{lem:tiltupper}. Therefore
\begin{equation}\label{eq:kernel-freebound}
\bar\kappa_m\le\frac{1}{m\,\sqrt{2\pi(m-1)}}\,
\exp\Bigl(\frac{3m}{2k}-1\Bigr)\qquad(2\le m\le k).
\end{equation}
By \eqref{eq:kernel-freebound} with $m=k$ we have
$\bar\kappa_k\le e^{1/2}/\bigl(k\sqrt{2\pi(k-1)}\bigr)$, so that it
suffices to show that $e^{1/2}\sqrt k\le0.662\,\sqrt{2\pi(k-1)}$, or, after
squaring, that $e\,k\le0.662^2\cdot2\pi\,(k-1)$. Since $e<2.7183$ and
$0.662^2\cdot2\pi>2.7534$, this follows from the inequality
$2.7183\,k\le2.7534\,(k-1)$, which holds because $k\ge82$.
\end{proof}

The second bound is for the number $\bar\kappa_{n-q+1}$, where $n\ge2k$ is a
row beyond the first layer and $q$ is the first positive row of the column;
this is the number with which Section~\ref{subsec:kernel-order} compares the
weights that the kernel gives to the rows. There $q\le k+27k/100$, so that
the index $L=n-q+1$ satisfies $L-1=n-q\ge73k/100$; this index may lie in the
free range or beyond it.

\begin{lemma}\label{lem:kernelfar}
Let $k\ge82$ and $L$ be integers with $100\,(L-1)\ge73\,k$. Then the prime
kernel $\bar\kappa$ of threshold $k$ at the critical tilt satisfies
\[
\bar\kappa_L\le\frac{0.73}{k\,\sqrt k} .
\]
\end{lemma}

\begin{proof}
Let $m_0$ be the least integer with $100\,(m_0-1)\ge73\,k$, so that
$73\,k\le100\,(m_0-1)\le73\,k+99$. Then $m_0\le L$ and $m_0\ge2$, and
$m_0\le k$ because $100\,m_0\le73\,k+199\le100\,k$ for $k\ge82$. By
\eqref{eq:kernel-monotone} and by \eqref{eq:kernel-freebound} with
$m=m_0$, which was shown in the proof of Lemma~\ref{lem:kernelatk},
\[
\bar\kappa_L\le\bar\kappa_{m_0}\le\frac{1}{m_0\,\sqrt{2\pi(m_0-1)}}\,
\exp\Bigl(\frac32\,\sigma-1\Bigr),\qquad\text{where }\sigma=\frac{m_0}{k} .
\]
The bounds for $m_0$ give $0.73\le\sigma\le0.73+1.99/k\le0.7543$, so that
the exponent $(3/2)\,\sigma-1$ lies between $0$ and $1.7$. By the
inequality $e^{v}\le1+v+v^2$ for $0\le v\le1.7$, which is recalled in
Section~\ref{sec:prelim},
\[
\exp\Bigl(\frac32\,\sigma-1\Bigr)\le1+\Bigl(\frac32\,\sigma-1\Bigr)
+\Bigl(\frac32\,\sigma-1\Bigr)^2=1-\frac32\,\sigma+\frac94\,\sigma^2 .
\]
The right-hand side is at most $1.563368\,\sigma$. Indeed, the difference
\[
\frac94\,\sigma^2-3.063368\,\sigma+1
\]
is a convex function of $\sigma$, and it is negative at $\sigma=0.73$ and
at $\sigma=0.7543$; hence it is negative on the interval between these two
points. Finally
\[
2\pi\,(m_0-1)\ge2\pi\cdot0.73\,k\ge2.1416^2\,k ,
\]
where the second inequality is checked by direct computation, with
$2\pi>6.283$. Altogether,
\[
\bar\kappa_L\le\frac{1.563368\,\sigma}{m_0\cdot2.1416\,\sqrt k}
=\frac{1.563368}{2.1416\,k\,\sqrt k}=\frac{0.73}{k\,\sqrt k} .\qedhere
\]
\end{proof}

\subsection{The weight $\Psi_k$ and the order of the kernel weights}
\label{subsec:kernel-order}

For integers $k\ge1$ and $N\ge0$ let
\[
\Psi_k(N)=\max(N,\,2k-1).
\]
The number $\Psi_k(N)$ is the \emph{weight} of row $N$ in the argument for
column $k$. It is equal to $2k-1$ up to row $2k-1$,
the last row of the first layer, and to the row number from there on.
Clearly $\Psi_k(N)\ge N$ for every $N$, and $\Psi_k(N)\le\Psi_k(N')$
whenever $N\le N'$.

The weight serves the following step of Section~\ref{sec:beyond}. Let
$k\ge2$, and let $\eta_j$ and $U_j$ be the tilted step of column $k$ and
the tilted column count that are defined in Section~\ref{sec:diagonal},
with \eqref{eq:diagonal-columncounts}: they are the normalised step of
column $k$ and the normalised count of column $k-1$, both tilted by the
critical tilt $x_k$.
Multiplied by $x_k^{\,n}$, the inequality of
Proposition~\ref{prop:columninequality} for a row $n\ge k+1$ reads
\[
n\,\eta_n\ \ge\ (k+1)\,x_k\,\eta_{n-1}+2\,x_k\,U_{n-1}
+n\sum_{j=k}^{n-1}\bar\kappa_{n-j}\,\eta_j .
\]
The sum is the feedback of rows $k\le j\le n-1$, in which row $j$ has the
\emph{kernel weight} $\bar\kappa_{n-j}$.
In the situation of Section~\ref{sec:beyond} the column has a first positive
row $q$, and row $n$ is reached by an induction on the row: the tilted
step $\eta_j$ is not positive for $k\le j<q$ and positive for
$q\le j\le n-1$. The feedback thus contains terms that are not positive, and
nothing is known about the size of the single steps in the rows before $q$.
What is known about these rows is a lower bound for one weighted sum:
Lemma~\ref{lem:debt} bounds the sum $\sum_{j=k}^{J}\Psi_k(j)\,\eta_j$ from
below by terms of rows $J$ and $J-1$. To use this bound for the
feedback, the kernel weights $\bar\kappa_{n-j}$ have to be exchanged for the
weights $\Psi_k(j)$. This is possible term by term, without any knowledge
of the sizes of the $\eta_j$, as soon as one number $\theta\ge0$ separates
the ratios of the two weights at row $q$:
\begin{equation}\label{eq:kernel-separation}
\frac{\bar\kappa_{n-j}}{\Psi_k(j)}\le\theta\quad(k\le j<q),\qquad
\frac{\bar\kappa_{n-j}}{\Psi_k(j)}\ge\theta\quad(q\le j\le n-1).
\end{equation}
For then $\bar\kappa_{n-j}\,\eta_j\ge\theta\,\Psi_k(j)\,\eta_j$ in every row
$k\le j\le n-1$, because $\eta_j$ is not positive in the rows of the first
kind and positive in those of the second, and the feedback is at least
$\theta$ times the sum $\sum_{j=k}^{n-1}\Psi_k(j)\,\eta_j$.
The two lemmas of this
subsection show that \eqref{eq:kernel-separation} holds with
\[
\theta=\frac{\bar\kappa_{n-q+1}}{2k-1},
\]
the value of the ratio at row $j=q-1$, for every row $n\ge2k$, provided
that $k+1\le q\le k+27k/100$. This is the range in which the first positive
row of a column $k\ge82$ lies (Proposition~\ref{prop:onechange}).

The weight has its form because it has to meet two requirements. One is
the separation \eqref{eq:kernel-separation}. The other comes from
Lemma~\ref{lem:debt}. That lemma is proved for columns $k\ge7$, and for such
a column it rests on a convexity along the diagonal of the tilted
normalised counts of three consecutive rows $N-1$, $N$ and $N+1$, in which
each count is multiplied by the weight of its row
(Corollary~\ref{cor:weightedall}).
With a constant weight this convexity is proved for $k\le N\le2k-2$
(Corollary~\ref{cor:convexlayer}), and with each count multiplied by its row
number it is proved for $N\ge2k-1$ (Proposition~\ref{prop:weightedbeyond}).
The weight $\Psi_k$ is chosen so that the two statements together give the
convexity for every $N\ge k$.

For the separation a constant weight is the convenient one, since
\eqref{eq:kernel-separation} then asks only that the kernel do not increase.
This is the situation on the rows before $q$.

\begin{lemma}\label{lem:orderbefore}
Let $k\ge2$, $q$, $n$ and $j$ be nonnegative integers with
$100\,(q-k)\le27\,k$, $n\ge2k$ and $j<q$. Then the prime kernel
$\bar\kappa$ of threshold $k$ at the critical tilt and the weight $\Psi_k$
satisfy
\[
\bar\kappa_{n-j}\,(2k-1)\le\bar\kappa_{n-q+1}\,\Psi_k(j) .
\]
\end{lemma}

\begin{proof}
Since $27\,k\le100\,(k-1)$ for $k\ge2$, the hypothesis on $q$ gives
$q\le2k-1$. Hence $j<2k-1$, so that $\Psi_k(j)=2k-1$, and $q<n$, so that
$1\le n-q+1\le n-j$. By \eqref{eq:kernel-monotone} we have
$\bar\kappa_{n-j}\le\bar\kappa_{n-q+1}$, and multiplying by $2k-1$ gives
the claim.
\end{proof}

On the rows from $q$ on the weight is constant only up to row $2k-1$,
and on these rows the comparison is once more the monotonicity of the
kernel. For a row $j\ge2k$ the weight is $j$, and the second half of
\eqref{eq:kernel-separation} says that, as $j$ increases and the index
$n-j$ of the kernel decreases, the kernel grows at least in proportion to
$j$. The proof of this half depends on where the indices $n-j$ and $n-q+1$ lie
relative to the free range, and it rests on two further properties of the
kernel, one beyond the free range and one inside it.

We can now prove the second half of \eqref{eq:kernel-separation}.

\begin{lemma}\label{lem:orderafter}
Let $k\ge2$, $q$, $n$ and $j$ be integers with $k+1\le q$,
$100\,(q-k)\le27\,k$, $n\ge2k$ and $q\le j\le n-1$. Then the prime kernel
$\bar\kappa$ of threshold $k$ at the critical tilt and the weight $\Psi_k$
satisfy
\[
\bar\kappa_{n-q+1}\,\Psi_k(j)\le\bar\kappa_{n-j}\,(2k-1) .
\]
\end{lemma}

\begin{proof}
Put $L=n-q+1$ and $\ell=n-j$. As in the proof of Lemma~\ref{lem:orderbefore}
we have $q\le2k-1<n$, and $q\le j\le n-1$ gives $1\le\ell<L$.

If $j\le2k-1$, then $\Psi_k(j)=2k-1$, and the claim is
$\bar\kappa_L\le\bar\kappa_\ell$, which holds by \eqref{eq:kernel-monotone}.

Let $j\ge2k$. Then $\Psi_k(j)=j$, and we have to show that
\begin{equation}\label{eq:kernel-claim}
j\,\bar\kappa_L\le(2k-1)\,\bar\kappa_\ell .
\end{equation}
We distinguish three cases, according to the position of indices $\ell$
and $L$ relative to the free range.

Suppose first that $\ell\ge k$. The kernel decays by a fixed factor beyond
the free range: for every integer $m\ge k$,
\begin{equation}\label{eq:kernel-factor}
(2k-1)\,\bar\kappa_{m+1}\le(2k-2)\,\bar\kappa_m .
\end{equation}
To prove it, let $t\ge1$ be the integer with
$k+(t-1)(k-1)\le m<k+t\,(k-1)$. By \eqref{eq:kernel-rate}, which was shown
in the proof of Lemma~\ref{lem:quadratic},
$(m+1)\,\bar\kappa_{m+1}\le(m+1-t)\,\bar\kappa_m$. Moreover
$m+1\le k+t\,(k-1)\le t\,(2k-1)$, and therefore
\[
(2k-1)\,(m+1-t)-(2k-2)\,(m+1)=m+1-t\,(2k-1)\le0 .
\]
Since $\bar\kappa_m\ge0$ by \eqref{eq:kernel-sign}, the two inequalities
give
\[
(2k-1)\,(m+1)\,\bar\kappa_{m+1}\le(2k-1)\,(m+1-t)\,\bar\kappa_m
\le(2k-2)\,(m+1)\,\bar\kappa_m ,
\]
and division by $m+1$ gives \eqref{eq:kernel-factor}.

We show by induction on $i$ that, for every
integer $i$ with $2k-1\le i\le n-k$,
\begin{equation}\label{eq:kernel-ratio}
i\,\bar\kappa_{n-2k+1}\le(2k-1)\,\bar\kappa_{n-i} .
\end{equation}
For $i=2k-1$ the two sides are equal. Suppose that \eqref{eq:kernel-ratio}
holds for $i$, and that $i+1\le n-k$. By \eqref{eq:kernel-factor} with
$m=n-i-1\ge k$ we have
$(2k-1)\,\bar\kappa_{n-i}\le(2k-2)\,\bar\kappa_{n-i-1}$, so that
$i\,\bar\kappa_{n-2k+1}\le(2k-2)\,\bar\kappa_{n-i-1}$. Since
$(2k-2)\,(i+1)\le(2k-1)\,i$ for $i\ge2k-2$, it follows that
$(i+1)\,\bar\kappa_{n-2k+1}\le(2k-1)\,\bar\kappa_{n-i-1}$, which is
\eqref{eq:kernel-ratio} for $i+1$. Now
$2k\le j=n-\ell\le n-k$, so that \eqref{eq:kernel-ratio} with $i=j$ gives
$j\,\bar\kappa_{n-2k+1}\le(2k-1)\,\bar\kappa_\ell$. Since $q\le2k$, we have
$L\ge n-2k+1\ge1$, hence $\bar\kappa_L\le\bar\kappa_{n-2k+1}$ by
\eqref{eq:kernel-monotone}, and \eqref{eq:kernel-claim} follows.

Suppose next that $\ell<k$ and $L\le k$. Let $1\le m\le k$ be an integer.
By \eqref{eq:kernel-free} we have $\bar\kappa_m=\varepsilon_m\,x_k^{\,m}$,
and
\[
m\,\varepsilon_m=\frac{(m-1)^{m-1}}{(m-1)!}=\nu_{m-1}=e^{m-1}\,P(m-1),
\]
where $\nu$ and the Poisson mass $P$ are as in Section~\ref{sec:prelim}.
Hence
\begin{equation}\label{eq:kernel-closed}
\bar\kappa_m=\frac{P(m-1)\,(e\,x_k)^m}{e\,m}\qquad(1\le m\le k).
\end{equation}
We show that the kernel drops across a long gap in the free range:
for all integers $i$ and $m$ with $1\le i<m\le k$ and
$100\,(m-i)\ge73\,k$,
\begin{equation}\label{eq:kernel-drop}
1.22\,\bar\kappa_m\le\bar\kappa_i .
\end{equation}
For $i=1$ we have $\bar\kappa_1=x_k$ and $\bar\kappa_2=x_k^{\,2}/2$ by
\eqref{eq:kernel-free}, and since $m\ge2$ and $x_k\le1$, the monotonicity
\eqref{eq:kernel-monotone} gives
$\bar\kappa_m\le\bar\kappa_2\le\bar\kappa_1/2$. Let $i\ge2$. By
\eqref{eq:kernel-closed} for indices $m$ and $i$,
\[
\bar\kappa_m=\frac{P(m-1)}{P(i-1)}\,(e\,x_k)^{m-i}\,\frac{i}{m}\;
\bar\kappa_i .
\]
Here $P(m-1)\le1/\sqrt{2\pi(m-1)}$ by the upper bound of
\eqref{eq:prelim-stirling}, and
$P(i-1)\ge(11/12)/\sqrt{2\pi(i-1)}$ by its lower bound,
because $1-1/(12\,(i-1))\ge11/12$. By Lemma~\ref{lem:tiltupper},
$(e\,x_k)^{m-i}\le e^{3(m-i)/(2k)}\le e^{3/2}$, because $m-i\le k$. The
gap gives $m-i\ge0.73\,k\ge0.73\,m$, hence $i\le0.27\,m$ and
$i-1\le0.27\,(m-1)$. Therefore
\[
\bar\kappa_m\le\frac{12}{11}\,\sqrt{\frac{i-1}{m-1}}\;e^{3/2}\,\frac{i}{m}\;
\bar\kappa_i\le\frac{12}{11}\,e^{3/2}\,0.27^{3/2}\;\bar\kappa_i .
\]
The number $1.22\cdot(12/11)\,e^{3/2}\,0.27^{3/2}$ is smaller than $1$,
because its square $1.22^2\cdot(144/121)\,e^3\cdot0.27^3$ is smaller than
$0.71$, as a direct computation shows. This proves \eqref{eq:kernel-drop}.

We have $L-\ell=j-q+1$, and
therefore
\[
100\,(L-\ell)\ge100\,(2k-q+1)=100\,(k+1)-100\,(q-k)\ge73\,k .
\]
Thus \eqref{eq:kernel-drop} applies to indices $\ell<L$ and gives
$1.22\,\bar\kappa_L\le\bar\kappa_\ell$. On the other hand, $L\le k$ says that
$n\le k+q-1$, so that $j\le n-1\le k+q-2$ and
\[
100\,j\le100\,(2k-2)+100\,(q-k)\le227\,k-200\le122\,(2k-1).
\]
Hence
\[
j\,\bar\kappa_L\le1.22\,(2k-1)\,\bar\kappa_L\le(2k-1)\,\bar\kappa_\ell .
\]

Suppose finally that $\ell<k<L$. By Corollary~\ref{cor:harmoniccritical} we
have $L\,\bar\kappa_L\le k\,\bar\kappa_k$, and it remains to compare
$\bar\kappa_k$ with $\bar\kappa_\ell$; how this is done depends on the size of
$\ell$. Note that $j=n-\ell=L+q-1-\ell$.

If $100\,\ell\le27\,k$, then \eqref{eq:kernel-drop} applies to
indices $\ell<k$ and gives $1.22\,\bar\kappa_k\le\bar\kappa_\ell$. Since
$\ell\ge1$ and $100\,(q-k)\le27\,k$, we have $100\,j\le100\,L+127\,k-200$.
Moreover, since $L\ge k+1$,
\[
127\,k^2-200\,k\le(144\,k-122)\,(k+1)\le(144\,k-122)\,L .
\]
Adding $100\,k\,L$, we obtain
\[
100\,k\,j\le100\,k\,L+127\,k^2-200\,k\le122\,(2k-1)\,L ,
\]
so that $k\,j\le1.22\,(2k-1)\,L$. Hence
\[
j\,L\,\bar\kappa_L\le k\,j\,\bar\kappa_k\le1.22\,(2k-1)\,L\,\bar\kappa_k
\le(2k-1)\,L\,\bar\kappa_\ell .
\]

If $100\,\ell>27\,k$, then $\bar\kappa_k\le\bar\kappa_\ell$ by
\eqref{eq:kernel-monotone}, because $\ell<k$. Moreover
\[
100\,(q-\ell)=100\,(q-k)+100\,(k-\ell)<100\,k ,
\]
so that $q-\ell\le k-1$ and $j\le L+k-2$. Since $L\ge k+1$, this gives
\[
k\,j\le k\,L+k^2-2k\le k\,L+(k-1)\,(k+1)\le(2k-1)\,L ,
\]
and hence
\[
j\,L\,\bar\kappa_L\le k\,j\,\bar\kappa_k\le(2k-1)\,L\,\bar\kappa_k
\le(2k-1)\,L\,\bar\kappa_\ell .
\]
For both sizes of $\ell$, division by $L$ gives \eqref{eq:kernel-claim}.
\end{proof}

\section{Property K and diagonal convexity in the first layer}\label{sec:propk}
In this section we prove that, for a column $k\ge7$ and at the critical tilt
of that column, the tilted normalised counts of thresholds $k-1$, $k$ and
$k+1$ are convex along the diagonal in the first layer
(Corollary~\ref{cor:convexlayer}), and that they remain so when they are
multiplied by the weights $c-1$, $c$ and $c+1$ for a real number $c\ge k$
(Corollary~\ref{cor:weightedlayer}). These two inequalities are what the
argument beyond the first layer, in Section~\ref{sec:beyond}, takes from the
first layer. We prove them by an induction along the last-car recursion,
which is carried out for a stronger property,
property K (Definition~\ref{def:propertyk}).

\subsection{Diagonal triples}\label{subsec:propk-triples}

Throughout the section the tilt is the critical tilt
$x_k=(k/(k+1))^{k-1}$ of a column $k\ge1$
(Section~\ref{sec:propagation}). As before, for integers $n\ge0$, $N\ge0$
and $y\ge1$,
\[
\omega_n=c_n\,x_k^{\,n}\qquad\text{and}\qquad
u^{(y)}_N=z^{(y)}_N\,x_k^{\,N}
\]
denote the tilted Cayley masses and the tilted normalised count of threshold
$y$ at this tilt; we say masses and tilted counts for short, and column $k$
is not shown in the notation. The index $N$ of a tilted count
is its \emph{size}, the length of the parking functions that the count
refers to. The
tilt is positive, so that $\omega_n>0$ and $u^{(y)}_N\ge0$.
By Proposition~\ref{prop:freerange} a
tilted count in the free range of its threshold is a mass:
\begin{equation}\label{eq:propk-free}
u^{(y)}_N=\omega_N\qquad(N\le y).
\end{equation}
The \emph{offset} of a tilted count $u^{(y)}_N$ is the integer $N-y$, so
that the counts of \eqref{eq:propk-free} are those whose offset is at most
$0$.

For integers $k\ge2$ and $\rho\ge0$ the \emph{diagonal triple of offset
$\rho$} of column $k$ is
\[
\bigl(u^{(k-1)}_{k-1+\rho},\ u^{(k)}_{k+\rho},\ u^{(k+1)}_{k+1+\rho}\bigr).
\]
Its members
are the tilted counts of the three thresholds $k-1$, $k$ and $k+1$, from
which the quantities of column $k$ are formed
(Section~\ref{sec:counts}), each at offset $\rho$ and all tilted by the
same number $x_k$. For $0\le\rho\le k-2$ the
largest of the three sizes, $k+1+\rho$, runs through rows $k+1,\dots,2k-1$ of
the first layer of column $k$.

The induction on the offset uses the last-car recursion in a form that
the proof of Proposition~\ref{prop:diagonalk} derives from
Lemma~\ref{lem:layersplit}. For integers $\rho$ and $y$ with
$1\le\rho\le y-1$ the \emph{direct term} is
\[
S_\rho(y)=x_k\,\omega_{\rho-1}\,\omega_y
+\frac{x_k}{2}\cdot\frac{y+\rho+1}{y+\rho}\,Q_\rho(y-\rho-1),
\]
where, for integers $\rho\ge0$ and $j\ge0$,
\begin{equation}\label{eq:propk-pairsums}
Q_\rho(j)=\sum_{i=0}^{j}\omega_{\rho+i}\,\omega_{\rho+j-i} .
\end{equation}
That proof writes the tilted count $u^{(y)}_{y+\rho}$ as the
sum of the direct term and of the \emph{chain part}, which comes from the
sum in Lemma~\ref{lem:layersplit} and contains the counts
$u^{(y)}_{y+\rho-1-i}$ of the same threshold at offsets $1,\dots,\rho-1$,
each multiplied by a number $x_k\,\omega_i$ and by a factor
$(y+i+1)/(y+\rho)$, for $0\le i\le\rho-2$.
Taken for the
three thresholds $y=k-1,k,k+1$, the proof thus builds the diagonal triple of
offset $\rho$ from the diagonal triples of smaller offset and from the
triple of the direct terms by three
operations: sums; multiplication by the numbers $x_k\,\omega_i\ge0$, which
are the same for the three thresholds; and multiplication by the factors
$(y+i+1)/(y+\rho)$, which depend on the threshold and increase with it.

\subsection{Property K and its stability under factors}
\label{subsec:propk-definition}

Property K concerns a triple of real numbers, whose members we write with
indices $-$, $0$ and $+$ in their order, and one further real number
$k$, which in the applications is the column.

\begin{definition}\label{def:propertyk}
Let $k\ne0$ be a real number. A triple $(t_-,t_0,t_+)$ of real numbers
\emph{has property K at $k$} if
\[
2\,e^{\mu/k}\,t_0\le e^{\mu}\,t_++e^{-\mu}\,t_-
\qquad\text{for every real number }\mu\ge0 .
\]
\end{definition}

In the applications
$t_-$, $t_0$, $t_+$ are the values at the three consecutive thresholds
$k-1$, $k$, $k+1$ of one quantity that depends on a threshold. The order of
the triple is the order of the thresholds, so that the value $t_-$ at the
lower threshold carries the factor $e^{-\mu}$ and the value $t_+$ at the
upper threshold the factor $e^{\mu}$. For $\mu=0$ the inequality is the
convexity $2\,t_0\le t_-+t_+$ of the three values. For $\mu>0$ it says that
the triple stays convex after the lower value has been multiplied by
$e^{-\mu}$ and the upper value by $e^{\mu}$, and that it does so with the
additional factor $e^{\mu/k}$ on the middle value. For each fixed $\mu$ the
inequality
is linear in the triple.
Hence the sum of two triples with property K at $k$ has property K at $k$,
and so has the multiple of such a triple by a nonnegative real number.

Convexity alone is not preserved when the three values are multiplied by
positive factors that increase with the threshold, as the factors of the
chain part do. The triple $(2,1,0)$ is convex.
Multiplied by factors $\theta_-<\theta_0<\theta_+$ it becomes
$(2\theta_-,\theta_0,0)$, which is not convex, because
$2\theta_0>2\theta_-$. For this reason the induction on the offset
cannot be carried out for convexity itself, and
we carry it out for property K. Its stability under factors rests on the
following computation.
Let $k>0$ be a real number and let $(t_-,t_0,t_+)$ be a triple with
property K at $k$. Then for all real numbers $L$, $M$ and $\mu$ with
$\mu+(L-M)/2\ge0$,
\begin{equation}\label{eq:propk-shift}
2\,e^{\mu/k}\exp\Bigl(\frac{(k+1)L+(k-1)M}{2k}\Bigr)\,t_0
\le e^{\mu}\,e^{L}\,t_++e^{-\mu}\,e^{M}\,t_- .
\end{equation}
Indeed, the number $\mu'=\mu+(L-M)/2$ is nonnegative, so that
$2e^{\mu'/k}\,t_0\le e^{\mu'}\,t_++e^{-\mu'}\,t_-$. Multiplying this
inequality by $e^{(L+M)/2}$ gives \eqref{eq:propk-shift}: the exponents on
the right become $\mu+L$ and $-\mu+M$, and the exponent on the left becomes
\[
\frac{L+M}{2}+\frac{\mu'}{k}=\frac{\mu}{k}+\frac{(k+1)L+(k-1)M}{2k} .
\]

In words, to multiply the upper value by $e^{L}$ and the lower value by
$e^{M}$ has the same effect on the inequality of property K as to move the
parameter from $\mu$ to $\mu+(L-M)/2$ and to multiply the middle value by
the exponential factor on the left-hand side of \eqref{eq:propk-shift}. If
$L\ge M$, the parameter moves upward, and \eqref{eq:propk-shift}
holds for every $\mu\ge0$. This is why the property is required for all
$\mu\ge0$: a hypothesis that held for $\mu=0$ only could not be used after a
single increasing factor. If $L<M$, the parameter moves downward, and
\eqref{eq:propk-shift} is available only for $\mu\ge(M-L)/2$, because
property K says nothing for negative parameters. We meet this case in
the proof of Lemma~\ref{lem:directk}.

The proofs of the two stability lemmas below use \eqref{eq:propk-shift} for
positive factors $\theta_-$, $\theta_0$, $\theta_+$ that do not decrease
from the lower to the upper threshold, and they verify that the number
\[
(k+1)\log\frac{\theta_+}{\theta_0}+(k-1)\log\frac{\theta_-}{\theta_0}
\]
is nonnegative. This condition on the factors explains the factor
$e^{\mu/k}$ in Definition~\ref{def:propertyk}. When $\theta_-$, $\theta_0$,
$\theta_+$ are the values at $k-1$, $k$, $k+1$ of a positive function
$\theta$ of the threshold, the number above is the
second difference of the function $y\mapsto y\log\theta(y)$ at these three
points. The weights $k+1$ and $k-1$ come from the exponent $\mu/k$. Without
the factor $e^{\mu/k}$, the same computation would lead to the condition
$\log\theta_++\log\theta_-\ge2\log\theta_0$, the log-convexity of $\theta$
at the three points. With it, the condition is the one that the proof of
Lemma~\ref{lem:quotientfactor} verifies for the factors of the chain part,
through the relative increments of the factors.

The first of the two stability lemmas treats the factors
$(y+i+1)/(y+\rho)$ that thresholds $y=k-1$, $k$ and $k+1$ carry in the
chain part.

\begin{lemma}\label{lem:quotientfactor}
Let $k>1$ be a real number, let $(t_-,t_0,t_+)$ be a triple of real numbers
that has property K at $k$ and satisfies $t_0\ge0$, and let $\alpha$ and
$\beta$ be real numbers with $0\le\alpha\le\beta$. Then the triple
\[
\Bigl(\frac{k-1+\alpha}{k-1+\beta}\,t_-,\ \frac{k+\alpha}{k+\beta}\,t_0,
\ \frac{k+1+\alpha}{k+1+\beta}\,t_+\Bigr)
\]
has property K at $k$.
\end{lemma}

\begin{proof}
Let $\theta_-$, $\theta_0$, $\theta_+$ be any positive real numbers such
that
\begin{equation}\label{eq:propk-factors}
\theta_-\le\theta_+\qquad\text{and}\qquad
(k+1)\log\frac{\theta_+}{\theta_0}+(k-1)\log\frac{\theta_-}{\theta_0}\ge0 .
\end{equation}
Then the triple $(\theta_-t_-,\theta_0t_0,\theta_+t_+)$ has property K at
$k$. To see this, put $L=\log(\theta_+/\theta_0)$ and
$M=\log(\theta_-/\theta_0)$, and let $E=(k+1)L+(k-1)M$ be the left-hand
side of the second condition.
Since $L\ge M$, every $\mu\ge0$ satisfies $\mu+(L-M)/2\ge0$, and
\eqref{eq:propk-shift} gives
\[
2\,e^{\mu/k}\,e^{E/(2k)}\,t_0\le e^{\mu}\,\frac{\theta_+}{\theta_0}\,t_+
+e^{-\mu}\,\frac{\theta_-}{\theta_0}\,t_- .
\]
As $e^{E/(2k)}\ge1$ and $t_0\ge0$, the left-hand side is at least
$2e^{\mu/k}\,t_0$. Multiplying by $\theta_0$, we obtain the inequality of
property K for the triple $(\theta_-t_-,\theta_0t_0,\theta_+t_+)$.

For such positive numbers $\theta_-$, $\theta_0$, $\theta_+$ put
\[
\sigma_+=\frac{\theta_+}{\theta_0}-1,\qquad
\sigma_-=\frac{\theta_0}{\theta_-}-1 .
\]
If
\begin{equation}\label{eq:propk-increments}
\sigma_+\ge0\qquad\text{and}\qquad
(k-1)\,\sigma_-\le(k+1)\,\frac{2\sigma_+}{\sigma_++2},
\end{equation}
then the second condition of \eqref{eq:propk-factors} holds. This follows
from the elementary inequalities $\log(1+u)\ge2u/(u+2)$ for $u\ge0$, which
is easily verified by differentiation, and $\log r\ge1-1/r$ for $r>0$. The
first, with $u=\sigma_+$, gives
$\log(\theta_+/\theta_0)\ge2\sigma_+/(\sigma_++2)$, and the second, with
$r=\theta_-/\theta_0$, gives $\log(\theta_-/\theta_0)\ge-\sigma_-$. Hence
\[
(k+1)\log\frac{\theta_+}{\theta_0}+(k-1)\log\frac{\theta_-}{\theta_0}
\ge(k+1)\,\frac{2\sigma_+}{\sigma_++2}-(k-1)\,\sigma_-\ge0 .
\]

Now let $\theta(y)=(y+\alpha)/(y+\beta)$, and put $\theta_-=\theta(k-1)$,
$\theta_0=\theta(k)$ and $\theta_+=\theta(k+1)$. These numbers are positive,
and the triple of the lemma is $(\theta_-t_-,\theta_0t_0,\theta_+t_+)$. So
it suffices to verify the two conditions of \eqref{eq:propk-factors}, and
for the second of them we verify \eqref{eq:propk-increments}.

With the denominators cleared, the inequality $\theta_-\le\theta_+$ reduces
to $2(\beta-\alpha)\ge0$. A direct computation of the two quotients
$\theta_+/\theta_0$ and $\theta_0/\theta_-$ gives
\[
\sigma_+=\frac{\beta-\alpha}{(k+\alpha)(k+1+\beta)},\qquad
\sigma_-=\frac{\beta-\alpha}{(k+\beta)(k-1+\alpha)} .
\]
Thus $\sigma_+\ge0$, and with the denominators cleared, the second
inequality of \eqref{eq:propk-increments} reads
\begin{multline*}
(k-1)(\beta-\alpha)\bigl[2(k+\alpha)(k+1+\beta)+(\beta-\alpha)\bigr]\\
\le2(k+1)(\beta-\alpha)(k+\beta)(k-1+\alpha),
\end{multline*}
and when both sides are expanded, the right-hand side minus the left-hand
side is found to be
\[
(\beta-\alpha)\bigl[\alpha\,(4\beta+3k+1)+\beta\,(k-1)\bigr],
\]
which is nonnegative.
\end{proof}

The second stability lemma treats the weights $c-1$, $c$ and $c+1$ of
Corollary~\ref{cor:weightedlayer}, with a real number
$c\ge k$. These weights are the values at $y=k-1$, $k$ and $k+1$ of the
factor $y+\alpha$ with $\alpha=c-k$. This factor increases with the
threshold, but it is not a quotient of the kind treated in
Lemma~\ref{lem:quotientfactor}. We state the lemma for the more general
factor $(y+\alpha)(y+\beta)/y$ with $\alpha\ge0$ and $\beta\ge0$, of which
the factor $y+\alpha$ is the case $\beta=0$.

\begin{lemma}\label{lem:productfactor}
Let $k>1$ be a real number, let $(t_-,t_0,t_+)$ be a triple of real numbers
that has property K at $k$ and satisfies $t_0\ge0$, and let $\alpha\ge0$
and $\beta\ge0$ be real numbers with $\alpha\beta\le k^2-1$. Then the
triple
\[
\Bigl(\frac{(k-1+\alpha)(k-1+\beta)}{k-1}\,t_-,
\ \frac{(k+\alpha)(k+\beta)}{k}\,t_0,
\ \frac{(k+1+\alpha)(k+1+\beta)}{k+1}\,t_+\Bigr)
\]
has property K at $k$.
\end{lemma}

\begin{proof}
Let $\theta(y)=(y+\alpha)(y+\beta)/y$ for $y>0$, and put
$\theta_-=\theta(k-1)$, $\theta_0=\theta(k)$ and $\theta_+=\theta(k+1)$.
These numbers are positive, and the triple of the lemma is
$(\theta_-t_-,\theta_0t_0,\theta_+t_+)$. By what was shown in the proof of
Lemma~\ref{lem:quotientfactor} for any positive numbers $\theta_-$,
$\theta_0$, $\theta_+$, it suffices to verify the first condition of
\eqref{eq:propk-factors} and the two conditions of
\eqref{eq:propk-increments}.

Since $\theta(y)=y+\alpha+\beta+\alpha\beta/y$, we have
\[
\theta(k+1)-\theta(k-1)=2-\frac{2\,\alpha\beta}{k^2-1}\ge0 ,
\]
because $\alpha\beta\le k^2-1$; this is the first condition. In the same way
$\theta(k+1)-\theta(k)=1-\alpha\beta/(k(k+1))$ and
$\theta(k)-\theta(k-1)=1-\alpha\beta/(k(k-1))$, and dividing by $\theta(k)$
and by $\theta(k-1)$ we obtain
\[
\sigma_+=\frac{k(k+1)-\alpha\beta}{(k+1)(k+\alpha)(k+\beta)},\qquad
\sigma_-=\frac{k(k-1)-\alpha\beta}{k\,(k-1+\alpha)(k-1+\beta)} .
\]
Here $\sigma_+\ge0$, because $\alpha\beta\le k^2-1\le k(k+1)$; the number
$\sigma_-$ may be negative. With the denominators cleared, the second
inequality of \eqref{eq:propk-increments} reads
\begin{multline*}
(k-1)\bigl(k(k-1)-\alpha\beta\bigr)
\bigl[2(k+1)(k+\alpha)(k+\beta)+k(k+1)-\alpha\beta\bigr]\\
\le2k\,(k+1)\bigl(k(k+1)-\alpha\beta\bigr)(k-1+\alpha)(k-1+\beta) .
\end{multline*}
To compare the two sides we write them in terms of $s=\alpha+\beta$ and
$p=\alpha\beta$. We have
\begin{align*}
(k+\alpha)(k+\beta)&=k^2+ks+p,\\
(k-1+\alpha)(k-1+\beta)&=(k-1)^2+(k-1)\,s+p .
\end{align*}
The left-hand side of the inequality is therefore $k-1$ times
\begin{align*}
&\bigl(k(k-1)-p\bigr)\bigl[2(k+1)(k^2+ks+p)+k(k+1)-p\bigr]\\
&\qquad=\bigl(k(k-1)-p\bigr)
\bigl[k(k+1)(2k+1)+2k(k+1)\,s+(2k+1)\,p\bigr]\\
&\qquad=k^2(k^2-1)(2k+1)+2k^2(k^2-1)\,s-2k(k+1)\,sp\\
&\qquad\quad-2k(2k+1)\,p-(2k+1)\,p^2 ,
\end{align*}
and the right-hand side is $2k(k+1)$ times
\begin{align*}
&\bigl(k(k+1)-p\bigr)\bigl[(k-1)^2+(k-1)\,s+p\bigr]\\
&\qquad=k(k+1)(k-1)^2+k(k^2-1)\,s-(k-1)\,sp+(3k-1)\,p-p^2 .
\end{align*}
In the first expansion the coefficient of $p$ is $2k+1$ times
$k(k-1)-k(k+1)=-2k$, and in the second it is $k(k+1)-(k-1)^2=3k-1$. The
terms with $sp$ are $-2k(k-1)(k+1)\,sp$ on both sides of the inequality,
so that they cancel in the difference of the two sides. For the monomials
$1$, $s$, $p$ and $p^2$, in this order, the coefficient on the right-hand
side minus the coefficient on the left-hand side is
\begin{align*}
2k^2(k+1)^2(k-1)^2-k^2(k-1)^2(k+1)(2k+1)&=k^2(k-1)^2(k+1),\\
2k^2(k+1)(k^2-1)-2k^2(k-1)(k^2-1)&=4k^2(k^2-1),\\
2k(k+1)(3k-1)+2k(k-1)(2k+1)&=10k^3+2k^2-4k,\\
-2k(k+1)+(k-1)(2k+1)&=-(3k+1) .
\end{align*}
With $s=\alpha+\beta$ and $p=\alpha\beta$ inserted again, the right-hand
side minus the left-hand side is therefore
\begin{multline*}
k^2(k-1)^2(k+1)+4k^2(k^2-1)(\alpha+\beta)\\
+\alpha\beta\,\bigl(10k^3+2k^2-4k-(3k+1)\,\alpha\beta\bigr).
\end{multline*}
The first two terms are nonnegative. So is the third, because
$\alpha\beta\ge0$ and $\alpha\beta\le k^2-1$ gives
\begin{align*}
10k^3+2k^2-4k-(3k+1)\,\alpha\beta
&\ge10k^3+2k^2-4k-(3k+1)(k^2-1)\\
&=7k^3+k^2-k+1>0 . \qedhere
\end{align*}
\end{proof}

\subsection{The tilted Cayley masses}\label{subsec:propk-masses}

In this subsection we prove that three consecutive tilted Cayley masses
with indices up to $k+1$ satisfy the inequality of property K at $k$. For
the last triple, $(\omega_{k-1},\omega_k,\omega_{k+1})$, this is the base of
the induction on the offset. For the other triples it is the hypothesis
under which Section~\ref{subsec:propk-convolution} treats the sums of
products of two masses that occur in the direct term.

Everything is read off from the ratio of consecutive masses. For a real
number $y>0$ let
\[
\varphi(y)=(y-1)\log\Bigl(1+\frac1y\Bigr).
\]
For an integer $n\ge1$ we have $e^{\varphi(n)}=(1+1/n)^{n-1}$. Hence
$x_n=e^{-\varphi(n)}$, and \eqref{eq:prelim-quotient}, taken at $j=n+1$ and
written for the numbers $c_n=a_{n+1}$, says that
$c_{n+1}=e^{\varphi(n+1)}\,c_n$ for every integer $n\ge0$. Multiplying by
$x_k^{\,n+1}$ and using $x_k=e^{-\varphi(k)}$, we obtain for every
column $k\ge1$
\begin{equation}\label{eq:propk-ratio}
\omega_{n+1}=\omega_n\,e^{\varphi(n+1)-\varphi(k)}\qquad(n\ge0).
\end{equation}
By Lemma~\ref{lem:powermonotone} the function $\varphi$ does not decrease on
the positive integers. So the ratio $\omega_{n+1}/\omega_n$ does not
decrease in $n$, and it equals $1$
for $n+1=k$:
\begin{equation}\label{eq:propk-flat}
\omega_k=\omega_{k-1}\qquad(k\ge1).
\end{equation}

The sequence $\omega$ is also log-convex, in the following form: for all
integers $\rho\ge0$, $m\ge0$ and $i$ with $0\le i\le m$,
\begin{equation}\label{eq:propk-logconvex}
\omega_{\rho+i}\,\omega_{\rho+m-i}\le\omega_\rho\,\omega_{\rho+m} .
\end{equation}
Indeed, since $c_n=a_{n+1}$, Lemma~\ref{lem:cayleyconvex}, applied to
indices $\rho+1\le\rho+m-i+1$ with the shift $i$, gives
$c_{\rho+i}\,c_{\rho+m-i}\le c_\rho\,c_{\rho+m}$, and the two sides of
\eqref{eq:propk-logconvex} carry the same power $x_k^{\,2\rho+m}$ of the
tilt. The inequality is used in the proof of Lemma~\ref{lem:directk}.

By \eqref{eq:propk-ratio}, the inequality of property K for three
consecutive masses, divided by the middle one, is an inequality between
exponentials of increments of $\varphi$. The proofs below bound these
increments with the help of the formula
\begin{equation}\label{eq:propk-integral}
\log\Bigl(1+\frac1y\Bigr)=\int_0^1\frac{dt}{y+t} ,
\end{equation}
which holds for every real number $y>0$.

They also use the inequality
between the arithmetic and the geometric mean, with the weights that the
exponent $\mu/k$ of property K dictates. For a real number $k>1$ let
\[
\chi_k=(k+1)\log\frac{k+1}{k}+(k-1)\log\frac{k-1}{k} ;
\]
this is the second difference of the function $t\mapsto t\log t$ at the
points $k-1$, $k$ and $k+1$.

We now prove the inequality of property K for the triples of
masses other than the last one.

\begin{lemma}\label{lem:massesmiddle}
Let $k\ge3$ be an integer, and let $\omega_n$ be the tilted Cayley masses
at the critical tilt $x_k$ of column $k$. Then for every integer $j\ge0$
with $j+2\le k$ and every real number $\mu$,
\[
2\,e^{\mu/k}\,\omega_{j+1}\le e^{\mu}\,\omega_{j+2}+e^{-\mu}\,\omega_j .
\]
\end{lemma}

\begin{proof}
For all real numbers $L$, $M$ and $\mu$,
\begin{equation}\label{eq:propk-mean}
e^{\mu+L}+e^{-\mu+M}\ge
2\,e^{\mu/k}\exp\Bigl(\frac{(k+1)L+(k-1)M-\chi_k}{2k}\Bigr).
\end{equation}
To prove this, we write
\[
e^{\mu+L}=\frac{k+1}{k}\,e^{p_+},\qquad
e^{-\mu+M}=\frac{k-1}{k}\,e^{p_-},
\]
with $p_+=\mu+L-\log\frac{k+1}{k}$ and $p_-=-\mu+M-\log\frac{k-1}{k}$. The
two coefficients are positive and add up to $2$. Let
\[
\bar p=\frac{k+1}{2k}\,p_++\frac{k-1}{2k}\,p_-
=\frac{\mu}{k}+\frac{(k+1)L+(k-1)M-\chi_k}{2k} ,
\]
so that the right-hand side of \eqref{eq:propk-mean} is $2e^{\bar p}$. As
$\bar p$ is the mean of $p_+$ and $p_-$ with these coefficients as weights,
Jensen's inequality for the exponential function, in the form of
Section~\ref{sec:prelim}, shows that the left-hand side is at least
$2e^{\bar p}$.

Now let $n=j+1$, so that $1\le n\le k-1$, and put
$L=\varphi(n+1)-\varphi(k)$ and $M=\varphi(k)-\varphi(n)$. By
\eqref{eq:propk-ratio} we have $\omega_{j+2}=\omega_{j+1}\,e^{L}$ and
$\omega_j=\omega_{j+1}\,e^{M}$. Since $\omega_{j+1}>0$, the claim is that
$e^{\mu+L}+e^{-\mu+M}\ge2e^{\mu/k}$, and by \eqref{eq:propk-mean} it
suffices to show that $(k+1)L+(k-1)M\ge\chi_k$. For $1\le m\le k-1$ let
\[
h_m=(k+1)\,\varphi(m+1)-(k-1)\,\varphi(m).
\]
Then $(k+1)L+(k-1)M=h_n-2\,\varphi(k)$, and what has to be shown is that
\begin{equation}\label{eq:propk-exponent}
h_n-2\,\varphi(k)\ge\chi_k .
\end{equation}

For every real number $y\ge2$,
\begin{align}
\varphi(y)-\varphi(y-1)&\ge\frac{3}{2\,y\,(y+1)},
\label{eq:propk-increment}\\
(y+2)\bigl(\varphi(y+1)-\varphi(y)\bigr)
&\le y\,\bigl(\varphi(y)-\varphi(y-1)\bigr).
\label{eq:propk-decay}
\end{align}
For \eqref{eq:propk-increment} we write the left-hand side as
$(y-1)\log(1+1/y)-(y-2)\log(1+1/(y-1))$. By \eqref{eq:propk-integral} it is
the integral over $[0,1]$ of
\[
\frac{y-1}{y+t}-\frac{y-2}{y-1+t}=\frac{1+t}{(y+t)(y-1+t)}
\ge\frac{1+t}{y\,(y+1)},
\]
where the inequality holds because $0<(y+t)(y-1+t)\le(y+1)\,y$ for
$0\le t\le1$; and the integral of $1+t$ over $[0,1]$ is $3/2$. For
\eqref{eq:propk-decay}, the right-hand side minus the left-hand side is
$(2y+2)\,\varphi(y)-y\,\varphi(y-1)-(y+2)\,\varphi(y+1)$. By the definition
of $\varphi$ and by \eqref{eq:propk-integral} this is the integral over
$[0,1]$ of
\begin{multline*}
\frac{2\,(y^2-1)}{y+t}-\frac{y\,(y-2)}{y-1+t}-\frac{y\,(y+2)}{y+1+t}\\
=\frac{2\,(1-t^2)}{(y-1+t)(y+t)(y+1+t)},
\end{multline*}
which is nonnegative for $0\le t\le1$.

We now show that $h_{m+1}\le h_m$ for $1\le m\le k-2$. By
\eqref{eq:propk-decay} with $y=m+1$,
\[
(m+3)\bigl(\varphi(m+2)-\varphi(m+1)\bigr)
\le(m+1)\bigl(\varphi(m+1)-\varphi(m)\bigr).
\]
Moreover $\varphi(m+1)-\varphi(m)\ge0$ by \eqref{eq:propk-increment}, and
$(k+1)(m+1)\le(k-1)(m+3)$, because $m\le k-2$. Hence
\[
(k+1)\bigl(\varphi(m+2)-\varphi(m+1)\bigr)
\le(k-1)\bigl(\varphi(m+1)-\varphi(m)\bigr),
\]
which is the inequality $h_{m+1}\le h_m$. Consequently $h_n\ge h_{k-1}$.

It remains to compare $h_{k-1}-2\,\varphi(k)$ with $\chi_k$. We have
$h_{k-1}-2\,\varphi(k)=(k-1)\bigl(\varphi(k)-\varphi(k-1)\bigr)$ and
$\chi_k=(k+1)\log(1+1/k)-(k-1)\log(1+1/(k-1))$, and inserting the
definition of $\varphi$ we find
\begin{multline*}
h_{k-1}-2\,\varphi(k)-\chi_k\\
=(k-3)\Bigl(k\log\Bigl(1+\frac1k\Bigr)
-(k-1)\log\Bigl(1+\frac1{k-1}\Bigr)\Bigr).
\end{multline*}
This is nonnegative, because $k\ge3$ and because $(1+1/i)^{i}$ does not
decrease in $i\ge1$, as is recalled in Section~\ref{sec:prelim}. This
proves \eqref{eq:propk-exponent}.
\end{proof}

For the last triple, $(\omega_{k-1},\omega_k,\omega_{k+1})$, the next lemma
bounds from below the amount by which the inequality of property K holds;
the bound is needed in the proof of
Lemma~\ref{lem:directk}.

\begin{lemma}\label{lem:topsurplus}
Let $k\ge3$ be an integer, and let $\omega_n$ be the tilted Cayley masses
at the critical tilt $x_k$ of column $k$. Then for every real number $\mu$,
\[
\frac{k^2-2k-1}{2k^2(k+1)(k+2)}\;e^{\mu/k}\,\omega_k
\le e^{\mu}\,\omega_{k+1}+e^{-\mu}\,\omega_{k-1}-2\,e^{\mu/k}\,\omega_k .
\]
\end{lemma}

\begin{proof}
Let $L=\varphi(k+1)-\varphi(k)$. By \eqref{eq:propk-ratio} and
\eqref{eq:propk-flat} we have $\omega_{k+1}=\omega_k\,e^{L}$ and
$\omega_{k-1}=\omega_k$, so that
\begin{equation}\label{eq:propk-surplus}
\begin{split}
e^{\mu}\,\omega_{k+1}+e^{-\mu}\,\omega_{k-1}&-2\,e^{\mu/k}\,\omega_k\\
&=\omega_k\,\bigl(e^{\mu+L}+e^{-\mu}-2\,e^{\mu/k}\bigr).
\end{split}
\end{equation}
The inequality \eqref{eq:propk-mean}, which was shown in the proof of
Lemma~\ref{lem:massesmiddle}, with $M=0$ gives
\[
e^{\mu+L}+e^{-\mu}\ge
2\,e^{\mu/k}\exp\Bigl(\frac{(k+1)L-\chi_k}{2k}\Bigr),
\]
and we bound the number $(k+1)L-\chi_k$ from below. By the definitions of
$\varphi$ and of $\chi_k$,
\begin{multline*}
(k+1)L-\chi_k=(k-1)\log\Bigl(1+\frac1{k-1}\Bigr)
-k\,(k+1)\log\Bigl(1+\frac1k\Bigr)\\
+k\,(k+1)\log\Bigl(1+\frac1{k+1}\Bigr).
\end{multline*}
By \eqref{eq:propk-integral} this is the integral over $[0,1]$ of
\begin{multline*}
\frac{k-1}{k-1+t}-\frac{k\,(k+1)}{k+t}+\frac{k\,(k+1)}{k+1+t}\\
=\frac{t\,\bigl(k^2-2k-1+(k-1)\,t\bigr)}{(k-1+t)(k+t)(k+1+t)} .
\end{multline*}
For $0\le t\le1$ the numerator is at least $(k^2-2k-1)\,t$, where
$k^2-2k-1>0$ because $k\ge3$, and the denominator is positive and at most
$k\,(k+1)(k+2)$. Since the integral of $t$ over $[0,1]$ is $1/2$,
\[
(k+1)L-\chi_k\ge\frac{k^2-2k-1}{2k\,(k+1)(k+2)} .
\]
With $e^{b}\ge1+b$ for $b=((k+1)L-\chi_k)/(2k)$ we conclude that
\begin{align*}
e^{\mu+L}+e^{-\mu}-2\,e^{\mu/k}
&\ge e^{\mu/k}\,\frac{(k+1)L-\chi_k}{k}\\
&\ge\frac{k^2-2k-1}{2k^2(k+1)(k+2)}\;e^{\mu/k} ,
\end{align*}
and the lemma follows from \eqref{eq:propk-surplus}, because $\omega_k>0$.
\end{proof}

For the base of the induction on the offset only the sign of this amount is
needed.

\begin{corollary}\label{cor:massestop}
Let $k\ge3$ be an integer, and let $\omega_n$ be the tilted Cayley masses
at the critical tilt $x_k$ of column $k$. Then the triple
$(\omega_{k-1},\omega_k,\omega_{k+1})$ has property K at $k$.
\end{corollary}

\begin{proof}
This is immediate from Lemma~\ref{lem:topsurplus}, since the left-hand side
of its inequality is nonnegative for $k\ge3$.
\end{proof}

\subsection{Self-convolutions and the pair sums of the masses}
\label{subsec:propk-convolution}

The direct term $S_\rho(y)$ of Section~\ref{subsec:propk-triples} contains
the number $Q_\rho(y-\rho-1)$, one of the numbers $Q_\rho(j)$ that
\eqref{eq:propk-pairsums} defines for integers $\rho\ge0$ and $j\ge0$.
We call $Q_\rho(j)$ a \emph{pair sum of the masses}, of length $j$.
In this section a pair sum is always one of
the numbers $Q_\rho(j)$, and not the pair sum $\mathcal P(k,R)$ of
Definition~\ref{def:markedpairs}.
When $y$ runs through $k-1$, $k$, $k+1$ and
$\rho$ is fixed, the length $j=y-\rho-1$ runs through three consecutive
values $n$, $n+1$, $n+2$ with $\rho+n+2=k$, and the largest index of a mass
that occurs in these three pair sums of the masses is $k$. The aim of this
subsection is to show that the three pair sums of the masses form a triple
with property K at $k$ (Corollary~\ref{cor:pairsumsk}).

The pair sum of the masses $Q_\rho(j)$ is a self-convolution. For a sequence
$(v_i)_{i\ge0}$ of real numbers and an integer $j\ge0$, the
\emph{self-convolution} of $v$ at length $j$ is the sum of the products
$v_i\,v_{j-i}$ over $0\le i\le j$, and $Q_\rho(j)$ is the self-convolution
at length $j$ of the shifted sequence $v_i=\omega_{\rho+i}$, whose last two
terms that occur are $v_{n+1}=\omega_{k-1}$ and $v_{n+2}=\omega_k$. The
argument uses of this sequence only what
Section~\ref{subsec:propk-masses} has provided, and the next lemma is
stated for an arbitrary sequence with these properties.

The proof of the lemma sorts the combination of the three
self-convolutions that property K asks about by one of the two indices of
the products: the inner terms are then nonnegative,
because three consecutive terms of $v$ satisfy the inequality of property
K at $k$, but end terms remain, which may be negative. For
$\cosh\mu\ge e^{\mu/k}$ the proof replaces each value of $v$ in the inner
terms by the smallest one, $v_{n+1}$; for the remaining parameters it
compares a positive
multiple of the sequence $(e^{\mu i}\,v_i)$ with the solution of the
three-term recurrence $\psi(t+1)+\psi(t-1)=2\,e^{\mu/k}\,\psi(t)$ that
agrees with this multiple at indices $n+1$ and $n+2$. This solution
cancels the end terms,
and what is left is an explicit expression in the two values $v_{n+1}$ and
$v_{n+2}$, whose sign is decided by the hypothesis $n+2\le k$.

\begin{lemma}\label{lem:convolutionk}
Let $k$ and $n$ be integers with $n\ge0$ and $n+2\le k$, and let
$(v_i)_{i\ge0}$ be a sequence of real numbers with the following four
properties.
\begin{enumerate}
\item One has $v_i>0$ for every $i$ with $0\le i\le n+2$.
\item One has $v_{n+1}\le v_i$ for every $i$ with $0\le i\le n$.
\item One has $v_{n+2}=v_{n+1}$.
\item For every real number $\mu\ge0$ and every $j$ with $0\le j\le n$,
\[
2\,e^{\mu/k}\,v_{j+1}\le e^{\mu}\,v_{j+2}+e^{-\mu}\,v_j .
\]
\end{enumerate}
Then the self-convolutions
\[
\sum_{i=0}^{j}v_i\,v_{j-i}\qquad(j=n,\ n+1,\ n+2)
\]
of $v$ form, in this order, a triple with property K at $k$.
\end{lemma}

\begin{proof}
Fix a real number $\mu\ge0$, and for an integer $j\ge0$ let $\Sigma_j$
denote the self-convolution of $v$ at length $j$. We have to show that
the number
\[
I=e^{\mu}\,\Sigma_{n+2}+e^{-\mu}\,\Sigma_n-2\,e^{\mu/k}\,\Sigma_{n+1}
\]
is nonnegative. In the three sums we collect the terms according to the
factor $v_{n-j}$ with $0\le j\le n$:
\begin{equation}\label{eq:propk-sorted}
\begin{split}
\Sigma_n&=\sum_{j=0}^{n}v_{n-j}\,v_j,\\
\Sigma_{n+1}&=\sum_{j=0}^{n}v_{n-j}\,v_{j+1}+v_{n+1}\,v_0,\\
\Sigma_{n+2}&=\sum_{j=0}^{n}v_{n-j}\,v_{j+2}+v_{n+1}\,v_1+v_{n+2}\,v_0 .
\end{split}
\end{equation}
These identities hold for every sequence of real numbers. We distinguish two
cases, according to whether $\cosh\mu\ge e^{\mu/k}$ or
$\cosh\mu<e^{\mu/k}$.

Suppose first that $\cosh\mu\ge e^{\mu/k}$. By \eqref{eq:propk-sorted},
\begin{align*}
I={}&\sum_{j=0}^{n}v_{n-j}\,
\bigl(e^{\mu}\,v_{j+2}+e^{-\mu}\,v_j-2\,e^{\mu/k}\,v_{j+1}\bigr)\\
&+v_{n+1}\,\bigl(e^{\mu}\,v_1-2\,e^{\mu/k}\,v_0\bigr)+e^{\mu}\,v_{n+2}\,v_0 .
\end{align*}
By hypothesis~(4) each expression in parentheses under the sum is
nonnegative, and by hypothesis~(2) we have $v_{n-j}\ge v_{n+1}$. So the sum
does not increase when each factor $v_{n-j}$ is replaced by $v_{n+1}$. Let
$\sigma=v_0+v_1+\dots+v_{n+1}$. Then
\begin{gather*}
\sum_{j=0}^{n}v_j=\sigma-v_{n+1},\qquad
\sum_{j=0}^{n}v_{j+1}=\sigma-v_0,\\
\sum_{j=0}^{n}v_{j+2}=\sigma-v_0-v_1+v_{n+2} .
\end{gather*}
Inserting these three sums and using $v_{n+2}=v_{n+1}$, which is
hypothesis~(3), we obtain after cancellation
\[
I\ge2\,v_{n+1}\,\Bigl(\bigl(\cosh\mu-e^{\mu/k}\bigr)\,\sigma
+\sinh\mu\cdot v_{n+1}\Bigr).
\]
The right-hand side is nonnegative, because $v_{n+1}$ and $\sigma$ are
positive by hypothesis~(1) and $\cosh\mu\ge e^{\mu/k}$ in the present case.

Suppose now that $\cosh\mu<e^{\mu/k}$. Then $\mu>0$, since both sides are
equal to $1$ for $\mu=0$. As $e^{\mu/k}>1$, there is a real number $\tau>0$
with $\cosh\tau=e^{\mu/k}$, and $\cosh\mu<\cosh\tau$ gives $\mu<\tau$. We
pass to the sequence $\tilde v_i=e^{\mu i}\,v_i$ and write
$\widetilde{\Sigma}_j$ for its self-convolution at length $j$, so that
$\widetilde{\Sigma}_j=e^{\mu j}\,\Sigma_j$. Hypothesis~(4) for the parameter
$\mu$, multiplied by $e^{\mu(j+1)}$, says that the numbers
\[
w_j=\tilde v_{j+2}+\tilde v_j-2\cosh\tau\cdot\tilde v_{j+1}
\qquad(0\le j\le n)
\]
are nonnegative, and the number
\[
J=\widetilde{\Sigma}_{n+2}+\widetilde{\Sigma}_n
-2\cosh\tau\cdot\widetilde{\Sigma}_{n+1}
\]
is $e^{\mu(n+1)}\,I$, so that we have to show that $J\ge0$. By
\eqref{eq:propk-sorted} for the sequence $\tilde v$,
\begin{equation}\label{eq:propk-tilted}
\begin{split}
J={}&\sum_{j=0}^{n}\tilde v_{n-j}\,w_j\\
&+\tilde v_{n+1}\,\bigl(\tilde v_1-2\cosh\tau\cdot\tilde v_0\bigr)
+\tilde v_{n+2}\,\tilde v_0 .
\end{split}
\end{equation}

For a real number $t$ let
\[
\psi(t)=\sinh\bigl((n+2-t)\,\tau\bigr)\,\tilde v_{n+1}
-\sinh\bigl((n+1-t)\,\tau\bigr)\,\tilde v_{n+2} .
\]
By the addition formula
$\sinh(y+\tau)+\sinh(y-\tau)=2\cosh\tau\cdot\sinh y$,
\begin{equation}\label{eq:propk-recurrence}
\psi(t+1)+\psi(t-1)=2\cosh\tau\cdot\psi(t)
\end{equation}
for every real number $t$, and from the definition,
$\psi(n+1)=\sinh\tau\cdot\tilde v_{n+1}$ and
$\psi(n+2)=\sinh\tau\cdot\tilde v_{n+2}$. For $0\le j\le n+2$ let
\[
r_j=\sinh\tau\cdot\tilde v_j-\psi(j) .
\]
Then $r_{n+1}=r_{n+2}=0$, and by \eqref{eq:propk-recurrence},
\[
r_{j+2}+r_j-2\cosh\tau\cdot r_{j+1}=\sinh\tau\cdot w_j\ge0
\qquad(0\le j\le n).
\]
We claim that $0\le r_{j+1}\le r_j$ for $0\le j\le n+1$, and prove this by
descending induction on $j$. For $j=n+1$ both numbers are $0$. If
$j\le n$ and $0\le r_{j+2}\le r_{j+1}$, then
\[
r_j\ge2\cosh\tau\cdot r_{j+1}-r_{j+2}\ge2\,r_{j+1}-r_{j+2}\ge r_{j+1},
\]
where we used $\cosh\tau\ge1$ and $r_{j+1}\ge0$ in the second step and
$r_{j+2}\le r_{j+1}$ in the third. In particular $r_j\ge0$ for
$0\le j\le n$.

We multiply \eqref{eq:propk-tilted} by $\sinh\tau>0$ and write
$\sinh\tau\cdot\tilde v_{n-j}=r_{n-j}+\psi(n-j)$ in the sum. The products
$r_{n-j}\,w_j$ are nonnegative. For the other part let
\[
Y_j=\psi(n+1-j)\,\tilde v_{j+1}-\psi(n-j)\,\tilde v_j
\qquad(0\le j\le n+1).
\]
By \eqref{eq:propk-recurrence} at $t=n-j$, we have for $0\le j\le n$
\begin{align*}
Y_{j+1}-Y_j&=\psi(n-j)\,\bigl(\tilde v_{j+2}+\tilde v_j\bigr)\\
&\quad-\bigl(\psi(n-j-1)+\psi(n-j+1)\bigr)\,\tilde v_{j+1}\\
&=\psi(n-j)\,w_j ,
\end{align*}
so that the sum of $\psi(n-j)\,w_j$ over $0\le j\le n$ is
$Y_{n+1}-Y_0$. Now \eqref{eq:propk-recurrence} at $t=n+1$ gives
\[
\psi(n)=2\cosh\tau\cdot\psi(n+1)-\psi(n+2)
=\sinh\tau\,\bigl(2\cosh\tau\cdot\tilde v_{n+1}-\tilde v_{n+2}\bigr),
\]
and hence
\begin{align*}
Y_0&=\psi(n+1)\,\tilde v_1-\psi(n)\,\tilde v_0\\
&=\sinh\tau\,\Bigl(\tilde v_{n+1}\,
\bigl(\tilde v_1-2\cosh\tau\cdot\tilde v_0\bigr)
+\tilde v_{n+2}\,\tilde v_0\Bigr),
\end{align*}
which is $\sinh\tau$ times the two end terms of
\eqref{eq:propk-tilted}. These terms therefore cancel against $-Y_0$, and
\[
\sinh\tau\cdot J=\sum_{j=0}^{n}r_{n-j}\,w_j+Y_{n+1}\ge Y_{n+1} .
\]

It remains to show that $Y_{n+1}\ge0$. Reading $\psi(0)$ and $\psi(-1)$
from the definition of $\psi$, we find
\begin{align*}
Y_{n+1}&=\psi(0)\,\tilde v_{n+2}-\psi(-1)\,\tilde v_{n+1}\\
&=2\,\tilde v_{n+1}\,\tilde v_{n+2}\,\sinh\bigl((n+2)\,\tau\bigr)
-\tilde v_{n+2}^{\,2}\,\sinh\bigl((n+1)\,\tau\bigr)\\
&\quad-\tilde v_{n+1}^{\,2}\,\sinh\bigl((n+3)\,\tau\bigr).
\end{align*}
By hypothesis~(3) we have $\tilde v_{n+2}=e^{\mu}\,\tilde v_{n+1}$. With
$m=n+1$ this gives $Y_{n+1}=\tilde v_{n+1}^{\,2}\,A$, where
\[
A=2\,e^{\mu}\sinh\bigl((m+1)\,\tau\bigr)-e^{2\mu}\sinh(m\tau)
-\sinh\bigl((m+2)\,\tau\bigr).
\]
Put $h=m\tau/2$. Then
\[
A=2\,\bigl(\sinh(h+\tau)-e^{\mu}\sinh h\bigr)\,
\bigl(e^{\mu}\cosh h-\cosh(h+\tau)\bigr);
\]
indeed, by the addition formula for $\sinh$ and the formula
$2\sinh y\cosh y=\sinh(2y)$ the product on the right is
$2\,e^{\mu}\sinh(2h+\tau)-\sinh(2h+2\tau)-e^{2\mu}\sinh(2h)$. Both factors
are nonnegative. For the first, we have $e^{\mu}\le e^{\tau}
=\cosh\tau+\sinh\tau$, because $\mu\le\tau$, and $0\le\sinh h\le\cosh h$, so
that
\[
e^{\mu}\sinh h\le\sinh h\cosh\tau+\cosh h\sinh\tau=\sinh(h+\tau).
\]
For the second,
\begin{align*}
\cosh(h+\tau)&=\cosh h\cosh\tau+\sinh h\sinh\tau\\
&\le\cosh h\cosh\tau\,\Bigl(1+\frac{m\,\tau^2}{2}\Bigr)\\
&\le\cosh h\cosh\tau\,\Bigl(1+\frac{\tau^2}{2}\Bigr)^{m}
\le\cosh h\,(\cosh\tau)^{m+1}\\
&\le\cosh h\,(\cosh\tau)^{k}=e^{\mu}\cosh h .
\end{align*}
Here the first inequality uses the elementary inequality
$\sinh y\le y\cosh y$ for $y\ge0$, applied to $y=h$ and to $y=\tau$,
together with $h\,\tau=m\,\tau^2/2$. The second is Bernoulli's inequality,
and the third is $1+\tau^2/2\le\cosh\tau$. The fourth holds because
$m+1=n+2\le k$ and $\cosh\tau\ge1$, and the last equality is the choice of
$\tau$. Hence $A\ge0$ and $Y_{n+1}\ge0$, which completes the proof.
\end{proof}

Applied to the shifted sequence of the masses, the lemma gives the result on
the pair sums of the masses.

\begin{corollary}\label{cor:pairsumsk}
Let $k\ge3$, $\rho\ge0$ and $n\ge0$ be integers with $\rho+n+2=k$. Then the
pair sums $Q_\rho(n)$, $Q_\rho(n+1)$ and $Q_\rho(n+2)$ of the tilted Cayley
masses at the critical tilt $x_k$ of column $k$ form, in this order, a
triple with property K at $k$.
\end{corollary}

\begin{proof}
We apply Lemma~\ref{lem:convolutionk} to the integers $k$ and $n$ and to
the sequence $v_i=\omega_{\rho+i}$, whose self-convolution at length $j$ is
$Q_\rho(j)$. The masses are positive, which is hypothesis~(1).
Hypothesis~(2) follows from Lemma~\ref{lem:massesdecrease}, since the
indices $\rho+i$ with $0\le i\le n+1$ are at most $k-1$, and
hypothesis~(3) is \eqref{eq:propk-flat}. Finally, for $0\le j\le n$ the
inequality of~(4) is that of Lemma~\ref{lem:massesmiddle} at index
$\rho+j$, which is admissible because $\rho+j+2\le k$.
\end{proof}

\subsection{The direct term}\label{subsec:propk-direct}

Recall that for integers $\rho$ and $y$ with $1\le\rho\le y-1$ the direct
term is
\[
S_\rho(y)=x_k\,\omega_{\rho-1}\,\omega_y
+\frac{x_k}{2}\cdot\frac{y+\rho+1}{y+\rho}\,Q_\rho(y-\rho-1) .
\]
In the induction on the offset it is the part of the recursion
that the induction hypothesis does not cover, and
its property K has to be proved directly.

Both ingredients of $S_\rho(y)$ have property K by themselves when the
threshold runs through $k-1$, $k$, $k+1$: the masses $\omega_y$ by
Corollary~\ref{cor:massestop}, and the pair sums of the masses
$Q_\rho(y-\rho-1)$ by Corollary~\ref{cor:pairsumsk}. But the pair sum of
the masses is multiplied by the factor
$\theta(y)=(y+\rho+1)/(y+\rho)$, which decreases as the threshold grows,
whereas the stability statements of
Section~\ref{subsec:propk-definition} require factors that do not decrease.
A decreasing factor is the case $L<M$ of \eqref{eq:propk-shift}, which is
available only for parameters $\mu\ge(M-L)/2$. The proof of the next lemma
therefore compares two amounts: the amount by which the triple
$(\omega_{k-1},\omega_k,\omega_{k+1})$ exceeds the inequality of property
K, which it bounds from below, and the amount by which the triple of the
terms $\theta(y)\,Q_\rho(y-\rho-1)$ may fail it, which it bounds from
above. It divides the parameters
at $\mu=1/(k(k+2))$: above this value the downward shift caused by the
factor stays within the nonnegative parameters, so that
\eqref{eq:propk-shift} applies, and below it the second amount is bounded
through the two end terms of the longest of the three pair sums of the
masses. In
both ranges the two bounds combine to a
nonnegative multiple of
$\omega_{\rho-1}-\omega_\rho$.

\begin{lemma}\label{lem:directk}
Let $k\ge7$ and $\rho$ be integers with $1\le\rho\le k-2$. Then the direct
terms $S_\rho(k-1)$, $S_\rho(k)$ and $S_\rho(k+1)$ at the critical tilt
$x_k$ of column $k$ form, in this order, a triple with property K at $k$.
\end{lemma}

\begin{proof}
Fix a real number $\mu\ge0$, and put $N=k+\rho$, so that
$k+1\le N\le2k-2$. We write
\[
\theta_-=\frac{N}{N-1},\qquad\theta_0=\frac{N+1}{N},\qquad
\theta_+=\frac{N+2}{N+1}
\]
for the values of the factor $\theta(y)=(y+\rho+1)/(y+\rho)$ at
$y=k-1,k,k+1$, and
\[
Q^{-}=Q_\rho(k-\rho-2),\qquad Q^{0}=Q_\rho(k-\rho-1),\qquad
Q^{+}=Q_\rho(k-\rho)
\]
for the three pair sums of the masses that occur in the direct terms. These
are positive, as sums of products of positive masses, and by
Corollary~\ref{cor:pairsumsk}, applied with $n=k-\rho-2$, the triple
$(Q^{-},Q^{0},Q^{+})$ has property K at $k$. Inserting the definition of
the direct terms $S_\rho(y)$ at $y=k-1$, $k$ and $k+1$, we find
\begin{equation}\label{eq:propk-surdef}
\begin{split}
e^{\mu}\,S_\rho(k+1)+e^{-\mu}\,S_\rho(k-1)&-2\,e^{\mu/k}\,S_\rho(k)\\
&=x_k\,\omega_{\rho-1}\,\mathrm{Sur}+\frac{x_k}{2}\,\mathrm{Def},
\end{split}
\end{equation}
where
\begin{align*}
\mathrm{Sur}&=e^{\mu}\,\omega_{k+1}+e^{-\mu}\,\omega_{k-1}
-2\,e^{\mu/k}\,\omega_k,\\
\mathrm{Def}&=e^{\mu}\,\theta_+\,Q^{+}+e^{-\mu}\,\theta_-\,Q^{-}
-2\,e^{\mu/k}\,\theta_0\,Q^{0}.
\end{align*}
The number $\mathrm{Sur}$ is the amount by which the triple
$(\omega_{k-1},\omega_k,\omega_{k+1})$ exceeds the inequality of property K
at the parameter $\mu$, and $\mathrm{Def}$, which may be negative, is the
corresponding amount for the pair sums of the masses with their factors.
The left-hand
side of \eqref{eq:propk-surdef} is this amount for the triple of the lemma,
and we have to show that it is nonnegative. We show that there is a real
number $\lambda\ge0$ such that
\begin{equation}\label{eq:propk-balance}
\mathrm{Sur}\ge\lambda\,\omega_k\qquad\text{and}\qquad
\mathrm{Def}\ge-2\,\lambda\,\omega_\rho\,\omega_k .
\end{equation}
This proves the lemma. Indeed, by \eqref{eq:propk-balance} the right-hand
side of \eqref{eq:propk-surdef} is at least
$x_k\,\lambda\,\omega_k\,(\omega_{\rho-1}-\omega_\rho)$, and this is
nonnegative, because $\omega_\rho\le\omega_{\rho-1}$ by
Lemma~\ref{lem:massesdecrease}, as $\rho\le k$. We distinguish two cases,
according to whether $\mu\le1/(k(k+2))$ or $\mu\ge1/(k(k+2))$.

Suppose first that $\mu\le1/(k(k+2))$. We take $\lambda=e^{\mu}/(N(N+1))$.
The differences of the factors are
\[
\theta_0-\theta_+=\frac{1}{N(N+1)},\qquad
\theta_--\theta_0=\frac{1}{N(N-1)}\ge\frac{1}{N(N+1)} ,
\]
and we write
\begin{multline*}
\mathrm{Def}=\theta_0\,\bigl(e^{\mu}\,Q^{+}+e^{-\mu}\,Q^{-}
-2\,e^{\mu/k}\,Q^{0}\bigr)\\
-(\theta_0-\theta_+)\,e^{\mu}\,Q^{+}+(\theta_--\theta_0)\,e^{-\mu}\,Q^{-}.
\end{multline*}
The first term is nonnegative by property K of the pair sums of the masses,
and in the
last term $\theta_--\theta_0$ can be replaced by the smaller number
$\theta_0-\theta_+$. Hence
\begin{equation}\label{eq:propk-defect}
\mathrm{Def}\ge-\frac{1}{N(N+1)}\,
\bigl(e^{\mu}\,Q^{+}-e^{-\mu}\,Q^{-}\bigr).
\end{equation}
We bound the difference in parentheses by the two end terms of $Q^{+}$. Let
$n=k-\rho-2$, so that $Q^{-}=Q_\rho(n)$ and $Q^{+}=Q_\rho(n+2)$. The sum
$Q^{+}$ consists of the products $\omega_{\rho+i}\,\omega_{\rho+n+2-i}$
with $0\le i\le n+2$. Its terms $i=0$ and $i=n+2$ are both equal to
$\omega_\rho\,\omega_k$, because $\rho+n+2=k$, and its other terms are the
products $\omega_{\rho+i+1}\,\omega_{\rho+n-i+1}$ with $0\le i\le n$.
Since $\mu\le1/(k(k+2))$, we have for every integer $j\ge0$ with $j+2\le k$
\begin{equation}\label{eq:propk-stepdown}
e^{\mu}\,\omega_{j+1}\le\omega_j .
\end{equation}
Indeed, by \eqref{eq:propk-ratio} we have
$\omega_j=\omega_{j+1}\,e^{\varphi(k)-\varphi(j+1)}$, so that it suffices to
show that $\mu\le\varphi(k)-\varphi(j+1)$. Since $j+1\le k-1$ and $\varphi$
does not decrease on the positive integers,
$\varphi(k)-\varphi(j+1)\ge\varphi(k)-\varphi(k-1)$. By
\eqref{eq:propk-increment}, which was shown in the proof of
Lemma~\ref{lem:massesmiddle}, with $y=k$ this is at least
$3/(2k(k+1))\ge1/(k(k+2))$. For
such $i$ the indices $\rho+i$ and $\rho+n-i$ are at most $k-2$, so that
\eqref{eq:propk-stepdown} applies to both of them:
$e^{\mu}\,\omega_{\rho+i+1}\le\omega_{\rho+i}$ and
$e^{\mu}\,\omega_{\rho+n-i+1}\le\omega_{\rho+n-i}$. Multiplying these two
inequalities and then multiplying by $e^{-\mu}$, we get
\[
e^{\mu}\,\omega_{\rho+i+1}\,\omega_{\rho+n-i+1}
\le e^{-\mu}\,\omega_{\rho+i}\,\omega_{\rho+n-i} .
\]
The right-hand side is $e^{-\mu}$ times the term of index $i$ of $Q^{-}$.
Summing over $0\le i\le n$, we obtain
$e^{\mu}\,(Q^{+}-2\,\omega_\rho\,\omega_k)\le e^{-\mu}\,Q^{-}$, and hence
\[
e^{\mu}\,Q^{+}-e^{-\mu}\,Q^{-}\le2\,e^{\mu}\,\omega_\rho\,\omega_k .
\]
With \eqref{eq:propk-defect} this is the second inequality of
\eqref{eq:propk-balance}.

For the first inequality, \eqref{eq:propk-ratio} and \eqref{eq:propk-flat}
give $\omega_{k+1}=\omega_k\,e^{\varphi(k+1)-\varphi(k)}$ and
$\omega_{k-1}=\omega_k$, so that
\[
\mathrm{Sur}=\omega_k\,\bigl(e^{\mu+\varphi(k+1)-\varphi(k)}+e^{-\mu}
-2\,e^{\mu/k}\bigr).
\]
Since $e^{t}\ge1+t$ for every real number $t$, the sum of the first two
exponentials is at least $2+\varphi(k+1)-\varphi(k)$, and
$\varphi(k+1)-\varphi(k)\ge3/(2(k+1)(k+2))$ by \eqref{eq:propk-increment}
at the point $k+1$. Hence
\[
\mathrm{Sur}\ge\omega_k\,\Bigl(\frac{3}{2\,(k+1)(k+2)}
-2\,\bigl(e^{\mu/k}-1\bigr)\Bigr),
\]
and it remains to see that the number in parentheses is at least
$\lambda$. Since $N\ge k+1$, we have $\lambda\le e^{\mu}/((k+1)(k+2))$. As
$k\ge7$, the parameter satisfies $\mu\le1/63$ and
$\mu/k\le1/(k^2(k+2))\le1/441$. By the inequality $e^{t}\le1/(1-t)$ for
$0\le t<1$ this gives $e^{\mu}\le63/62$ and
\begin{align*}
e^{\mu/k}-1\le\frac{\mu/k}{1-\mu/k}
&\le\frac{441}{440}\cdot\frac{1}{k^2(k+2)}\\
&\le\frac{441}{440}\cdot\frac{8}{49}\cdot\frac{1}{(k+1)(k+2)} ,
\end{align*}
where the last step holds because $8k^2-49\,(k+1)=(k-7)(8k+7)\ge0$.
Therefore the number in parentheses is at least
\[
\Bigl(\frac32-2\cdot\frac{441}{440}\cdot\frac{8}{49}\Bigr)
\frac{1}{(k+1)(k+2)} ,
\]
and this is at least $\lambda$, because
$3/2-2\cdot(441/440)\cdot(8/49)\ge63/62\ge e^{\mu}$.

Suppose now that $\mu\ge1/(k(k+2))$. We take
\[
\lambda=e^{\mu/k}\,\theta_0\,\frac{(2k-N)(N-k+1)}{N^2\,(N+2)\,k} ,
\]
which is positive, because $k+1\le N\le2k-2$. Let
$L=\log(\theta_+/\theta_0)$ and $M=\log(\theta_-/\theta_0)$. We have
$\theta_+/\theta_0=N(N+2)/(N+1)^2$ and $\theta_-/\theta_0=N^2/(N^2-1)$, and
the inequalities $1-1/r\le\log r\le r-1$ for $r>0$ give
\[
L\ge-\frac{1}{N(N+2)},\qquad
\frac{1}{N^2}\le M\le\frac{1}{(N-1)(N+1)} .
\]
Hence
\begin{align*}
\frac{M-L}{2}
&\le\frac12\,\Bigl(\frac{1}{(N-1)(N+1)}+\frac{1}{N(N+2)}\Bigr)\\
&\le\frac{1}{(N-1)(N+1)}\le\frac{1}{k\,(k+2)}\le\mu ,
\end{align*}
where the third step uses $N\ge k+1$. So $\mu+(L-M)/2\ge0$, and
\eqref{eq:propk-shift}, applied to the triple $(Q^{-},Q^{0},Q^{+})$, gives
\[
2\,e^{\mu/k}\,e^{E/(2k)}\,Q^{0}\le e^{\mu}\,\frac{\theta_+}{\theta_0}\,
Q^{+}+e^{-\mu}\,\frac{\theta_-}{\theta_0}\,Q^{-}
\]
with $E=(k+1)L+(k-1)M$. Multiplying by $\theta_0$ and subtracting
$2\,e^{\mu/k}\,\theta_0\,Q^{0}$ on both sides, we get
$\mathrm{Def}\ge2\,e^{\mu/k}\,\theta_0\,Q^{0}\,(e^{E/(2k)}-1)$. By the
bounds for $L$ and $M$,
\[
E\ge-\frac{k+1}{N(N+2)}+\frac{k-1}{N^2}
=-\frac{2\,(N-k+1)}{N^2\,(N+2)} ,
\]
so that $e^{E/(2k)}-1\ge E/(2k)\ge-(N-k+1)/(N^2(N+2)\,k)$. As
$Q^{0}>0$, this gives
\[
\mathrm{Def}\ge-2\,e^{\mu/k}\,\theta_0\,
\frac{N-k+1}{N^2\,(N+2)\,k}\,Q^{0} .
\]
Finally, the sum $Q^{0}$ consists of the $k-\rho$ products
$\omega_{\rho+i}\,\omega_{k-1-i}$ with $0\le i\le k-\rho-1$, and by
\eqref{eq:propk-logconvex} with $m=k-\rho-1$ each of them is at most
$\omega_\rho\,\omega_{k-1}$, which is equal to $\omega_\rho\,\omega_k$
because $\omega_{k-1}=\omega_k$ by \eqref{eq:propk-flat}. So
$Q^{0}\le(k-\rho)\,\omega_\rho\,\omega_k=(2k-N)\,\omega_\rho\,\omega_k$,
and the second inequality of \eqref{eq:propk-balance} follows.

For the first inequality, Lemma~\ref{lem:topsurplus} gives
\[
\mathrm{Sur}\ge\frac{k^2-2k-1}{2k^2(k+1)(k+2)}\;e^{\mu/k}\,\omega_k .
\]
So it suffices to show that $\lambda$ is at most the coefficient of
$\omega_k$ on the right. After division by $e^{\mu/k}$ and
multiplication by $k$, and with $\theta_0=(N+1)/N$, this is the inequality
\[
\frac{(N+1)(2k-N)(N-k+1)}{N^3\,(N+2)}
\le\frac{k^2-2k-1}{2k\,(k+1)(k+2)} .
\]
The numbers $2k-N$ and $N-k+1$ are positive and add up to $k+1$, so that
their product is at most $(k+1)^2/4$; and
$(N+1)(k+1)^3\le(N+2)\,N^3$, because $k+1\le N$. Hence the left-hand side
is at most $1/(4(k+1))$. And the inequality
$1/(4(k+1))\le(k^2-2k-1)/(2k(k+1)(k+2))$ is equivalent to
$k\,(k+2)\le2\,(k^2-2k-1)$, or $k^2-6k-2\ge0$, which holds because
$k\ge7$ (and fails for $k=6$).
\end{proof}

\subsection{The recursion in the first layer and the diagonal triples}
\label{subsec:propk-diagonal}

We now sort the terms of the tilted last-car recursion
\eqref{eq:counts-tilted} by the sizes of the two tilted counts in each
term. For integers $y\ge1$ and $p$ with $0\le p\le y$ let
\[
\mathsf{D}(y,p)=\sum_{i=0}^{y-p}\min(p+i+1,\,y)\;\omega_{p+i}\;\omega_{y-i},
\]
the \emph{direct part} of the recursion at size $y+p+1$ and threshold $y$.
The next lemma separates it from the terms of the recursion that contain a
tilted count of a size larger than the threshold.

\begin{lemma}\label{lem:layersplit}
Let $k\ge1$, $y\ge1$ and $p\ge0$ be integers with $p\le y$. Then the tilted
normalised counts $u^{(y)}_N$ and the tilted Cayley masses $\omega_n$ at
the critical tilt $x_k$ of column $k$ satisfy
\[
(y+p+1)\,u^{(y)}_{y+p+1}
=x_k\,\Bigl(\sum_{i=0}^{p-1}(y+i+1)\,\omega_i\,u^{(y)}_{y+p-i}
+\mathsf{D}(y,p)\Bigr),
\]
where $\mathsf{D}(y,p)$ is the direct part of the recursion at size $y+p+1$
and threshold $y$.
\end{lemma}

\begin{proof}
The recursion \eqref{eq:counts-tilted}, for threshold $y$, the tilt $x_k$ and
$N=y+p$, reads
\[
(y+p+1)\,u^{(y)}_{y+p+1}
=x_k\sum_{j=0}^{y+p}\min(j+1,y)\,u^{(y)}_j\,u^{(y)}_{y+p-j} .
\]
We split the range of $j$ into the three parts $0\le j<p$, $p\le j\le y$
and $y<j\le y+p$. For $0\le i<p$ we have $i<y$, because $p\le y$. By
\eqref{eq:propk-free}, the term of the first part with $j=i$ is therefore
$(i+1)\,\omega_i\,u^{(y)}_{y+p-i}$, and the term of the third part with
$j=y+p-i$ is $y\,\omega_i\,u^{(y)}_{y+p-i}$.
The two parts together give the first sum of the lemma. In the second part
we write $j=p+i$ with $0\le i\le y-p$. Both $p+i$ and $y+p-j=y-i$ are at
most $y$, so that both factors are Cayley masses by \eqref{eq:propk-free},
and the terms are those of $\mathsf{D}(y,p)$.
\end{proof}

The
next proposition uses the lemma for an induction on the offset, which is
carried out for the three thresholds $k-1$, $k$ and $k+1$ at once.

\begin{proposition}\label{prop:diagonalk}
Let $k\ge7$ and $\rho\ge0$ be integers with $\rho+2\le k$. Then the
diagonal triple of offset $\rho$ of column $k$,
\[
\bigl(u^{(k-1)}_{k-1+\rho},\ u^{(k)}_{k+\rho},\ u^{(k+1)}_{k+1+\rho}\bigr),
\]
has property K at $k$.
\end{proposition}

\begin{proof}
We first derive the recursion \eqref{eq:propk-recursion}. For integers
$p\ge0$ and $y$ with $p+2\le y$ the direct part is
\begin{equation}\label{eq:propk-direct}
\mathsf{D}(y,p)=(y+p+1)\,\omega_p\,\omega_y
+\frac{y+p+2}{2}\,Q_{p+1}(y-p-2) .
\end{equation}
Indeed, the terms $i=0$ and $i=y-p$ of the sum that defines
$\mathsf{D}(y,p)$ are $(p+1)\,\omega_p\,\omega_y$ and
$y\,\omega_y\,\omega_p$, and their sum is the first term of
\eqref{eq:propk-direct}. Let $m=y-p-2\ge0$. For $0\le j\le m$ the term
$i=j+1$ is $(p+2+j)\,\pi_j$, where
$\pi_j=\omega_{p+1+j}\,\omega_{p+1+m-j}$. Since $\pi_{m-j}=\pi_j$, the sum
of $(p+2+j)\,\pi_j$ over $0\le j\le m$ is equal to the sum of
$(p+2+m-j)\,\pi_j$, and hence to $(2p+4+m)/2$ times the sum of
the $\pi_j$. Here $2p+4+m=y+p+2$, and the sum of the $\pi_j$ is
$Q_{p+1}(m)$, which gives \eqref{eq:propk-direct}. Now let $y$ and $\rho$
be integers with $1\le\rho\le y-1$. We apply Lemma~\ref{lem:layersplit} and
\eqref{eq:propk-direct} with $p=\rho-1$, for which $p+2\le y$, and divide
by $y+\rho$. Since
\begin{align*}
\frac{x_k}{y+\rho}\,\mathsf{D}(y,\rho-1)
&=x_k\,\omega_{\rho-1}\,\omega_y
+\frac{x_k}{2}\cdot\frac{y+\rho+1}{y+\rho}\,Q_\rho(y-\rho-1)\\
&=S_\rho(y),
\end{align*}
the result is
\begin{equation}\label{eq:propk-recursion}
u^{(y)}_{y+\rho}=\sum_{i=0}^{\rho-2}x_k\,\omega_i\,\frac{y+i+1}{y+\rho}\,
u^{(y)}_{y+\rho-1-i}+S_\rho(y).
\end{equation}

We now prove the proposition by induction on $\rho$, assuming it for all
smaller offsets. For $\rho=0$ the sizes $k-1$, $k$, $k+1$ do not exceed
thresholds $k-1$, $k$, $k+1$, so that by \eqref{eq:propk-free} the
triple is $(\omega_{k-1},\omega_k,\omega_{k+1})$, which has property K at
$k$ by Corollary~\ref{cor:massestop}.

Let $1\le\rho\le k-2$. Then \eqref{eq:propk-recursion} holds for each of
the thresholds $y=k-1$, $k$ and $k+1$, because $\rho\le(k-1)-1$. The range
of $i$ and the numbers $x_k\,\omega_i$ in it are the same for the three
thresholds, so that the triple of the proposition is the sum of the
following triples. For each $i$ with $0\le i\le\rho-2$ there is the triple
whose members are the numbers
\[
x_k\,\omega_i\,\frac{y+i+1}{y+\rho}\,u^{(y)}_{y+\rho'}
\qquad(y=k-1,\,k,\,k+1),
\]
where $\rho'=\rho-1-i$ satisfies $1\le\rho'\le\rho-1$. By the induction
hypothesis the diagonal triple of offset $\rho'$ has property K at $k$, and
its middle member $u^{(k)}_{k+\rho'}$ is nonnegative. By
Lemma~\ref{lem:quotientfactor} with $\alpha=i+1$ and $\beta=\rho$ it keeps
the property when its members are multiplied by the values of
$(y+i+1)/(y+\rho)$ at $y=k-1$, $k$, $k+1$, and it keeps it again when all
three members are multiplied by the nonnegative number $x_k\,\omega_i$.
The remaining triple consists of the direct terms $S_\rho(k-1)$,
$S_\rho(k)$ and $S_\rho(k+1)$, and it has property K at $k$ by
Lemma~\ref{lem:directk}. Since a sum of triples with property K at $k$ has
property K at $k$, the proposition follows.
\end{proof}

The two inequalities that Section~\ref{sec:beyond} uses are now read off.
In the first the three counts of a diagonal triple carry the weights $c-1$,
$c$ and $c+1$, where $c$ is a real number, not to be confused with the
numbers $c_n$.

\begin{corollary}\label{cor:weightedlayer}
Let $k\ge7$ and $N$ be integers with $k\le N\le2k-2$, and let $c\ge k$ be a
real number. Then the tilted normalised counts $u^{(y)}_n$ at the critical
tilt $x_k$ of column $k$ satisfy
\[
(c+1)\,u^{(k+1)}_{N+1}-2\,c\,u^{(k)}_N+(c-1)\,u^{(k-1)}_{N-1}\ge0 .
\]
\end{corollary}

\begin{proof}
Let $\rho=N-k$, so that $0\le\rho\le k-2$. By
Proposition~\ref{prop:diagonalk} the triple with the members
$u^{(k-1)}_{N-1}$, $u^{(k)}_N$ and $u^{(k+1)}_{N+1}$ has property K at
$k$, and its middle member is nonnegative. We apply
Lemma~\ref{lem:productfactor} with $\alpha=c-k\ge0$ and $\beta=0$; its
hypothesis $\alpha\beta\le k^2-1$ holds because $\alpha\beta=0$. For
$\beta=0$ the three factors of that lemma are $k-1+\alpha$, $k+\alpha$ and
$k+1+\alpha$, that is, $c-1$, $c$ and $c+1$, so that the triple
\[
\bigl((c-1)\,u^{(k-1)}_{N-1},\ c\,u^{(k)}_N,\ (c+1)\,u^{(k+1)}_{N+1}\bigr)
\]
has property K at $k$. For $\mu=0$ the inequality of property K for this
triple is the claim.
\end{proof}

The second inequality is the convexity of the diagonal triple itself.

\begin{corollary}\label{cor:convexlayer}
Let $k\ge7$ and $N$ be integers with $k\le N\le2k-2$. Then the tilted
normalised counts $u^{(y)}_n$ at the critical tilt $x_k$ of column $k$
satisfy
\[
u^{(k+1)}_{N+1}-2\,u^{(k)}_N+u^{(k-1)}_{N-1}\ge0 .
\]
\end{corollary}

\begin{proof}
This is the inequality of property K at the parameter $\mu=0$ for the triple
of Proposition~\ref{prop:diagonalk} with $\rho=N-k$, for which
$0\le\rho\le k-2$.
\end{proof}

\section{The signs beyond the first layer}\label{sec:beyond}
In this section we determine the signs of the steps beyond the first layer:
in every column $k\ge1$ and every row $n\ge2k$ the step $d_{n,k}$ is positive, unless
$(n,k)$ is one of the eleven exceptional positions, which are the positions
$(2k,k)$ with $1\le k\le7$ and $(2k+1,k)$ with $2\le k\le5$
(Proposition~\ref{prop:beyondlayer}). From this statement and from the
computed negative signs of Lemma~\ref{lem:fortyfive} we derive the column
rule and the diagonal rule beyond the first layer
(Corollaries~\ref{cor:columnbeyond} and~\ref{cor:diagonalbeyond}); in
Section~\ref{sec:staircase} these are joined with the corresponding
statements for the first layer. Columns
$k\le81$ need no new estimate. Most
of the section treats columns $k\ge82$, by an induction on the row
(Proposition~\ref{prop:largecolumns}) in which the rows before the first
positive row of the column contribute terms that are not positive and whose
size is not known; we bound a weighted sum of these terms by a summation by
parts (Lemma~\ref{lem:debt}), which rests on a convexity of the tilted
normalised counts along the diagonal (Corollary~\ref{cor:weightedall}).

\subsection{Convexity along the diagonal}\label{subsec:beyond-convexity}

The arguments for columns $k\ge82$ are carried out at the critical tilt
of Section~\ref{sec:propagation}. For an integer $k\ge1$ this is the number
$x_k=(k/(k+1))^{k-1}$; it is positive and satisfies
\eqref{eq:propagation-critical}. As in that section, the tilted Cayley
masses at this tilt are written $\omega_n=c_n\,x_k^{\,n}$.
They are positive, and \eqref{eq:propagation-ratio}, taken at the tilt
$x=x_k$, gives the ratio of two consecutive masses: for every integer
$j\ge0$ the mass $\omega_{j+1}$ is the mass $\omega_j$ multiplied by
$x_k\,(1+1/(j+1))^{j}$.
By Lemma~\ref{lem:massesdecrease} the masses do not increase up to index
$k$. We begin with the counterpart of that lemma: from index $k-1$ on
they do not decrease.

\begin{lemma}\label{lem:massesincrease}
Let $k\ge1$ and $n\ge0$ be integers with $k\le n+1$. Then, at the critical
tilt $x_k$ of column $k$,
\[
c_n\,x_k^{\,n}\le c_{n+1}\,x_k^{\,n+1} .
\]
\end{lemma}

\begin{proof}
By \eqref{eq:propagation-ratio} for $x=x_k$ and $j=n$ the right-hand side
is the left-hand side multiplied by $x_k\,(1+1/(n+1))^{n}$. This factor is
at least
$x_k\,(1+1/k)^{k-1}$ by Lemma~\ref{lem:powermonotone}, applied to the
integers $k\le n+1$, and $x_k\,(1+1/k)^{k-1}=1$ by
\eqref{eq:propagation-critical}.
\end{proof}

Now let $k\ge2$ be an integer. In the rest of this section we use the
notation of the earlier sections at the critical tilt of column $k$; the
column is not shown in it. For a threshold $y\ge1$ and an integer $N\ge0$,
$u^{(y)}_N$ is the tilted normalised count of threshold $y$
(Section~\ref{sec:propagation}), and $U_N$, $U^{+}_N$ and $\eta_N$ are the
tilted column counts and the tilted step of column $k$
(Section~\ref{sec:diagonal}). As in Section~\ref{sec:propk}, the index $N$
of $u^{(y)}_N$ is called its size, and a tilted normalised count is also
called a tilted count for short. The tilted normalised counts are
nonnegative and do not decrease in the threshold,
\begin{equation}\label{eq:beyond-order}
0\le u^{(y)}_N\le u^{(y')}_N\qquad(1\le y\le y'),
\end{equation}
and in the free range $N\le y$ of a threshold $y$ they are the tilted
Cayley masses $\omega_N$, by \eqref{eq:propk-free}. By
\eqref{eq:diagonal-columncounts} the tilted column counts are the
normalised counts of columns $k-1$ and $k$, tilted by $x_k$. Both are
nonnegative. By their definition, $U_N$ and $U^{+}_N$ are the differences
$u^{(k)}_N-u^{(k-1)}_N$ and $u^{(k+1)}_N-u^{(k)}_N$ of the tilted counts of
consecutive thresholds, and the tilted step is the difference of the two
tilted column counts:
\begin{equation}\label{eq:beyond-step}
\eta_N=U^{+}_N-U_N=u^{(k+1)}_N-2\,u^{(k)}_N+u^{(k-1)}_N .
\end{equation}
The tilted step $\eta_N$ has the sign of the step $d_{N,k}$. Finally, at
the tilt $x_k$ the tilted last-car recursion \eqref{eq:counts-tilted}
reads
\begin{equation}\label{eq:beyond-recursion}
(N+1)\,u^{(y)}_{N+1}=x_k\sum_{i=0}^{N}\min(i+1,y)\,u^{(y)}_i\,u^{(y)}_{N-i}
\qquad(y\ge1,\ N\ge0).
\end{equation}

The convexity that we need concerns three consecutive points of a diagonal.
For an integer $M\ge1$ and a real number $c$ let
\begin{equation}\label{eq:beyond-weighted}
\nabla_c(M)=(c+1)\,u^{(k+1)}_{M+1}-2c\,u^{(k)}_M+(c-1)\,u^{(k-1)}_{M-1} .
\end{equation}
This is a second difference along the diagonal with weights, and for $c=M$
each tilted count is
multiplied by its own size. Section~\ref{sec:propk} treats the sizes of the
first layer. For $k\ge7$ and $k\le M\le2k-2$ the second difference without
weights, $u^{(k+1)}_{M+1}-2\,u^{(k)}_M+u^{(k-1)}_{M-1}$, is nonnegative
(Corollary~\ref{cor:convexlayer}), and so is $\nabla_c(M)$ for every real
number $c\ge k$ (Corollary~\ref{cor:weightedlayer}). The following
proposition treats all larger sizes, for the same columns $k\ge7$ and with
the weight equal to the size. This is the weight that the last-car recursion
produces, since the left-hand side of \eqref{eq:beyond-recursion} is a
tilted count multiplied by its size.

The proof of the proposition is an induction on the size. It writes the three
terms of $\nabla_N(N)$ as convolutions by \eqref{eq:beyond-recursion} and
splits each convolution according to whether one of its two indices is at
most $k-2$. The terms with such an index combine to positive multiples of
numbers $\nabla_c(M)$ at smaller sizes $M$, which are nonnegative by
Corollary~\ref{cor:weightedlayer} or by the induction hypothesis, and the
remaining terms are controlled by Corollary~\ref{cor:diagonaltilted}.

\begin{proposition}\label{prop:weightedbeyond}
Let $k\ge7$ and $N\ge2k-1$ be integers. Then the tilted normalised counts
$u^{(y)}_M$ at the critical tilt $x_k$ of column $k$ satisfy
\[
(N+1)\,u^{(k+1)}_{N+1}-2N\,u^{(k)}_N+(N-1)\,u^{(k-1)}_{N-1}\ge0 .
\]
\end{proposition}

\begin{proof}
In the notation \eqref{eq:beyond-weighted} the claim is $\nabla_N(N)\ge0$. We
prove it by induction on $N$. Let $N\ge2k-1$, suppose that $\nabla_M(M)\ge0$
for every integer $M$ with $2k-1\le M<N$, and put $r=N-(2k-1)$, so that
$r\ge0$.

By Proposition~\ref{prop:diagonalmonotone} we have
$z^{(y)}_i/c_i\le z^{(y+1)}_{i+1}/c_{i+1}$ for all integers $y\ge1$ and
$i\ge0$. Multiplying this inequality by $\omega_i$ and using that
$\omega_i\le\omega_{i+1}$ for $i\ge k-1$, by
Lemma~\ref{lem:massesincrease}, we see that at the critical tilt the tilted
normalised counts do not decrease along a diagonal step, on which the
threshold and the size are both raised by one, as soon as the size is at
least $k-1$:
\begin{equation}\label{eq:beyond-diagonal}
u^{(y)}_i\le u^{(y+1)}_{i+1}\qquad(y\ge1,\ i\ge k-1).
\end{equation}

We now show that the numbers $\nabla_c(M)$ at smaller sizes are
nonnegative: for every integer $M$ and every real number $c$ in the ranges
indicated,
\begin{equation}\label{eq:beyond-smaller}
\nabla_c(M)\ge0\qquad(k\le M\le N-1,\ \ k\le c\le2k-1).
\end{equation}
If $M\le2k-2$, this is Corollary~\ref{cor:weightedlayer}, and it is here
that the hypothesis $k\ge7$ is used. If $2k-1\le M\le N-1$, then
$\nabla_M(M)\ge0$ by the induction hypothesis, and a comparison of the
coefficients of the three tilted counts shows that
\[
M\,\nabla_c(M)=c\,\nabla_M(M)
+(M-c)\,\bigl(u^{(k+1)}_{M+1}-u^{(k-1)}_{M-1}\bigr) .
\]
Here $c>0$ and $M-c\ge0$, and
$u^{(k-1)}_{M-1}\le u^{(k)}_M\le u^{(k+1)}_{M+1}$ by
\eqref{eq:beyond-diagonal}; hence $M\,\nabla_c(M)\ge0$. We also need the
weight $k$ at size $k-1$, which lies below the range of
\eqref{eq:beyond-smaller}:
\begin{equation}\label{eq:beyond-bottom}
\nabla_k(k-1)=(k+1)\,\omega_k-2k\,\omega_{k-1}+(k-1)\,\omega_{k-2}\ge0 .
\end{equation}
Here the equality holds by \eqref{eq:propk-free}, and the inequality
because $\omega_k\ge\omega_{k-1}$ by Lemma~\ref{lem:massesincrease} and
$\omega_{k-2}\ge\omega_{k-1}$ by Lemma~\ref{lem:massesdecrease}.

Next we apply \eqref{eq:beyond-recursion} to
threshold $k+1$ at size $N+1$, to threshold $k$ at size $N$
and to threshold $k-1$ at size $N-1$. This gives
$(N+1)\,u^{(k+1)}_{N+1}=x_k\,\Sigma^{(k+1)}$,
$N\,u^{(k)}_N=x_k\,\Sigma^{(k)}$ and
$(N-1)\,u^{(k-1)}_{N-1}=x_k\,\Sigma^{(k-1)}$, where
\begin{align*}
\Sigma^{(k+1)}&=\sum_{i=0}^{N}\min(i+1,k+1)\,u^{(k+1)}_i\,u^{(k+1)}_{N-i},\\
\Sigma^{(k)}&=\sum_{i=0}^{N-1}\min(i+1,k)\,u^{(k)}_i\,u^{(k)}_{N-1-i},\\
\Sigma^{(k-1)}&=\sum_{i=0}^{N-2}\min(i+1,k-1)\,u^{(k-1)}_i\,u^{(k-1)}_{N-2-i}.
\end{align*}
Hence $\nabla_N(N)$ is $x_k$ times
$\Sigma^{(k+1)}-2\,\Sigma^{(k)}+\Sigma^{(k-1)}$, and it suffices to show
that
\[
\Sigma^{(k+1)}-2\,\Sigma^{(k)}+\Sigma^{(k-1)}\ge0 .
\]

To this end we sort the terms of the three sums. Every term
is a weight times a product of two tilted counts. The indices of the
two counts, the first and the second index of the term, add up to $N$, to
$N-1$ and to $N-2$ in the three sums. In each sum we separate the
terms whose first index is at most $k-2$, the terms whose second index is at
most $k-2$, and the terms in which both indices are at least $k-1$. No term
has both indices at most $k-2$, because the two indices add up to at least
$N-2\ge2k-3$.

Let $0\le\ell\le k-2$. By \eqref{eq:propk-free} a tilted count with index
$\ell$ is equal to $\omega_\ell$ for each of the three thresholds. In
the term with the first index $\ell$ the weight $\min(\ell+1,y)$ is $\ell+1$
for each of the thresholds $y=k+1$, $k$, $k-1$. In the term with the second
index $\ell$ the first index is $N-\ell$, $N-1-\ell$ or $N-2-\ell$, and
since $N-2-\ell\ge k-1$, the weight is the threshold itself. The two terms
with index $\ell$ therefore contribute to $\Sigma^{(k+1)}$, to
$\Sigma^{(k)}$ and to $\Sigma^{(k-1)}$ the amounts
\[
(\ell+k+2)\,\omega_\ell\,u^{(k+1)}_{N-\ell},\qquad
(\ell+k+1)\,\omega_\ell\,u^{(k)}_{N-1-\ell},\qquad
(\ell+k)\,\omega_\ell\,u^{(k-1)}_{N-2-\ell},
\]
and to $\Sigma^{(k+1)}-2\,\Sigma^{(k)}+\Sigma^{(k-1)}$ the amount
$\omega_\ell\,\nabla_{\ell+k+1}(N-1-\ell)$.

The terms in which both indices are at least $k-1$ are sorted by their
second index, which we write as $m_t=N-k-t$ with an integer $t$. In
$\Sigma^{(k)}$ these are the terms with $0\le t\le r$; their first index is
$k-1+t$, and their weight is $k$. In $\Sigma^{(k+1)}$ the terms with
$0\le t\le r$ have the first index $k+t$ and the weight $k+1$, and there is
one further term, with the first index $k-1$ and the second index $N-k+1$;
its weight is $k$, and by \eqref{eq:propk-free} it is equal to
$k\,\omega_{k-1}\,u^{(k+1)}_{N-k+1}$. In $\Sigma^{(k-1)}$ the first index is
$k-2+t$ and the weight is $k-1$, and the terms in question are those with
$1\le t\le r$, since for $t=0$ the first index is $k-2$. In order that the
three sums run over the same values of $t$, we add to the terms of
$\Sigma^{(k-1)}$ the product for $t=0$, which is
\[
(k-1)\,u^{(k-1)}_{k-2}\,u^{(k-1)}_{m_0}
=(k-1)\,\omega_{k-2}\,u^{(k-1)}_{N-k} ,
\]
and subtract it again. For $0\le t\le r$ let
\[
\Upsilon_t=(k+1)\,u^{(k+1)}_{k+t}\,u^{(k+1)}_{m_t}
-2k\,u^{(k)}_{k-1+t}\,u^{(k)}_{m_t}
+(k-1)\,u^{(k-1)}_{k-2+t}\,u^{(k-1)}_{m_t} .
\]
Then we have found that
\begin{equation}\label{eq:beyond-split}
\begin{split}
\Sigma^{(k+1)}-2\,\Sigma^{(k)}+\Sigma^{(k-1)}
={}&\sum_{\ell=0}^{k-2}\omega_\ell\,\nabla_{\ell+k+1}(N-1-\ell)
+\sum_{t=0}^{r}\Upsilon_t\\
&+k\,\omega_{k-1}\,u^{(k+1)}_{N-k+1}
-(k-1)\,\omega_{k-2}\,u^{(k-1)}_{N-k} .
\end{split}
\end{equation}

It remains to show that the right-hand side of \eqref{eq:beyond-split} is
nonnegative. For
$0\le\ell\le k-2$ the weight $c=\ell+k+1$ satisfies $k+1\le c\le2k-1$, and
the size $M=N-1-\ell$ satisfies $k\le N-k+1\le M\le N-1$. Hence
$\nabla_{\ell+k+1}(N-1-\ell)\ge0$ by \eqref{eq:beyond-smaller}, and the
first sum on the right of \eqref{eq:beyond-split} is nonnegative.

We turn to the sum of the $\Upsilon_t$. Let $0\le t\le r$. We insert
$u^{(k+1)}_{m_t}=u^{(k)}_{m_t}+U^{+}_{m_t}$ and
$u^{(k-1)}_{m_t}=u^{(k)}_{m_t}-U_{m_t}$ into $\Upsilon_t$ and obtain
\[
\Upsilon_t=u^{(k)}_{m_t}\,\nabla_k(k-1+t)
+(k+1)\,u^{(k+1)}_{k+t}\,U^{+}_{m_t}
-(k-1)\,u^{(k-1)}_{k-2+t}\,U_{m_t} .
\]
Here $\nabla_k(k-1+t)\ge0$: for $t=0$ this is \eqref{eq:beyond-bottom}, and
for $t\ge1$ it is \eqref{eq:beyond-smaller} with the weight $k$ and the size
$k-1+t$, which lies between $k$ and $k-1+r=N-k$. Moreover
$(k+1)\,U^{+}_{m_t}\ge k\,U_{m_t-1}$ by Corollary~\ref{cor:diagonaltilted},
applied at row $m_t\ge1$. Since $u^{(k)}_{m_t}$ and $u^{(k+1)}_{k+t}$ are
nonnegative by \eqref{eq:beyond-order}, it follows that
\[
\Upsilon_t\ge k\,u^{(k+1)}_{k+t}\,U_{m_t-1}
-(k-1)\,u^{(k-1)}_{k-2+t}\,U_{m_t} .
\]
We sum over $0\le t\le r$ and combine the positive term of $t$ with the
negative term of $t+1$, which is $(k-1)\,u^{(k-1)}_{k-1+t}\,U_{m_t-1}$
because $m_{t+1}=m_t-1$. The positive term of $t=r$, in which $m_r-1=k-2$,
and the negative term of $t=0$, in which $m_0=N-k$, are left over, so that
\begin{multline*}
\sum_{t=0}^{r}\Upsilon_t\ge\sum_{t=0}^{r-1}U_{m_t-1}\,
\bigl(k\,u^{(k+1)}_{k+t}-(k-1)\,u^{(k-1)}_{k-1+t}\bigr)\\
+k\,u^{(k+1)}_{k+r}\,U_{k-2}-(k-1)\,\omega_{k-2}\,U_{N-k} .
\end{multline*}
Every term of the sum on the right is nonnegative, because
$U_{m_t-1}\ge0$ and
$u^{(k-1)}_{k-1+t}\le u^{(k)}_{k+t}\le u^{(k+1)}_{k+t}$ by
\eqref{eq:beyond-diagonal} and \eqref{eq:beyond-order}. Further,
$U_{k-2}=0$ by \eqref{eq:propk-free}. Hence the sum of the
$\Upsilon_t$ is at least $-(k-1)\,\omega_{k-2}\,U_{N-k}$, and the last three
terms on the right of \eqref{eq:beyond-split} together are at least
\begin{align*}
&k\,\omega_{k-1}\,u^{(k+1)}_{N-k+1}
-(k-1)\,\omega_{k-2}\,\bigl(U_{N-k}+u^{(k-1)}_{N-k}\bigr)\\
&\qquad=k\,\omega_{k-1}\,u^{(k+1)}_{N-k+1}
-(k-1)\,\omega_{k-2}\,u^{(k)}_{N-k} .
\end{align*}
This number is nonnegative. Indeed,
$u^{(k)}_{N-k}\le u^{(k+1)}_{N-k+1}$ by \eqref{eq:beyond-diagonal}, which
applies because $N-k\ge k-1$, and
$(k-1)\,\omega_{k-2}\le k\,\omega_{k-1}$: by \eqref{eq:propagation-ratio}
for $x=x_k$ and $j=k-2$,
\[
k\,\omega_{k-1}=k\,x_k\,\Bigl(1+\frac{1}{k-1}\Bigr)^{k-2}\,\omega_{k-2}
=(k-1)\,x_k\,\Bigl(1+\frac{1}{k-1}\Bigr)^{k-1}\,\omega_{k-2},
\]
and the factor $x_k\,(1+1/(k-1))^{k-1}$ is at least
$x_k\,(1+1/k)^{k-1}$, which is $1$ by \eqref{eq:propagation-critical}. Since
the first sum on the right of \eqref{eq:beyond-split} is nonnegative as
well, the right-hand side of \eqref{eq:beyond-split} is nonnegative, which
completes the induction.
\end{proof}

The proposition and the convexity of the first layer
are joined into one statement by the weight $\Psi_k(N)=\max(N,2k-1)$ of
Section~\ref{sec:kernel}. It is constant on the sizes $N\le2k-1$, where the
convexity without weights is available, and equal to the size from there
on, where only the weighted form of Proposition~\ref{prop:weightedbeyond}
is. One weight for all sizes is needed because the sum of the next
subsection runs through the first layer and beyond it.

\begin{corollary}\label{cor:weightedall}
Let $k\ge7$ and $N\ge k$ be integers. Then the tilted normalised counts
$u^{(y)}_M$ at the critical tilt $x_k$ of column $k$ and the weight $\Psi_k$
satisfy
\[
\Psi_k(N+1)\,u^{(k+1)}_{N+1}-2\,\Psi_k(N)\,u^{(k)}_N
+\Psi_k(N-1)\,u^{(k-1)}_{N-1}\ge0 .
\]
\end{corollary}

\begin{proof}
If $k\le N\le2k-2$, the three weights are equal to $2k-1$, and the
left-hand side is $2k-1$ times
$u^{(k+1)}_{N+1}-2\,u^{(k)}_N+u^{(k-1)}_{N-1}$, which is nonnegative by
Corollary~\ref{cor:convexlayer}. If $N=2k-1$, the weights are
$\Psi_k(N+1)=N+1$, $\Psi_k(N)=N$ and $\Psi_k(N-1)=2k-1=N$, and the left-hand
side is
\[
\bigl((N+1)\,u^{(k+1)}_{N+1}-2N\,u^{(k)}_N+(N-1)\,u^{(k-1)}_{N-1}\bigr)
+u^{(k-1)}_{N-1} ;
\]
this is nonnegative by Proposition~\ref{prop:weightedbeyond} and because
$u^{(k-1)}_{N-1}\ge0$. If $N\ge2k$, the weights are $N+1$, $N$ and $N-1$,
and the claim is Proposition~\ref{prop:weightedbeyond}.
\end{proof}

\subsection{The weighted debt of a column}\label{subsec:beyond-debt}

Let $k\ge7$, and let $J\ge k$ be an integer. We now sum the tilted steps
$\eta_n$ of column $k$ over rows $k\le n\le J$, with the weights
$\Psi_k(n)$. In the rows before the first positive row of the column the
steps are not positive, so that such a sum can be negative, and the
following lemma says how negative it can be: the sum is at least minus a
number that consists of quantities of the last two rows $J-1$ and $J$
alone. We call this number, which stands in the large parentheses on the
right-hand side of the lemma, the \emph{weighted debt} of column $k$ at row
$J$. Nothing about the
single rows enters it, and nothing about the row in which the sign changes.

The reason is a summation by parts. By \eqref{eq:beyond-step} the tilted step
$\eta_n$ is a second difference in the threshold, at the fixed size $n$.
When the terms $\Psi_k(n)\,\eta_n$ of rows $k,\dots,J$ are added and
regrouped along the diagonals, they become, up to terms of rows $J-1$ and
$J$, the left-hand sides of Corollary~\ref{cor:weightedall} at sizes
$k,\dots,J-1$ and one term of rows $k-1$ and $k$, the number
$(2k-1)\,(\omega_k-\omega_{k-1})$. Each of these left-hand sides is
nonnegative, and so is this number, by Lemma~\ref{lem:massesincrease}.

\begin{lemma}\label{lem:debt}
Let $k\ge7$ and $J\ge k$ be integers. Then, at the critical tilt $x_k$ of
column $k$, the tilted steps $\eta_n$ of the column, its tilted column
counts $U_N$, the tilted normalised counts $u^{(k)}_N$ of threshold $k$ and
the weight $\Psi_k$ satisfy
\[
\sum_{n=k}^{J}\Psi_k(n)\,\eta_n\ge-\Bigl(\Psi_k(J)\,U_J
+\Psi_k(J-1)\,U_{J-1}
+\bigl(\Psi_k(J)-\Psi_k(J-1)\bigr)\,u^{(k)}_J\Bigr).
\]
\end{lemma}

\begin{proof}
For an integer $J\ge k$ let
\begin{align*}
Y_J={}&\sum_{n=k}^{J}\Psi_k(n)\,\eta_n+\Psi_k(J)\,U_J
+\Psi_k(J-1)\,U_{J-1}\\
&-\Psi_k(J-1)\,u^{(k)}_{J-1}+\Psi_k(J)\,u^{(k)}_J .
\end{align*}
We show by induction on $J$ that $Y_J\ge0$.

For $J=k$ the weights $\Psi_k(k)$ and $\Psi_k(k-1)$ are both equal to
$2k-1$. Inserting \eqref{eq:beyond-step} for $\eta_k$ and the definitions of
$U_k$ and $U_{k-1}$, we find by a direct computation that all terms cancel
except two, and by \eqref{eq:propk-free},
\[
Y_k=(2k-1)\,\bigl(u^{(k+1)}_k-u^{(k-1)}_{k-1}\bigr)
=(2k-1)\,(\omega_k-\omega_{k-1}) .
\]
This is nonnegative by Lemma~\ref{lem:massesincrease}.

Let $J\ge k$. In the difference $Y_{J+1}-Y_J$ the terms
$\Psi_k(J)\,U_J$ cancel, and
\begin{align*}
Y_{J+1}-Y_J={}&\Psi_k(J+1)\,
\bigl(\eta_{J+1}+U_{J+1}+u^{(k)}_{J+1}\bigr)-2\,\Psi_k(J)\,u^{(k)}_J\\
&+\Psi_k(J-1)\,\bigl(u^{(k)}_{J-1}-U_{J-1}\bigr).
\end{align*}
By \eqref{eq:beyond-step} and the definition of $U_N$, the first parenthesis
is $u^{(k+1)}_{J+1}$ and the second is $u^{(k-1)}_{J-1}$. Hence
\[
Y_{J+1}-Y_J=\Psi_k(J+1)\,u^{(k+1)}_{J+1}-2\,\Psi_k(J)\,u^{(k)}_J
+\Psi_k(J-1)\,u^{(k-1)}_{J-1},
\]
which is nonnegative by Corollary~\ref{cor:weightedall}, applied with
$N=J$. This completes the induction.

From $Y_J\ge0$ we get
\[
\sum_{n=k}^{J}\Psi_k(n)\,\eta_n\ge-\Psi_k(J)\,U_J-\Psi_k(J-1)\,U_{J-1}
+\Psi_k(J-1)\,u^{(k)}_{J-1}-\Psi_k(J)\,u^{(k)}_J .
\]
By Proposition~\ref{prop:tilteddecrease} we have
$u^{(k)}_J\le u^{(k)}_{J-1}$. Therefore
\[
\Psi_k(J-1)\,u^{(k)}_{J-1}-\Psi_k(J)\,u^{(k)}_J
\ge-\bigl(\Psi_k(J)-\Psi_k(J-1)\bigr)\,u^{(k)}_J ,
\]
and the lemma follows.
\end{proof}

\subsection{Columns $k\ge82$}\label{subsec:beyond-large}

Let $k\ge82$. By Proposition~\ref{prop:onechange} column $k$ changes its
sign once between rows $k$ and $2k-1$: there is a row $q$ with
$k+1\le q\le k+27k/100$ such that the step of the column is not positive in
rows $k\le j<q$ and positive in rows $q\le j\le2k-1$. Row $q$ is
the first positive row of the column. We now show that the step stays
positive in every later row.

The proof is an induction on the row: for a row $n\ge2k$ it starts from the
inequality of Proposition~\ref{prop:columninequality} at the critical tilt,
to which the rows before row $q$ contribute terms that are not positive and
whose size is not known. Lemmas~\ref{lem:orderbefore}
and~\ref{lem:orderafter} compare the weights that the prime kernel gives to
the rows before row $n$ in this inequality with the weights $\Psi_k$, so
that by Lemma~\ref{lem:debt} the contribution of all these rows is at least
minus a nonnegative multiple of the weighted debt of the column at row
$n-1$. The proof then bounds the debt by the quantities $\eta_{n-1}$ and
$U_{n-1}$ of the other two terms of the inequality; the debt enters with
coefficients that are at most a fixed multiple of $1/\sqrt k$, whereas the
coefficients of those two terms are bounded below by a positive constant.

\begin{proposition}\label{prop:largecolumns}
Let $k\ge82$ and $n\ge2k$ be integers. Then $d_{n,k}>0$.
\end{proposition}

\begin{proof}
We use the notation of Section~\ref{subsec:beyond-convexity} at the critical
tilt $x_k$ of column $k$, and the prime kernel
$\bar\kappa_m=\kappa^{(k)}_m\,x_k^{\,m}$ of threshold $k$ at this tilt; the
numbers $\bar\kappa_m$ are nonnegative by Lemma~\ref{lem:kernelsign}. Since
$\eta_j$ has the sign of $d_{j,k}$, the claim is $\eta_n>0$. Let $q$ be a
row as in Proposition~\ref{prop:onechange}, so that $k+1\le q$ and
$100\,(q-k)\le27\,k$, that $\eta_j\le0$ for $k\le j<q$, and that $\eta_j>0$
for $q\le j\le2k-1$. It suffices to prove the following claim: if $n\ge2k$
is an integer and $\eta_j>0$ for every $j$ with $q\le j\le n-1$, then
$\eta_n>0$. Indeed, since $q<2k$ and $\eta_m>0$ for $q\le m\le2k-1$ by the
choice of $q$, the claim gives $\eta_m>0$ for every integer $m\ge q$ by
induction on $m$.

To prove the claim, let $n\ge2k$, and suppose that $\eta_j>0$ for
$q\le j\le n-1$. Since $q<2k\le n$, row $n-1$ is one of these rows, so
that $\eta_{n-1}>0$. We multiply the inequality of
Proposition~\ref{prop:columninequality}, which applies because $n\ge k+1$,
by $x_k^{\,n}$ and index its sum by row $j=n-m$. Since
$\kappa^{(k)}_m\,\delta^{(k)}_{n-m}\,x_k^{\,n}=\bar\kappa_m\,\eta_{n-m}$ and
$\beta^{(k-1)}_{n-1}\,x_k^{\,n}=x_k\,U_{n-1}$, it becomes
\begin{equation}\label{eq:beyond-column}
n\,\eta_n\ge(k+1)\,x_k\,\eta_{n-1}+2\,x_k\,U_{n-1}
+n\sum_{j=k}^{n-1}\bar\kappa_{n-j}\,\eta_j .
\end{equation}
The first term on the right is a multiple of the tilted step $\eta_{n-1}$
of the row before, which is positive; the second is the source, which comes
from column $k-1$ and is nonnegative; and the sum is the feedback of rows
$k\le j\le n-1$, in which row $j$ has the kernel weight
$\bar\kappa_{n-j}$. Let
\[
L=n-q+1,\qquad\theta=\frac{\bar\kappa_L}{2k-1},\qquad
\mathfrak K=n(n-1)\,\bar\kappa_L .
\]
Then $\theta\ge0$ and $n(n-1)\,\theta=\mathfrak K/(2k-1)$.

We first bound the feedback from below in terms of the weighted debt of the
column. Let $\mathfrak S$ be the feedback in \eqref{eq:beyond-column}, the
sum of $\bar\kappa_{n-j}\,\eta_j$ over $k\le j\le n-1$.
Lemmas~\ref{lem:orderbefore} and~\ref{lem:orderafter} apply to the rows $q$
and $n$, because $k+1\le q$, $100\,(q-k)\le27\,k$ and $n\ge2k$. For
$k\le j<q$ we have $\eta_j\le0$, and Lemma~\ref{lem:orderbefore} gives
$\bar\kappa_{n-j}\,(2k-1)\le\bar\kappa_L\,\Psi_k(j)$, so that
$\bar\kappa_{n-j}\le\theta\,\Psi_k(j)$. For $q\le j\le n-1$ we have
$\eta_j>0$, and Lemma~\ref{lem:orderafter} gives
$\bar\kappa_L\,\Psi_k(j)\le\bar\kappa_{n-j}\,(2k-1)$, so that
$\bar\kappa_{n-j}\ge\theta\,\Psi_k(j)$. In both cases
$\bar\kappa_{n-j}\,\eta_j\ge\theta\,\Psi_k(j)\,\eta_j$. We sum over $j$ and
apply Lemma~\ref{lem:debt} with $J=n-1$:
\begin{equation}\label{eq:beyond-feedback}
\mathfrak S\ge\theta\sum_{j=k}^{n-1}\Psi_k(j)\,\eta_j\ge-\theta\,\mathfrak D,
\end{equation}
where
\[
\mathfrak D=\Psi_k(n-1)\,U_{n-1}+\Psi_k(n-2)\,U_{n-2}
+\bigl(\Psi_k(n-1)-\Psi_k(n-2)\bigr)\,u^{(k)}_{n-1}
\]
is the weighted debt of the column at row $n-1$. Since $n-1\ge2k-1$, we
have $\Psi_k(n-1)=n-1$ and $n-2\le\Psi_k(n-2)\le n-1$. As $U_{n-2}\ge0$ and
$u^{(k)}_{n-1}\ge0$, it follows that
\begin{equation}\label{eq:beyond-debt-rows}
\mathfrak D\le(n-1)\,(U_{n-1}+U_{n-2})+u^{(k)}_{n-1} .
\end{equation}

Next we bound $U_{n-2}$ and
$u^{(k)}_{n-1}$ by the quantities $\eta_{n-1}$ and $U_{n-1}$ of the first
two terms on the right of \eqref{eq:beyond-column}.
Corollary~\ref{cor:diagonaltilted}, applied at row $n-1$, states that
$k\,U_{n-2}\le(k+1)\,U^{+}_{n-1}$, and $U^{+}_{n-1}=U_{n-1}+\eta_{n-1}$ by
\eqref{eq:beyond-step}. Hence
\begin{equation}\label{eq:beyond-previous}
U_{n-2}\le\Bigl(1+\frac1k\Bigr)\,(U_{n-1}+\eta_{n-1}) .
\end{equation}
For $u^{(k)}_{n-1}$ we show that
\begin{equation}\label{eq:beyond-domination}
x_k\,\Omega\,u^{(k)}_{n-1}\le(n-1)\,U_{n-1},
\qquad\text{where }\Omega=\sum_{j=0}^{8}\omega_j .
\end{equation}
We subtract the recursion \eqref{eq:beyond-recursion} of threshold $k-1$
at size $n-1$ from that of threshold $k$ at the same size, and
obtain
\[
(n-1)\,U_{n-1}=x_k\sum_{i=0}^{n-2}\mathfrak T_i ,
\]
where
\[
\mathfrak T_i=\min(i+1,k)\,u^{(k)}_i\,u^{(k)}_{n-2-i}
-\min(i+1,k-1)\,u^{(k-1)}_i\,u^{(k-1)}_{n-2-i} .
\]
Every $\mathfrak T_i$ is nonnegative, by \eqref{eq:beyond-order} and because
$\min(i+1,k-1)\le\min(i+1,k)$. Let $0\le j\le8$ and $i=n-2-j$. Then
$i+1\ge2k-9\ge k$, so that the two minima are $k$ and $k-1$, and
\[
\mathfrak T_{n-2-j}=k\,u^{(k)}_{n-2-j}\,u^{(k)}_j
-(k-1)\,u^{(k-1)}_{n-2-j}\,u^{(k-1)}_j
\ge u^{(k)}_{n-2-j}\,u^{(k)}_j\ge u^{(k)}_{n-1}\,\omega_j .
\]
Here the first inequality is \eqref{eq:beyond-order}. The second holds
because $u^{(k)}_j=\omega_j$ by \eqref{eq:propk-free}, as $j\le8\le k$,
and because $u^{(k)}_{n-2-j}\ge u^{(k)}_{n-1}$ by
Proposition~\ref{prop:tilteddecrease}. Keeping only these nine terms of the
sum, we obtain \eqref{eq:beyond-domination}. Now we insert
\eqref{eq:beyond-previous} and \eqref{eq:beyond-domination} into
\eqref{eq:beyond-debt-rows}. With the number
\[
\sigma=2+\frac1k+\frac{1}{x_k\,\Omega}
\]
we find
\begin{equation}\label{eq:beyond-debt-source}
\mathfrak D\le(n-1)\,\Bigl(\sigma\,U_{n-1}
+\Bigl(1+\frac1k\Bigr)\,\eta_{n-1}\Bigr).
\end{equation}

By \eqref{eq:beyond-column},
\eqref{eq:beyond-feedback} and \eqref{eq:beyond-debt-source}, and since
$n(n-1)\,\theta=\mathfrak K/(2k-1)$,
\begin{align*}
n\,\eta_n&\ge(k+1)\,x_k\,\eta_{n-1}+2\,x_k\,U_{n-1}-n\,\theta\,\mathfrak D\\
&\ge U_{n-1}\,\Bigl(2\,x_k-\frac{\mathfrak K\,\sigma}{2k-1}\Bigr)
+(k+1)\,\eta_{n-1}\,\Bigl(x_k-\frac{\mathfrak K}{(2k-1)\,k}\Bigr).
\end{align*}
Since $U_{n-1}\ge0$ and $\eta_{n-1}>0$, we obtain $\eta_n>0$ as soon as
\begin{equation}\label{eq:beyond-conditions}
\mathfrak K\,\sigma\le2\,x_k\,(2k-1)\qquad\text{and}\qquad
\mathfrak K<x_k\,k\,(2k-1).
\end{equation}

It remains to verify the two conditions \eqref{eq:beyond-conditions}. To
this end we first show that
\begin{equation}\label{eq:beyond-kernel}
\mathfrak K=n(n-1)\,\bar\kappa_L\le3.762\,\sqrt k .
\end{equation}
Since $q\le1.27\,k$, we have $n=L+q-1\le L+1.27\,k-1$. We distinguish three
cases according to the size of $L$.

If $L\le k$, then $n\le2.27\,k$. Moreover
$100\,(L-1)=100\,(n-q)\ge73\,k$, because $n\ge2k$ and $100\,q\le127\,k$, so
that $\bar\kappa_L\le0.73/(k\sqrt k)$ by Lemma~\ref{lem:kernelfar}. Hence
\[
\mathfrak K\le n^2\,\bar\kappa_L\le2.27^2\cdot0.73\,\sqrt k\le3.762\,\sqrt k .
\]

If $k<L\le2k$, Corollary~\ref{cor:harmoniccritical} gives
$L\,\bar\kappa_L\le k\,\bar\kappa_k$, and
$\bar\kappa_k\le0.662/(k\sqrt k)$ by Lemma~\ref{lem:kernelatk}; hence
$L\,\bar\kappa_L\le0.662/\sqrt k$. Furthermore
$n^2\le(L+1.27\,k)^2\le5.35\,k\,L$. With $t=L/k$ the second of these
inequalities reads $t^2-2.81\,t+1.6129\le0$; the polynomial on the left is
convex and negative at $t=1$ and at $t=2$, so that it is negative on the
interval $1\le t\le2$. Hence
\[
\mathfrak K\le n^2\,\bar\kappa_L\le5.35\,k\,L\,\bar\kappa_L
\le5.35\cdot0.662\,\sqrt k\le3.762\,\sqrt k .
\]

If $L>2k$, Lemma~\ref{lem:quadratic} gives
$(L-1)\,L\,\bar\kappa_L\le2k(k-1)\,\bar\kappa_k$, and with
Lemma~\ref{lem:kernelatk} we get
$(L-1)\,L\,\bar\kappa_L\le2\cdot0.662\,\sqrt k$. In this case
$k\le(L-1)/2$, so that $n\le(L-1)+1.27\,k\le1.635\,(L-1)$, and therefore
\[
n(n-1)\le1.635^2\,(L-1)\,L\le2.675\,(L-1)\,L .
\]
Hence $\mathfrak K\le2.675\cdot2\cdot0.662\,\sqrt k\le3.762\,\sqrt k$.

To derive \eqref{eq:beyond-conditions} from \eqref{eq:beyond-kernel}, we
bound $x_k$ and $\Omega$ from below. By \eqref{eq:propagation-critical},
$1/x_k=(1+1/k)^{k-1}\le e^{(k-1)/k}\le e$, so that $x_k\ge1/e\ge0.3678$.
Hence $\Omega$, which is the sum of $c_j\,x_k^{\,j}$ over
$0\le j\le8$, is at least the value of this sum with $x_k$ replaced by
$0.3678$, and a direct computation with the numbers $c_j=(j+1)^{j-1}/j!$
shows that this value is at least $2.01$. Therefore
$\Omega\ge2.01$. It follows that
\[
\sigma\le2+\frac{1}{82}+\frac{1}{0.3678\cdot2.01}\le3.3652 .
\]
By \eqref{eq:beyond-kernel}, the first condition of
\eqref{eq:beyond-conditions} follows from
\[
3.762\cdot3.3652\,\sqrt k\le12.66\,\sqrt k\le2\cdot0.3678\,(2k-1).
\]
With $s=\sqrt k$ the second of these inequalities reads
\[
1.4712\,s^2-12.66\,s-0.7356\ge0 .
\]
The left-hand side increases for $s\ge4.31$ and is positive at $s=9.05$,
and $\sqrt k\ge9.05$, because $9.05^2\le82$. The second
condition of \eqref{eq:beyond-conditions} holds as well: since
$2k-1\ge163$,
\[
\mathfrak K\le3.762\,k<0.3678\,k\,(2k-1)\le x_k\,k\,(2k-1).
\]
This proves the claim, and with it the proposition.
\end{proof}

\subsection{All columns, and the two rules beyond the first layer}
\label{subsec:beyond-all}

We now assemble the statement for all columns. The positions that the
hypothesis of the following proposition excludes are the eleven exceptional
positions.

\begin{proposition}\label{prop:beyondlayer}
Let $k\ge1$ and $n\ge2k$ be integers, and
suppose that neither $n=2k$ and $k\le7$, nor $n=2k+1$ and $2\le k\le5$. Then
$d_{n,k}>0$.
\end{proposition}

\begin{proof}
We distinguish column $k=1$, columns $2\le k\le81$ and columns
$k\ge82$.

Let first $k=1$. Then the supposition excludes $n=2$, so that $n\ge3$.
Since $b_{n,0}=A_{n,0}$ and $b_{n,1}=A_{n,1}-A_{n,0}$
(Section~\ref{sec:counts}), the claim $d_{n,1}>0$ is the inequality
\begin{equation}\label{eq:beyond-first-column}
2\,A_{n,0}<A_{n,1} ,
\end{equation}
which we prove by induction on $n\ge3$. Starting from $A_{0,\ell}=1$, the
last-car recursion (Proposition~\ref{prop:lastcar}) gives the values $1$,
$2$, $6$ of $A_{n,0}$ and the values $1$, $3$, $13$ of $A_{n,1}$ for
$n=1,2,3$. Thus $2\,A_{3,0}=12<13=A_{3,1}$. Now let $n\ge3$, and suppose
that \eqref{eq:beyond-first-column} holds for $n$. For $\ell=0$ and for
$\ell=1$ the last-car recursion reads
\[
A_{n+1,0}=\sum_{j=0}^{n}\binom nj\,A_{j,0}\,A_{n-j,0},\qquad
A_{n+1,1}=\sum_{j=0}^{n}\min(j+1,2)\binom nj\,A_{j,1}\,A_{n-j,1} .
\]
We compare twice the first sum with the second sum term by term. For
$j\ge1$ the term of the second sum carries the factor $\min(j+1,2)=2$, and
$A_{j,0}\,A_{n-j,0}\le A_{j,1}\,A_{n-j,1}$, because the cumulative counts do
not decrease in the bound on the displacement. For $j=0$ the two terms are
$A_{n,0}$ and $A_{n,1}$, and $2\,A_{n,0}<A_{n,1}$ by the induction
hypothesis. Adding up, we obtain $2\,A_{n+1,0}<A_{n+1,1}$.

Now let $2\le k\le81$, and let $n_0=n_0(k)$ be the row of
Definition~\ref{def:junctionrows}. By Lemma~\ref{lem:junction} we have
$2k+3\le n_0$, and $d_{m,k}>0$ for every row $m$ with
$2k\le m<n_0$ other than $m=2k$ when $k\le7$ and $m=2k+1$ when $k\le5$. If
$n<n_0$, row $n$ is one of these rows. Indeed, for $k\ge2$ the supposition
of the proposition says that $n\ne2k$ when $k\le7$ and that $n\ne2k+1$ when
$k\le5$, so that the rows which it excludes are the rows which the lemma
excepts. Hence the claim is part of
the lemma.

Let $n\ge n_0$, and put $m=n_0-1$. Then $2k+2\le m<n_0$, so that $m$ is one
of the rows of Lemma~\ref{lem:junction} and is not excepted. Hence
$d_{m,k}>0$. We apply
Proposition~\ref{prop:propagation} to column $k$, to row $m\ge k$
and to the tilt $x=x_k$, which satisfies $0<x_k<1$ because $k\ge2$. Its
hypotheses hold. The tilted normalised counts
$z^{(k)}_N\,x_k^{\,N}$ of threshold $k$ do not increase in $N$, by
Proposition~\ref{prop:tilteddecrease}. The sums of
$\xi^{(k)}_t\,x_k^{\,t}$ over $s\le t\le m+1$, where $\xi^{(k)}$ is the tree
curvature of column $k$ (Definition~\ref{def:treecurvature}), are
nonnegative for every $s$ with $k\le s\le m+1$; as $m+1=n_0$, this is part
of Lemma~\ref{lem:smallcolumns}. And $\delta^{(k)}_m>0$, because
$d_{m,k}>0$. The proposition gives $\delta^{(k)}_j>0$
for every $j\ge m$, and in particular $d_{n,k}>0$.

Finally, for $k\ge82$ the claim is
Proposition~\ref{prop:largecolumns}.
\end{proof}

At the ten exceptional positions with
$k\ge2$ the step is negative, by Lemma~\ref{lem:fortyfive}. The eleventh is
$(2,1)$, the one position of row $2$.

We turn to the two rules. Beyond the first layer both follow from
Proposition~\ref{prop:beyondlayer} and Lemma~\ref{lem:fortyfive}. The first
corollary is the column rule in the case in which the lower position
$(n+1,k)$ lies beyond the first layer.

\begin{corollary}\label{cor:columnbeyond}
Let $k\ge1$ and $n$ be integers with
$n\ge k+1$ and $n+1\ge2k$. If $d_{n,k}>0$, then $d_{n+1,k}>0$.
\end{corollary}

\begin{proof}
Suppose first that $(n+1,k)$ is not an exceptional position. Then
Proposition~\ref{prop:beyondlayer}, applied to column $k$ and to row
$n+1\ge2k$, gives $d_{n+1,k}>0$.

Suppose now that $(n+1,k)$ is exceptional. Then $d_{n,k}<0$ by
Lemma~\ref{lem:fortyfive}, so that the hypothesis $d_{n,k}>0$ fails and
there is nothing to prove. Indeed, if $n+1=2k$ and $k\le7$, then $n=2k-1$,
and $n\ge k+1$ gives $k\ge2$, so that the lemma applies with its condition
that $n\le2k-1$ and $k\le8$. If $n+1=2k+1$ and $2\le k\le5$, then $n=2k$,
and the lemma applies with its condition that $n\le2k$ and $k\le7$.
\end{proof}

The second corollary is the diagonal rule in the case in which the upper
position $(n,k)$ lies beyond the first layer.

\begin{corollary}\label{cor:diagonalbeyond}
Let $k\ge1$ and $n\ge2k$ be integers. If $d_{n,k}\le0$, then
$d_{n+1,k+1}<0$.
\end{corollary}

\begin{proof}
If neither $n=2k$ and $k\le7$, nor $n=2k+1$ and $2\le k\le5$, then
Proposition~\ref{prop:beyondlayer} gives $d_{n,k}>0$, against the
hypothesis. Hence $n=2k$ and $k\le7$, or $n=2k+1$ and $2\le k\le5$.

Put $k'=k+1$ and $n'=n+1$. Then $k'\ge2$, and $n'\ge k'+1$ because
$n\ge2k\ge k+1$. If $n=2k$ and $k\le7$, then $n'=2k'-1$ and $k'\le8$, and
Lemma~\ref{lem:fortyfive} applies with its condition that $n'\le2k'-1$ and
$k'\le8$. If $n=2k+1$ and $2\le k\le5$, then $n'=2k'$ and $k'\le6$, and the
lemma applies with its condition that $n'\le2k'$ and $k'\le7$. In both
cases $d_{n+1,k+1}=d_{n',k'}<0$.
\end{proof}

\section{The gain of the boundary term along the diagonal}\label{sec:gain}
Proposition~\ref{prop:diagonalstep}, the diagonal rule at the positions
behind the curve in columns $k\ge265$, is proved in Section~\ref{sec:step}
from three comparisons along the diagonal, and we begin with the one for the
boundary term. At a position $(k+r,k)$ of the first layer the normalised
step is a curvature term minus the boundary term
$(a_k/k)\,\Lambda_{r+1}(k-1;k+r+1)$
(Corollary~\ref{cor:curvatureboundary}), and at the diagonal neighbour
$(k+r+1,k+1)$ the boundary term is
$(a_{k+1}/(k+1))\,\Lambda_{r+1}(k;k+r+2)$. For
Proposition~\ref{prop:diagonalstep} the boundary term has to grow along the
diagonal by more than the terms of the curvature term do, and
Proposition~\ref{prop:boundarygain} bounds the growth of its chain sum from
below by an explicit factor.

The two chain sums run over the compositions of the same integer $q=r+1$,
with the same weights of the parts, and they differ only in the weights of
the cuts. The proof of Proposition~\ref{prop:boundarygain} compares the
weights of the compositions in the first chain sum with their weights when
all cuts have the same weight, by an instance of the correlation inequality
of Fortuin, Kasteleyn and Ginibre \cite[Proposition~1]{fortuin1971} on the
sets of cut points. Section~\ref{subsec:gain-homogeneous} treats the
compositions whose cuts all have the same weight and bounds their mean
number of cuts from below (Lemma~\ref{lem:meancuts}), and
Section~\ref{subsec:gain-cutsets} writes the weighted sums over
compositions as sums over sets of cut points and recalls the correlation
inequality, whose hypothesis the proof derives
from the log-convexity of the Cayley coefficients
(Lemma~\ref{lem:cayleyconvex}).

Throughout the section $q$ denotes the integer of
Proposition~\ref{prop:boundarygain}, whose compositions the two chain sums
run over. The statements of Sections~\ref{subsec:gain-homogeneous}
and~\ref{subsec:gain-cutsets} are about the compositions of an arbitrary
integer $m$, and we write $Q$ for the number of parts of a composition.

\subsection{The homogeneous law and its mean number of cuts}
\label{subsec:gain-homogeneous}

We use the weighted sums over compositions of
Section~\ref{sec:compositions}: the weight $V_Q(m;f)$ of the compositions of
$m$ with exactly $Q$ parts for a cut profile $f$, the weight
$\Lambda_m[f]=\sum_{Q=1}^{m}V_Q(m;f)$ of all compositions of $m$, and the
numbers $T_{m,Q}=V_Q(m;\mathbf 1)$, which are nonnegative and satisfy
$T_{m,1}=a_m$.

For an integer $m\ge0$ and a real number $h$ let
\[
\tau_m(h)=\Lambda_m[h]
\]
be the total weight of the compositions of $m$ for the constant cut profile
with the value $h$. Then
\begin{equation}\label{eq:gain-total}
\tau_m(h)=\sum_{Q=1}^{m}T_{m,Q}\,h^{Q-1}.
\end{equation}
For $m\ge1$ and $h\ge0$ we have $\tau_m(h)>0$.

Let $m\ge1$ and $h\ge0$. The \emph{homogeneous law with the parameter $h$} on
the compositions of $m$ gives every composition a probability proportional to
its weight for the constant cut profile $h$. Under this law the number of
parts takes the value
$Q$ with probability $T_{m,Q}\,h^{Q-1}/\tau_m(h)$. Accordingly, the mean
number of cuts of a composition of $m$ under the homogeneous law with the
parameter $h$ is
\[
\mu_m(h)=\frac{\sum_{Q=1}^{m}(Q-1)\,T_{m,Q}\,h^{Q-1}}{\tau_m(h)}.
\]
Clearly $\mu_m(h)\ge0$. We use the definition in the form
\begin{equation}\label{eq:gain-meandef}
\mu_m(h)\,\tau_m(h)=\sum_{Q=1}^{m}(Q-1)\,T_{m,Q}\,h^{Q-1}.
\end{equation}

The mean $\mu_m(h)$ is a quotient of two sums of $m$ terms, and
Proposition~\ref{prop:boundarygain} needs a lower bound for it that is
explicit in $m$ and $h$. The bound is expressed by the function
\[
\mu^0_m(X)=\frac{\sqrt{(4+X)^2+12\,(m-1)}-(4+X)}{3},
\]
defined for an integer $m\ge1$ and a real number $X$. For $X\ge0$ the square
root is at least $4+X$, so that $\mu^0_m(X)\ge0$.

By Proposition~\ref{prop:partslaw} the number $J$ of cuts under the
homogeneous law takes the value $j$ with a probability proportional to
$(j+1)\,h^{j}B_{m,j+1}$, where $B_{m,j+1}$ is the birthday tail of
Section~\ref{sec:compositions}, and the proof of Lemma~\ref{lem:meancuts}
shows that the sequence of the numbers $h^{j}B_{m,j+1}$ is log-concave. The
proof sums the relation between consecutive terms of this sequence over
$j$, which gives an exact relation between the first two moments of $J$, and
combines it with the inequality
$\mathbb E[J(J-1)]\le(3/2)\,\mathbb E[J]^2$; the result is the inequality
$2\,(m-1)\le\tfrac32\,\mu^2+(4+X)\,\mu$ for $\mu=\mu_m(h)=\mathbb E[J]$ and
$X=m\,(1/h-1)$, and the bound of the lemma is its solution for $\mu$. For the
inequality between the moments the factor $j+1$ is read as the number of
pairs $(i,\ell)$ of nonnegative integers with $i+\ell=j$: then $J$ has the
law of the sum of the two coordinates of a random pair, and the proof
derives from the log-concavity that the two coordinates are negatively
correlated.

\begin{lemma}\label{lem:meancuts}
For every integer $m\ge1$ and every real number $h$ with $0<h\le1$, the mean
number of cuts $\mu_m(h)$ of a composition of $m$ under the homogeneous law
with the parameter $h$ satisfies
\[
\mu^0_m\bigl(m\,(1/h-1)\bigr)\le\mu_m(h).
\]
\end{lemma}

\begin{proof}
Fix $m$ and $h$, and write $\mu=\mu_m(h)$ and $X=m\,(1/h-1)$. For an integer
$j\ge0$ let
\[
v_j=h^{j}\,B_{m,j+1}=h^{j}\prod_{i=0}^{j}\Bigl(1-\frac im\Bigr),
\]
where $B_{m,j}$ is the birthday tail of Section~\ref{sec:compositions}. Since
$h>0$, we have $v_j>0$ for $0\le j\le m-1$, while $v_j=0$ for $j\ge m$.
Moreover,
\begin{equation}\label{eq:gain-ratio}
m\,v_{j+1}=h\,(m-1-j)\,v_j\qquad\text{for every integer }j\ge0.
\end{equation}
We put
\[
M_0=\sum_{j=0}^{m-1}(j+1)\,v_j,\quad
M_1=\sum_{j=0}^{m-1}j\,(j+1)\,v_j,\quad
M_2=\sum_{j=0}^{m-1}j^2\,(j+1)\,v_j.
\]
By Proposition~\ref{prop:partslaw} we have $T_{m,Q}=a_m\,Q\,B_{m,Q}$ for
$1\le Q\le m$. With $Q=j+1$, the formulas \eqref{eq:gain-total} and
\eqref{eq:gain-meandef} become $\tau_m(h)=a_m\,M_0$ and
$\mu\,\tau_m(h)=a_m\,M_1$. Since $a_m>0$ and $\tau_m(h)>0$, it follows that
\begin{equation}\label{eq:gain-moments}
M_0>0\qquad\text{and}\qquad M_1=\mu\,M_0.
\end{equation}

We first show that
\begin{equation}\label{eq:gain-identity}
m\,M_1=h\,\bigl(2\,(m-1)\,M_0+(m-3)\,M_1-M_2\bigr).
\end{equation}
The term $j=0$ of $M_1$ vanishes, and $v_m=0$. Shifting the index of
summation by one, we therefore have
\[
M_1=\sum_{j=0}^{m-1}(j+1)(j+2)\,v_{j+1}.
\]
By \eqref{eq:gain-ratio}, for every $j$,
\begin{align*}
m\,(j+1)(j+2)\,v_{j+1}&=h\,(j+1)(j+2)(m-1-j)\,v_j\\
&=h\,(j+1)\,v_j\,\bigl(2\,(m-1)+(m-3)\,j-j^2\bigr),
\end{align*}
and the sum of these equations over $0\le j\le m-1$ is
\eqref{eq:gain-identity}.

Next we show that
\begin{equation}\label{eq:gain-correlation}
\tfrac16\,(M_2-M_1)\,M_0\le\tfrac14\,M_1^2.
\end{equation}
The sequence $(v_j)$ is log-concave in the following form: for all integers
$s$, $t$ and $d$ with $0\le s\le t$ and $d\ge0$,
\begin{equation}\label{eq:gain-logconcave}
v_s\,v_{t+d}\le v_{s+d}\,v_t.
\end{equation}
Indeed, if $t+d\ge m$, the left-hand side is $0$ and the right-hand side is
nonnegative. If $t+d\le m-1$, the four numbers are positive, and
\[
\frac{v_{t+d}}{v_t}=h^{d}\prod_{i=1}^{d}\Bigl(1-\frac{t+i}{m}\Bigr)
\le h^{d}\prod_{i=1}^{d}\Bigl(1-\frac{s+i}{m}\Bigr)=\frac{v_{s+d}}{v_s},
\]
because the factors of both products are positive and $s\le t$.

We use a weighted form of Chebyshev's sum inequality. Let $\chi$, $\psi$ and
$\varphi$ be real functions on the set $\{0,\dots,m-1\}$ such that
\[
\chi(j')\,\psi(j)\le\chi(j)\,\psi(j')\quad\text{and}\quad
\varphi(j)\le\varphi(j')\qquad\text{whenever }j\le j'.
\]
Then
\begin{equation}\label{eq:gain-chebyshev}
\Bigl(\sum_{j}\chi(j)\,\varphi(j)\Bigr)\Bigl(\sum_{j}\psi(j)\Bigr)
\le\Bigl(\sum_{j}\psi(j)\,\varphi(j)\Bigr)\Bigl(\sum_{j}\chi(j)\Bigr),
\end{equation}
all sums running over $0\le j\le m-1$. When $\chi$ and $\psi$ are positive,
the hypothesis says that the quotient $\chi/\psi$ does not increase, and the
conclusion says that the mean of the nondecreasing function $\varphi$ for the
weights $\chi$ is at most its mean for the weights $\psi$; but no assumption
on the signs is needed. The inequality \eqref{eq:gain-chebyshev} follows by
the usual symmetrisation: if $D$ denotes its left-hand side minus its
right-hand side, then
\[
2D=\sum_{j}\sum_{j'}\bigl(\chi(j)\,\psi(j')-\chi(j')\,\psi(j)\bigr)
\bigl(\varphi(j)-\varphi(j')\bigr),
\]
and by the two hypotheses no term of this sum is positive.

We apply \eqref{eq:gain-chebyshev} twice. First let $i$ and $i'$ be integers
with $0\le i\le i'\le m-1$, and take $\chi(\ell)=v_{i'+\ell}$,
$\psi(\ell)=v_{i+\ell}$ and $\varphi(\ell)=\ell$. For $\ell\le\ell'$ the
inequality \eqref{eq:gain-logconcave} with $s=i+\ell$, $t=i'+\ell$ and
$d=\ell'-\ell$ reads
$v_{i+\ell}\,v_{i'+\ell'}\le v_{i+\ell'}\,v_{i'+\ell}$, and this is the
hypothesis on $\chi$ and $\psi$. Hence
\[
\Bigl(\sum_{\ell}\ell\,v_{i'+\ell}\Bigr)\Bigl(\sum_{\ell}v_{i+\ell}\Bigr)
\le\Bigl(\sum_{\ell}\ell\,v_{i+\ell}\Bigr)\Bigl(\sum_{\ell}v_{i'+\ell}\Bigr)
\qquad(0\le i\le i'\le m-1),
\]
the sums running over $0\le\ell\le m-1$. This in turn is the hypothesis of
\eqref{eq:gain-chebyshev} for the functions
$\chi(i)=\sum_{\ell}\ell\,v_{i+\ell}$ and $\psi(i)=\sum_{\ell}v_{i+\ell}$,
and with $\varphi(i)=i$ we obtain
\begin{equation}\label{eq:gain-double}
\Bigl(\sum_{i,\ell}i\,\ell\,v_{i+\ell}\Bigr)\Bigl(\sum_{i,\ell}v_{i+\ell}\Bigr)
\le\Bigl(\sum_{i,\ell}i\,v_{i+\ell}\Bigr)
\Bigl(\sum_{i,\ell}\ell\,v_{i+\ell}\Bigr),
\end{equation}
where the double sums run over $0\le i\le m-1$ and $0\le\ell\le m-1$.

It remains to express the four double sums by $M_0$, $M_1$ and $M_2$. Since
$v_j=0$ for $j\ge m$, only the pairs $(i,\ell)$ with $i+\ell\le m-1$
contribute, and for every $j\le m-1$ the pairs with $i+\ell=j$ are the $j+1$
pairs $(i,j-i)$ with $0\le i\le j$. Since $\sum_{i=0}^{j}i=j(j+1)/2$ and
$\sum_{i=0}^{j}i\,(j-i)=(j-1)j(j+1)/6$, we find
\begin{gather*}
\sum_{i,\ell}v_{i+\ell}=M_0,\qquad
\sum_{i,\ell}i\,v_{i+\ell}=\sum_{i,\ell}\ell\,v_{i+\ell}=\tfrac12\,M_1,\\
\sum_{i,\ell}i\,\ell\,v_{i+\ell}=\tfrac16\,(M_2-M_1).
\end{gather*}
With these values \eqref{eq:gain-double} is \eqref{eq:gain-correlation}.

We now multiply \eqref{eq:gain-correlation} by $6h$ and eliminate $h\,M_2$
by means of \eqref{eq:gain-identity}. The result is
\[
0\le\tfrac32\,h\,M_1^2+(4h+m-mh)\,M_1\,M_0-2\,(m-1)\,h\,M_0^2.
\]
We divide by $h\,M_0^2$, which is positive by \eqref{eq:gain-moments} and by
the hypothesis $h>0$, and use $M_1=\mu\,M_0$ and $(4h+m-mh)/h=4+X$. This
gives
\[
2\,(m-1)\le\tfrac32\,\mu^2+(4+X)\,\mu.
\]
Completing the square, we obtain
$(4+X)^2+12\,(m-1)\le(3\mu+4+X)^2$. Since $X\ge0$ by the hypothesis $h\le1$,
and $\mu\ge0$, taking square roots gives
\[
\sqrt{(4+X)^2+12\,(m-1)}\le3\mu+4+X,
\]
which is the inequality $\mu^0_m(X)\le\mu$.
\end{proof}

\subsection{Sets of cut points and the correlation inequality}
\label{subsec:gain-cutsets}

Let $m\ge1$ be an integer. The cut points of a composition of $m$ are the
partial sums of its parts that lie strictly between $0$ and $m$. Every set
$S$ of integers between $1$ and $m-1$ is the set of cut points of exactly one
composition of $m$: if $0=s_0<s_1<\dots<s_Q=m$ are the elements of
$S\cup\{0,m\}$, this composition has the $Q$ parts $s_i-s_{i-1}$ with
$1\le i\le Q$. Let
\[
\Pi_m(S)=\prod_{i=1}^{Q}a_{s_i-s_{i-1}}
\]
be the product of the Cayley coefficients of its parts, which is positive.
The empty set belongs to the composition with the single part $m$, and
$\Pi_m(\emptyset)=a_m$.

For a cut profile $f$, the weight of the composition with the set $S$ of
cut points is $\Pi_m(S)$ times the product of $f(s)$ over $s\in S$. So
$V_Q(m;f)$ is the contribution of the sets with $Q-1$ elements to the sum
over all sets, and for every cut profile $f$
\begin{equation}\label{eq:gain-cutsets}
\Lambda_m[f]=\sum_{S\subseteq\{1,\dots,m-1\}}\Pi_m(S)\prod_{s\in S}f(s).
\end{equation}
In particular $T_{m,Q}$ is the sum of $\Pi_m(S)$ over the sets
$S$ with $Q-1$ elements, and \eqref{eq:gain-total} and
\eqref{eq:gain-meandef} take the form
\begin{equation}\label{eq:gain-cutcount}
\tau_m(h)=\sum_{S}\Pi_m(S)\,h^{|S|},\qquad
\mu_m(h)\,\tau_m(h)=\sum_{S}|S|\,\Pi_m(S)\,h^{|S|}
\end{equation}
for every real number $h\ge0$, where $S$ runs over the subsets of
$\{1,\dots,m-1\}$ and $|S|$ is the number of elements of $S$. Thus the
homogeneous law with the parameter $h$ gives the set $S$ the probability
$\Pi_m(S)\,h^{|S|}/\tau_m(h)$, and $\mu_m(h)$ is the mean of $|S|$.

The proof of Proposition~\ref{prop:boundarygain} applies the correlation
inequality to weights of this kind.
A function $\theta$ with positive values on the set of all subsets of a
finite set $\Omega$ is called \emph{log-supermodular} if
$\theta(S)\,\theta(S')\le\theta(S\cap S')\,\theta(S\cup S')$ for all subsets
$S$ and $S'$ of $\Omega$.

For log-supermodular functions we use the correlation inequality of Fortuin,
Kasteleyn and Ginibre \cite[Proposition~1]{fortuin1971}. Their result
concerns a positive function on a finite distributive lattice that satisfies
the inequality above with the intersection and the union replaced by the
meet and the join of two elements. It states that two nondecreasing real
functions on the lattice are nonnegatively correlated when this function is
taken as the weight: the mean of their product is at least the product of
their means. The subsets of a finite set $\Omega$, ordered by inclusion, form
a distributive lattice whose meet and join are the intersection and the
union. For this lattice, and after multiplication by the square of the total
weight, the result reads as follows. Let $\theta$ be a log-supermodular
function with positive values on the subsets of $\Omega$, and let $\varphi_1$
and $\varphi_2$ be real functions on the subsets of $\Omega$ that do not
decrease under inclusion. Then
\begin{multline}\label{eq:gain-fkg}
\Bigl(\sum_{S}\theta(S)\,\varphi_1(S)\,\varphi_2(S)\Bigr)
\Bigl(\sum_{S}\theta(S)\Bigr)\\
\ge\Bigl(\sum_{S}\theta(S)\,\varphi_1(S)\Bigr)
\Bigl(\sum_{S}\theta(S)\,\varphi_2(S)\Bigr),
\end{multline}
where all sums run over the subsets $S$ of $\Omega$.

\subsection{The gain of the chain sum}\label{subsec:gain-proof}

The proof of Proposition~\ref{prop:boundarygain} first shows that the
weight of the composition with the set $S$ of cut points grows, from the
first chain sum to the second, by at least the factor
$(1+x)^{\sigma(S)/q}$, where $x$ is the number of the proposition and
$\sigma(S)$ is the sum of $q+1-s$ over the cut points $s$, so that by
Jensen's inequality the chain sum grows by at least $1+x$ raised to the mean
of $\sigma/q$ for the weights of the first chain sum. It then bounds the
weight of every cut of the first chain sum from below by the weight $h$ of
the cut at position $1$, and it uses the correlation inequality to show that
this mean of $\sigma$ is at least the mean of $\sigma$ under the homogeneous
law with the parameter $h$. It computes the mean of $\sigma$ under
the homogeneous law as $(q+2)/2$ times the mean number of cuts, by reading
the compositions backwards, and Lemma~\ref{lem:meancuts} bounds the mean
number of cuts from below.

\begin{proposition}\label{prop:boundarygain}
Let $k\ge2$ and $q\ge1$ be integers, and put
\[
x=\frac{q}{(k-1)(k+q)},\qquad
E_0=\frac{q+2}{2q}\;\mu^0_q\Bigl(\frac{q^2}{k-1}\Bigr),
\]
with the function $\mu^0_q$ of Section~\ref{subsec:gain-homogeneous}. Then
the chain sums satisfy
\[
(1+x)^{E_0}\,\Lambda_q(k-1;k+q)\le\Lambda_q(k;k+q+1).
\]
\end{proposition}

\begin{proof}
For $q=1$ both chain sums are equal to $a_1=1$, and $E_0=0$ because
$\mu^0_1(X)=0$ for $X\ge0$; the two sides are then equal. Let $q\ge2$. For
an integer $s$ with $0<s<q$ we write
\[
\lambda_s=\frac{k-1}{k+q-s},\qquad\lambda'_s=\frac{k}{k+q+1-s}
\]
for the values at $s$ of the threshold profiles $\lambda^{k-1,k+q}$ and
$\lambda^{k,k+q+1}$ of Section~\ref{sec:compositions}; these numbers are
positive, since $k\ge2$. In this
proof $S$ denotes a subset of $\{1,\dots,q-1\}$, a sum over $S$ runs over all
these subsets, and
\[
\lambda_S=\prod_{s\in S}\lambda_s,\qquad
\lambda'_S=\prod_{s\in S}\lambda'_s,\qquad
\sigma(S)=\sum_{s\in S}(q+1-s).
\]
We write $L$ and $L'$ for the two chain sums. By \eqref{eq:gain-cutsets}
with $m=q$,
\begin{equation}\label{eq:gain-twosums}
\begin{split}
L&=\Lambda_q(k-1;k+q)=\sum_{S}\Pi_q(S)\,\lambda_S,\\
L'&=\Lambda_q(k;k+q+1)=\sum_{S}\Pi_q(S)\,\lambda'_S,
\end{split}
\end{equation}
and all terms of both sums are positive.

We first compare the weights of one cut in the two sums. For $0<s<q$ we have
\[
\frac{\lambda'_s}{\lambda_s}=\frac{k\,(k+q-s)}{(k-1)(k+q+1-s)}
=1+\frac{q+1-s}{(k-1)(k+q+1-s)}.
\]
Because $s\ge1$, the denominator of the last fraction is at most
$(k-1)(k+q)$, so that $\lambda'_s/\lambda_s\ge1+x\,(q+1-s)/q$. The number
$(q+1-s)/q$ lies between $0$ and $1$, and for an exponent in this range
Bernoulli's inequality gives $(1+x)^{(q+1-s)/q}\le1+x\,(q+1-s)/q$. Hence
\[
\lambda_s\,(1+x)^{(q+1-s)/q}\le\lambda'_s\qquad(0<s<q).
\]
We take the product of these inequalities over the elements $s$ of a set $S$
and multiply by $\Pi_q(S)$. This gives
\[
\Pi_q(S)\,\lambda_S\,(1+x)^{\sigma(S)/q}\le\Pi_q(S)\,\lambda'_S,
\]
and by \eqref{eq:gain-twosums} the sum over $S$ is
\begin{equation}\label{eq:gain-setgain}
\sum_{S}\Pi_q(S)\,\lambda_S\,(1+x)^{\sigma(S)/q}\le L'.
\end{equation}

Now let $\bar\sigma$ be the mean of $\sigma$ for the weights of the first
chain sum, so that
\[
\bar\sigma\,L=\sum_{S}\Pi_q(S)\,\lambda_S\,\sigma(S).
\]
The function $t\mapsto(1+x)^{t}$ is convex. We apply Jensen's inequality for
it, in the form recalled in Section~\ref{sec:prelim} in which the weights
need not be normalised, to the weights $\Pi_q(S)\,\lambda_S$ and the
exponents $\sigma(S)/q$, whose mean is $\bar\sigma/q$, and obtain
\begin{equation}\label{eq:gain-jensen}
(1+x)^{\bar\sigma/q}\,L\le\sum_{S}\Pi_q(S)\,\lambda_S\,(1+x)^{\sigma(S)/q}.
\end{equation}

We next compute the mean of $\sigma$ when all cuts have the same weight. For
every real number $h\ge0$,
\begin{equation}\label{eq:gain-reflection}
\sum_{S}\Pi_q(S)\,h^{|S|}\,\sigma(S)=\frac{q+2}{2}\,\mu_q(h)\,\tau_q(h).
\end{equation}
To see this, let $S^{*}=\{q-s:s\in S\}$, so that $S\mapsto S^{*}$ is an
involution of the set of subsets of $\{1,\dots,q-1\}$. The
composition with the cut points $S^{*}$ has the parts of the composition
with the cut points $S$ in the reverse order, so that
$\Pi_q(S^{*})=\Pi_q(S)$; moreover $|S^{*}|=|S|$ and
$\sigma(S^{*})=\sum_{s\in S}(s+1)$. Replacing $S$ by $S^{*}$ in the sum on
the left of \eqref{eq:gain-reflection}, we see that this sum is equal to the
sum of $\Pi_q(S)\,h^{|S|}\sum_{s\in S}(s+1)$ over $S$. Since
$(q+1-s)+(s+1)=q+2$, half the sum of the two expressions is
\[
\frac{q+2}{2}\sum_{S}|S|\,\Pi_q(S)\,h^{|S|},
\]
and this is the right-hand side of \eqref{eq:gain-reflection} by
\eqref{eq:gain-cutcount}.

Next we show that, for all subsets $S$ and $S'$ of $\{1,\dots,q-1\}$,
\begin{equation}\label{eq:gain-supermodular}
\Pi_q(S)\,\Pi_q(S')\le\Pi_q(S\cap S')\,\Pi_q(S\cup S').
\end{equation}
To prove it, we first consider the effect of one additional cut point. Let
$S\subseteq\{1,\dots,q-1\}$, and let $u$ be an integer with $0<u<q$ that does
not belong to $S$. Among the elements of $S\cup\{0,q\}$ let $u-p$ be the
largest one below $u$ and $u+p'$ the smallest one above $u$, so that $p\ge1$
and $p'\ge1$. The composition with the cut points $S\cup\{u\}$ arises from the
composition with the cut points $S$ when the part $p+p'$ is cut into the two
parts $p$ and $p'$, and therefore
\begin{equation}\label{eq:gain-onecut}
\Pi_q(S\cup\{u\})=\frac{a_p\,a_{p'}}{a_{p+p'}}\,\Pi_q(S).
\end{equation}
The factor depends on $p$ and $p'$ only, and it does not increase when one
of them grows: for all integers $p\ge1$, $p'\ge1$ and $d\ge0$,
\begin{equation}\label{eq:gain-pieces}
\frac{a_{p+d}\,a_{p'}}{a_{p+d+p'}}\le\frac{a_p\,a_{p'}}{a_{p+p'}}.
\end{equation}
Indeed, after clearing denominators and cancelling $a_{p'}$ this is the
inequality
\[
a_{p+d}\,a_{p+p'}\le a_p\,a_{p+p'+d},
\]
which is Lemma~\ref{lem:cayleyconvex} for indices $p\le p+p'$ and the shift
$d$. The factor is symmetric in $p$ and $p'$, so that the same holds when $p'$
grows.

Now let $S_0\subseteq S_1\subseteq\{1,\dots,q-1\}$, and let $u$ be an integer
with $0<u<q$ that does not belong to $S_1$. Let $u-p_0$ and $u+p_0'$ be the
nearest elements of $S_0\cup\{0,q\}$ below and above $u$, and let $u-p_1$
and $u+p_1'$ be the nearest elements of $S_1\cup\{0,q\}$ below and above
$u$. Since $S_0\cup\{0,q\}$ is contained in
$S_1\cup\{0,q\}$, we have $p_1\le p_0$ and $p_1'\le p_0'$. By
\eqref{eq:gain-pieces}, applied in the first argument with the shift
$p_0-p_1$ and then in the second with the shift $p_0'-p_1'$,
\[
\frac{a_{p_0}\,a_{p_0'}}{a_{p_0+p_0'}}
\le\frac{a_{p_1}\,a_{p_0'}}{a_{p_1+p_0'}}
\le\frac{a_{p_1}\,a_{p_1'}}{a_{p_1+p_1'}},
\]
and by \eqref{eq:gain-onecut} this says that
\begin{equation}\label{eq:gain-increment}
\frac{\Pi_q(S_0\cup\{u\})}{\Pi_q(S_0)}\le\frac{\Pi_q(S_1\cup\{u\})}{\Pi_q(S_1)}
\qquad(S_0\subseteq S_1,\ u\notin S_1).
\end{equation}
In words, the factor by which a new cut point multiplies $\Pi_q$ does not
decrease when further cut points are present.

We prove \eqref{eq:gain-supermodular} by induction on the number of elements
of $S\setminus S'$. If $S\subseteq S'$, the two sides are equal. Otherwise
choose $u\in S\setminus S'$ and put $S_u=S\setminus\{u\}$. Then
$S_u\cap S'=S\cap S'$, and the induction hypothesis for the sets $S_u$ and
$S'$ reads
\[
\Pi_q(S_u)\,\Pi_q(S')\le\Pi_q(S\cap S')\,\Pi_q(S_u\cup S').
\]
Since $u$ does not belong to $S_u\cup S'$, we can apply
\eqref{eq:gain-increment} with $S_0=S_u$ and $S_1=S_u\cup S'$, and we obtain
\begin{align*}
\Pi_q(S)\,\Pi_q(S')
&=\frac{\Pi_q(S_u\cup\{u\})}{\Pi_q(S_u)}\cdot\Pi_q(S_u)\,\Pi_q(S')\\
&\le\frac{\Pi_q(S\cup S')}{\Pi_q(S_u\cup S')}\cdot
\Pi_q(S\cap S')\,\Pi_q(S_u\cup S')\\
&=\Pi_q(S\cap S')\,\Pi_q(S\cup S').
\end{align*}
This proves \eqref{eq:gain-supermodular}.

We now compare $\bar\sigma$ with the mean of $\sigma$ for a constant weight.
From now on let $h=(k-1)/(k-1+q)$, so that $0<h<1$. For $0<s<q$ we have
$\lambda_s\ge h$, because $k+q-s\le k-1+q$. We apply the correlation
inequality \eqref{eq:gain-fkg} on the set $\Omega=\{1,\dots,q-1\}$ to
\[
\theta(S)=\Pi_q(S)\,h^{|S|},\qquad\varphi_1(S)=\sigma(S),\qquad
\varphi_2(S)=\prod_{s\in S}\frac{\lambda_s}{h}.
\]
The function $\theta$ has positive values, and it is log-supermodular by
\eqref{eq:gain-supermodular}, because $|S|+|S'|=|S\cap S'|+|S\cup S'|$.
Neither $\sigma$ nor $\varphi_2$
decreases under inclusion, because the summands $q+1-s$ are positive and the
factors $\lambda_s/h$ are at least $1$. Since
$\theta(S)\,\varphi_2(S)=\Pi_q(S)\,\lambda_S$, the
inequality \eqref{eq:gain-fkg} reads
\[
\Bigl(\sum_{S}\Pi_q(S)\,\lambda_S\,\sigma(S)\Bigr)\,\tau_q(h)
\ge\Bigl(\sum_{S}\Pi_q(S)\,h^{|S|}\,\sigma(S)\Bigr)\,L,
\]
where the two sums without the factor $\sigma$ have been evaluated by
\eqref{eq:gain-cutcount} and \eqref{eq:gain-twosums}. The sum on the left is
$\bar\sigma\,L$ by the definition of $\bar\sigma$, and the sum on the right
is given by \eqref{eq:gain-reflection}. Dividing by $\tau_q(h)\,L>0$, we
obtain
\begin{equation}\label{eq:gain-meanbound}
\bar\sigma\ge\frac{q+2}{2}\,\mu_q(h).
\end{equation}

Finally, we apply Lemma~\ref{lem:meancuts} with $m=q$ and with the number
$h$ just chosen. Since $1/h-1=q/(k-1)$, it gives
$\mu^0_q\bigl(q^2/(k-1)\bigr)\le\mu_q(h)$. Together with
\eqref{eq:gain-meanbound} and the definition of $E_0$ this shows that
$E_0\le\bar\sigma/q$, and hence $(1+x)^{E_0}\le(1+x)^{\bar\sigma/q}$, because
$1+x\ge1$. By \eqref{eq:gain-jensen} and \eqref{eq:gain-setgain} we conclude
that
\[
(1+x)^{E_0}\,L\le(1+x)^{\bar\sigma/q}\,L
\le\sum_{S}\Pi_q(S)\,\lambda_S\,(1+x)^{\sigma(S)/q}\le L'.\qedhere
\]
\end{proof}

\section{The growth of the curvature sum along the diagonal}\label{sec:growth}
We turn to the second of the three comparisons along the diagonal
from which Proposition~\ref{prop:diagonalstep} is obtained in
Section~\ref{sec:step}. By Corollary~\ref{cor:curvatureboundary}, the
curvature term of a position $(k+r,k)$ of the first layer is a combination
of the curvature sums $\mathsf C_m(k;N)$ with $1\le m\le r$, where
$N=k+r+1$, and the curvature term of the diagonal neighbour $(k+r+1,k+1)$ is
the corresponding combination of the sums $\mathsf C_m(k+1;N+1)$. Here we
compare the two curvature sums for one $m\ge2$ with $401\,m^2\le12\,k^2$:
by Proposition~\ref{prop:curvaturegrowth}, a diagonal step, from a
position to its diagonal neighbour, multiplies
$\mathsf C_m$ by at most $k/(k+1)$ times an exponential factor. The
curvature sum is a sum over the compositions of $m$ in which every cut
carries a factor of the threshold profile, and the exponent of the
exponential factor contains a mean number of cuts of these compositions;
we bound this mean through the birthday law of $m$
(Lemma~\ref{lem:momentquotient}).

\subsection{Reflection and log-convexity in the cut profile}
\label{subsec:growth-profile}

We work with the weights $V_Q(m;f)$ of Section~\ref{sec:compositions}, the
total weight of the compositions of $m$ with exactly $Q$ parts for a cut
profile $f\colon\mathbb N\to\mathbb R$; we write $Q$ for the number of
parts. Recall that $V_0(m;f)=0$ and that
for integers $Q\ge1$ and $m\ge0$
\begin{equation}\label{eq:growth-chains}
V_Q(m;f)=\sum_{0=s_0<\dots<s_Q=m}\ \prod_{i=1}^{Q}a_{s_i-s_{i-1}}\
\prod_{i=1}^{Q-1}f(s_i),
\end{equation}
where the sum runs over the sequences of integers
$0=s_0<s_1<\dots<s_Q=m$. Such a sequence is a composition of $m$ with the
parts $s_i-s_{i-1}$ and the cut points $s_1,\dots,s_{Q-1}$. The cut points
of a composition of $m$ lie strictly
between $0$ and $m$, and the integers $s$ with $0<s<m$ are the
\emph{interior points} of $m$.

We use throughout the section the following elementary properties of these
weights, which Section~\ref{sec:compositions} records. The weight
$V_Q(m;f)$ depends only on the values of $f$ at the interior points of
$m$. It is nonnegative for a nonnegative profile: if $f(s)\ge0$ at every
interior point $s$ of $m$, then $V_Q(m;f)\ge0$. It is homogeneous in the
threshold of a threshold profile: for all real numbers $y$ and $N$ and
every $Q\ge1$
\begin{equation}\label{eq:growth-homogeneous}
V_Q(m;\lambda^{y,N})=y^{\,Q-1}\,V_Q(m;\lambda^{1,N}).
\end{equation}
For the constant profile
$\mathbf 1$ the weight is $T_{m,Q}=V_Q(m;\mathbf 1)$, so that
$T_{m,Q}\ge0$, and we write $\rho\,\mathbf 1$ for the constant profile with
the value $\rho$.

The comparison of two profiles in Section~\ref{subsec:growth-comparison}
rests on two properties of the weights as functions of the profile. The
first is a symmetry. Read from its right end, a
composition of $m$ has the same parts in the reverse order, and a cut point
$s$ has become the cut point $m-s$; the weights therefore do not change when
the cut profile is reflected in the same way.

\begin{lemma}\label{lem:reflection}
Let $f$ be a cut profile, let $Q\ge0$ and $m\ge0$ be integers, and let
$\tilde f$ be a cut profile with $\tilde f(s)=f(m-s)$ for every integer $s$
with $0<s<m$. Then
\[
V_Q(m;\tilde f)=V_Q(m;f).
\]
\end{lemma}

\begin{proof}
For $Q=0$ both sides are $0$. Let $Q\ge1$. The map that sends a sequence
$0=s_0<s_1<\dots<s_Q=m$ to the sequence $s'_i=m-s_{Q-i}$, $0\le i\le Q$, is
an involution of the set of these sequences, and it reverses the order of
the parts. The cut points of the image are the integers $s'_i$ with
$1\le i\le Q-1$. They satisfy $0<s'_i<m$, so that
$\tilde f(s'_i)=f(m-s'_i)=f(s_{Q-i})$ by the assumption on $\tilde f$.
Hence the term of the image in the sum \eqref{eq:growth-chains} for
$\tilde f$ is equal to the term of the original sequence in the sum for
$f$, and the claim follows.
\end{proof}

The second property is the log-convexity of the weight in the profile. By
\eqref{eq:growth-chains}, $V_Q(m;f)$ is a polynomial in the values of $f$ at
the interior points of $m$, and the coefficients of this polynomial are
nonnegative. For such a function H\"older's inequality gives log-convexity
along geometric interpolation. We use H\"older's inequality for finite sums
in the following form: if $\theta_1>0$ and $\theta_2>0$ are real numbers
with $\theta_1+\theta_2=1$, and $x_i\ge0$ and $y_i\ge0$ are real numbers
indexed by a finite set $I$, then
\begin{equation}\label{eq:growth-hoelder}
\sum_{i\in I}x_i^{\theta_1}\,y_i^{\theta_2}
\le\Bigl(\sum_{i\in I}x_i\Bigr)^{\theta_1}
\Bigl(\sum_{i\in I}y_i\Bigr)^{\theta_2}.
\end{equation}
This is the usual inequality with the conjugate exponents $1/\theta_1$ and
$1/\theta_2$, applied to the numbers $x_i^{\theta_1}$ and $y_i^{\theta_2}$.

\begin{lemma}\label{lem:profileconvex}
Let $\theta_1>0$ and $\theta_2>0$ be real numbers with
$\theta_1+\theta_2=1$, let $Q\ge0$ and $m\ge0$ be integers, and let $f$ and
$g$ be cut profiles such that $f(s)\ge0$ and $g(s)\ge0$ for every integer
$s$ with $0<s<m$. Then
\[
V_Q\bigl(m;f^{\theta_1}g^{\theta_2}\bigr)
\le V_Q(m;f)^{\theta_1}\,V_Q(m;g)^{\theta_2},
\]
where $f^{\theta_1}g^{\theta_2}$ denotes a cut profile that takes the value
$f(s)^{\theta_1}\,g(s)^{\theta_2}$ at every integer $s$ with $0<s<m$.
\end{lemma}

\begin{proof}
For $Q=0$ both sides are $0$. Let $Q\ge1$, and for a sequence
$c=(s_0,\dots,s_Q)$ as in \eqref{eq:growth-chains} let $x_c$ and $y_c$ be
its terms in the sums for the profiles $f$ and $g$. They are nonnegative,
because $f$ and $g$ are nonnegative at the cut points of $c$, and since
$\theta_1+\theta_2=1$, the term of $c$ in the sum for the profile
$f^{\theta_1}g^{\theta_2}$ is $x_c^{\theta_1}\,y_c^{\theta_2}$. The claim
is therefore \eqref{eq:growth-hoelder} for the numbers $x_c$ and $y_c$,
indexed by the sequences $c$.
\end{proof}

\subsection{The comparison of two profiles}\label{subsec:growth-comparison}

We now compare one weighted sum over the compositions of $m$ for two cut
profiles. Let a weight $v_Q\ge0$ be given for each number $Q$ of parts with
$1\le Q\le m$, and for a cut profile $\phi$ that is nonnegative at the
interior points of $m$ put
\begin{equation}\label{eq:growth-sum}
S(\phi)=\sum_{Q=1}^{m}v_Q\,V_Q(m;\phi).
\end{equation}
Since every $V_Q(m;\phi)$ is nonnegative, this is a nonnegative number.
Lemma~\ref{lem:profilecomparison} bounds $S(f)$ by a multiple of $S(g)$ for
two profiles $f$ and $g$, without assuming that $f\le g$. In
Section~\ref{subsec:growth-step} it is applied with $f$ the profile of the
diagonal neighbour of a position and $g$ the profile of the position.
The proof of the lemma needs a lower bound for the profile $f$;
hypothesis~(3) asks for one only for the products $f(s)\,f(m-s)$. A lower
bound of these products by the square of a constant serves in the proof, by
Lemmas~\ref{lem:reflection} and~\ref{lem:profileconvex}, as a constant
lower bound of $f$ would; the proof of
Proposition~\ref{prop:curvaturegrowth} verifies it for its profiles.

\begin{lemma}\label{lem:profilecomparison}
Let $k\ge1$ and $m\ge0$ be integers, let $v_Q$, $1\le Q\le m$, be real
numbers, let $f$ and $g$ be cut profiles, and let $\rho$ and $\mu$ be real
numbers. Assume the following.
\begin{enumerate}
\item One has $v_Q\ge0$ for $1\le Q\le m$.
\item For every integer $s$ with $0<s<m$ one has $f(s)\ge0$, $g(s)\ge0$ and
$f(s)^{k+1}\le g(s)^{k}$.
\item One has $0<\rho\le1$, and $\rho^2\le f(s)\,f(m-s)$ for every integer
$s$ with $0<s<m$.
\item The sum of $v_Q\,T_{m,Q}$ over $1\le Q\le m$ is positive,
and
\[
\sum_{Q=1}^{m}v_Q\,T_{m,Q}\,(Q-1)\le\mu\sum_{Q=1}^{m}v_Q\,T_{m,Q}.
\]
\end{enumerate}
Then
\[
\Bigl(\sum_{Q=1}^{m}v_Q\,V_Q(m;f)\Bigr)^{k}
\le\rho^{-\mu}\,\Bigl(\sum_{Q=1}^{m}v_Q\,V_Q(m;g)\Bigr)^{k}.
\]
\end{lemma}

\begin{proof}
By hypotheses~(1) and~(2) the numbers $S(f)$, $S(g)$ and $S(\mathbf 1)$ of
\eqref{eq:growth-sum} are defined and nonnegative, and $S(\mathbf 1)>0$ by
hypothesis~(4). We first show that
\begin{equation}\label{eq:growth-chord}
S(f)^{k+1}\le S(\mathbf 1)\,S(g)^{k}.
\end{equation}
Let $\theta=k/(k+1)$, so that $0<\theta<1$. At an interior point $s$ of
$m$, hypothesis~(2) gives $f(s)\le g(s)^{\theta}=1^{1-\theta}\,g(s)^{\theta}$.
By \eqref{eq:growth-chains}, $V_Q(m;\phi)$ is a polynomial with
nonnegative coefficients in the values of a cut profile $\phi$ at the
interior points of $m$. Hence, for cut profiles $\phi$ and $\psi$ with
$0\le\phi(s)\le\psi(s)$ at every interior point $s$ of $m$,
\begin{equation}\label{eq:growth-monotone}
0\le V_Q(m;\phi)\le V_Q(m;\psi).
\end{equation}
By
\eqref{eq:growth-monotone}, and by Lemma~\ref{lem:profileconvex} for the
profiles $\mathbf 1$ and $g$ with the exponents $1-\theta$ and $\theta$,
\[
V_Q(m;f)\le V_Q\bigl(m;\mathbf 1^{1-\theta}g^{\theta}\bigr)
\le T_{m,Q}^{\,1-\theta}\,V_Q(m;g)^{\theta}\qquad(1\le Q\le m).
\]
We multiply by $v_Q=v_Q^{1-\theta}\,v_Q^{\theta}$, sum over $Q$, and apply
\eqref{eq:growth-hoelder} to the numbers $x_Q=v_Q\,T_{m,Q}$ and
$y_Q=v_Q\,V_Q(m;g)$:
\[
S(f)\le\sum_{Q=1}^{m}x_Q^{1-\theta}\,y_Q^{\theta}
\le S(\mathbf 1)^{1-\theta}\,S(g)^{\theta}.
\]
Raising this to the power $k+1$ gives \eqref{eq:growth-chord}.

Next we show that
\begin{equation}\label{eq:growth-unit}
\rho^{\mu}\,S(\mathbf 1)\le S(f).
\end{equation}
Let $\tilde f$ be a cut profile with
$\tilde f(s)=f(m-s)$ at every interior point $s$ of $m$. Then $\tilde f$ is
nonnegative at the interior points of $m$ by hypothesis~(2), and
hypothesis~(3) gives $\rho\le f(s)^{1/2}\,\tilde f(s)^{1/2}$ there. By
\eqref{eq:growth-monotone}, by Lemma~\ref{lem:profileconvex} for the
profiles $f$ and $\tilde f$ with the exponents $1/2$ and $1/2$, and by
Lemma~\ref{lem:reflection},
\[
V_Q(m;\rho\,\mathbf 1)\le V_Q\bigl(m;f^{1/2}\tilde f^{1/2}\bigr)
\le V_Q(m;f)^{1/2}\,V_Q(m;\tilde f)^{1/2}=V_Q(m;f).
\]
The left side is $\rho^{\,Q-1}\,T_{m,Q}$, because by
\eqref{eq:growth-chains} each of the $Q-1$ cut points of a composition with
$Q$ parts contributes the factor $\rho$ to its term.
Multiplying by $v_Q\ge0$ and summing over $Q$, we get
\begin{equation}\label{eq:growth-constant}
\sum_{Q=1}^{m}v_Q\,T_{m,Q}\,\rho^{\,Q-1}\le S(f).
\end{equation}
Since $S(\mathbf 1)>0$, there is a real number $\bar\mu$ with
\[
\bar\mu\,S(\mathbf 1)=\sum_{Q=1}^{m}v_Q\,T_{m,Q}\,(Q-1),
\]
the mean of the number $Q-1$ of cuts with respect to the weights
$v_Q\,T_{m,Q}\ge0$. Jensen's inequality for the exponential function
$t\mapsto\rho^{t}$, in the form given in Section~\ref{sec:prelim} and
applied with these weights and with the exponents $Q-1$, shows that the left
side of \eqref{eq:growth-constant} is at least
$\rho^{\bar\mu}\,S(\mathbf 1)$. By hypothesis~(4) we have $\bar\mu\le\mu$,
and since $0<\rho\le1$ this implies $\rho^{\mu}\le\rho^{\bar\mu}$. Together
with \eqref{eq:growth-constant} this proves \eqref{eq:growth-unit}. In
particular $S(f)>0$.

By \eqref{eq:growth-chord} and \eqref{eq:growth-unit},
\[
\rho^{\mu}\,S(f)^{k+1}\le\rho^{\mu}\,S(\mathbf 1)\,S(g)^{k}
\le S(f)\,S(g)^{k},
\]
and dividing by $\rho^{\mu}\,S(f)>0$ we obtain
$S(f)^{k}\le\rho^{-\mu}\,S(g)^{k}$.
\end{proof}

\subsection{The moment ratio of the birthday law}\label{subsec:growth-moment}

The mean number of cuts that the comparison needs in
Section~\ref{subsec:growth-step} is expressed by a number formed from the
mean of the birthday law. Recall the birthday tail $B_{m,j}$, the product of
$1-i/m$ over $0\le i<j$ (Section~\ref{sec:compositions}), and
the mean
\[
s_m=\sum_{j=1}^{m}B_{m,j}
\]
of the birthday law of $m$ (Section~\ref{sec:birthday}); the symbol $s_m$
is not related to the cut points $s_0<s_1<\dots<s_Q$ of a composition.
Since
$B_{m,j}\le1$ for every $j$ and $B_{m,2}<1$, we have $s_m<m$ for $m\ge2$.
For an integer $m\ge2$ the \emph{moment ratio} of the birthday law of $m$ is
therefore well defined by
\begin{equation}\label{eq:growth-moment}
D_m=\frac{3\bigl((m+2)\,s_m-3m\bigr)}{2\,(m-s_m)}.
\end{equation}

By Proposition~\ref{prop:partslaw} the weights
$T_{m,Q}=a_m\,Q\,B_{m,Q}$ are proportional to the point masses
$(Q/m)\,B_{m,Q}$ of the birthday law of $m$, and $D_m$ is the ratio of
the third to the second factorial moment of this law: the sum of
$Q(Q-1)(Q-2)\,T_{m,Q}$ over $1\le Q\le m$ is $D_m$ times the sum of
$Q(Q-1)\,T_{m,Q}$, as is verified in the proof of
Proposition~\ref{prop:curvaturegrowth}. A composition with $Q$ parts has
$Q-1=(Q-2)+1$ cuts. Hence
$D_m+1$ is the mean number of cuts of a composition of $m$ when the weight
of every composition with $Q$ parts is multiplied by $Q(Q-1)$. We define
$D_m$ by the closed formula
in $m$ and $s_m$, and not as a quotient of two sums, because a bound for
$s_m$ then turns directly into a bound for $D_m$.
Lemma~\ref{lem:momentquotient} is obtained in this way from the upper bound
of Lemma~\ref{lem:birthdayupper}.

\begin{lemma}\label{lem:momentquotient}
For every integer $m\ge2$, the moment ratio $D_m$ of the birthday law of
$m$ satisfies
\[
D_m+1\le\frac32\sqrt{\frac{\pi m}{2}}-\frac85 .
\]
\end{lemma}

\begin{proof}
Let $p=\sqrt{\pi m/2}$. The claim is $D_m\le3p/2-13/5$. Multiplying by the
denominator $2\,(m-s_m)$ of $D_m$, which is positive, and collecting the
terms with $s_m$ on the left side, we see that this is equivalent to
\begin{equation}\label{eq:growth-linear}
s_m\,\Bigl(3m+3p+\frac45\Bigr)\le m\,\Bigl(3p+\frac{19}5\Bigr).
\end{equation}

For $m=2$ and $m=3$ we check \eqref{eq:growth-linear} directly. By the
definition of the birthday tails, $B_{2,1}=B_{3,1}=1$, $B_{2,2}=1/2$,
$B_{3,2}=2/3$ and $B_{3,3}=(2/3)(1/3)=2/9$, so that the definition of $s_m$
gives
\[
s_2=1+\frac12=\frac32,\qquad s_3=1+\frac23+\frac29=\frac{17}9 .
\]
For $m=2$ the inequality \eqref{eq:growth-linear} rearranges to
$p\ge26/15$, which
holds because $p^2=\pi>3.14$ and $(26/15)^2<3.01$. For $m=3$ it is
equivalent to $p\ge32/15$, which holds because $p^2=3\pi/2>4.71$ and
$(32/15)^2<4.56$.

Let $m\ge4$. By Lemma~\ref{lem:birthdayupper},
$s_m\le p\,e^{1/(12m)}-1/3$. The number $y=1/(12m)$ satisfies
$0\le y\le1/48$, so that the inequality $e^{y}\le1+y+y^2$ of
Section~\ref{sec:prelim} gives
\[
e^{1/(12m)}\le1+\frac{49}{48}\,y\le1+\frac{48}{47}\,y=1+\frac{4}{47m},
\]
and hence
\[
s_m\le p\,\Bigl(1+\frac{4}{47m}\Bigr)-\frac13 .
\]
Since the factor of $s_m$ in \eqref{eq:growth-linear} is positive, it
suffices to prove \eqref{eq:growth-linear} with $s_m$ replaced by this upper
bound, and after multiplication by $m$ this is the inequality
\[
\Bigl(pm+\frac{4p}{47}-\frac m3\Bigr)\Bigl(3m+3p+\frac45\Bigr)
\le3pm^2+\frac{19}5\,m^2 .
\]
We expand the product and replace $p^2$ by $\pi m/2$. The right side minus
the left side is then
\begin{equation}\label{eq:growth-difference}
\Bigl(\frac{24}5-\frac{3\pi}2\Bigr)m^2-\frac{13}{235}\,pm
-\frac{6\pi}{47}\,m-\frac{16}{235}\,p+\frac4{15}\,m,
\end{equation}
and we have to show that this number is nonnegative. In
\eqref{eq:growth-difference} the numbers $\pi$ and $p$ occur with negative
coefficients, so that they may be replaced by larger numbers. We use
$\pi<3.1416$ and
\[
p\le\frac{p^2}5+\frac54\le0.31416\,m+1.25,
\]
where the first inequality is $(p-5/2)^2\ge0$ and the second holds because
$p^2=\pi m/2$. A direct computation with these two bounds shows that the
number \eqref{eq:growth-difference} is at least
\[
0.0702\,m^2-0.225\,m-0.0852=m\,(0.0702\,m-0.225)-0.0852 .
\]
For $m\ge4$ we have $0.0702\,m-0.225\ge0.0558$, so that this is at least
$4\cdot0.0558-0.0852>0$.
\end{proof}

\subsection{The comparison of the curvature sums along the diagonal}
\label{subsec:growth-step}

We turn to the curvature sum $\mathsf C_m(y;N)$ of
Definition~\ref{def:curvaturesum}, the second difference in the threshold
$y$ of $y\,\Lambda_m(y;N)$. By Lemma~\ref{lem:curvatureparts} it is a sum
over the number $Q$ of parts, in which the weight of the compositions of $m$
with $Q$ parts is multiplied by the second difference
$(y+1)^Q-2\,y^Q+(y-1)^Q$ of the power $y^Q$. Formula
\eqref{eq:frozen-binomial} expands this second difference in powers of $y$,
for every real number $y$ and every integer exponent $q\ge0$; we use it with
the letters $Q$ and $j$ in the place of $q$ and $i$, and we write its right
side as a power of $y$ times a polynomial in $1/y$. For an integer $Q\ge0$
and a real number $x$
let
\begin{equation}\label{eq:growth-sigma}
\sigma_Q(x)=2\sum_{j\ge1}\binom{Q}{2j}\,x^{2j-2}.
\end{equation}
The sum is finite, because the binomial coefficient vanishes for $2j>Q$.
The proof of Proposition~\ref{prop:curvaturegrowth} bounds $\sigma_Q$ from
both sides.

Now let $k\ge1$, $m\ge2$ and $q\ge m+1$ be integers. We compare the
curvature sums $\mathsf C_m(k;k+q)$ and $\mathsf C_m(k+1;k+q+1)$. For
$q=r+1$ these are the curvature sums of the positions $(k+r,k)$ and
$(k+r+1,k+1)$ that were described at the beginning of the section: on a
diagonal step the threshold $y$ and the number $N$ of the curvature sum
$\mathsf C_m(y;N)$ are both raised by one, while the integer $m$ whose
compositions are summed stays the same.
The threshold profiles $\lambda^{k,k+q}$ and $\lambda^{k+1,k+q+1}$ of the
two sums, for which we write $\lambda$ and $\lambda'$, are
\[
\lambda(s)=\frac{k}{k+q-s}\qquad\text{and}\qquad
\lambda'(s)=\frac{k+1}{k+q+1-s}.
\]
For an interior point $s$ of $m$ the number $h=q-s$ is positive, and the
factors that the two profiles give to a cut at $s$
are $\lambda(s)=1/(1+h/k)$ and $\lambda'(s)=1/(1+h/(k+1))$. They are the
values at $u=1/k$ and at $u=1/(k+1)$ of the function $u\mapsto1/(1+hu)$,
whose value at $u=0$ is $1$. A diagonal step thus moves $u$ from $1/k$ to
$1/(k+1)$, towards $0$.
The function $u\mapsto1/(1+hu)$ is log-convex, and $1/(k+1)$ is the convex
combination of $0$ and $1/k$ with the coefficients $1/(k+1)$ and $k/(k+1)$.
This log-convexity is the reason behind hypothesis~(2) of
Lemma~\ref{lem:profilecomparison} for $f=\lambda'$ and $g=\lambda$, which
the proof of Proposition~\ref{prop:curvaturegrowth} verifies; the constant
profile $\mathbf 1$ of that lemma belongs to $u=0$.

\begin{proposition}\label{prop:curvaturegrowth}
Let $k$, $q$ and $m$ be integers with $k\ge1$, $m\ge2$, $m+1\le q$ and
$401\,m^2\le12\,k^2$. Then the curvature sum of
Definition~\ref{def:curvaturesum} and the moment ratio $D_m$ of the
birthday law of $m$ satisfy
\begin{multline*}
\mathsf C_m(k+1;\,k+q+1)\\
\le\frac{k}{k+1}\,
\exp\biggl(\frac{\bigl(1+\frac1{400}\bigr)(D_m+1)\bigl(q-\frac m2\bigr)}
{k\,(k+1)}\biggr)\,\mathsf C_m(k;\,k+q).
\end{multline*}
\end{proposition}

\begin{proof}
Let $\lambda=\lambda^{k,k+q}$ and $\lambda'=\lambda^{k+1,k+q+1}$ be the
threshold profiles of the two curvature sums, so that
$\lambda(s)=k/(k+q-s)$ and $\lambda'(s)=(k+1)/(k+q+1-s)$, and put
\begin{gather*}
v_Q=\sigma_Q(1/k)\quad(1\le Q\le m),\qquad\bar h=q-\frac m2,\\
\rho=\frac{k+1}{k+1+\bar h},\qquad
\mu=\Bigl(1+\frac1{400}\Bigr)(D_m+1).
\end{gather*}
All sums over $Q$ run over $1\le Q\le m$. We first express the two curvature
sums by the weights of their own profiles. Let $(y,N)$ be one of the pairs
$(k,k+q)$ and $(k+1,k+q+1)$. By Lemma~\ref{lem:curvatureparts},
$\mathsf C_m(y;N)$ is the sum over $Q$ of
\[
\bigl((y+1)^Q-2\,y^Q+(y-1)^Q\bigr)\,V_Q\bigl(m;\lambda^{1,N}\bigr).
\]
Here $y\ne0$, and the term $j$ of the sum \eqref{eq:growth-sigma} at
$x=1/y$, multiplied by $y^{Q-2}$, is $\binom{Q}{2j}\,y^{\,Q-2j}$, the term
$i=j$ of the sum in \eqref{eq:frozen-binomial} for $q=Q$. Hence
\eqref{eq:frozen-binomial} says that the first factor is
\[
(y+1)^Q-2\,y^Q+(y-1)^Q=y^{Q-2}\,\sigma_Q(1/y).
\]
The threshold profile $\lambda^{y,N}$
is $y$ times the profile $\lambda^{1,N}$, so that the second factor is
$y^{-(Q-1)}\,V_Q(m;\lambda^{y,N})$ by \eqref{eq:growth-homogeneous}. Hence
\[
\mathsf C_m(y;N)=\frac1y\sum_{Q}\sigma_Q(1/y)\,V_Q\bigl(m;\lambda^{y,N}\bigr).
\]
We write $X$ and $Y$ for the sums of $v_Q\,V_Q(m;\lambda')$ and
$v_Q\,V_Q(m;\lambda)$ over $Q$, so that $\mathsf C_m(k;k+q)=Y/k$. At an
interior point $s$ of $m$ we have $q-s>q-m>0$, so that both profiles are
positive at the interior points of $m$, and the weights $V_Q(m;\lambda)$
and $V_Q(m;\lambda')$ are nonnegative, as recalled in
Section~\ref{subsec:growth-profile}. The polynomial $\sigma_Q$ contains only
even powers of $x$, its coefficients are nonnegative, and its constant
term, the term $j=1$ of \eqref{eq:growth-sigma}, is $Q(Q-1)$. Consequently,
for all real numbers $x$ and $x'$ with $0\le x\le x'$,
\begin{equation}\label{eq:growth-sigmalower}
Q(Q-1)\le\sigma_Q(x)\le\sigma_Q(x').
\end{equation}
Since
$0\le\sigma_Q(1/(k+1))\le v_Q$ by \eqref{eq:growth-sigmalower}, the formula
for the pair $(k+1,k+q+1)$ gives $\mathsf C_m(k+1;k+q+1)\le X/(k+1)$. As
$\mu\,\bar h/(k\,(k+1))$ is the exponent in the statement, it therefore
suffices to show that
\begin{equation}\label{eq:growth-goal}
X\le\exp\Bigl(\frac{\mu\,\bar h}{k\,(k+1)}\Bigr)\,Y.
\end{equation}
We derive \eqref{eq:growth-goal} from Lemma~\ref{lem:profilecomparison},
applied with the exponent $k$, the weights $v_Q$, the profiles $f=\lambda'$
and $g=\lambda$, and the numbers $\rho$ and $\mu$.

Hypothesis~(1) holds by \eqref{eq:growth-sigmalower}. For hypothesis~(2),
let $s$ be an interior point of $m$ and $h=q-s$, so that
$\lambda(s)=1/(1+h/k)$ and $\lambda'(s)=1/(1+h/(k+1))$. The inequality
between the geometric and the arithmetic mean of the $k+1$ numbers that
consist of $k$ copies of $1+h/k$ and of the number $1$ gives
\[
\Bigl(1+\frac hk\Bigr)^{k}\le\Bigl(1+\frac h{k+1}\Bigr)^{k+1},
\]
and taking reciprocals we obtain $\lambda'(s)^{k+1}\le\lambda(s)^{k}$. For
hypothesis~(3), we have
$\bar h>0$ because $q>m$, and hence $0<\rho<1$. For an interior point $s$
of $m$,
\[
\lambda'(s)\,\lambda'(m-s)
=\frac{(k+1)^2}{(k+1+q-s)\,(k+1+q-m+s)}.
\]
The two factors of the denominator are positive and their sum is
$2\,(k+1+\bar h)$, so that their product is at most $(k+1+\bar h)^2$. Hence
$\lambda'(s)\,\lambda'(m-s)\ge\rho^2$.

For hypothesis~(4), let $M_2$
and $M_3$ be the sums over $Q$ of $Q(Q-1)\,T_{m,Q}$ and
$Q(Q-1)(Q-2)\,T_{m,Q}$. By Proposition~\ref{prop:partslaw} we have
$T_{m,Q}=a_m\,Q\,B_{m,Q}$, so that $M_2$ and $M_3$ are $a_m$ times the sums
of $Q^2(Q-1)\,B_{m,Q}$ and $Q^2(Q-1)(Q-2)\,B_{m,Q}$.
Corollaries~\ref{cor:secondfactorial} and~\ref{cor:thirdfactorial} evaluate
these two sums, and we obtain
\[
M_2=2m\,a_m\,(m-s_m),\qquad M_3=3m\,a_m\,\bigl((m+2)\,s_m-3m\bigr).
\]
Since $a_m>0$, and since $s_m<m$ as was noted before
\eqref{eq:growth-moment}, the number $M_2$ is positive, and by the
definition \eqref{eq:growth-moment} of $D_m$,
\begin{equation}\label{eq:growth-ratio}
M_3=D_m\,M_2.
\end{equation}
Clearly $M_3\ge0$, so that \eqref{eq:growth-ratio} shows that $D_m\ge0$; in
particular $\mu>0$. By
\eqref{eq:growth-sigmalower} we have $v_Q\ge Q(Q-1)$, and therefore
\[
\sum_{Q}v_Q\,T_{m,Q}\ge M_2>0 .
\]
In the other direction, if $1\le Q\le m$ and $x$ is a real number with
$m^2x^2<12$, then
\begin{equation}\label{eq:growth-sigmaupper}
\sigma_Q(x)\le\frac{Q(Q-1)}{1-m^2x^2/12}.
\end{equation}
To see this, we compare consecutive terms of \eqref{eq:growth-sigma}. For
$j\ge1$ with $2j\le Q$,
\[
\binom{Q}{2j+2}\,x^{2j}
=\frac{(Q-2j)(Q-2j-1)}{(2j+1)(2j+2)}\,x^2\cdot\binom{Q}{2j}\,x^{2j-2}
\le\frac{m^2x^2}{12}\,\binom{Q}{2j}\,x^{2j-2},
\]
since $0\le(Q-2j)(Q-2j-1)\le Q^2\le m^2$ and $(2j+1)(2j+2)\ge12$; for $2j>Q$
both terms vanish. The terms of \eqref{eq:growth-sigma} are thus at most
those of the geometric series with first term $Q(Q-1)$ and ratio
$m^2x^2/12$, whose sum is the right side of \eqref{eq:growth-sigmaupper}.
The hypothesis $401\,m^2\le12\,k^2$ says that
$m^2/(12k^2)\le1/401$, so that \eqref{eq:growth-sigmaupper} with $x=1/k$
gives $v_Q\le(1+1/400)\,Q(Q-1)$ for $1\le Q\le m$. Writing $Q-1=(Q-2)+1$
and using \eqref{eq:growth-ratio}, we obtain
\begin{align*}
\sum_{Q}v_Q\,T_{m,Q}\,(Q-1)
&\le\Bigl(1+\frac1{400}\Bigr)\sum_{Q}Q(Q-1)\bigl((Q-2)+1\bigr)\,T_{m,Q}\\
&=\Bigl(1+\frac1{400}\Bigr)(M_3+M_2)=\mu\,M_2
\le\mu\sum_{Q}v_Q\,T_{m,Q},
\end{align*}
where the last step uses $\mu\ge0$. This is hypothesis~(4).

Lemma~\ref{lem:profilecomparison} now gives $X^k\le\rho^{-\mu}\,Y^k$. Since
$\rho^{-1}=1+\bar h/(k+1)\le\exp\bigl(\bar h/(k+1)\bigr)$ and $\mu\ge0$,
\[
X^k\le\exp\Bigl(\frac{\mu\,\bar h}{k+1}\Bigr)\,Y^k
=\biggl(\exp\Bigl(\frac{\mu\,\bar h}{k\,(k+1)}\Bigr)\,Y\biggr)^{k}.
\]
Since $X$ and $Y$ are nonnegative, \eqref{eq:growth-goal} follows by taking
$k$-th roots.
\end{proof}

\section{The diagonal step behind the curve}\label{sec:step}
We can now prove Proposition~\ref{prop:diagonalstep}: for $k\ge265$,
$r\ge1$ and $r^3\le k^2$, a step $d_{k+r,k}$ that is not positive is
followed by a negative step $d_{k+r+1,k+1}$ at the diagonal neighbour.
Together with Proposition~\ref{prop:curveincrease} and
Corollary~\ref{cor:layerpersistenceall}, this gives the diagonal rule in
the first layer of columns $k\ge265$ in
Section~\ref{sec:staircase}. The proof compares the two positions term by
term, through the comparisons of Propositions~\ref{prop:boundarygain}
and~\ref{prop:curvaturegrowth} and a third one, for the kernel deficit
(Proposition~\ref{prop:deficitgrowth}).

\subsection{The two positions}\label{subsec:step-positions}

The two positions are compared through
Corollary~\ref{cor:curvatureboundary}, which writes the step at a position of
the first layer as a curvature term minus a boundary term. Let $k\ge2$ and
$r$ be integers with $1\le r\le k-1$, and put $N=k+r+1$. For $1\le m\le r$
let
\begin{equation}\label{eq:step-terms}
\chi_m=\frac{\varkappa_k(N-m)}{N-m}\;\mathsf C_m(k;N),\qquad
B=\frac{a_k}{k}\;\Lambda_{r+1}(k-1;N),
\end{equation}
and let $\chi'_m$ and $B'$ be the same expressions for the diagonal neighbour
$(k+r+1,k+1)$, that is, with $k$ and $N$ replaced by $k+1$ and $N+1$:
\begin{equation}\label{eq:step-neighbour}
\begin{aligned}
\chi'_m&=\frac{\varkappa_{k+1}(N+1-m)}{N+1-m}\;\mathsf C_m(k+1;N+1),\\
B'&=\frac{a_{k+1}}{k+1}\;\Lambda_{r+1}(k;N+1).
\end{aligned}
\end{equation}
Here $\varkappa$ is the kernel deficit of Section~\ref{sec:frozen},
$\mathsf C_n(y;L)$ is the curvature sum of
Definition~\ref{def:curvaturesum}, and $\Lambda_n(y;L)$ is the chain sum of
Section~\ref{sec:compositions}. Let $C=\chi_1+\dots+\chi_r$ and
$C'=\chi'_1+\dots+\chi'_r$. Corollary~\ref{cor:curvatureboundary}, applied to
column $k$ and to column $k+1$ with the same $r$, states that
\begin{equation}\label{eq:step-twopositions}
\frac{d_{k+r,k}}{(k+r)!}=C-B,\qquad
\frac{d_{k+r+1,k+1}}{(k+r+1)!}=C'-B'.
\end{equation}
The numbers $C$ and $B$ are the curvature term and the boundary term of
position $(k+r,k)$, and $C'$ and $B'$ are those of its diagonal neighbour.
The two positions have the same deficit $r$. Their curvature terms are
therefore sums over the same indices $m$, and the chain sums of their
boundary terms run over the compositions of the same number $r+1$.

The chain sums of the boundary terms are positive. Indeed, by the properties
of the threshold profile noted in Section~\ref{sec:compositions}, for a real
number $y$ and integers $n$ and $L$ we have
\begin{equation}\label{eq:step-chainpositive}
\Lambda_n(y;L)>0\qquad(y>0,\ 1\le n<L).
\end{equation}

By \eqref{eq:step-twopositions}, the hypothesis of
Proposition~\ref{prop:diagonalstep} is $C\le B$, and its claim is
$C'<B'$. The terms with $m=1$ do not contribute to
the two curvature terms: a curvature sum $\mathsf C_1(y;L)$ is zero, as
noted after Lemma~\ref{lem:curvatureparts} in Section~\ref{sec:frozen}. Hence
\begin{equation}\label{eq:step-firstterm}
\chi_1=\chi'_1=0 .
\end{equation}
For $r\ge2$, and for $k\ge265$ and $r^3\le k^2$ as in the proposition, we
shall derive the claim from the \emph{cross inequality}
\begin{equation}\label{eq:step-cross}
C'\,B<C\,B' ,
\end{equation}
which says that the quotient of the curvature term by the boundary term
decreases strictly from the position to its diagonal neighbour. The cross
inequality does not involve the hypothesis $C\le B$, and it is the reason
why the conclusion is strict although the hypothesis admits a step equal to
zero. It is restricted to $r\ge2$, because for $r=1$ both of its sides are
zero.

The cross inequality is proved term by term: we show that
$\chi'_m\,B<\chi_m\,B'$ for $2\le m\le r$
(Proposition~\ref{prop:oneterm}). Each side of the inequality
for one term is a product of three ingredients, namely a coefficient formed
from a kernel deficit, a curvature sum, and the chain sum of a boundary
term, and each ingredient can be followed from the position to its diagonal
neighbour by itself. For the chain sum of the boundary term this was done in
Proposition~\ref{prop:boundarygain}, by which it grows by at least an
explicit factor, and for the curvature sum in
Proposition~\ref{prop:curvaturegrowth}, by which it grows by at most an
explicit factor; the comparison of the coefficients is
Proposition~\ref{prop:deficitgrowth} below. What the three
comparisons leave to be shown is an inequality between their factors
(Lemma~\ref{lem:scalar}). After taking logarithms and inserting closed bounds
for the factors, it becomes an inequality between explicit functions of $k$,
$r$ and $m$ (Lemma~\ref{lem:closedforms}), and this is where the hypothesis
$r^3\le k^2$ is needed.

\subsection{The kernel deficit along the diagonal}\label{subsec:step-deficit}

The coefficient of the term $\chi_m$ in \eqref{eq:step-terms} is
$\varkappa_k(t)/t$ with $t=N-m$. We write $t=k+a$, so that $a=r+1-m$; here
and below $a$ without an index denotes a nonnegative integer, while $a_j$
continues to denote the Cayley coefficients. In $\chi'_m$ the
coefficient is $\varkappa_{k+1}(t+1)/(t+1)$: the index of the kernel deficit
and its argument have grown by one each, and their difference $a$ is the
same. We compare the two coefficients after dividing each kernel deficit by
the Cayley coefficient of its argument. When the proof of
Proposition~\ref{prop:oneterm} returns from these quotients to the
coefficients, a diagonal step brings the factor $\gamma_{k+a}$ into a term of
the curvature term and the factor $\gamma_k$ into the boundary term, where
$\gamma_t=(1+1/t)^{t-2}$ is the diagonal factor of
Section~\ref{sec:diagonal}. We need to know by how much the first exceeds
the second, and the following lemma bounds the difference
$\log\gamma_i-\log\gamma_k$ from above by $(5/2)(1/k-1/i)$.

\begin{lemma}\label{lem:factorlog}
For all integers $k$ and $i$ with $2\le k\le i$, the diagonal factors
satisfy
\[
\log\gamma_i-\log\gamma_k\le\frac52\Bigl(\frac1k-\frac1i\Bigr).
\]
\end{lemma}

\begin{proof}
It suffices to show that
\begin{equation}\label{eq:step-onestep}
\log\gamma_{m+1}-\log\gamma_m\le\frac{5}{2\,m\,(m+1)}
=\frac52\Bigl(\frac1m-\frac1{m+1}\Bigr)
\end{equation}
for every integer $m\ge2$; the lemma follows by adding these inequalities
for $m=k,\dots,i-1$. By the definition of the diagonal factors,
\[
\log\gamma_{m+1}=(m-1)\,\log\Bigl(1+\frac{1}{m+1}\Bigr),\qquad
\log\gamma_m=(m-2)\,\log\Bigl(1+\frac1m\Bigr),
\]
and the factors $m-1$ and $m-2$ are nonnegative. We bound the first
logarithm from above and the second from below by the Taylor bounds of
Section~\ref{sec:prelim}, namely $\log(1+u)\le u-u^2/2+u^3/3$ at
$u=1/(m+1)$ and $\log(1+u)\ge u-u^2/2+u^3/3-u^4/4$ at $u=1/m$. This gives
\begin{align*}
\log\gamma_{m+1}-\log\gamma_m\le{}&(m-1)\Bigl(\frac{1}{m+1}
-\frac{1}{2\,(m+1)^2}+\frac{1}{3\,(m+1)^3}\Bigr)\\
&-(m-2)\Bigl(\frac1m-\frac{1}{2\,m^2}+\frac{1}{3\,m^3}-\frac{1}{4\,m^4}\Bigr).
\end{align*}
A straightforward computation shows that $5/(2m(m+1))$ minus the right-hand
side equals
\[
\frac{29\,(m-2)^4+253\,(m-2)^3+823\,(m-2)^2+1191\,(m-2)+656}
{12\,m^4\,(m+1)^3},
\]
which is positive for $m\ge2$. This proves \eqref{eq:step-onestep}.
\end{proof}

We turn to the kernel deficits $\varkappa_k(t)$ of
Section~\ref{sec:frozen}. Let $k\ge0$ and $a\ge0$ be integers such that
$t=k+a\ge1$. The quotient $\varkappa_k(t)/a_t$ is the \emph{normalised
kernel deficit}. On a diagonal step it is replaced by
$\varkappa_{k+1}(t+1)/a_{t+1}$, and we bound the increase by a multiple of
the normalised kernel deficit itself.

The exact bound is expressed by the following closed form. For integers
$k\ge1$ and $a\ge0$ with $2a\le k$ we put $\xi=a/k$ and
\begin{equation}\label{eq:step-theta}
\Theta(k,a)=\frac32\,\xi\,\Bigl(1+\frac{2}{3\,(k+2)}\Bigr)\cdot
\frac{1+\frac32\,\xi+\frac38\,\xi^2}
{\frac{23}{24}-\frac32\,\xi\,\bigl(1+\frac{\xi}{4}\bigr)}.
\end{equation}
The ratio $\xi$ carries no index and is unrelated to the tree curvature of
Definition~\ref{def:treecurvature}. Since $0\le\xi\le1/2$, we have
$(3/2)\,\xi\,(1+\xi/4)\le27/32<23/24$.
The denominator in \eqref{eq:step-theta} is therefore positive, and
$\Theta(k,a)\ge0$.

The following proposition is the relative bound in its exact form, and the
form of $\Theta$ is the one that its proof produces. The proof bounds the
normalised kernel deficit from below and its increase from above, both in
terms of the sum of the Poisson masses $P(j)$ over $0<j<a$. It then compares
this sum with the tail of the Borel distribution by the two bounds
of~\eqref{eq:prelim-stirling}, and this comparison turns the bound for the
increase into a multiple of the normalised kernel deficit.

\begin{proposition}\label{prop:deficitgrowth}
Let $k\ge1$ and $a\ge0$ be integers with $2a\le k$. Then the normalised
kernel deficits satisfy
\[
\frac{\varkappa_{k+1}(k+a+1)}{a_{k+a+1}}\le
\Bigl(1+\frac{\Theta(k,a)}{k}\Bigr)\,\frac{\varkappa_k(k+a)}{a_{k+a}},
\]
where $\Theta(k,a)$ is the closed form \eqref{eq:step-theta}.
\end{proposition}

\begin{proof}
Let $t=k+a$ and $\xi=a/k$, so that $0\le\xi\le1/2$, and write
\begin{gather*}
\sigma=\frac{\varkappa_k(t)}{a_t},\qquad
\sigma'=\frac{\varkappa_{k+1}(t+1)}{a_{t+1}},\\
\Sigma=\sum_{0<j<a}P(j),\qquad
\phi=1-\sum_{0<j<a}\frac{P(j)}{j}.
\end{gather*}
Further, let
\[
\theta=\frac32\,\Bigl(1+\frac{2}{3\,(k+2)}\Bigr)
\Bigl(1+\frac32\,\xi+\frac38\,\xi^2\Bigr),\qquad
\rho=\frac{23}{24}-\frac32\,\xi\,\Bigl(1+\frac{\xi}{4}\Bigr),
\]
so that $\rho>0$ and $\Theta(k,a)=\theta\,\xi/\rho$. The assertion is
$\sigma'\le(1+\Theta(k,a)/k)\,\sigma$.

We begin with two comparisons of Cayley coefficients. First, for all
integers $i\ge1$ and $j\ge1$,
\begin{equation}\label{eq:step-borel}
\frac{a_j\,a_i}{a_{i+j}}\le\frac{P(j)}{j}\,
\Bigl(1+\frac32\cdot\frac{j}{i}+\frac38\cdot\frac{j^2}{i^2}\Bigr).
\end{equation}
Indeed, $a_j=e^{j}P(j)/j$, so that the left-hand side is $P(j)/j$ times
$e^{j}a_i/a_{i+j}$, and Lemma~\ref{lem:borelquotient}, applied to indices
$i\le i+j$, bounds $e^{j}a_i/a_{i+j}$ by the polynomial in $j/i$ on
the right. Secondly, for all integers $i\ge2$ and $j\ge0$,
\begin{equation}\label{eq:step-shift}
\frac{a_i}{a_{i+j}}-\frac{a_{i+1}}{a_{i+j+1}}\le\frac{a_i}{a_{i+j}}\cdot
\frac{j}{i+j}\,\Bigl(\frac{3}{2\,i}+\frac{1}{i\,(i+1)}\Bigr).
\end{equation}
To see this, recall that $a_{n+1}=\gamma_n\,(1+1/n)\,a_n$ for $n\ge2$, by
\eqref{eq:prelim-quotient} and \eqref{eq:diagonal-factor}. The left-hand
side of \eqref{eq:step-shift} is therefore $(a_i/a_{i+j})\,(1-1/u)$ with
\[
u=\frac{\gamma_{i+j}}{\gamma_i}\cdot\frac{1+1/(i+j)}{1+1/i}>0 .
\]
Since $\log y\le y-1$ for every $y>0$, we have $1-1/u\le\log u$, and
$\log u$ is the sum of $\log\gamma_{i+j}-\log\gamma_i$ and of the logarithm
of the second factor of $u$. The first of these is at most
$(5/2)(1/i-1/(i+j))=5j/(2i(i+j))$ by Lemma~\ref{lem:factorlog}, and the
second, by the same inequality for the logarithm, is at most
\[
\frac{1+1/(i+j)}{1+1/i}-1=-\frac{j}{(i+j)\,(i+1)} .
\]
The two bounds add up to
\[
\frac{j}{i+j}\,\Bigl(\frac{5}{2\,i}-\frac{1}{i+1}\Bigr)
=\frac{j}{i+j}\,\Bigl(\frac{3}{2\,i}+\frac{1}{i\,(i+1)}\Bigr),
\]
and \eqref{eq:step-shift} follows, because $a_i/a_{i+j}>0$.

Next we bound the normalised kernel deficit $\sigma$ from below: we show
that
\begin{equation}\label{eq:step-lower}
\sigma\ge\phi-\frac32\,\Bigl(1+\frac{\xi}{4}\Bigr)\,\frac{\Sigma}{k+1}.
\end{equation}
On replacing the index of summation $i$ in the definition of
$\varkappa_k(t)$ by $j=t-i$, which runs over $0<j<a$, and dividing by
$a_t$, we obtain
\begin{equation}\label{eq:step-normalised}
\frac{\varkappa_k(t)}{a_t}=1-\sum_{0<j<a}\frac{a_j\,a_{t-j}}{a_t}.
\end{equation}
Let $0<j<a$ and $i=t-j$, so that $i\ge k+1$. Then $j/i\le j/(k+1)$ and
$j/i\le a/k=\xi$, and \eqref{eq:step-borel} gives
\begin{align*}
\frac{a_j\,a_{t-j}}{a_t}
&\le\frac{P(j)}{j}\,
\Bigl(1+\frac32\,\Bigl(1+\frac{j}{4\,i}\Bigr)\frac{j}{i}\Bigr)\\
&\le\frac{P(j)}{j}+\frac32\,\Bigl(1+\frac{\xi}{4}\Bigr)\,\frac{P(j)}{k+1}.
\end{align*}
Summing over $j$ and inserting the result into \eqref{eq:step-normalised},
we obtain \eqref{eq:step-lower}.

For the increase $\sigma'-\sigma$ we show that
\begin{equation}\label{eq:step-increase}
\sigma'-\sigma\le\theta\cdot\frac{\Sigma}{(k+a)\,(k+1)}.
\end{equation}
The identity \eqref{eq:step-normalised} holds for the pair $(k,t)$ and for
the pair $(k+1,t+1)$, with the same range $0<j<a$ of summation. Subtracting
the two identities, we get
\[
\sigma'-\sigma=\sum_{0<j<a}a_j\,
\Bigl(\frac{a_{t-j}}{a_t}-\frac{a_{t+1-j}}{a_{t+1}}\Bigr).
\]
Let $0<j<a$ and $i=t-j$, so that $i\ge k+1\ge2$ and $i+j=k+a$. By
\eqref{eq:step-shift} and then \eqref{eq:step-borel}, the term with index
$j$ is at most
\[
\frac{P(j)}{j}\,
\Bigl(1+\frac32\cdot\frac{j}{i}+\frac38\cdot\frac{j^2}{i^2}\Bigr)\cdot
\frac{j}{k+a}\,\Bigl(\frac{3}{2\,i}+\frac{1}{i\,(i+1)}\Bigr).
\]
Here $j/i\le\xi$, and since $i\ge k+1$,
\[
\frac{3}{2\,i}+\frac{1}{i\,(i+1)}\le\frac{3}{2\,(k+1)}+\frac{1}{(k+1)(k+2)}
=\frac{3}{2\,(k+1)}\,\Bigl(1+\frac{2}{3\,(k+2)}\Bigr).
\]
The factors $j$ cancel, so that this term is at most
$\theta\,P(j)/((k+a)(k+1))$, and the sum over $j$ gives
\eqref{eq:step-increase}.

We now compare the sum $\Sigma$ of the Poisson masses with $\phi$, and show
that
\begin{equation}\label{eq:step-tail}
\frac{23}{24}\,\Sigma\le a\,\phi .
\end{equation}
For $a\le1$ the two sums over $j$ are empty, so that $\Sigma=0$ and
$\phi=1$, and \eqref{eq:step-tail} holds. Let $a\ge2$. By
the upper bound of \eqref{eq:prelim-stirling} we have
$P(j)\le c_{\mathrm S}/\sqrt j$ for $j\ge1$, and the sum of
$1/\sqrt j$ over $1\le j\le a-1$ is at most $2\sqrt{a-1}-1$
(Section~\ref{sec:prelim}). Hence
$\Sigma\le c_{\mathrm S}\,(2\sqrt{a-1}-1)$. For the tail, let $n\ge a$.
The lower bound of \eqref{eq:prelim-stirling} gives
\[
P(n)\ge\frac{c_{\mathrm S}}{\sqrt n}\,\Bigl(1-\frac{1}{12\,n}\Bigr)
\ge\frac{c_{\mathrm S}}{\sqrt n}\,\Bigl(1-\frac{1}{12\,a}\Bigr),
\]
and $1/(n\sqrt n)\ge2\,(1/\sqrt n-1/\sqrt{n+1})$
(Section~\ref{sec:prelim}). Hence
\[
\frac{P(n)}{n}\ge2\,c_{\mathrm S}\,\Bigl(1-\frac{1}{12\,a}\Bigr)
\Bigl(\frac{1}{\sqrt n}-\frac{1}{\sqrt{n+1}}\Bigr)\qquad(n\ge a).
\]
The numbers $P(j)/j=e^{-j}a_j$, $j\ge1$, are nonnegative, and their series
has the sum $1$ by Lemma~\ref{lem:borelmass}; so $\phi$ is the sum of
$P(n)/n$ over all $n\ge a$. The differences on the right of the last
display telescope to $1/\sqrt a$, and we obtain
\[
\phi\ge\frac{2\,c_{\mathrm S}}{\sqrt a}\,\Bigl(1-\frac{1}{12\,a}\Bigr)
=\frac{c_{\mathrm S}\,(2a-\frac16)}{a\,\sqrt a}.
\]
Now let $u=\sqrt{a-1}$ and $v=\sqrt a$. Then $2uv\le u^2+v^2=2a-1$, hence
$(2u-1)\,v\le2a-1$, and
\[
\frac{23}{24}\,\Sigma\le\frac{23}{24}\,c_{\mathrm S}\,(2u-1)
\le\frac{23\,c_{\mathrm S}\,(2a-1)}{24\,v}
\le\frac{c_{\mathrm S}\,(2a-\frac16)}{v}\le a\,\phi ,
\]
where the third inequality holds because $23\,(2a-1)\le24\,(2a-1/6)$. This
is \eqref{eq:step-tail}.

It remains to combine the three bounds. By \eqref{eq:step-tail},
\[
\frac{23}{24}\cdot\frac{\Sigma}{k+1}\le\frac{23}{24}\cdot\frac{\Sigma}{k}
\le\frac{a\,\phi}{k}=\xi\,\phi .
\]
We use this and \eqref{eq:step-lower}, multiplied by $\xi\ge0$, and get
\begin{align*}
\rho\cdot\frac{\Sigma}{k+1}
&=\frac{23}{24}\cdot\frac{\Sigma}{k+1}
-\frac32\,\xi\,\Bigl(1+\frac{\xi}{4}\Bigr)\,\frac{\Sigma}{k+1}\\
&\le\xi\,\phi-\frac32\,\xi\,\Bigl(1+\frac{\xi}{4}\Bigr)\,\frac{\Sigma}{k+1}
\le\xi\,\sigma ,
\end{align*}
and hence $\Sigma/(k+1)\le\xi\,\sigma/\rho$, because $\rho>0$. In
\eqref{eq:step-increase} we bound $1/(k+a)$ by $1/k$ and insert this
inequality:
\[
\sigma'-\sigma\le\frac{\theta}{k}\cdot\frac{\Sigma}{k+1}
\le\frac{\theta}{k}\cdot\frac{\xi\,\sigma}{\rho}
=\frac{\Theta(k,a)}{k}\,\sigma . \qedhere
\]
\end{proof}

\subsection{The inequality between closed forms}\label{subsec:step-closed}

We now collect what the three comparisons along the diagonal give for one
term. Consider a position $(k+r,k)$ as in \eqref{eq:step-terms} and an index
$m$ with $2\le m\le r$, and put
\[
q=r+1,\qquad a=q-m,
\]
so that $1\le a\le q-2$ and $k+q=N$. The proof of
Proposition~\ref{prop:oneterm} bounds the growth of the coefficient of the
term $\chi_m$ on a diagonal step by the factor
$\gamma_{k+a}\,(1+\Theta(k,a)/k)$, by means of
Proposition~\ref{prop:deficitgrowth}. By
Proposition~\ref{prop:curvaturegrowth} its curvature sum grows by at most the
factor $(k/(k+1))\,e^{\Upsilon_m}$, where we write
\begin{equation}\label{eq:step-growthexponent}
\Upsilon_m=\frac{\bigl(1+\frac{1}{400}\bigr)\,(D_m+1)\,\bigl(q-\frac{m}{2}\bigr)}
{k\,(k+1)}
\end{equation}
for the exponent of that proposition; here $D_m$ is the moment ratio of the
birthday law (Section~\ref{sec:growth}). The same proof bounds the growth of
the boundary term from below by the factor $\gamma_k\,(1+x)^{E_0}$, by means
of Proposition~\ref{prop:boundarygain}, where $x$ and $E_0$ are the two
numbers of that proposition for the integers $k$ and $q$,
\begin{equation}\label{eq:step-gainexponent}
x=\frac{q}{(k-1)(k+q)},\qquad
E_0=\frac{q+2}{2\,q}\;\mu^0_q\Bigl(\frac{q^2}{k-1}\Bigr),
\end{equation}
with the explicit function
\[
\mu^0_q(X)=\frac{\sqrt{(4+X)^2+12\,(q-1)}-(4+X)}{3}
\]
of Section~\ref{sec:gain}. The hypotheses of the three statements are
verified in the proof of Proposition~\ref{prop:oneterm}. That proof
multiplies the three comparisons and so reduces $\chi'_m\,B<\chi_m\,B'$ to
the inequality
\begin{equation}\label{eq:step-factors}
\gamma_{k+a}\,\Bigl(1+\frac{\Theta(k,a)}{k}\Bigr)\cdot\frac{k}{k+1}\cdot
e^{\Upsilon_m}<\gamma_k\,(1+x)^{E_0}.
\end{equation}
This is Lemma~\ref{lem:scalar} below, and it is the source of the strict
sign in the cross inequality.

We prove \eqref{eq:step-factors} in its logarithmic form, with each factor
replaced by a closed bound, and this form is the following lemma. Its
proof first bounds $k^2$ times the left side from above by an explicit
expression in $a$, $q$ and $k$, and $k^2$ times the right side from below by
$k-1/2$ plus a nonnegative term that comes from the gain. It then compares
the two bounds, for $q\le36$ without the gain and for $q>36$ with it.

\begin{lemma}\label{lem:closedforms}
Let $k$, $q$ and $a$ be integers with $k\ge265$, $(q-1)^3\le k^2$, $a\ge1$
and $a+2\le q$, and let $\Theta(k,a)$, $x$ and $E_0$ be given by
\eqref{eq:step-theta} and \eqref{eq:step-gainexponent}. Then
\begin{multline*}
\frac52\Bigl(\frac1k-\frac1{k+a}\Bigr)+\frac{\Theta(k,a)}{k}
+\frac{401}{400}\Bigl(\frac32\sqrt{\frac{\pi\,(q-a)}{2}}-\frac85\Bigr)\cdot
\frac{(q+a)/2}{k\,(k+1)}\\
<\log\Bigl(1+\frac1k\Bigr)+E_0\,\log(1+x).
\end{multline*}
\end{lemma}

\begin{proof}
Let $m=q-a$ and $\xi=a/k$, so that $m\ge2$ and $q\ge3$, and put
\begin{gather*}
\Omega=4.0705\,a+6.675\,\frac{a^2}{k}
+(q+a)\,\bigl(0.94235\,\sqrt m-0.8\bigr),\\
J=k^2\,E_0\,\Bigl(x-\frac{x^2}{2}\Bigr).
\end{gather*}

We first show that
\begin{equation}\label{eq:step-range}
q^3\le1.036^2\,k^2\qquad\text{and}\qquad q-2\le0.151\,k .
\end{equation}
If $q\le42$, then $q^3\le42^3<1.036^2\cdot265^2$ and
$q-2\le40<0.151\cdot265$, so that both inequalities follow from $k\ge265$.
Let $q\ge43$. Then $q/(q-1)\le43/42$, hence
$q^3\le(43/42)^3\,(q-1)^3\le1.036^2\,k^2$, because $(43/42)^3<1.036^2$.
Moreover the quotient $(q-2)^2/(q-1)^3$ decreases for $q\ge4$, and for
$q=43$ it equals $41^2/42^3<0.151^2$. Hence
$(q-2)^2\le0.151^2\,(q-1)^3\le(0.151\,k)^2$. This proves
\eqref{eq:step-range}. Since $a\le q-2$, the second inequality gives
$\xi\le0.151$; in particular $2a\le k$, so that $\Theta(k,a)$ is defined.

Next we show that the left side of the asserted inequality is at most
$\Omega/k^2$, bounding its three terms in turn. The first is
\[
\frac52\Bigl(\frac1k-\frac1{k+a}\Bigr)=\frac{5\,a}{2\,k\,(k+a)}
\le\frac52\cdot\frac{a}{k^2}.
\]
For the second, let $\rho=23/24-(3/2)\,\xi\,(1+\xi/4)$ be the denominator
in \eqref{eq:step-theta}, which is positive. Since $k\ge265$, we have
$1+2/(3(k+2))\le1.0025$. We claim that
\begin{equation}\label{eq:step-chord}
1.0025\,\Bigl(1+\frac32\,\xi+\frac38\,\xi^2\Bigr)
\le(1.047+4.45\,\xi)\,\rho .
\end{equation}
Indeed, the right-hand side minus the left-hand side is
\[
\frac{7}{8000}+\frac{3571}{3000}\,\xi-\frac{119097}{16000}\,\xi^2
-\frac{267}{160}\,\xi^3 .
\]
Since $0\le\xi\le0.151$, we have $\xi^2\le0.151\,\xi$ and
$\xi^3\le0.151^2\,\xi$, so that this number is at least
\[
\frac{7}{8000}+(1.1903-1.1240-0.0381)\,\xi\ge0 .
\]
Dividing \eqref{eq:step-chord} by $\rho$ and multiplying by
$(3/2)\,\xi$, we get $\Theta(k,a)\le(3/2)\,\xi\,(1.047+4.45\,\xi)$, and with
$\xi=a/k$
\[
\frac{\Theta(k,a)}{k}
\le\frac{1}{k^2}\,\Bigl(1.5705\,a+6.675\,\frac{a^2}{k}\Bigr).
\]
For the third term we use $\pi<3.1416$, so that $\pi/2<1.5708<1.25332^2$ and
$\sqrt{\pi m/2}\le1.25332\,\sqrt m$. Since
$(401/400)\cdot(3/2)\cdot1.25332<1.8847$ and $(401/400)\cdot(8/5)\ge1.6$,
\[
\frac{401}{400}\Bigl(\frac32\sqrt{\frac{\pi\,m}{2}}-\frac85\Bigr)
\le1.8847\,\sqrt m-1.6=2\,\bigl(0.94235\,\sqrt m-0.8\bigr).
\]
The right-hand side is positive, because $\sqrt m\ge1.4$. As $q+a\ge0$ and
$k\,(k+1)\ge k^2$, the third term is at most
$(q+a)\,(0.94235\,\sqrt m-0.8)/k^2$. The three bounds add up to
$\Omega/k^2$.

We turn to the right side of the asserted inequality. Let $X=q^2/(k-1)$
be the number at which the function $\mu^0_q$ is taken in
\eqref{eq:step-gainexponent}, so that $E_0$ is $(q+2)/(2q)$ times
$\mu^0_q(X)$. Since $q\ge1$ and $X\ge0$, the
square root in $\mu^0_q(X)$ is at least $4+X$, so that $\mu^0_q(X)\ge0$ and
$E_0\ge0$. Moreover $0\le x\le1/(k-1)$. By the inequality
$\log(1+u)\ge u-u^2/2$ for $u\ge0$ (Section~\ref{sec:prelim}), applied at
$u=1/k$ and at $u=x$, the right side of the asserted inequality is at least
\[
\frac1k-\frac{1}{2\,k^2}+E_0\,\Bigl(x-\frac{x^2}{2}\Bigr)
=\frac{k-\frac12+J}{k^2},
\]
and $J\ge0$. It remains to show that $\Omega<k-1/2+J$, and for this we
distinguish the cases $q\le36$ and $q\ge37$.

Suppose first that $q\le36$. Here we show that $\Omega<k-1/2$, which
suffices because $J\ge0$. The inequality
$(\sqrt m-7/2)^2\ge0$ gives $\sqrt m\le(m+12.25)/7$, and hence
\[
0.94235\,\sqrt m-0.8\le0.134622\,m+0.8491125 .
\]
We insert this bound, with $m=q-a$, and the bound $a^2/k\le a^2/265$ into
the definition of $\Omega$. Since $6.675/265<0.025189$, expanding gives
\[
\Omega\le0.134622\,q^2+0.8491125\,q+4.9196125\,a-0.109433\,a^2 .
\]
The first two terms increase with $q$, and for $q=36$ their sum is less than
$205.1$. The last two terms are together at most
$4.9196125^2/(4\cdot0.109433)<55.3$. Hence
$\Omega<260.4<264.5\le k-1/2$.

Now suppose that $q\ge37$. Then $\sqrt q>6$, and
$q^{3/2}\le1.036\,k$ by \eqref{eq:step-range}. We first bound $J$ from
below. Let $\varsigma=4.657\,\sqrt q+5.25$ and $u=4\,(q-1)/\varsigma$. We
have
\[
(2.577\,\sqrt q-2.75)\,\varsigma-12\,(q-1)
=0.001089\,q+0.7225\,\sqrt q-2.4375 ,
\]
and this is positive because $\sqrt q>6$. As $3u\,\varsigma=12\,(q-1)$ and
$\varsigma>0$, it follows that $3u<2.577\,\sqrt q-2.75$. Since
$1.036\,k\le1.04\,(k-1)$ for $k\ge260$, we have
$q^{3/2}\le1.04\,(k-1)$. As $q^2=\sqrt q\cdot q^{3/2}$, this gives
$X\le1.04\,\sqrt q$. Hence
\[
2\,(4+X)+3u\le8+2.08\,\sqrt q+2.577\,\sqrt q-2.75=\varsigma ,
\]
and therefore
\begin{align*}
(4+X+3u)^2&=(4+X)^2+3u\,\bigl(2\,(4+X)+3u\bigr)\\
&\le(4+X)^2+3u\,\varsigma=(4+X)^2+12\,(q-1).
\end{align*}
Taking square roots, we get $u\le\mu^0_q(X)$, and so
\[
E_0\ge\frac{q+2}{2\,q}\cdot\frac{4\,(q-1)}{\varsigma}
\ge\frac{2\,q}{\varsigma},
\]
because $(q+2)(q-1)\ge q^2$. Next, $x-x^2/2\ge0.998\,x$, because
$x\le1/(k-1)\le1/264$. Finally, since $q^{3/2}\le1.036\,k$,
\[
\sqrt q\,(k-1)(k+q)\le\sqrt q\,k^2+k\,q^{3/2}\le(\sqrt q+1.036)\,k^2 ,
\]
and hence
\[
k^2\,x=\frac{k^2\,q}{(k-1)(k+q)}\ge\frac{q^{3/2}}{\sqrt q+1.036}.
\]
The product of the three bounds is
\[
J\ge0.998\cdot\frac{q^{3/2}}{\sqrt q+1.036}\cdot\frac{2\,q}{\varsigma}
=\frac{1.996\,q^{5/2}}{(\sqrt q+1.036)\,\varsigma}\ge0.305\,q^{3/2} .
\]
The last inequality holds because
\[
1.996\,q-0.305\,(\sqrt q+1.036)\,\varsigma
=0.575615\,q-3.07276886\,\sqrt q-1.658895 ,
\]
and the right-hand side, which is a polynomial of degree two in $\sqrt q$,
is positive for $\sqrt q=6$ and increases for $\sqrt q\ge6$.

We now bound $\Omega$ from above. The inequality
$(\sqrt m-3\sqrt q/5)^2\ge0$ gives
$\sqrt m\le(m+0.36\,q)/(1.2\,\sqrt q)$, and $1/k\le1.036/q^{3/2}$. Let
$h=1/\sqrt q$ and $v=m/q$, so that $0<h<1/6$ and $0<v<1$, and
\[
a=q\,(1-v),\qquad q+a=q\,(2-v),\qquad m=q\,v .
\]
We insert these expressions and the two bounds into the definition of
$\Omega$. Since $6.675\cdot1.036=6.9153$ and $0.94235/1.2=18847/24000$, the
result is $\Omega+1/2\le q^{3/2}\,\mathcal M(h,v)$, where
\begin{multline*}
\mathcal M(h,v)=4.0705\,(1-v)\,h+6.9153\,(1-v)^2\,h^2\\
+(2-v)\,\Bigl(\frac{18847}{24000}\,(v+0.36)-0.8\,h\Bigr)+\frac{h^3}{2};
\end{multline*}
the term $h^3/2$ accounts for $1/2=q^{3/2}h^3/2$. We compare
$\mathcal M(h,v)$ with its values for $h=0$ and for $h=1/6$. Replacing
$h^2$ by $h/6$ and $h^3$ by $h/36$ in $\mathcal M(h,v)$ gives an expression
that is affine in $h$ and that agrees with $\mathcal M(h,v)$ for $h=0$ and
for $h=1/6$. Since the coefficients of $h^2$ and of $h^3$ in
$\mathcal M(h,v)$ are not negative, and since $h^2\le h/6$ and
$h^3\le h/36$, the number $\mathcal M(h,v)$ is at most the value of this
expression, that is,
\[
\mathcal M(h,v)\le(1-6h)\,\mathcal M(0,v)+6h\,\mathcal M(1/6,v),
\]
where
\begin{align*}
\mathcal M(0,v)&=\frac{18847}{24000}\,(2-v)\,(v+0.36)
=\frac{18847}{24000}\,\bigl(0.72+1.64\,v-v^2\bigr),\\
\mathcal M(1/6,v)&=\mathcal M(0,v)+\frac16\,\Bigl(4.0705\,(1-v)
+\frac{6.9153}{6}\,(1-v)^2\\
&\qquad\qquad\qquad-0.8\,(2-v)+\frac{1}{72}\Bigr).
\end{align*}
Both are polynomials of degree two in $v$. The first is largest at
$v=0.82$, where its value is less than $1.094$. Expanding the second, one
finds that, for $v\ge0$,
\[
\mathcal M(1/6,v)\le1.1716+0.3587\,v-0.5932\,v^2 ,
\]
and the right-hand side of this inequality is at most
$1.1716+0.3587^2/(4\cdot0.5932)<1.226$. As $0\le6h\le1$, it follows that
$\mathcal M(h,v)\le1.226$, and therefore
\[
\Omega+\frac12\le1.226\,q^{3/2}
<\Bigl(\frac{250}{259}+0.305\Bigr)\,q^{3/2}\le k+J .
\]
Here $250/259=1/1.036$, so that $(250/259)\,q^{3/2}\le k$, and
$0.305\,q^{3/2}\le J$. Thus $\Omega<k-1/2+J$ in this case as well.
\end{proof}

With the bounds for the single factors, the lemma gives the inequality
\eqref{eq:step-factors} between the factors of the three comparisons.

\begin{lemma}\label{lem:scalar}
Let $k$, $q$ and $a$ be integers with $k\ge265$, $(q-1)^3\le k^2$, $a\ge1$
and $a+2\le q$, and put $m=q-a$. Let $\Theta(k,a)$ be given by
\eqref{eq:step-theta}, the exponent $\Upsilon_m$ of the curvature sum by
\eqref{eq:step-growthexponent}, and the numbers $x$ and $E_0$ of the gain of
the boundary term by \eqref{eq:step-gainexponent}. Then the
diagonal factors $\gamma_{k+a}$ and $\gamma_k$ satisfy
\[
\gamma_{k+a}\,\Bigl(1+\frac{\Theta(k,a)}{k}\Bigr)\cdot\frac{k}{k+1}\cdot
e^{\Upsilon_m}<\gamma_k\,(1+x)^{E_0}.
\]
\end{lemma}

\begin{proof}
Since $(2\,(q-1))^3\le8\,k^2\le k^3$, we have $2a\le2\,(q-1)\le k$, so that
$\Theta(k,a)$ is defined and nonnegative. Both sides of the asserted
inequality are positive, and it suffices to compare their logarithms, that
is, to show that
\begin{multline*}
\log\gamma_{k+a}-\log\gamma_k+\log\Bigl(1+\frac{\Theta(k,a)}{k}\Bigr)
+\Upsilon_m\\
<\log\Bigl(1+\frac1k\Bigr)+E_0\,\log(1+x).
\end{multline*}
By Lemma~\ref{lem:factorlog}, applied to indices $2\le k\le k+a$, the
difference $\log\gamma_{k+a}-\log\gamma_k$ is at most
$(5/2)(1/k-1/(k+a))$. The logarithm of $1+\Theta(k,a)/k$ is at most
$\Theta(k,a)/k$, because $\log y\le y-1$ for $y>0$. By
Lemma~\ref{lem:momentquotient}, which applies because $m\ge2$, we have
$D_m+1\le(3/2)\sqrt{\pi m/2}-8/5$, and $q-m/2=(q+a)/2\ge0$, so that
\[
\Upsilon_m\le\frac{401}{400}\Bigl(\frac32\sqrt{\frac{\pi\,m}{2}}-\frac85\Bigr)
\cdot
\frac{(q+a)/2}{k\,(k+1)}.
\]
Since $m=q-a$, the sum of the three bounds is the left side of the
inequality of Lemma~\ref{lem:closedforms}, which has the same hypotheses,
and the right side of that inequality is the right side of the inequality
to be shown. The assertion follows.
\end{proof}

\subsection{The cross inequality}\label{subsec:step-cross}

We can now compare one term of the curvature term with the boundary term, at
a position behind the curve and at its diagonal neighbour. In the notation
of \eqref{eq:step-terms} and \eqref{eq:step-neighbour}, the following
proposition states that $\chi'_m\,B<\chi_m\,B'$ for $2\le m\le r$.

\begin{proposition}\label{prop:oneterm}
Let $k$, $r$ and $m$ be integers with $k\ge265$, $r^3\le k^2$ and
$2\le m\le r$, and put $N=k+r+1$. Then
\begin{multline*}
\frac{\varkappa_{k+1}(N+1-m)}{N+1-m}\;\mathsf C_m(k+1;N+1)\cdot
\frac{a_k}{k}\;\Lambda_{r+1}(k-1;N)\\
<\frac{\varkappa_k(N-m)}{N-m}\;\mathsf C_m(k;N)\cdot
\frac{a_{k+1}}{k+1}\;\Lambda_{r+1}(k;N+1).
\end{multline*}
\end{proposition}

\begin{proof}
Let $q=r+1$, $a=q-m$ and $t=k+a$. Then $1\le a\le r-1$, $a+2\le q$,
$t=N-m$ and $t+1=N+1-m$, and $N=k+q$. Let $\Upsilon_m$ be the exponent
\eqref{eq:step-growthexponent}, and let $x$ and $E_0$ be the numbers
\eqref{eq:step-gainexponent}.

We have $2r\le k$, because $(2r)^3\le8\,k^2\le k^3$. Hence
$r\le k-1$, so that the notation of \eqref{eq:step-terms} and
\eqref{eq:step-neighbour} is available, and $2a\le2\,(r-1)\le k$. Moreover
$401\,m^2\le12\,k^2$. Indeed, if $r\le41$, then $m\le41$ and
$401\cdot41^2\le12\cdot265^2$; and if $r\ge42$, then
$m^2\le r^2\le r^3/42\le k^2/42$ and $401\le12\cdot42$.

We first compare the coefficients $\varkappa_k(t)/t$ and
$\varkappa_{k+1}(t+1)/(t+1)$ of the terms $\chi_m$ and $\chi'_m$.
Proposition~\ref{prop:deficitgrowth} applies because
$2a\le k$. By \eqref{eq:prelim-quotient} and \eqref{eq:diagonal-factor} we
have $a_{n+1}=\gamma_n\,(1+1/n)\,a_n$ for $n\ge2$, and on dividing by $n+1$
we obtain
\begin{equation}\label{eq:step-cayleystep}
\frac{a_{n+1}}{n+1}=\gamma_n\,\frac{a_n}{n}\qquad(n\ge2).
\end{equation}
We multiply the inequality of Proposition~\ref{prop:deficitgrowth} by
$a_{t+1}/(t+1)$, which equals $\gamma_t\,a_t/t$ by
\eqref{eq:step-cayleystep}, and obtain
\begin{equation}\label{eq:step-coefficient}
\frac{\varkappa_{k+1}(t+1)}{t+1}\le\gamma_t\,
\Bigl(1+\frac{\Theta(k,a)}{k}\Bigr)\,\frac{\varkappa_k(t)}{t}.
\end{equation}

Next, Proposition~\ref{prop:curvaturegrowth} compares
the curvature sums $\mathsf C_m(k+1;k+q+1)$ and $\mathsf C_m(k;k+q)$ under
the hypotheses $k\ge1$, $m\ge2$, $m+1\le q$ and $401\,m^2\le12\,k^2$. These
hold for $q=r+1$, and since $k+q=N$ the proposition gives
\begin{equation}\label{eq:step-curvature}
\mathsf C_m(k+1;N+1)\le\frac{k}{k+1}\,e^{\Upsilon_m}\;\mathsf C_m(k;N).
\end{equation}

For the boundary terms, Proposition~\ref{prop:boundarygain}, for the
integers $k\ge2$ and $q=r+1\ge1$, states that
\[
(1+x)^{E_0}\,\Lambda_{r+1}(k-1;N)\le\Lambda_{r+1}(k;N+1).
\]
We multiply by $\gamma_k\,a_k/k$, which equals $a_{k+1}/(k+1)$ by
\eqref{eq:step-cayleystep}, and obtain
\begin{equation}\label{eq:step-boundary}
\gamma_k\,(1+x)^{E_0}\,B\le B' .
\end{equation}

Before we multiply these inequalities, we check the signs of their terms.
The kernel deficits $\varkappa_k(t)$ and
$\varkappa_{k+1}(t+1)$ are positive by Lemma~\ref{lem:deficitpositive},
since $1\le t\le2k+1$ and $1\le t+1\le2\,(k+1)+1$. The boundary term $B$ is
positive by \eqref{eq:step-chainpositive}. The curvature sum
$\mathsf C_m(k;N)$ is positive as well. Indeed, by
Lemma~\ref{lem:curvatureparts},
\[
\mathsf C_m(k;N)=\sum_{Q=1}^{m}\bigl((k+1)^Q-2\,k^Q+(k-1)^Q\bigr)\,
V_Q\bigl(m;\lambda^{1,N}\bigr).
\]
Here $Q$ is the number of parts, which Lemma~\ref{lem:curvatureparts} writes
$q$; in this proof $q$ denotes $r+1$.
No term of this sum is negative: the second differences
$(k+1)^Q-2\,k^Q+(k-1)^Q$ are nonnegative by convexity, and so are the
weights $V_Q(m;\lambda^{1,N})$, whose cut factors $1/(N-s)$ are positive for
$0<s<m$. The term with $Q=2$ is $2\,V_2(m;\lambda^{1,N})$, twice the sum of
$a_s\,a_{m-s}/(N-s)$ over $0<s<m$, which is positive because
$m\ge2$. Hence $\mathsf C_m(k;N)>0$, and $\chi_m>0$.

Multiplying \eqref{eq:step-curvature} by
$\varkappa_{k+1}(t+1)/(t+1)$ and using \eqref{eq:step-coefficient}, we now
get
\begin{equation}\label{eq:step-termgrowth}
\chi'_m\le\gamma_t\,\Bigl(1+\frac{\Theta(k,a)}{k}\Bigr)\,\frac{k}{k+1}\,
e^{\Upsilon_m}\cdot\chi_m .
\end{equation}
By Lemma~\ref{lem:scalar} for $k$, $q$ and $a$, whose hypothesis
$(q-1)^3\le k^2$ is $r^3\le k^2$ and in which $q-a=m$ and $k+a=t$, the
factor of $\chi_m$ in \eqref{eq:step-termgrowth} is smaller than
$\gamma_k\,(1+x)^{E_0}$. As $\chi_m\,B>0$, we conclude with
\eqref{eq:step-boundary} that
\begin{align*}
\chi'_m\,B&\le\gamma_t\,\Bigl(1+\frac{\Theta(k,a)}{k}\Bigr)\,\frac{k}{k+1}\,
e^{\Upsilon_m}\cdot\chi_m\,B\\
&<\gamma_k\,(1+x)^{E_0}\cdot\chi_m\,B\le\chi_m\,B' ,
\end{align*}
which is the assertion.
\end{proof}

The proof of the following proposition sums the inequalities of
Proposition~\ref{prop:oneterm} over $m$ and obtains the cross inequality
\eqref{eq:step-cross} for $r\ge2$.

\begin{proposition}\label{prop:diagonalstep}
Let $k$ and $r$ be integers with $k\ge265$,
$r\ge1$ and $r^3\le k^2$. If $d_{k+r,k}\le0$, then $d_{k+r+1,k+1}<0$.
\end{proposition}

\begin{proof}
Since $(2r)^3\le8\,k^2\le k^3$, we have $2r\le k$, and hence
$1\le r\le k-1$. Both positions therefore lie in the first layer of their
columns, and \eqref{eq:step-twopositions} holds with the notation of
\eqref{eq:step-terms} and \eqref{eq:step-neighbour}. The hypothesis
$d_{k+r,k}\le0$ is $C\le B$, and the claim $d_{k+r+1,k+1}<0$ is $C'<B'$. The
boundary terms $B$ and $B'$ are positive by \eqref{eq:step-chainpositive},
and the terms with $m=1$ vanish by \eqref{eq:step-firstterm}.

If $r=1$, then $C'=\chi'_1=0$, and $C'<B'$ holds because $B'>0$. Let
$r\ge2$. For $2\le m\le r$, Proposition~\ref{prop:oneterm} gives
$\chi'_m\,B<\chi_m\,B'$, and $\chi'_1\,B=\chi_1\,B'$. The sum of these
relations over $1\le m\le r$, in which the one for $m=2$ is strict, is the
cross inequality $C'\,B<C\,B'$. Since $C\le B$ and $B'>0$, it follows that
$C'\,B<C\,B'\le B\,B'$, and hence $C'<B'$.
\end{proof}

The restriction of Proposition~\ref{prop:diagonalstep} to the positions
behind the curve comes from Lemma~\ref{lem:closedforms}, whose inequality
is proved in the range $(q-1)^3\le k^2$. For
columns $k\le264$ the diagonal rule in the first layer is
obtained with the help of the computed lemmas of
Section~\ref{sec:tables}, in Proposition~\ref{prop:rulesmall}.

\section{The two rules and the proofs of the theorems}\label{sec:staircase}
We now prove Theorems~\ref{thm:peaks} and~\ref{thm:main}. The column
rule and the diagonal rule for the array $(b_{n,k})$ are
assembled from the statements of the earlier sections, region by region
(Propositions~\ref{prop:rulecolumn} and~\ref{prop:rulediagonal}), and an
induction on the row, for an arbitrary array of natural numbers, derives
from them the conclusion of Theorem~\ref{thm:peaks}
(Proposition~\ref{prop:rowinduction}). We pass between an inequality for a step
$d_{n,k}$ and the corresponding inequality between two entries of the array
without comment.

\subsection{The two rules}\label{subsec:staircase-rules}

We begin with the column rule, which has been proved in two parts, according
to the region in which the lower position $(n+1,k)$ lies.

\begin{proposition}\label{prop:rulecolumn}
Let $n\ge2$ and $k$ be integers with $1\le k\le n-1$. If $d_{n,k}>0$, then
$d_{n+1,k}>0$.
\end{proposition}

\begin{proof}
Suppose that $d_{n,k}>0$. If $n+1\le2k-1$, then $k+1\le n\le2k-2$, and
Proposition~\ref{prop:layerpersistence} gives $d_{n+1,k}>0$. If
$n+1\ge2k$, the claim is Corollary~\ref{cor:columnbeyond}.
\end{proof}

We turn to the diagonal rule, and
begin with the first layer of columns $2\le k\le264$, which is as far as the
computations of
Section~\ref{sec:tables} reach.

\begin{proposition}\label{prop:rulesmall}
Let $k$ and $n$ be integers with $2\le k\le264$ and $k+1\le n\le2k-1$. If
$d_{n,k}\le0$, then $d_{n+1,k+1}<0$.
\end{proposition}

\begin{proof}
Let first $2\le k\le8$. Then $3\le k+1\le9$ and
$(k+1)+1\le n+1\le2(k+1)-2$, so that Lemma~\ref{lem:fortyfive}, applied to
column $k+1$ and to row $n+1$ with its fourth condition, gives
$d_{n+1,k+1}<0$.

Now let $9\le k\le264$, and put $j=k+\zeta(k)+1$, where $\zeta(k)$ is the
number of Definition~\ref{def:zeta}. By Lemma~\ref{lem:signpair} we have
$d_{j,k}>0$. If $n\ge j$, then $k+1\le j\le n\le2k-1$, and
Corollary~\ref{cor:layerpersistenceall} gives $d_{n,k}>0$, contrary to the
hypothesis. Hence $n\le j-1$, so that $k+2\le n+1\le k+\zeta(k)+1$, and
Lemma~\ref{lem:nextcolumn} gives $d_{n+1,k+1}<0$.
\end{proof}

We can now prove the diagonal rule at all positions.

\begin{proposition}\label{prop:rulediagonal}
Let $n\ge2$ and $k$ be integers with $1\le k\le n-1$. If $d_{n,k}\le0$, then
$d_{n+1,k+1}<0$.
\end{proposition}

\begin{proof}
Suppose that $d_{n,k}\le0$. If $n\ge2k$, the claim is
Corollary~\ref{cor:diagonalbeyond}. Let $n\le2k-1$, so that $k\ge2$ and
$k+1\le n\le2k-1$. For $k\le264$ the claim is
Proposition~\ref{prop:rulesmall}.

It remains to consider $k\ge265$, with $k+1\le n\le2k-1$ as before. By
Proposition~\ref{prop:curveincrease} there is an integer $j>k$ with
$(j-k-1)^3\le k^2$ and $d_{j,k}>0$. If $n\ge j$, then
$k+1\le j\le n\le2k-1$, and Corollary~\ref{cor:layerpersistenceall} gives
$d_{n,k}>0$, contrary to the hypothesis. Hence $n\le j-1$, so that the
deficit $r=n-k$ satisfies $1\le r\le j-k-1$ and
\[
r^3\le(j-k-1)^3\le k^2 .
\]
Since
$d_{k+r,k}=d_{n,k}\le0$, Proposition~\ref{prop:diagonalstep} gives
$d_{n+1,k+1}<0$.
\end{proof}

\subsection{The induction on the row}\label{subsec:staircase-induction}

It remains to pass from the two rules to the shape of the rows. We do this
for an arbitrary array of
natural numbers, with the notions of Section~\ref{sec:intro}.

\begin{proposition}\label{prop:rowinduction}
Let $f(n,k)$, for integers $n\ge0$ and $k\ge0$, be an array of natural
numbers. Suppose that for all integers $n\ge2$ and $k$ with $1\le k\le n-1$
the following two implications hold: the column rule,
\[
f(n,k-1)<f(n,k)\quad\text{implies}\quad f(n+1,k-1)<f(n+1,k),
\]
and the diagonal rule,
\[
f(n,k)\le f(n,k-1)\quad\text{implies}\quad f(n+1,k+1)<f(n+1,k).
\]
Then for every integer $n\ge2$ there is an integer $p\ge0$ such that row
$n$ of $f$ has its peak at position $p$, and row $n+1$ of $f$ has
its peak at position $p$ or a strict peak at position $p+1$.
\end{proposition}

\begin{proof}
Let $n\ge2$, and suppose that row $n$ of $f$ has its peak at position $p$.
Then $p\le n-1$, and we have $f(n,k-1)<f(n,k)$ for $1\le k\le p$ and
$f(n,k)\le f(n,k-1)$ for $p+1\le k\le n-1$. The column rule at the
positions $(n,k)$ with $k\le p$ and the diagonal rule at those with
$k\ge p+1$ therefore give
\[
f(n+1,0)<\dots<f(n+1,p)\quad\text{and}\quad
f(n+1,p+1)>\dots>f(n+1,n).
\]
Hence row $n+1$ has a strict peak at position $p+1$ if
$f(n+1,p)<f(n+1,p+1)$, and its peak at position $p$ otherwise; in either
case it has its peak at some position. Since row $2$ consists of two
entries, it has its peak at position $0$ or $1$, and it follows by
induction on $n$ that every row $n\ge2$ of $f$ has its peak at some
position $p$. By what we have shown, row $n+1$ then has its peak at
position $p$ or a strict peak at position $p+1$.
\end{proof}

\subsection{Proofs of the two theorems}\label{subsec:staircase-theorems}

\begin{proof}[Proof of Theorem~\ref{thm:peaks}]
By Propositions~\ref{prop:rulecolumn} and~\ref{prop:rulediagonal}, the array
$(b_{n,k})$ satisfies the column rule and the diagonal rule.
The theorem therefore follows from Proposition~\ref{prop:rowinduction},
applied to this array.
\end{proof}

\begin{proof}[Proof of Theorem~\ref{thm:main}]
For $n=1$ the sequence has a single term and is unimodal. For $n\ge2$,
Theorem~\ref{thm:peaks} gives a position $p$ at which row $n$ of the array
$(b_{n,k})$ has its peak, and the definition of a unimodal
sequence is satisfied with $j^*=p$.
\end{proof}

\bibliographystyle{amsplain}
\bibliography{references}
\end{document}
